\input{preamble}

% OK, start here.
%
\begin{document}

\title{Cohomologie des faisceaux}


\maketitle

\phantomsection
\label{section-phantom}

\tableofcontents

\section{Introduction}

\label{section-introduction}

\noindent
Dans ce document, nous étudions quelques thèmes de cohomologie des faisceaux sur les espaces topologiques. Nous travaillons principalement dans le cadre général des modules sur un faisceau d'anneaux et des morphismes d'espaces annelés. Pour le cas des faisceaux sur les sites, voir le chapitre Cohomologie sur les sites, section \ref{sites-cohomology-section-introduction}. Les références de base sont \cite{Godement} et \cite{Iversen}.





\section{Cohomologie des faisceaux}

\label{section-cohomology-sheaves}

\noindent

Soit $X$ un espace topologique. Soit $\mathcal{F}$ un faisceau abélien. La catégorie des faisceaux abéliens sur $X$ possède suffisamment d'objets injectifs, voir Injectifs, lemme \ref{injectives-lemma-abelian-sheaves-space}. Nous pouvons donc choisir une résolution injective
$\mathcal{F}[0] \to \mathcal{I}^\bullet$. Comme il est d'usage, nous définissons

\begin{equation}
\label{equation-cohomology}
H^i(X, \mathcal{F}) = H^i(\Gamma(X, \mathcal{I}^\bullet))
\end{equation}

le membre de gauche comme le {\it $i$-ième groupe de cohomologie du faisceau abélien $\mathcal{F}$}. La famille des foncteurs $H^i(X, -)$ forme un $\delta$-foncteur universel de $\textit{Ab}(X) \to \textit{Ab}$.

\medskip\noindent

Soit $f : X \to Y$ une application continue d'espaces topologiques. Avec
$\mathcal{F}[0] \to \mathcal{I}^\bullet$ comme ci-dessus, nous définissons

\begin{equation}
\label{equation-higher-direct-image}
R^if_*\mathcal{F} = H^i(f_*\mathcal{I}^\bullet)
\end{equation}

le membre de gauche comme la {\it $i$-ième image directe supérieure de $\mathcal{F}$}. La famille des foncteurs $R^if_*$ forme un $\delta$-foncteur universel de $\textit{Ab}(X) \to \textit{Ab}(Y)$.

\medskip\noindent

Soit $(X, \mathcal{O}_X)$ un espace annelé. Soit $\mathcal{F}$ un
$\mathcal{O}_X$-module. La catégorie des $\mathcal{O}_X$-modules sur $X$ possède suffisamment d'objets injectifs, voir Injectifs, lemme \ref{injectives-lemma-sheaves-modules-space}. Nous pouvons donc choisir une résolution injective
$\mathcal{F}[0] \to \mathcal{I}^\bullet$. Comme il est d'usage, nous définissons

\begin{equation}
\label{equation-cohomology-modules}
H^i(X, \mathcal{F}) = H^i(\Gamma(X, \mathcal{I}^\bullet))
\end{equation}

le membre de gauche comme le {\it $i$-ième groupe de cohomologie de $\mathcal{F}$}. La famille des foncteurs $H^i(X, -)$ forme un $\delta$-foncteur universel de $\textit{Mod}(\mathcal{O}_X) \to \text{Mod}_{\mathcal{O}_X(X)}$.

\medskip\noindent

Soit $f : (X, \mathcal{O}_X) \to (Y, \mathcal{O}_Y)$ un morphisme d'espaces annelés. Avec $\mathcal{F}[0] \to \mathcal{I}^\bullet$ comme ci-dessus, nous définissons

\begin{equation}
\label{equation-higher-direct-image-modules}
R^if_*\mathcal{F} = H^i(f_*\mathcal{I}^\bullet)
\end{equation}

le membre de gauche comme la {\it $i$-ième image directe supérieure de $\mathcal{F}$}. La famille des foncteurs $R^if_*$ forme un $\delta$-foncteur universel de $\textit{Mod}(\mathcal{O}_X) \to \textit{Mod}(\mathcal{O}_Y)$.





\section{Foncteurs dérivés}

\label{section-derived-functors}

\noindent
Nous expliquons brièvement comment obtenir les foncteurs dérivés à droite à l'aide de foncteurs de résolution. Pour les foncteurs dérivés non bornés, voir la section \ref{section-unbounded}.

\medskip\noindent

Soit $(X, \mathcal{O}_X)$ un espace annelé. La catégorie
$\textit{Mod}(\mathcal{O}_X)$ est abélienne, voir Faisceaux de modules, lemme \ref{modules-lemma-abelian}. Dans ce chapitre, nous écrirons
$$
K(\mathcal{O}_X) = K(\textit{Mod}(\mathcal{O}_X))
\quad
\text{et}
\quad
D(\mathcal{O}_X) = D(\textit{Mod}(\mathcal{O}_X)).
$$
et de même pour les versions bornées des catégories triangulées introduites dans Catégories dérivées, définition \ref{derived-definition-complexes-notation} et définition \ref{derived-definition-unbounded-derived-category}. D'après Catégories dérivées, remarque \ref{derived-remark-big-abelian-category},
il existe un foncteur de résolution
$$
j = j_X :
K^{+}(\textit{Mod}(\mathcal{O}_X))
\longrightarrow
K^{+}(\mathcal{I})
$$
où $\mathcal{I}$ est la sous-catégorie additive strictement pleine de
$\textit{Mod}(\mathcal{O}_X)$ formée des faisceaux injectifs. Pour tout foncteur exact à gauche
$F : \textit{Mod}(\mathcal{O}_X) \to \mathcal{B}$
à valeurs dans une catégorie abélienne $\mathcal{B}$, nous noterons $RF$ le foncteur dérivé à droite décrit dans Catégories dérivées, section \ref{derived-section-right-derived-functor},
et construit à l'aide du foncteur de résolution $j_X$ que nous venons de décrire :

\begin{equation}
\label{equation-RF}
RF = F \circ j_X' : D^{+}(X) \longrightarrow D^{+}(\mathcal{B})
\end{equation}

Voir Catégories dérivées, lemme \ref{derived-lemma-right-derived-functor},
pour cette notation. Remarquons que l'on peut considérer $RF$ comme défini sur
$\textit{Mod}(\mathcal{O}_X)$,
$\text{Comp}^{+}(\textit{Mod}(\mathcal{O}_X))$,
$K^{+}(X)$ ou $D^{+}(X)$,
selon la situation. D'après Catégories dérivées, définition \ref{derived-definition-higher-derived-functors},
nous obtenons le $i$-ième foncteur dérivé à droite

\begin{equation}
\label{equation-RFi}
R^iF = H^i \circ RF : \textit{Mod}(\mathcal{O}_X) \longrightarrow \mathcal{B}
\end{equation}

de sorte que $R^0F = F$ et que $\{R^iF, \delta\}_{i \geq 0}$ est un
$\delta$-foncteur universel, voir Catégories dérivées, lemme \ref{derived-lemma-higher-derived-functors}.

\medskip\noindent

Voici deux cas particuliers de cette construction. Pour un anneau $R$, nous écrivons $K(R) = K(\text{Mod}_R)$ et
$D(R) = D(\text{Mod}_R)$, et de même pour les versions bornées. Pour tout ouvert $U \subset X$, nous avons un foncteur exact à gauche
$
\Gamma(U, -) :
\textit{Mod}(\mathcal{O}_X)
\longrightarrow
\text{Mod}_{\mathcal{O}_X(U)}
$
qui donne lieu à

\begin{equation}
\label{equation-total-derived-cohomology}
R\Gamma(U, -) :
D^{+}(X)
\longrightarrow
D^{+}(\mathcal{O}_X(U))
\end{equation}

d'après la discussion précédente. Nous posons $H^i(U, -) = R^i\Gamma(U, -)$. Si $U = X$, nous retrouvons (\ref{equation-cohomology-modules}). Si $f : X \to Y$ est un morphisme d'espaces annelés, nous avons le foncteur exact à gauche
$
f_* :
\textit{Mod}(\mathcal{O}_X)
\longrightarrow
\textit{Mod}(\mathcal{O}_Y)
$
qui donne lieu à l'{\it image directe dérivée}

\begin{equation}
\label{equation-total-derived-direct-image}
Rf_* :
D^{+}(X)
\longrightarrow
D^{+}(Y)
\end{equation}

Le $i$-ième faisceau de cohomologie de $Rf_*\mathcal{F}^\bullet$ est noté
$R^if_*\mathcal{F}^\bullet$ et appelé la $i$-ième {\it image directe supérieure},
conformément à (\ref{equation-higher-direct-image-modules}). Les deux foncteurs affichés ci-dessus sont des foncteurs exacts de catégories dérivées.

\medskip\noindent
{\bf Abus de notation:} Lorsque le foncteur $Rf_*$, ou tout autre foncteur dérivé, est appliqué à un faisceau $\mathcal{F}$ sur $X$ ou à un complexe de faisceaux, il est entendu que $\mathcal{F}$ a été remplacé par une résolution appropriée de $\mathcal{F}$. Afin de faciliter ce type d'opération, nous dirons, étant donné un objet $\mathcal{F}^\bullet \in D(\mathcal{O}_X)$, qu'un complexe borné inférieurement $\mathcal{I}^\bullet$ d'objets injectifs de $\textit{Mod}(\mathcal{O}_X)$ {\it représente $\mathcal{F}^\bullet$ dans la catégorie dérivée} s'il existe un quasi-isomorphisme $\mathcal{F}^\bullet \to \mathcal{I}^\bullet$. Dans le même esprit, la phrase « soit $\alpha : \mathcal{F}^\bullet \to \mathcal{G}^\bullet$ un morphisme de $D(\mathcal{O}_X)$ » ne signifie pas que $\alpha$ est représenté par un morphisme de complexes. S'il s'agit effectivement d'un morphisme de complexes, nous le préciserons.









\section{Première cohomologie et torseurs}

\label{section-h1-torsors}

\begin{definition}

\label{definition-torsor}

Soit $X$ un espace topologique. Soit $\mathcal{G}$ un faisceau de groupes (éventuellement non commutatifs) sur $X$. Un {\it pseudo-torseur}, ou plus précisément un {\it pseudo-$\mathcal{G}$-torseur}, est un faisceau d'ensembles $\mathcal{F}$ sur $X$ muni d'une action
$\mathcal{G} \times \mathcal{F} \to \mathcal{F}$ telle que

\begin{enumerate}

\item dès que $\mathcal{F}(U)$ est non vide, l'action
$\mathcal{G}(U) \times \mathcal{F}(U) \to \mathcal{F}(U)$
est simplement transitive

\end{enumerate}

Un {\it morphisme de pseudo-$\mathcal{G}$-torseurs}  $\mathcal{F} \to \mathcal{F}'$
est un morphisme de faisceaux d'ensembles compatible avec les
actions de $\mathcal{G}$. Un {\it torseur}, ou plus précisément un
{\it $\mathcal{G}$-torseur}, est un pseudo-$\mathcal{G}$-torseur tel que, de plus,

\begin{enumerate}

\item[(2)] pour tout $x \in X$, la fibre $\mathcal{F}_x$ est non vide.

\end{enumerate}

Un {\it morphisme de $\mathcal{G}$-torseurs} est un morphisme de pseudo-$\mathcal{G}$-torseurs. Le {\it $\mathcal{G}$-torseur trivial}
est le faisceau $\mathcal{G}$ muni de l'action à gauche évidente de
$\mathcal{G}$.

\end{definition}

\noindent
Il est clair qu'un morphisme de torseurs est automatiquement un isomorphisme.

\begin{lemma}

\label{lemma-trivial-torsor}
Soit $X$ un espace topologique. Soit $\mathcal{G}$ un faisceau de groupes (éventuellement non commutatifs) sur $X$. Un $\mathcal{G}$-torseur $\mathcal{F}$ est trivial si et seulement si $\mathcal{F}(X) \not = \emptyset$.
\end{lemma}

\begin{proof}
Démonstration omise.
\end{proof}

\begin{lemma}

\label{lemma-torsors-h1}
Soit $X$ un espace topologique. Soit $\mathcal{H}$ un faisceau abélien sur $X$. Il existe une bijection canonique entre l'ensemble des classes d'isomorphisme de $\mathcal{H}$-torseurs et $H^1(X, \mathcal{H})$.
\end{lemma}

\begin{proof}

Soit $\mathcal{F}$ un $\mathcal{H}$-torseur. Considérons le faisceau abélien libre $\mathbf{Z}[\mathcal{F}]$
sur $\mathcal{F}$. C'est la faisceautisation du préfaisceau qui associe à tout ouvert $U \subset X$ l'ensemble des sommes formelles finies $\sum n_i[s_i]$ avec $n_i \in \mathbf{Z}$
et $s_i \in \mathcal{F}(U)$. Il existe un morphisme naturel
$$
\sigma : \mathbf{Z}[\mathcal{F}] \longrightarrow \underline{\mathbf{Z}}
$$
qui, à une section locale $\sum n_i[s_i]$, associe $\sum n_i$. Le noyau de $\sigma$ est engendré par les sections locales de la forme
$[s] - [s']$. Il y a une application canonique
$a : \Ker(\sigma) \to \mathcal{H}$
qui envoie $[s] - [s'] \mapsto h$, où $h$ est la section locale de
$\mathcal{H}$ telle que $h \cdot s' = s$. Considérons le diagramme cocartésien
$$
\xymatrix{
0 \ar[r] &
\Ker(\sigma) \ar[r] \ar[d]^a &
\mathbf{Z}[\mathcal{F}] \ar[r] \ar[d] &
\underline{\mathbf{Z}} \ar[r] \ar[d] &
0 \\
0 \ar[r] &
\mathcal{H} \ar[r] &
\mathcal{E} \ar[r] &
\underline{\mathbf{Z}} \ar[r] &
0
}
$$
Ici, $\mathcal{E}$ est l'extension obtenue par somme amalgamée. La longue suite exacte de cohomologie associée à la suite exacte courte inférieure fournit un élément
$\xi = \xi_\mathcal{F} \in H^1(X, \mathcal{H})$
en appliquant le morphisme de bord à $1 \in H^0(X, \underline{\mathbf{Z}})$.

\medskip\noindent

Réciproquement, étant donné $\xi \in H^1(X, \mathcal{H})$, nous pouvons associer à $\xi$
un torseur comme suit. Choisissons un plongement $\mathcal{H} \to \mathcal{I}$
de $\mathcal{H}$ dans un faisceau abélien injectif $\mathcal{I}$. Posons
$\mathcal{Q} = \mathcal{I}/\mathcal{H}$, de sorte que nous avons une suite exacte courte
$$
\xymatrix{
0 \ar[r] &
\mathcal{H} \ar[r] &
\mathcal{I} \ar[r] &
\mathcal{Q} \ar[r] &
0
}
$$
L'élément $\xi$ est l'image d'une section globale $q \in H^0(X, \mathcal{Q})$
car $H^1(X, \mathcal{I}) = 0$ (voir Catégories dérivées, lemme \ref{derived-lemma-higher-derived-functors}). Soit $\mathcal{F} \subset \mathcal{I}$ le sous-faisceau (d'ensembles) des sections qui s'envoient sur $q$ dans le faisceau $\mathcal{Q}$. Il est facile de vérifier que
$\mathcal{F}$ est un torseur.

\medskip\noindent

Nous omettons de vérifier que les deux constructions données ci-dessus sont inverses l'une de l'autre.

\end{proof}







\section{Première cohomologie et extensions}

\label{section-h1-extensions}

\begin{lemma}

\label{lemma-h1-extensions}
Soit $(X, \mathcal{O}_X)$ un espace annelé. Soit $\mathcal{F}$ un faisceau de $\mathcal{O}_X$-modules. Il existe une bijection canonique $$
\Ext^1_{\textit{Mod}(\mathcal{O}_X)}(\mathcal{O}_X, \mathcal{F})
\longrightarrow
H^1(X, \mathcal{F})
$$ qui associe à l'extension $$
0 \to \mathcal{F} \to \mathcal{E} \to \mathcal{O}_X \to 0
$$ l'image de $1 \in \Gamma(X, \mathcal{O}_X)$ dans $H^1(X, \mathcal{F})$.
\end{lemma}

\begin{proof}
Construisons l'inverse de l'application donnée dans le lemme. Soit $\xi \in H^1(X, \mathcal{F})$. Choisissons une injection $\mathcal{F} \subset \mathcal{I}$, où $\mathcal{I}$ est injectif dans $\textit{Mod}(\mathcal{O}_X)$. Posons $\mathcal{Q} = \mathcal{I}/\mathcal{F}$. Par la longue suite exacte de cohomologie, $\xi$ est l'image d'une section $\tilde \xi \in \Gamma(X, \mathcal{Q}) =
\Hom_{\mathcal{O}_X}(\mathcal{O}_X, \mathcal{Q})$. Formons alors le produit fibré $$
\xymatrix{
0 \ar[r] &
\mathcal{F} \ar[r] \ar@{=}[d] &
\mathcal{E} \ar[r] \ar[d] &
\mathcal{O}_X \ar[r] \ar[d]^{\tilde \xi} &
0 \\
0 \ar[r] &
\mathcal{F} \ar[r] &
\mathcal{I} \ar[r] &
\mathcal{Q} \ar[r] &
0
}
$$ voir Homologie, section \ref{homology-section-extensions}.
\end{proof}







\section{Première cohomologie et faisceaux inversibles}

\label{section-invertible-sheaves}

\noindent
Le groupe de Picard d'un espace annelé est défini dans Faisceaux de modules, section \ref{modules-section-invertible}.

\begin{lemma}

\label{lemma-h1-invertible}

Soit $(X, \mathcal{O}_X)$ un espace annelé. Si tous les anneaux locaux
$\mathcal{O}_{X, x}$ sont locaux, alors il existe un isomorphisme canonique
$$
H^1(X, \mathcal{O}_X^*) = \Pic(X).
$$
de groupes abéliens.

\end{lemma}

\begin{proof}

Soit $\mathcal{L}$ un $\mathcal{O}_X$-module inversible. Considérons le préfaisceau $\mathcal{L}^*$ défini par la règle
$$
U \longmapsto \{s \in \mathcal{L}(U)
\text{ tel que } \mathcal{O}_U \xrightarrow{s \cdot -} \mathcal{L}_U
\text{ est un isomorphisme}\}
$$
Ce préfaisceau satisfait à la condition de faisceau. De plus, si
$f \in \mathcal{O}_X^*(U)$ et $s \in \mathcal{L}^*(U)$, alors il est clair que
$fs \in \mathcal{L}^*(U)$. De la même façon, si $s, s' \in \mathcal{L}^*(U)$
alors il existe un unique $f \in \mathcal{O}_X^*(U)$ tel que
$fs = s'$. De plus, le faisceau $\mathcal{L}^*$ possède localement des sections d'après Faisceaux de modules, lemme \ref{modules-lemma-invertible-is-locally-free-rank-1}. Ainsi, $\mathcal{L}^*$ est un $\mathcal{O}_X^*$-torseur. Nous obtenons donc une application
$$
\begin{matrix}
\text{faisceaux inversibles sur }(X, \mathcal{O}_X) \\
\text{ à isomorphisme près}
\end{matrix}
\longrightarrow
\begin{matrix}
\mathcal{O}_X^*\text{-torsors} \\
\text{ à isomorphisme près}
\end{matrix}
$$
Nous omettons de vérifier qu'il s'agit d'un homomorphisme de groupes abéliens. D'après le lemme \ref{lemma-torsors-h1},
le membre de droite est canoniquement en bijection avec $H^1(X, \mathcal{O}_X^*)$. Il reste donc à montrer que cette application est injective et surjective.

\medskip\noindent
Injectivité. Si le torseur $\mathcal{L}^*$ est trivial, le lemme \ref{lemma-trivial-torsor} montre que $\mathcal{L}^*$ possède une section globale. Cela signifie précisément que $\mathcal{L} \cong \mathcal{O}_X$, qui est l'élément neutre de $\Pic(X)$.

\medskip\noindent

Surjectivité. Soit $\mathcal{F}$ un $\mathcal{O}_X^*$-torseur. Considérons le préfaisceau d'ensembles
$$
\mathcal{L}_1 : U \longmapsto
(\mathcal{F}(U) \times \mathcal{O}_X(U))/\mathcal{O}_X^*(U)
$$
où l'action de $f \in \mathcal{O}_X^*(U)$ sur
$(s, g)$ est donnée par $(fs, f^{-1}g)$. Alors $\mathcal{L}_1$ est un préfaisceau de $\mathcal{O}_X$-modules si l'on pose
$(s, g) + (s', g') = (s, g + (s'/s)g')$, où $s'/s$ est la section locale $f$ de $\mathcal{O}_X^*$ telle que $fs = s'$,
$h(s, g) = (s, hg)$ pour toute section locale $h$ de $\mathcal{O}_X$. Nous omettons de vérifier que la faisceautisation
$\mathcal{L} = \mathcal{L}_1^\#$ est un $\mathcal{O}_X$-module inversible dont le $\mathcal{O}_X^*$-torseur associé $\mathcal{L}^*$ est isomorphe à $\mathcal{F}$.

\end{proof}













\section{Caractère local de la cohomologie}

\label{section-locality}

\noindent
Le lemme suivant montre qu'il n'y a aucune ambiguïté dans la définition de la cohomologie d'un faisceau $\mathcal{F}$ sur un ouvert.

\begin{lemma}

\label{lemma-cohomology-of-open}
Soit $X$ un espace annelé. Soit $U \subset X$ un sous-espace ouvert.
\begin{enumerate}
\item Si $\mathcal{I}$ est un $\mathcal{O}_X$-module injectif, alors $\mathcal{I}|_U$ est un $\mathcal{O}_U$-module injectif.
\item Pour tout faisceau de $\mathcal{O}_X$-modules $\mathcal{F}$, nous avons $H^p(U, \mathcal{F}) = H^p(U, \mathcal{F}|_U)$.
\end{enumerate}

\end{lemma}

\begin{proof}
Notons $j : U \to X$ l'immersion ouverte. Rappelons que le foncteur de restriction $j^{-1}$ à $U$ est adjoint à droite du foncteur $j_!$ d'extension par $0$, voir Faisceaux, lemme \ref{sheaves-lemma-j-shriek-modules}. De plus, $j_!$ est exact. Par conséquent, (1) découle d'Homologie, lemme \ref{homology-lemma-adjoint-preserve-injectives}.

\medskip\noindent
Par définition, $H^p(U, \mathcal{F}) = H^p(\Gamma(U, \mathcal{I}^\bullet))$, où $\mathcal{F} \to \mathcal{I}^\bullet$ est une résolution injective dans $\textit{Mod}(\mathcal{O}_X)$. D'après ce qui précède, $\mathcal{F}|_U \to \mathcal{I}^\bullet|_U$ est une résolution injective dans $\textit{Mod}(\mathcal{O}_U)$. Par conséquent, $H^p(U, \mathcal{F}|_U)$ est égal à $H^p(\Gamma(U, \mathcal{I}^\bullet|_U))$. Bien entendu, $\Gamma(U, \mathcal{F}) = \Gamma(U, \mathcal{F}|_U)$ pour tout faisceau $\mathcal{F}$ sur $X$. On obtient ainsi l'égalité de (2).
\end{proof}

\noindent

Soit $X$ un espace annelé. Soit $\mathcal{F}$ un faisceau de $\mathcal{O}_X$-modules. Soient $U \subset V \subset X$ des ouverts. Il existe alors une {\it application de restriction} canonique

\begin{equation}
\label{equation-restriction-mapping}
H^n(V, \mathcal{F})
\longrightarrow
H^n(U, \mathcal{F}), \quad
\xi \longmapsto \xi|_U
\end{equation}

fonctorielle en $\mathcal{F}$. En effet, choisissons une résolution injective $\mathcal{F} \to \mathcal{I}^\bullet$. Les applications de restriction des faisceaux $\mathcal{I}^p$ donnent un morphisme de complexes
$$
\Gamma(V, \mathcal{I}^\bullet)
\longrightarrow
\Gamma(U, \mathcal{I}^\bullet)
$$
Le membre de gauche est un complexe représentant $R\Gamma(V, \mathcal{F})$,
et celui de droite un complexe représentant $R\Gamma(U, \mathcal{F})$. En appliquant le foncteur $H^n$, nous obtenons l'application sur les groupes de cohomologie. Comme indiqué, nous la noterons $\xi \mapsto \xi|_U$. Ainsi, la règle $U \mapsto H^n(U, \mathcal{F})$ définit un préfaisceau de
$\mathcal{O}_X$-modules. Ce préfaisceau est habituellement noté
$\underline{H}^n(\mathcal{F})$. Nous en donnerons une autre interprétation dans le lemme \ref{lemma-include}.

\begin{lemma}

\label{lemma-kill-cohomology-class-on-covering}
Soit $X$ un espace annelé. Soit $\mathcal{F}$ un faisceau de $\mathcal{O}_X$-modules. Soit $U \subset X$ un sous-espace ouvert. Soient $n > 0$ et $\xi \in H^n(U, \mathcal{F})$. Il existe alors un recouvrement ouvert $U = \bigcup_{i\in I} U_i$ tel que $\xi|_{U_i} = 0$ pour tout $i \in I$.
\end{lemma}

\begin{proof}

Soit $\mathcal{F} \to \mathcal{I}^\bullet$ une résolution injective. Alors
$$
H^n(U, \mathcal{F}) =
\frac{\Ker(\mathcal{I}^n(U) \to \mathcal{I}^{n + 1}(U))}
{\Im(\mathcal{I}^{n - 1}(U) \to \mathcal{I}^n(U))}.
$$
Choisissons un élément $\tilde \xi \in \mathcal{I}^n(U)$ représentant la classe de cohomologie dans la présentation ci-dessus. Puisque $\mathcal{I}^\bullet$
est une résolution injective de $\mathcal{F}$ et que $n > 0$, le complexe $\mathcal{I}^\bullet$ est exact en degré $n$. Par conséquent,
$\Im(\mathcal{I}^{n - 1} \to \mathcal{I}^n) =
\Ker(\mathcal{I}^n \to \mathcal{I}^{n + 1})$ comme faisceaux. Puisque $\tilde \xi$ est une section du faisceau noyau sur $U$,
il existe un recouvrement ouvert $U = \bigcup_{i \in I} U_i$
tel que $\tilde \xi|_{U_i}$ soit l'image par $d$ d'une section
$\xi_i \in \mathcal{I}^{n - 1}(U_i)$. Par définition de la restriction, $\xi|_{U_i}$ correspond à la classe de
$\tilde \xi|_{U_i}$, ce qui conclut la preuve.

\end{proof}

\begin{lemma}

\label{lemma-describe-higher-direct-images}

Soit $f : X \to Y$ un morphisme d'espaces annelés. Soit $\mathcal{F}$ un $\mathcal{O}_X$-module. Les faisceaux $R^if_*\mathcal{F}$ sont les faisceaux associés aux préfaisceaux
$$
V \longmapsto H^i(f^{-1}(V), \mathcal{F})
$$
munis des applications de restriction de l'équation (\ref{equation-restriction-mapping}). Il existe un énoncé analogue pour $R^if_*$ appliqué à un complexe borné inférieurement $\mathcal{F}^\bullet$.

\end{lemma}

\begin{proof}

Soit $\mathcal{F} \to \mathcal{I}^\bullet$ une résolution injective. Par définition, $R^if_*\mathcal{F}$ est le $i$-ième faisceau de cohomologie du complexe
$$
f_*\mathcal{I}^0 \to f_*\mathcal{I}^1 \to f_*\mathcal{I}^2 \to \ldots
$$
Par définition de la structure de catégorie abélienne sur les $\mathcal{O}_Y$-modules, ce faisceau de cohomologie est le faisceau associé au préfaisceau
$$
V
\longmapsto
\frac{\Ker(f_*\mathcal{I}^i(V) \to f_*\mathcal{I}^{i + 1}(V))}
{\Im(f_*\mathcal{I}^{i - 1}(V) \to f_*\mathcal{I}^i(V))}
$$
et cela est évidemment égal à
$$
\frac{\Ker(\mathcal{I}^i(f^{-1}(V)) \to \mathcal{I}^{i + 1}(f^{-1}(V)))}
{\Im(\mathcal{I}^{i - 1}(f^{-1}(V)) \to \mathcal{I}^i(f^{-1}(V)))}
$$
qui est égal à $H^i(f^{-1}(V), \mathcal{F})$,
ce qui démontre le résultat.

\end{proof}

\begin{lemma}

\label{lemma-localize-higher-direct-images}
Soit $f : X \to Y$ un morphisme d'espaces annelés. Soit $\mathcal{F}$ un $\mathcal{O}_X$-module. Soit $V \subset Y$ un sous-espace ouvert. Notons $g : f^{-1}(V) \to V$ la restriction de $f$. Alors $$
R^pg_*(\mathcal{F}|_{f^{-1}(V)}) = (R^pf_*\mathcal{F})|_V
$$ Il existe un énoncé analogue pour l'image dérivée $Rf_*\mathcal{F}^\bullet$, où $\mathcal{F}^\bullet$ est un complexe borné inférieurement de $\mathcal{O}_X$-modules.
\end{lemma}

\begin{proof}
Première démonstration. Appliquer les lemmes \ref{lemma-describe-higher-direct-images} et \ref{lemma-cohomology-of-open} donne l'égalité affichée. Deuxième démonstration. Choisissons une résolution injective $\mathcal{F} \to \mathcal{I}^\bullet$ et utilisons le fait que $\mathcal{F}|_{f^{-1}(V)} \to \mathcal{I}^\bullet|_{f^{-1}(V)}$ est également une résolution injective.
\end{proof}

\begin{remark}

\label{remark-daniel}

Voici une autre approche des démonstrations des lemmes \ref{lemma-kill-cohomology-class-on-covering} et
\ref{lemma-describe-higher-direct-images} ci-dessus. Soit $(X, \mathcal{O}_X)$ un espace annelé. Soit $i_X : \textit{Mod}(\mathcal{O}_X) \to \textit{PMod}(\mathcal{O}_X)$
le foncteur d'inclusion, et soit $\#$ le foncteur de faisceautisation. Rappelons que $i_X$ est exact à gauche et que $\#$ est exact.

\begin{enumerate}

\item Démontrons d'abord le lemme \ref{lemma-include} ci-dessous, qui affirme que les foncteurs dérivés à droite de $i_X$ sont donnés par
$R^pi_X\mathcal{F} = \underline{H}^p(\mathcal{F})$. En voici une autre démonstration : l'égalité est claire pour $p = 0$. Les deux familles $(R^pi_X)_{p \geq 0}$ et $(\underline{H}^p)_{p \geq 0}$
sont des foncteurs delta qui s'annulent sur les objets injectifs ; ils sont donc universels, et par conséquent isomorphes. Voir Homologie, section \ref{homology-section-cohomological-delta-functor}.
\item Une reformulation du lemme \ref{lemma-kill-cohomology-class-on-covering}
est que $(\underline{H}^p(\mathcal{F}))^\# = 0$, $p > 0$, pour tout faisceau
de $\mathcal{O}_X$-modules $\mathcal{F}$. Pour le vérifier, utilisons l'exactitude de ${}^\#$ :
$$
(\underline{H}^p(\mathcal{F}))^\# =
(R^pi_X\mathcal{F})^\# =
R^p(\# \circ i_X)(\mathcal{F}) = 0
$$
car $\# \circ i_X$ est le foncteur identité.
\item Soit $f : X \to Y$ un morphisme d'espaces annelés. Soit $\mathcal{F}$ un $\mathcal{O}_X$-module. Le préfaisceau
$V \mapsto H^p(f^{-1}V, \mathcal{F})$ est égal à
$R^p (i_Y \circ f_*)\mathcal{F}$. On peut le démontrer en remarquant que les deux constructions donnent des foncteurs delta universels, comme dans l'argument de (1). Le lemme \ref{lemma-describe-higher-direct-images}
affirme que $R^p f_* \mathcal{F}= (R^p (i_Y \circ f_*)\mathcal{F})^\#$. En utilisant à nouveau l'exactitude de $\#$ et le fait que $\# \circ i_Y$ est le foncteur identité, on obtient
$$
R^p f_* \mathcal{F} =
R^p(\# \circ i_Y \circ f_*)\mathcal{F} =
(R^p (i_Y \circ f_*)\mathcal{F})^\#
$$
comme souhaité.

\end{enumerate}

\end{remark}





\section{Mayer--Vietoris}

\label{section-mayer-vietoris}

\noindent
Nous construirons plus loin la suite spectrale reliant la cohomologie de {\v C}ech à la cohomologie des faisceaux, voir le lemme \ref{lemma-cech-spectral-sequence}. La longue suite exacte de Mayer--Vietoris en est un cas particulier. Comme il s'agit d'une variante fondamentale, utile et facile à comprendre de cette suite spectrale, nous la traitons ici séparément.

\begin{lemma}

\label{lemma-injective-restriction-surjective}

\begin{slogan}
Les injectifs sont flasques.
\end{slogan}
Soit $X$ un espace annelé. Soient $U' \subset U \subset X$ des sous-espaces ouverts. Pour tout $\mathcal{O}_X$-module injectif $\mathcal{I}$, l'application de restriction $\mathcal{I}(U) \to \mathcal{I}(U')$ est surjective.
\end{lemma}

\begin{proof}

Soient $j : U \to X$ et $j' : U' \to X$ les immersions ouvertes. Rappelons que $j_!\mathcal{O}_U$ est l'extension par zéro de
$\mathcal{O}_U = \mathcal{O}_X|_U$, voir Faisceaux, section \ref{sheaves-section-open-immersions}. Puisque $j_!$ est adjoint à gauche de la restriction, pour tout faisceau $\mathcal{F}$ de $\mathcal{O}_X$-modules nous avons
$$
\Hom_{\mathcal{O}_X}(j_!\mathcal{O}_U, \mathcal{F})
=
\Hom_{\mathcal{O}_U}(\mathcal{O}_U, \mathcal{F}|_U)
=
\mathcal{F}(U)
$$
voir Faisceaux, lemme \ref{sheaves-lemma-j-shriek-modules}. De même, le faisceau $j'_!\mathcal{O}_{U'}$ représente le foncteur $\mathcal{F} \mapsto \mathcal{F}(U')$. En outre, il existe une application canonique évidente de $\mathcal{O}_X$-modules
$$
j'_!\mathcal{O}_{U'} \longrightarrow j_!\mathcal{O}_U
$$
qui correspond à l'application de restriction
$\mathcal{F}(U) \to \mathcal{F}(U')$ par le lemme de Yoneda (Catégories, lemme \ref{categories-lemma-yoneda}). D'après la description des fibres de
$j'_!\mathcal{O}_{U'}$, $j_!\mathcal{O}_U$
nous voyons que l'application affichée ci-dessus est injective (voir le lemme cité ci-dessus). Si $\mathcal{I}$ est un $\mathcal{O}_X$-module injectif, alors l'application
$$
\Hom_{\mathcal{O}_X}(j_!\mathcal{O}_U, \mathcal{I})
\longrightarrow
\Hom_{\mathcal{O}_X}(j'_!\mathcal{O}_{U'}, \mathcal{I})
$$
est surjective, voir Homologie, lemme \ref{homology-lemma-characterize-injectives}. En réunissant ces observations, nous obtenons le lemme.

\end{proof}

\begin{lemma}
[Mayer--Vietoris]
\label{lemma-mayer-vietoris}
Soit $X$ un espace annelé. Supposons que $X = U \cup V$ soit la réunion de deux ouverts. Pour tout $\mathcal{O}_X$-module $\mathcal{F}$, il existe une longue suite exacte de cohomologie $$
0 \to
H^0(X, \mathcal{F}) \to
H^0(U, \mathcal{F}) \oplus H^0(V, \mathcal{F}) \to
H^0(U \cap V, \mathcal{F}) \to
H^1(X, \mathcal{F}) \to \ldots
$$ Cette longue suite exacte est fonctorielle en $\mathcal{F}$.
\end{lemma}

\begin{proof}

La condition de faisceau affirme que le noyau de
$(1, -1) : \mathcal{F}(U) \oplus \mathcal{F}(V) \to \mathcal{F}(U \cap V)$
est égal à l'image de $\mathcal{F}(X)$ par la première application, pour tout faisceau abélien $\mathcal{F}$. Le lemme \ref{lemma-injective-restriction-surjective} ci-dessus implique que l'application
$(1, -1) : \mathcal{I}(U) \oplus \mathcal{I}(V) \to \mathcal{I}(U \cap V)$
est surjective dès que $\mathcal{I}$ est un $\mathcal{O}_X$-module injectif. Ainsi, si $\mathcal{F} \to \mathcal{I}^\bullet$ est une résolution injective de $\mathcal{F}$, nous obtenons une suite exacte courte de complexes
$$
0 \to
\mathcal{I}^\bullet(X) \to
\mathcal{I}^\bullet(U) \oplus \mathcal{I}^\bullet(V) \to
\mathcal{I}^\bullet(U \cap V) \to
0.
$$
Le passage à la cohomologie donne le résultat (utiliser Homologie, lemme \ref{homology-lemma-long-exact-sequence-cochain}). Nous omettons la vérification de la fonctorialité de la suite.

\end{proof}

\begin{lemma}
[Relative Mayer--Vietoris]
\label{lemma-relative-mayer-vietoris}
Soit $f : X \to Y$ un morphisme d'espaces annelés. Supposons que $X = U \cup V$ soit la réunion de deux ouverts. Notons $a = f|_U : U \to Y$, $b = f|_V : V \to Y$ et $c = f|_{U \cap V} : U \cap V \to Y$. Pour tout $\mathcal{O}_X$-module $\mathcal{F}$, il existe une longue suite exacte $$
0 \to
f_*\mathcal{F} \to
a_*(\mathcal{F}|_U) \oplus b_*(\mathcal{F}|_V) \to
c_*(\mathcal{F}|_{U \cap V}) \to
R^1f_*\mathcal{F} \to \ldots
$$ Cette longue suite exacte est fonctorielle en $\mathcal{F}$.
\end{lemma}

\begin{proof}

Soit $\mathcal{F} \to \mathcal{I}^\bullet$ une résolution injective de $\mathcal{F}$. Nous affirmons que nous obtenons une suite exacte courte de complexes
$$
0 \to
f_*\mathcal{I}^\bullet \to
a_*\mathcal{I}^\bullet|_U \oplus b_*\mathcal{I}^\bullet|_V \to
c_*\mathcal{I}^\bullet|_{U \cap V} \to
0.
$$
En effet, pour tout ouvert $W \subset Y$ et tout $n \geq 0$, la suite correspondante des groupes de sections sur $W$

$$
0 \to
\mathcal{I}^n(f^{-1}(W)) \to
\mathcal{I}^n(U \cap f^{-1}(W))
\oplus \mathcal{I}^n(V \cap f^{-1}(W)) \to
\mathcal{I}^n(U \cap V \cap f^{-1}(W)) \to
0
$$
est exacte courte d'après la démonstration du lemme \ref{lemma-mayer-vietoris}. Le lemme s'obtient en prenant les faisceaux de cohomologie et en utilisant le fait que
$\mathcal{I}^\bullet|_U$ est une résolution injective de $\mathcal{F}|_U$
et de même pour $\mathcal{I}^\bullet|_V$ et $\mathcal{I}^\bullet|_{U \cap V}$,
voir le lemme \ref{lemma-cohomology-of-open}.

\end{proof}



















\section{Le complexe de {\v C}
ech et la cohomologie de {\v C}ech}
\label{section-cech}

\noindent

Soit $X$ un espace topologique. Soit $\mathcal{U} : U = \bigcup_{i \in I} U_i$ un recouvrement ouvert, voir Topologie, notion de base (\ref{topology-item-covering}). Comme il est d'usage, nous notons
$U_{i_0\ldots i_p} = U_{i_0} \cap \ldots \cap U_{i_p}$ pour
l'intersection de $(p + 1)$ membres de $\mathcal{U}$. Soit $\mathcal{F}$ un préfaisceau abélien sur $X$. Posons
$$
\check{\mathcal{C}}^p(\mathcal{U}, \mathcal{F})
=
\prod\nolimits_{(i_0, \ldots, i_p) \in I^{p + 1}}
\mathcal{F}(U_{i_0\ldots i_p}).
$$
C'est un groupe abélien. Pour
$s \in \check{\mathcal{C}}^p(\mathcal{U}, \mathcal{F})$, nous notons
$s_{i_0\ldots i_p}$ sa valeur dans $\mathcal{F}(U_{i_0\ldots i_p})$. Remarquons que si $s \in \check{\mathcal{C}}^1(\mathcal{U}, \mathcal{F})$
et $i, j \in I$, alors $s_{ij}$ et $s_{ji}$ appartiennent tous deux à $\mathcal{F}(U_i \cap U_j)$, mais aucune relation n'est imposée entre $s_{ij}$ et $s_{ji}$. Autrement dit, nous ne travaillons {\it pas}
avec des cochaînes alternées (celles-ci seront définies dans la section \ref{section-alternating-cech}). Nous définissons
$$
d : \check{\mathcal{C}}^p(\mathcal{U}, \mathcal{F})
\longrightarrow
\check{\mathcal{C}}^{p + 1}(\mathcal{U}, \mathcal{F})
$$
par la formule

\begin{equation}
\label{equation-d-cech}
d(s)_{i_0\ldots i_{p + 1}}
=
\sum\nolimits_{j = 0}^{p + 1}
(-1)^j
s_{i_0\ldots \hat i_j \ldots i_{p + 1}}|_{U_{i_0\ldots i_{p + 1}}}
\end{equation}

On vérifie immédiatement que $d \circ d = 0$. Autrement dit,
$\check{\mathcal{C}}^\bullet(\mathcal{U}, \mathcal{F})$ est un complexe.

\begin{definition}

\label{definition-cech-complex}
Soit $X$ un espace topologique. Soit $\mathcal{U} : U = \bigcup_{i \in I} U_i$ un recouvrement ouvert. Soit $\mathcal{F}$ un préfaisceau abélien sur $X$. Le complexe $\check{\mathcal{C}}^\bullet(\mathcal{U}, \mathcal{F})$ est le {\it complexe de {\v C}ech} associé à $\mathcal{F}$ et au recouvrement ouvert $\mathcal{U}$. Ses groupes de cohomologie $H^i(\check{\mathcal{C}}^\bullet(\mathcal{U}, \mathcal{F}))$ sont appelés les {\it groupes de cohomologie de {\v C}ech} associés à $\mathcal{F}$ et au recouvrement $\mathcal{U}$. Ils sont notés $\check H^i(\mathcal{U}, \mathcal{F})$.
\end{definition}

\begin{lemma}

\label{lemma-cech-h0}

Soit $X$ un espace topologique. Soit $\mathcal{F}$ un préfaisceau abélien sur $X$. Les assertions suivantes sont équivalentes :

\begin{enumerate}

\item  $\mathcal{F}$ est un faisceau abélien, et
\item pour chaque recouvrement ouvert $\mathcal{U} : U = \bigcup_{i \in I} U_i$
l'application naturelle
$$
\mathcal{F}(U) \to \check{H}^0(\mathcal{U}, \mathcal{F})
$$
est bijective.

\end{enumerate}

\end{lemma}

\begin{proof}

Cela résulte du fait que la condition de faisceau affirme précisément que
$\mathcal{F}(U) \to \check{H}^0(\mathcal{U}, \mathcal{F})$
est bijectif pour chaque recouvrement ouvert.

\end{proof}

\begin{lemma}

\label{lemma-cech-trivial}
Soit $X$ un espace topologique. Soit $\mathcal{F}$ un préfaisceau abélien sur $X$. Soit $\mathcal{U} : U = \bigcup_{i \in I} U_i$ un recouvrement ouvert. Si $U_i = U$ pour un certain $i \in I$, alors le complexe de {\v C}ech augmenté $$
\mathcal{F}(U) \to \check{\mathcal{C}}^\bullet(\mathcal{U}, \mathcal{F})
$$ obtenu en plaçant $\mathcal{F}(U)$ en degré $-1$, avec pour différentielle l'application canonique de $\mathcal{F}(U)$ dans $\check{\mathcal{C}}^0(\mathcal{U}, \mathcal{F})$, est homotopiquement équivalent à $0$.
\end{lemma}

\begin{proof}

Fixons un élément $i \in I$ tel que $U = U_i$. Remarquons que
$U_{i_0 \ldots i_p} = U_{i_0 \ldots \hat i_j \ldots i_p}$ si $i_j = i$. Définissons une homotopie
$$
h :
\prod\nolimits_{i_0 \ldots i_{p + 1}} \mathcal{F}(U_{i_0 \ldots i_{p + 1}})
\longrightarrow
\prod\nolimits_{i_0 \ldots i_p} \mathcal{F}(U_{i_0 \ldots i_p})
$$
par la règle
$$
h(s)_{i_0 \ldots i_p} = s_{i i_0 \ldots i_p}
$$
En d'autres termes, $h : \prod_{i_0} \mathcal{F}(U_{i_0}) \to \mathcal{F}(U)$
est la projection sur le facteur $\mathcal{F}(U_i) = \mathcal{F}(U)$
et en général l'application $h$ égale la projection sur les facteurs
$\mathcal{F}(U_{i i_1 \ldots i_{p + 1}}) =
\mathcal{F}(U_{i_1 \ldots i_{p + 1}})$. Nous calculons

\begin{align*}
(dh + hd)(s)_{i_0 \ldots i_p}
& =
\sum\nolimits_{j = 0}^p
(-1)^j
h(s)_{i_0 \ldots \hat i_j \ldots i_p}
+
d(s)_{i i_0 \ldots i_p}\\
& =
\sum\nolimits_{j = 0}^p
(-1)^j
s_{i i_0 \ldots \hat i_j \ldots i_p}
+
s_{i_0 \ldots i_p}
+
\sum\nolimits_{j = 0}^p
(-1)^{j + 1}
s_{i i_0 \ldots \hat i_j \ldots i_p} \\
& =
s_{i_0 \ldots i_p}
\end{align*}

Ceci prouve que l'application identité est homotope à zéro, comme souhaité.

\end{proof}





\section{La cohomologie de {\v C}
ech comme foncteur sur les préfaisceaux}
\sectionmark{Cohomologie de {\v C}ech des préfaisceaux}
\label{section-cech-functor}

\noindent
Attention : dans cette section, nous travaillons presque exclusivement avec des préfaisceaux et des catégories de préfaisceaux ; les résultats sont entièrement faux dans le cadre des faisceaux et des catégories de faisceaux !

\medskip\noindent

Soit $X$ un espace annelé. Soit $\mathcal{U} : U = \bigcup_{i \in I} U_i$ un recouvrement ouvert. Soit $\mathcal{F}$ un préfaisceau de $\mathcal{O}_X$-modules. Nous disposons du complexe de {\v C}ech
$\check{\mathcal{C}}^\bullet(\mathcal{U}, \mathcal{F})$
de $\mathcal{F}$ en considérant simplement $\mathcal{F}$
comme un préfaisceau de groupes abéliens. Cependant, chaque terme
$\check{\mathcal{C}}^p(\mathcal{U}, \mathcal{F})$ est naturellement muni d'une structure de $\mathcal{O}_X(U)$-module, et la différentielle est donnée par des
applications $\mathcal{O}_X(U)$-linéaires. De plus, il est clair que la construction
$$
\mathcal{F} \longmapsto \check{\mathcal{C}}^\bullet(\mathcal{U}, \mathcal{F})
$$
est fonctorielle en $\mathcal{F}$. Il s'agit en fait d'un foncteur

\begin{equation}
\label{equation-cech-functor}
\check{\mathcal{C}}^\bullet(\mathcal{U}, -) :
\textit{PMod}(\mathcal{O}_X)
\longrightarrow
\text{Comp}^{+}(\text{Mod}_{\mathcal{O}_X(U)})
\end{equation}

Voir Catégories dérivées, définition \ref{derived-definition-complexes-notation},
pour la notation. Rappelons que la catégorie des complexes bornés inférieurement d'une catégorie abélienne est elle-même abélienne, voir Homologie, lemme \ref{homology-lemma-cat-cochain-abelian}.

\begin{lemma}

\label{lemma-cech-exact-presheaves}

Le foncteur donné par l'équation (\ref{equation-cech-functor}) est exact (voir Homologie, \ref{homology-lemma-exact-functor}).

\end{lemma}

\begin{proof}

Pour tout ouvert $W \subset U$, le foncteur
$\mathcal{F} \mapsto \mathcal{F}(W)$ est un foncteur additif exact de $\textit{PMod}(\mathcal{O}_X)$ vers $\text{Mod}_{\mathcal{O}_X(U)}$. Les termes
$\check{\mathcal{C}}^p(\mathcal{U}, \mathcal{F})$
du complexe sont des produits de ces foncteurs exacts, donc sont exacts. En outre, une suite de complexes est exacte si et seulement si la suite des termes de chaque degré est exacte.

\end{proof}

\begin{lemma}

\label{lemma-cech-cohomology-delta-functor-presheaves}

Soit $X$ un espace annelé. Soit $\mathcal{U} : U = \bigcup_{i \in I} U_i$ un recouvrement ouvert. Les foncteurs $\mathcal{F} \mapsto \check{H}^n(\mathcal{U}, \mathcal{F})$
forment un $\delta$-foncteur de la catégorie abélienne des préfaisceaux de $\mathcal{O}_X$-modules vers la catégorie des $\mathcal{O}_X(U)$-modules (voir Homologie, définition \ref{homology-definition-cohomological-delta-functor}).

\end{lemma}

\begin{proof}
Par le lemme \ref{lemma-cech-exact-presheaves}, une suite exacte courte de préfaisceaux de $\mathcal{O}_X$-modules $0 \to \mathcal{F}_1 \to \mathcal{F}_2 \to \mathcal{F}_3 \to 0$ donne une suite exacte courte de complexes de $\mathcal{O}_X(U)$-modules. Nous pouvons donc appliquer Homologie, lemme \ref{homology-lemma-long-exact-sequence-cochain}, pour obtenir les morphismes de bord $\delta_{\mathcal{F}_1 \to \mathcal{F}_2 \to \mathcal{F}_3} :
\check{H}^n(\mathcal{U}, \mathcal{F}_3) \to
\check{H}^{n + 1}(\mathcal{U}, \mathcal{F}_1)$ et la longue suite exacte correspondante. Nous omettons de vérifier que ces applications sont compatibles avec les morphismes entre suites exactes courtes de préfaisceaux.
\end{proof}

\noindent

Dans la formulation du lemme suivant, nous utilisons le foncteur $j_{p!}$ d'extension par $0$ pour les préfaisceaux de modules relativement à une immersion ouverte $j : U \to X$. Voir Faisceaux, section \ref{sheaves-section-open-immersions}. Pour tout ouvert
$W \subset X$ et tout préfaisceau $\mathcal{G}$ de $\mathcal{O}_X|_U$-modules, nous avons
$$
(j_{p!}\mathcal{G})(W) =
\left\{
\begin{matrix}
\mathcal{G}(W) & \text{si } W \subset U \\
0 & \text{sinon.}
\end{matrix}
\right.
$$
De plus, le foncteur $j_{p!}$ est adjoint à gauche du foncteur de restriction, voir Faisceaux, lemme \ref{sheaves-lemma-j-shriek-modules}. En particulier, nous avons la formule
$$
\Hom_{\mathcal{O}_X}(j_{p!}\mathcal{O}_U, \mathcal{F})
=
\Hom_{\mathcal{O}_U}(\mathcal{O}_U, \mathcal{F}|_U)
=
\mathcal{F}(U).
$$
Puisque le foncteur $\mathcal{F} \mapsto \mathcal{F}(U)$ est un foncteur exact sur la catégorie des préfaisceaux, nous concluons que le préfaisceau
$j_{p!}\mathcal{O}_U$ est un objet projectif dans la catégorie
$\textit{PMod}(\mathcal{O}_X)$, voir Homologie, lemme \ref{homology-lemma-characterize-projectives}.

\medskip\noindent
Notons que, si l'on se donne des ouverts $U \subset V \subset X$ et les immersions ouvertes associées $j_U, j_V$, alors on dispose d'une application canonique $(j_U)_{p!}\mathcal{O}_U \to (j_V)_{p!}\mathcal{O}_V$. Sur les sections au-dessus de tout ouvert $W \subset U$, c'est l'identité, et c'est $0$ sinon. Sous l'identification $\Hom_{\mathcal{O}_X}((j_U)_{p!}\mathcal{O}_U, (j_V)_{p!}\mathcal{O}_V) =
(j_V)_{p!}\mathcal{O}_V(U) = \mathcal{O}_V(U)$, elle correspond à l'élément $1 \in \mathcal{O}_V(U)$.

\begin{lemma}

\label{lemma-cech-map-into}
Soit $X$ un espace annelé. Soit $\mathcal{U} : U = \bigcup_{i \in I} U_i$ un recouvrement ouvert. Notons $j_{i_0\ldots i_p} : U_{i_0 \ldots i_p} \to X$ l'immersion ouverte. Considérons le complexe de chaînes $K(\mathcal{U})_\bullet$ de préfaisceaux de $\mathcal{O}_X$-modules $$
\ldots
\to
\bigoplus_{i_0i_1i_2} (j_{i_0i_1i_2})_{p!}\mathcal{O}_{U_{i_0i_1i_2}}
\to
\bigoplus_{i_0i_1} (j_{i_0i_1})_{p!}\mathcal{O}_{U_{i_0i_1}}
\to
\bigoplus_{i_0} (j_{i_0})_{p!}\mathcal{O}_{U_{i_0}}
\to 0 \to \ldots
$$ où le dernier terme non nul est placé en degré $0$ et où l'application $$
(j_{i_0\ldots i_{p + 1}})_{p!}\mathcal{O}_{U_{i_0\ldots i_{p + 1}}}
\longrightarrow
(j_{i_0\ldots \hat i_j \ldots i_{p + 1}})_{p!}
\mathcal{O}_{U_{i_0\ldots \hat i_j \ldots i_{p + 1}}}
$$ est donnée par $(-1)^j$ fois l'application canonique. Alors il existe un isomorphisme $$
\Hom_{\mathcal{O}_X}(K(\mathcal{U})_\bullet, \mathcal{F})
=
\check{\mathcal{C}}^\bullet(\mathcal{U}, \mathcal{F})
$$ fonctoriel en $\mathcal{F} \in \Ob(\textit{PMod}(\mathcal{O}_X))$.
\end{lemma}

\begin{proof}

Nous avons vu dans la discussion qui précède le lemme que
$$
\Hom_{\mathcal{O}_X}(
(j_{i_0\ldots i_p})_{p!}\mathcal{O}_{U_{i_0\ldots i_p}},
\mathcal{F})
=
\mathcal{F}(U_{i_0\ldots i_p}).
$$
Par conséquent, on voit bien que la somme directe
$$
\bigoplus\nolimits_{i_0 \ldots i_p}
(j_{i_0 \ldots i_p})_{p!}\mathcal{O}_{U_{i_0 \ldots i_p}}
$$
représente le foncteur
$$
\mathcal{F}
\longmapsto
\prod\nolimits_{i_0\ldots i_p} \mathcal{F}(U_{i_0\ldots i_p}).
$$
Par conséquent, d'après Catégories, lemme de Yoneda \ref{categories-lemma-yoneda},
nous voyons qu'il existe un complexe $K(\mathcal{U})_\bullet$ dont les termes sont ceux indiqués. Il est immédiat que les applications sont celles données dans le lemme.

\end{proof}

\begin{lemma}

\label{lemma-homology-complex}
Soit $X$ un espace annelé. Soit $\mathcal{U} : U = \bigcup_{i \in I} U_i$ un recouvrement ouvert. Soit $\mathcal{O}_\mathcal{U} \subset \mathcal{O}_X$ le préfaisceau image de l'application $\bigoplus j_{p!}\mathcal{O}_{U_i} \to \mathcal{O}_X$. Le complexe de chaînes $K(\mathcal{U})_\bullet$ de préfaisceaux du lemme \ref{lemma-cech-map-into} ci-dessus a pour préfaisceaux d'homologie $$
H_i(K(\mathcal{U})_\bullet) =
\left\{
\begin{matrix}
0 & \text{si} & i \not = 0 \\
\mathcal{O}_\mathcal{U} & \text{si} & i = 0
\end{matrix}
\right.
$$
\end{lemma}

\begin{proof}
Considérons le complexe étendu $K^{ext}_\bullet$ obtenu en plaçant $\mathcal{O}_\mathcal{U}$ en degré $-1$, muni de l'application évidente $K(\mathcal{U})_0 =
\bigoplus_{i_0} (j_{i_0})_{p!}\mathcal{O}_{U_{i_0}} \to
\mathcal{O}_\mathcal{U}$. Il suffit de montrer que le complexe obtenu en prenant les sections de ce complexe étendu sur tout ouvert $W \subset X$ est acyclique. En fait, nous affirmons que, pour chaque $W \subset X$, le complexe $K^{ext}_\bullet(W)$ est homotopiquement équivalent au complexe nul. Écrivons $I = I_1 \amalg I_2$, où $W \subset U_i$ si et seulement si $i \in I_1$.

\medskip\noindent
Si $I_1 = \emptyset$, alors le complexe $K^{ext}_\bullet(W) = 0$, donc il n'y a rien à prouver.

\medskip\noindent

Si $I_1 \not = \emptyset$, alors
$\mathcal{O}_\mathcal{U}(W) = \mathcal{O}_X(W)$
et
$$
K^{ext}_p(W) =
\bigoplus\nolimits_{i_0 \ldots i_p \in I_1} \mathcal{O}_X(W).
$$
Cela résulte de la description simple des préfaisceaux
$(j_{i_0 \ldots i_p})_{p!}\mathcal{O}_{U_{i_0 \ldots i_p}}$. En outre, le différentiel du complexe $K^{ext}_\bullet(W)$
est donné par
$$
d(s)_{i_0 \ldots i_p} =
\sum\nolimits_{j = 0, \ldots, p + 1} \sum\nolimits_{i \in I_1}
(-1)^j s_{i_0 \ldots i_{j - 1} i i_j \ldots i_p}.
$$
La somme est finie, car l'élément $s$ est à support fini. Fixons un élément $i_{\text{fix}} \in I_1$. Définissons une application
$$
h : K^{ext}_p(W) \longrightarrow K^{ext}_{p + 1}(W)
$$
par la règle
$$
h(s)_{i_0 \ldots i_{p + 1}} =
\left\{
\begin{matrix}
0 & \text{si} & i_0 \not = i_{\text{fix}} \\
s_{i_1 \ldots i_{p + 1}} & \text{si} & i_0 = i_{\text{fix}}
\end{matrix}
\right.
$$
Nous utiliserons la notation abrégée
$h(s)_{i_0 \ldots i_{p + 1}} = (i_0 = i_{\text{fix}}) s_{i_1 \ldots i_p}$
pour cela. Alors nous calculons

\begin{eqnarray*}
& & (dh + hd)(s)_{i_0 \ldots i_p} \\
& = &
\sum_j \sum_{i \in I_1} (-1)^j h(s)_{i_0 \ldots i_{j - 1} i i_j \ldots i_p}
+
(i = i_0) d(s)_{i_1 \ldots i_p} \\
& = &
s_{i_0 \ldots i_p} +
\sum_{j \geq 1}\sum_{i \in I_1}
(-1)^j (i_0 = i_{\text{fix}}) s_{i_1 \ldots i_{j - 1} i i_j \ldots i_p}
+
(i_0 = i_{\text{fix}}) d(s)_{i_1 \ldots i_p}
\end{eqnarray*}

qui est égal à $s_{i_0 \ldots i_p}$, comme voulu.

\end{proof}

\begin{lemma}

\label{lemma-cech-cohomology-derived-presheaves}
Soit $X$ un espace annelé. Soit $\mathcal{U} : U = \bigcup_{i \in I} U_i$ un recouvrement ouvert de $U \subset X$. Les foncteurs de cohomologie de {\v C}ech $\check{H}^p(\mathcal{U}, -)$ sont canoniquement isomorphes, en tant que $\delta$-foncteur, aux foncteurs dérivés à droite du foncteur $$
\check{H}^0(\mathcal{U}, -) :
\textit{PMod}(\mathcal{O}_X)
\longrightarrow
\text{Mod}_{\mathcal{O}_X(U)}.
$$ En outre, il existe un quasi-isomorphisme fonctoriel $$
\check{\mathcal{C}}^\bullet(\mathcal{U}, \mathcal{F})
\longrightarrow
R\check{H}^0(\mathcal{U}, \mathcal{F})
$$ où le membre de droite désigne le foncteur dérivé à droite $$
R\check{H}^0(\mathcal{U}, -) :
D^{+}(\textit{PMod}(\mathcal{O}_X))
\longrightarrow
D^{+}(\mathcal{O}_X(U))
$$ du foncteur exact à gauche $\check{H}^0(\mathcal{U}, -)$.
\end{lemma}


\begin{proof}
Notons que la catégorie des préfaisceaux de $\mathcal{O}_X$-modules a suffisamment d'injectifs, voir Injectifs, proposition \ref{injectives-proposition-presheaves-modules}. Notons aussi que $\check{H}^0(\mathcal{U}, -)$ est un foncteur exact à gauche de la catégorie des préfaisceaux de $\mathcal{O}_X$-modules vers celle des $\mathcal{O}_X(U)$-modules. Le foncteur dérivé et le foncteur dérivé à droite existent donc, voir Catégories dérivées, section \ref{derived-section-right-derived-functor}.

\medskip\noindent

Soit $\mathcal{I}$ un préfaisceau injectif de $\mathcal{O}_X$-modules. Dans ce cas, le foncteur $\Hom_{\mathcal{O}_X}(-, \mathcal{I})$
est exact sur $\textit{PMod}(\mathcal{O}_X)$. D'après le lemme \ref{lemma-cech-map-into}, nous avons
$$
\Hom_{\mathcal{O}_X}(K(\mathcal{U})_\bullet, \mathcal{I})
=
\check{\mathcal{C}}^\bullet(\mathcal{U}, \mathcal{I}).
$$
D'après le lemme \ref{lemma-homology-complex}, nous savons que $K(\mathcal{U})_\bullet$ est quasi-isomorphe à $\mathcal{O}_\mathcal{U}[0]$. Par l'exactitude du foncteur Hom à valeurs dans $\mathcal{I}$ mentionnée ci-dessus, nous obtenons donc $\check{H}^i(\mathcal{U}, \mathcal{I}) = 0$ pour tous
$i > 0$. Ainsi, le $\delta$-foncteur $(\check{H}^n, \delta)$
(voir lemme \ref{lemma-cech-cohomology-delta-functor-presheaves}) satisfait aux hypothèses de Homologie, lemme \ref{homology-lemma-efface-implies-universal}, et est donc un $\delta$-foncteur universel.

\medskip\noindent
D'après Catégories dérivées, lemme \ref{derived-lemma-higher-derived-functors}, la suite $R^i\check{H}^0(\mathcal{U}, -)$ forme également un $\delta$-foncteur universel. Par l'unicité des $\delta$-foncteurs universels, voir Homologie, lemme \ref{homology-lemma-uniqueness-universal-delta-functor}, nous concluons que $R^i\check{H}^0(\mathcal{U}, -) = \check{H}^i(\mathcal{U}, -)$. Cela suffit pour la plupart des applications, et l'on suggère au lecteur de sauter le reste de la démonstration.

\medskip\noindent

Soit $\mathcal{F}$ un préfaisceau quelconque de $\mathcal{O}_X$-modules. Choisissons une résolution injective $\mathcal{F} \to \mathcal{I}^\bullet$
dans la catégorie $\textit{PMod}(\mathcal{O}_X)$. Considérons le double complexe
$\check{\mathcal{C}}^\bullet(\mathcal{U}, \mathcal{I}^\bullet)$ avec des termes
$\check{\mathcal{C}}^p(\mathcal{U}, \mathcal{I}^q)$. Considérons le complexe total associé
$\text{Tot}(\check{\mathcal{C}}^\bullet(\mathcal{U}, \mathcal{I}^\bullet))$, voir Homologie, définition \ref{homology-definition-associated-simple-complex}. Il existe un morphisme de complexes
$$
\check{\mathcal{C}}^\bullet(\mathcal{U}, \mathcal{F})
\longrightarrow
\text{Tot}(\check{\mathcal{C}}^\bullet(\mathcal{U}, \mathcal{I}^\bullet))
$$
provenant des applications
$\check{\mathcal{C}}^p(\mathcal{U}, \mathcal{F})
\to \check{\mathcal{C}}^p(\mathcal{U}, \mathcal{I}^0)$
et il existe un morphisme de complexes
$$
\check{H}^0(\mathcal{U}, \mathcal{I}^\bullet)
\longrightarrow
\text{Tot}(\check{\mathcal{C}}^\bullet(\mathcal{U}, \mathcal{I}^\bullet))
$$
provenant des applications
$\check{H}^0(\mathcal{U}, \mathcal{I}^q) \to
\check{\mathcal{C}}^0(\mathcal{U}, \mathcal{I}^q)$. Ces deux applications sont des quasi-isomorphismes d'après Homologie, lemme \ref{homology-lemma-double-complex-gives-resolution}. En effet, les colonnes du double complexe sont exactes en degrés positifs parce que le complexe de {\v C}ech est exact en tant que foncteur (lemme \ref{lemma-cech-exact-presheaves}), et les lignes du double complexe sont exactes en degrés positifs puisque, comme nous venons de le voir, les groupes de cohomologie supérieure de {\v C}ech des préfaisceaux injectifs $\mathcal{I}^q$ sont nuls. Comme les quasi-isomorphismes deviennent inversibles dans $D^{+}(\mathcal{O}_X(U))$, on obtient le dernier morphisme affiché du lemme. Nous omettons de vérifier que ce morphisme est fonctoriel.

\end{proof}





\section{Cohomologie de {\v C}
ech et cohomologie}
\label{section-cech-cohomology-cohomology}

\begin{lemma}

\label{lemma-injective-trivial-cech}
Soit $X$ un espace annelé. Soit $\mathcal{U} : U = \bigcup_{i \in I} U_i$ un recouvrement ouvert. Soit $\mathcal{I}$ un $\mathcal{O}_X$-module injectif. Alors $$
\check{H}^p(\mathcal{U}, \mathcal{I}) =
\left\{
\begin{matrix}
\mathcal{I}(U) & \text{si} & p = 0 \\
0 & \text{si} & p > 0
\end{matrix}
\right.
$$
\end{lemma}

\begin{proof}

Un $\mathcal{O}_X$-module injectif est également injectif comme objet de la catégorie $\textit{PMod}(\mathcal{O}_X)$ (par exemple, puisque la faisceautisation est un adjoint à gauche exact du foncteur d'inclusion, en utilisant Homologie, lemme \ref{homology-lemma-adjoint-preserve-injectives}). Nous pouvons donc appliquer le lemme \ref{lemma-cech-cohomology-derived-presheaves}
(ou sa preuve) pour voir le résultat.

\end{proof}

\begin{lemma}

\label{lemma-cech-cohomology}
Soit $X$ un espace annelé. Soit $\mathcal{U} : U = \bigcup_{i \in I} U_i$ un recouvrement ouvert. Il existe une transformation $$
\check{\mathcal{C}}^\bullet(\mathcal{U}, -)
\longrightarrow
R\Gamma(U, -)
$$ des foncteurs $\textit{Mod}(\mathcal{O}_X) \to D^{+}(\mathcal{O}_X(U))$. En particulier, cela fournit des applications canoniques $\check{H}^p(\mathcal{U}, \mathcal{F}) \to H^p(U, \mathcal{F})$ lorsque $\mathcal{F}$ parcourt $\textit{Mod}(\mathcal{O}_X)$.
\end{lemma}

\begin{proof}

Soit $\mathcal{F}$ un $\mathcal{O}_X$-module. Choisissons une résolution injective
$\mathcal{F} \to \mathcal{I}^\bullet$. Considérons le double complexe
$\check{\mathcal{C}}^\bullet(\mathcal{U}, \mathcal{I}^\bullet)$ avec des termes
$\check{\mathcal{C}}^p(\mathcal{U}, \mathcal{I}^q)$. Il existe un morphisme de complexes
$$
\alpha :
\Gamma(U, \mathcal{I}^\bullet)
\longrightarrow
\text{Tot}(\check{\mathcal{C}}^\bullet(\mathcal{U}, \mathcal{I}^\bullet))
$$
provenant des applications
$\mathcal{I}^q(U) \to \check{H}^0(\mathcal{U}, \mathcal{I}^q)$
et un morphisme de complexes
$$
\beta :
\check{\mathcal{C}}^\bullet(\mathcal{U}, \mathcal{F})
\longrightarrow
\text{Tot}(\check{\mathcal{C}}^\bullet(\mathcal{U}, \mathcal{I}^\bullet))
$$
provenant de l'application $\mathcal{F} \to \mathcal{I}^0$. Le lemme \ref{homology-lemma-double-complex-gives-resolution} d'Homologie montre que
$\alpha$ est un quasi-isomorphisme. En effet, le lemme \ref{lemma-injective-trivial-cech} implique que la $q$-ième ligne du double complexe
$\check{\mathcal{C}}^\bullet(\mathcal{U}, \mathcal{I}^\bullet)$ est une résolution de $\Gamma(U, \mathcal{I}^q)$. Ainsi, $\alpha$ devient inversible dans $D^{+}(\mathcal{O}_X(U))$ et la transformation du lemme est la composition de $\beta$
suivie de l'inverse de $\alpha$. Nous omettons de vérifier que cette construction est fonctorielle.

\end{proof}

\begin{lemma}

\label{lemma-cech-h1}
Soit $X$ un espace topologique. Soit $\mathcal{H}$ un faisceau abélien sur $X$. Soit $\mathcal{U} : X = \bigcup_{i \in I} U_i$ un recouvrement ouvert. L'application $$
\check{H}^1(\mathcal{U}, \mathcal{H}) \longrightarrow H^1(X, \mathcal{H})
$$ est injective et identifie $\check{H}^1(\mathcal{U}, \mathcal{H})$, au moyen de la bijection du lemme \ref{lemma-torsors-h1}, à l'ensemble des classes d'isomorphisme de $\mathcal{H}$-torseurs dont la restriction à chaque $U_i$ est triviale.
\end{lemma}

\begin{proof}

Pour le voir, construisons une application inverse. Soit $\mathcal{F}$ un
$\mathcal{H}$-torseur dont la restriction à $U_i$ est triviale. D'après le lemme \ref{lemma-trivial-torsor}, cela signifie qu'il existe une section $s_i \in \mathcal{F}(U_i)$. Sur $U_{i_0} \cap U_{i_1}$
il existe une section unique $s_{i_0i_1}$ de $\mathcal{H}$ telle que
$s_{i_0i_1} \cdot s_{i_0}|_{U_{i_0} \cap U_{i_1}} =
s_{i_1}|_{U_{i_0} \cap U_{i_1}}$. Un calcul montre que $s_{i_0i_1}$ est un cocycle de {\v C}ech et que sa classe est bien définie (c.-à-d. ne dépend pas du choix des sections $s_i$). L'application inverse envoie la classe d'isomorphisme de $\mathcal{F}$ sur la classe de cohomologie du cocycle $(s_{i_0i_1})$. Nous omettons de vérifier que cette application est bien un inverse.

\end{proof}

\begin{lemma}

\label{lemma-include}
Soit $X$ un espace annelé. Considérons le foncteur $i : \textit{Mod}(\mathcal{O}_X) \to \textit{PMod}(\mathcal{O}_X)$. C'est un foncteur exact à gauche dont les foncteurs dérivés à droite sont donnés par $$
R^pi(\mathcal{F}) = \underline{H}^p(\mathcal{F}) :
U \longmapsto H^p(U, \mathcal{F})
$$ voir la discussion dans la section \ref{section-locality}.
\end{lemma}

\begin{proof}
Il est clair que $i$ est exact à gauche. Choisissons une résolution injective $\mathcal{F} \to \mathcal{I}^\bullet$. Par définition, $R^pi$ est le $p$-ième {\it préfaisceau} de cohomologie du complexe $\mathcal{I}^\bullet$. En d'autres termes, les sections de $R^pi(\mathcal{F})$ sur un ouvert $U$ sont données par $$
\frac{\Ker(\mathcal{I}^p(U) \to \mathcal{I}^{p + 1}(U))}
{\Im(\mathcal{I}^{p - 1}(U) \to \mathcal{I}^p(U))}.
$$ ce qui est la définition de $H^p(U, \mathcal{F})$.
\end{proof}

\begin{lemma}

\label{lemma-cech-spectral-sequence}
Soit $X$ un espace annelé. Soit $\mathcal{U} : U = \bigcup_{i \in I} U_i$ un recouvrement ouvert. Pour tout faisceau de $\mathcal{O}_X$-modules $\mathcal{F}$, il existe une suite spectrale $(E_r, d_r)_{r \geq 0}$ avec $$
E_2^{p, q} = \check{H}^p(\mathcal{U}, \underline{H}^q(\mathcal{F}))
$$ qui converge vers $H^{p + q}(U, \mathcal{F})$. Cette suite spectrale est fonctorielle en $\mathcal{F}$.
\end{lemma}

\begin{proof}
C'est une suite spectrale de Grothendieck (voir Catégories dérivées, lemme \ref{derived-lemma-grothendieck-spectral-sequence}) pour les foncteurs $$
i :  \textit{Mod}(\mathcal{O}_X) \to \textit{PMod}(\mathcal{O}_X)
\quad\text{et}\quad
\check{H}^0(\mathcal{U}, - ) : \textit{PMod}(\mathcal{O}_X)
\to \text{Mod}_{\mathcal{O}_X(U)}.
$$ Nous avons $\check{H}^0(\mathcal{U}, i(\mathcal{F})) = \mathcal{F}(U)$ d'après le lemme \ref{lemma-cech-h0}. Le préfaisceau $i(\mathcal{I})$ est acyclique pour la cohomologie de {\v C}ech d'après le lemme \ref{lemma-injective-trivial-cech}. Enfin, $\check{H}^p(\mathcal{U}, -) = R^p\check{H}^0(\mathcal{U}, -)$ comme foncteurs sur $\textit{PMod}(\mathcal{O}_X)$ d'après le lemme \ref{lemma-cech-cohomology-derived-presheaves}. L'assemblage de ces faits donne le lemme.
\end{proof}

\begin{lemma}

\label{lemma-cech-spectral-sequence-application}
Soit $X$ un espace annelé. Soit $\mathcal{U} : U = \bigcup_{i \in I} U_i$ un recouvrement ouvert. Soit $\mathcal{F}$ un $\mathcal{O}_X$-module. Supposons que $H^i(U_{i_0 \ldots i_p}, \mathcal{F}) = 0$ pour tout $i > 0$, tout $p \geq 0$ et tous $i_0, \ldots, i_p \in I$. Alors $\check{H}^p(\mathcal{U}, \mathcal{F}) = H^p(U, \mathcal{F})$ comme $\mathcal{O}_X(U)$-modules.
\end{lemma}

\begin{proof}

Nous utiliserons la suite spectrale du lemme \ref{lemma-cech-spectral-sequence}. Les hypothèses signifient que $E_2^{p, q} = 0$ pour tous $(p, q)$ avec
$q \not = 0$. La suite spectrale dégénère donc en $E_2$,
et le résultat s'ensuit.

\end{proof}

\begin{lemma}

\label{lemma-ses-cech-h1}
Soit $X$ un espace annelé. Soit $$
0 \to \mathcal{F} \to \mathcal{G} \to \mathcal{H} \to 0
$$ une suite exacte courte de $\mathcal{O}_X$-modules. Soit $U \subset X$ un ouvert. S'il existe un système cofinal de recouvrements ouverts $\mathcal{U}$ de $U$ tel que $\check{H}^1(\mathcal{U}, \mathcal{F}) = 0$, alors l'application $\mathcal{G}(U) \to \mathcal{H}(U)$ est surjective.
\end{lemma}

\begin{proof}

Prenons un élément $s \in \mathcal{H}(U)$. Choisissons un recouvrement ouvert
$\mathcal{U} : U = \bigcup_{i \in I} U_i$ tel que a) $\check{H}^1(\mathcal{U}, \mathcal{F}) = 0$ et b)
$s|_{U_i}$ soit l'image d'une section $s_i \in \mathcal{G}(U_i)$. Puisque nous pouvons certainement trouver $\mathcal{U}$ satisfaisant à (b), les hypothèses du lemme montrent que nous pouvons trouver
$\mathcal{U}$ satisfaisant à la fois à a) et à b). Considérons les sections
$$
s_{i_0i_1} = s_{i_1}|_{U_{i_0i_1}} - s_{i_0}|_{U_{i_0i_1}}.
$$
Puisque $s_i$ relève $s$, nous voyons que $s_{i_0i_1} \in \mathcal{F}(U_{i_0i_1})$. Par l'annulation de $\check{H}^1(\mathcal{U}, \mathcal{F})$, nous pouvons trouver des sections $t_i \in \mathcal{F}(U_i)$ telles que
$$
s_{i_0i_1} = t_{i_1}|_{U_{i_0i_1}} - t_{i_0}|_{U_{i_0i_1}}.
$$
Alors les sections $s_i - t_i$ satisfont clairement à la condition de recollement et se recollent en une section de $\mathcal{G}$ sur $U$ dont l'image est $s$. D'où le résultat.

\end{proof}

\begin{lemma}

\label{lemma-cech-vanish}

\begin{slogan}
Si la cohomologie supérieure de {\v C}ech d'un faisceau abélien s'annule pour tous les recouvrements ouverts, alors la cohomologie supérieure s'annule.
\end{slogan}
Soit $X$ un espace annelé. Soit $\mathcal{F}$ un $\mathcal{O}_X$-module tel que $$
\check{H}^p(\mathcal{U}, \mathcal{F}) = 0
$$ pour tout $p > 0$ et tout recouvrement ouvert $\mathcal{U} : U = \bigcup_{i \in I} U_i$ d'un ouvert de $X$. Alors $H^p(U, \mathcal{F}) = 0$ pour tout $p > 0$ et tout ouvert $U \subset X$.
\end{lemma}

\begin{proof}

Soit $\mathcal{F}$ un faisceau satisfaisant à l'hypothèse du lemme. Nous exprimerons cette propriété en disant que « la cohomologie supérieure de {\v C}ech de $\mathcal{F}$ s'annule pour tout recouvrement ouvert ». Choisissons un plongement $\mathcal{F} \to \mathcal{I}$ dans un $\mathcal{O}_X$-module injectif. D'après le lemme \ref{lemma-injective-trivial-cech}, la cohomologie supérieure de {\v C}ech de $\mathcal{I}$ s'annule pour tout recouvrement ouvert. Posons $\mathcal{Q} = \mathcal{I}/\mathcal{F}$
de sorte que nous avons une suite exacte courte
$$
0 \to \mathcal{F} \to \mathcal{I} \to \mathcal{Q} \to 0.
$$
D'après le lemme \ref{lemma-ses-cech-h1} et nos hypothèses, cette suite est même exacte comme suite de préfaisceaux ! En particulier, nous disposons d'une suite exacte longue de groupes de cohomologie de {\v C}ech pour tout recouvrement ouvert $\mathcal{U}$ ; voir, par exemple, le lemme \ref{lemma-cech-cohomology-delta-functor-presheaves}.
Il en résulte que $\mathcal{Q}$ est lui aussi un $\mathcal{O}_X$-module dont la cohomologie supérieure de {\v C}ech s'annule pour tout recouvrement ouvert.

\medskip\noindent

Ensuite, nous examinons la suite exacte longue de cohomologie
$$
\xymatrix{
0 \ar[r] &
H^0(U, \mathcal{F}) \ar[r] &
H^0(U, \mathcal{I}) \ar[r] &
H^0(U, \mathcal{Q}) \ar[lld] \\
&
H^1(U, \mathcal{F}) \ar[r] &
H^1(U, \mathcal{I}) \ar[r] &
H^1(U, \mathcal{Q}) \ar[lld] \\
&
\ldots & \ldots & \ldots \\
}
$$
pour tout ouvert $U \subset X$. Puisque $\mathcal{I}$ est injectif, nous avons $H^n(U, \mathcal{I}) = 0$ pour $n > 0$ (voir Catégories dérivées, lemme \ref{derived-lemma-higher-derived-functors}). D'après ce qui précède, l'application $H^0(U, \mathcal{I}) \to H^0(U, \mathcal{Q})$
est surjective, donc $H^1(U, \mathcal{F}) = 0$. Comme $\mathcal{F}$ était un $\mathcal{O}_X$-module arbitraire dont la cohomologie supérieure de {\v C}ech s'annule, nous en déduisons également que
$H^1(U, \mathcal{Q}) = 0$, puisque $\mathcal{Q}$ est un autre faisceau de ce type (voir ci-dessus). La suite exacte longue entraîne alors $H^2(U, \mathcal{F}) = 0$, et l'on poursuit de même.

\end{proof}

\begin{lemma}

\label{lemma-cech-vanish-basis}

(Variante du lemme \ref{lemma-cech-vanish}.) Soit $X$ un espace annelé. Soit $\mathcal{B}$ une base de la topologie de $X$. Soit $\mathcal{F}$ un $\mathcal{O}_X$-module. Supposons qu'il existe un ensemble de recouvrements ouverts $\text{Cov}$
ayant les propriétés suivantes :

\begin{enumerate}

\item Pour chaque $\mathcal{U} \in \text{Cov}$
avec $\mathcal{U} : U = \bigcup_{i \in I} U_i$, nous avons
$U, U_i \in \mathcal{B}$ et chaque $U_{i_0 \ldots i_p} \in \mathcal{B}$.
\item Pour chaque $U \in \mathcal{B}$, les recouvrements ouverts de $U$
qui appartiennent à $\text{Cov}$ forment un système cofinal de recouvrements ouverts de $U$.
\item Pour chaque $\mathcal{U} \in \text{Cov}$, nous avons
$\check{H}^p(\mathcal{U}, \mathcal{F}) = 0$ pour tout $p > 0$.

\end{enumerate}

Alors $H^p(U, \mathcal{F}) = 0$ pour tout $p > 0$ et tout $U \in \mathcal{B}$.

\end{lemma}

\begin{proof}

Soient $\mathcal{F}$ et $\text{Cov}$ comme dans le lemme. Nous exprimerons cette propriété en disant que « la cohomologie supérieure de {\v C}ech de $\mathcal{F}$ s'annule pour tout $\mathcal{U} \in \text{Cov}$ ». Choisissons un plongement $\mathcal{F} \to \mathcal{I}$ dans un $\mathcal{O}_X$-module injectif. D'après le lemme \ref{lemma-injective-trivial-cech}, la cohomologie supérieure de {\v C}ech de $\mathcal{I}$
s'annule pour tout $\mathcal{U} \in \text{Cov}$. Posons $\mathcal{Q} = \mathcal{I}/\mathcal{F}$
de sorte que nous avons une suite exacte courte
$$
0 \to \mathcal{F} \to \mathcal{I} \to \mathcal{Q} \to 0.
$$
D'après le lemme \ref{lemma-ses-cech-h1} et notre hypothèse (2), cette suite donne lieu à une suite exacte
$$
0 \to \mathcal{F}(U) \to \mathcal{I}(U) \to \mathcal{Q}(U) \to 0.
$$
pour chaque $U \in \mathcal{B}$. Par conséquent, pour tout $\mathcal{U} \in \text{Cov}$,
nous obtenons une suite exacte courte de complexes de {\v C}ech
$$
0 \to
\check{\mathcal{C}}^\bullet(\mathcal{U}, \mathcal{F}) \to
\check{\mathcal{C}}^\bullet(\mathcal{U}, \mathcal{I}) \to
\check{\mathcal{C}}^\bullet(\mathcal{U}, \mathcal{Q}) \to 0
$$
car chaque terme du complexe de {\v C}ech est constitué d'un produit de valeurs sur des éléments de $\mathcal{B}$, d'après l'hypothèse (1). En particulier, nous avons une suite exacte longue de groupes de cohomologie de {\v C}ech pour tout recouvrement ouvert $\mathcal{U} \in \text{Cov}$. Cela implique que $\mathcal{Q}$ est également un $\mathcal{O}_X$-module dont la cohomologie supérieure de {\v C}ech s'annule pour tous
$\mathcal{U} \in \text{Cov}$.

\medskip\noindent

Ensuite, nous examinons la suite exacte longue de cohomologie
$$
\xymatrix{
0 \ar[r] &
H^0(U, \mathcal{F}) \ar[r] &
H^0(U, \mathcal{I}) \ar[r] &
H^0(U, \mathcal{Q}) \ar[lld] \\
&
H^1(U, \mathcal{F}) \ar[r] &
H^1(U, \mathcal{I}) \ar[r] &
H^1(U, \mathcal{Q}) \ar[lld] \\
&
\ldots & \ldots & \ldots \\
}
$$
pour tout $U \in \mathcal{B}$. Puisque $\mathcal{I}$ est injectif, nous avons $H^n(U, \mathcal{I}) = 0$ pour $n > 0$ (voir Catégories dérivées, lemme \ref{derived-lemma-higher-derived-functors}). D'après ce qui précède, l'application $H^0(U, \mathcal{I}) \to H^0(U, \mathcal{Q})$
est surjective, donc $H^1(U, \mathcal{F}) = 0$. Comme $\mathcal{F}$ était un $\mathcal{O}_X$-module arbitraire dont la cohomologie supérieure de {\v C}ech s'annule pour tout $\mathcal{U} \in \text{Cov}$,
nous en déduisons que $H^1(U, \mathcal{Q}) = 0$, puisque $\mathcal{Q}$ est un autre faisceau de ce type (voir ci-dessus). La suite exacte longue entraîne alors $H^2(U, \mathcal{F}) = 0$, et l'on poursuit de même.

\end{proof}

\begin{lemma}

\label{lemma-pushforward-injective}

Soit $f : X \to Y$ un morphisme d'espaces annelés. Soit $\mathcal{I}$ un $\mathcal{O}_X$-module injectif. Alors :

\begin{enumerate}

\item  $\check{H}^p(\mathcal{V}, f_*\mathcal{I}) = 0$
pour tout $p > 0$ et tout recouvrement ouvert
$\mathcal{V} : V = \bigcup_{j \in J} V_j$ de $Y$.
\item  $H^p(V, f_*\mathcal{I}) = 0$ pour tout $p > 0$ et tout ouvert $V \subset Y$.

\end{enumerate}

En d'autres termes, $f_*\mathcal{I}$ est acyclique à droite pour $\Gamma(V, -)$
(voir Catégories dérivées, définition \ref{derived-definition-derived-functor}) pour tout ouvert $V \subset Y$.

\end{lemma}

\begin{proof}
Posons $\mathcal{U} : f^{-1}(V) = \bigcup_{j \in J} f^{-1}(V_j)$. Il s'agit d'un recouvrement par des ouverts de $X$ et $$
\check{\mathcal{C}}^\bullet(\mathcal{V}, f_*\mathcal{I}) =
\check{\mathcal{C}}^\bullet(\mathcal{U}, \mathcal{I}).
$$ Ceci est vrai parce que $$
f_*\mathcal{I}(V_{j_0 \ldots j_p})
= \mathcal{I}(f^{-1}(V_{j_0 \ldots j_p})) =
\mathcal{I}(f^{-1}(V_{j_0}) \cap \ldots \cap f^{-1}(V_{j_p}))
= \mathcal{I}(U_{j_0 \ldots j_p}).
$$ Ainsi, la première affirmation du lemme résulte du lemme \ref{lemma-injective-trivial-cech}. La seconde résulte de la première et du lemme \ref{lemma-cech-vanish}.
\end{proof}

\noindent
Le lemme suivant implique en particulier que $f_* : \textit{Ab}(X) \to \textit{Ab}(Y)$ transforme les faisceaux abéliens injectifs en faisceaux abéliens injectifs.

\begin{lemma}

\label{lemma-pushforward-injective-flat}
Soit $f : X \to Y$ un morphisme d'espaces annelés. Supposons que $f$ soit plat. Alors $f_*\mathcal{I}$ est un $\mathcal{O}_Y$-module injectif pour tout $\mathcal{O}_X$-module injectif $\mathcal{I}$.
\end{lemma}

\begin{proof}

Dans ce cas, le foncteur $f^*$ transforme les monomorphismes en monomorphismes (Faisceaux de modules, lemme \ref{modules-lemma-pullback-flat}). Le résultat découle alors d'Homologie, lemme \ref{homology-lemma-adjoint-preserve-injectives}.

\end{proof}

\begin{lemma}

\label{lemma-cohomology-products}
Soit $(X, \mathcal{O}_X)$ un espace annelé et soit $I$ un ensemble. Pour $i \in I$, soit $\mathcal{F}_i$ un $\mathcal{O}_X$-module. Soit $U \subset X$ un ouvert. L'application canonique $$
H^p(U, \prod\nolimits_{i \in I} \mathcal{F}_i)
\longrightarrow
\prod\nolimits_{i \in I} H^p(U, \mathcal{F}_i)
$$ est un isomorphisme pour $p = 0$ et injective pour $p = 1$.
\end{lemma}

\begin{proof}

L'assertion pour $p = 0$ est vraie parce que le produit des faisceaux est égal au produit des préfaisceaux sous-jacents ; voir Faisceaux, section \ref{sheaves-section-limits-sheaves}. Prouvons l'assertion pour $p = 1$. Posons $\mathcal{F} = \prod \mathcal{F}_i$. Soit $\xi \in H^1(U, \mathcal{F})$ dont l'image est nulle dans
$\prod H^1(U, \mathcal{F}_i)$. Par le caractère local de la cohomologie, voir lemme \ref{lemma-kill-cohomology-class-on-covering}, il existe un recouvrement ouvert $\mathcal{U} : U = \bigcup U_j$ tel que
$\xi|_{U_j} = 0$ pour tout $j$. D'après le lemme \ref{lemma-cech-h1}, cela signifie que
$\xi$ provient d'un élément
$\check \xi \in \check H^1(\mathcal{U}, \mathcal{F})$. Puisque les applications
$\check H^1(\mathcal{U}, \mathcal{F}_i) \to H^1(U, \mathcal{F}_i)$
sont injectives pour tout $i$ (par le lemme \ref{lemma-cech-h1}) et, puisque l'image de $\xi$ est nulle dans $\prod H^1(U, \mathcal{F}_i)$, nous voyons que l'image
$\check \xi_i = 0$ dans $\check H^1(\mathcal{U}, \mathcal{F}_i)$. Cependant, puisque $\mathcal{F} = \prod \mathcal{F}_i$, nous voyons que
$\check{\mathcal{C}}^\bullet(\mathcal{U}, \mathcal{F})$ est le produit des complexes
$\check{\mathcal{C}}^\bullet(\mathcal{U}, \mathcal{F}_i)$, donc par Homologie, lemme \ref{homology-lemma-product-abelian-groups-exact}
nous concluons que $\check \xi = 0$, comme voulu.

\end{proof}







\section{Faisceaux flasques}

\label{section-flasque}

\noindent
Voici la définition.

\begin{definition}

\label{definition-flasque}
Soit $X$ un espace topologique. On dit qu'un préfaisceau d'ensembles $\mathcal{F}$ est {\it flasque}, ou {\it mou}, si, pour tous ouverts $U \subset V$ de $X$, l'application de restriction $\mathcal{F}(V) \to \mathcal{F}(U)$ est surjective.
\end{definition}

\noindent
Nous utiliserons également cette terminologie pour les faisceaux abéliens et les faisceaux de modules lorsque $X$ est un espace annelé. Il suffit clairement de supposer que les applications de restriction $\mathcal{F}(X) \to \mathcal{F}(U)$ sont surjectives pour chaque ouvert $U \subset X$.

\begin{lemma}

\label{lemma-injective-flasque}
Soit $(X, \mathcal{O}_X)$ un espace annelé. Alors tout $\mathcal{O}_X$-module injectif est flasque.
\end{lemma}

\begin{proof}
C'est une reformulation du lemme \ref{lemma-injective-restriction-surjective}.
\end{proof}

\begin{lemma}

\label{lemma-flasque-acyclic}
Soit $(X, \mathcal{O}_X)$ un espace annelé. Tout $\mathcal{O}_X$-module flasque est acyclique pour $R\Gamma(X, -)$ ainsi que pour $R\Gamma(U, -)$, pour tout ouvert $U$ de $X$.
\end{lemma}

\begin{proof}

Nous utiliserons Catégories dérivées, lemme \ref{derived-lemma-subcategory-right-acyclics}. Comme tout module injectif est flasque, chaque $\mathcal{O}_X$-module se plonge dans un module flasque ; voir Injectifs, lemme \ref{injectives-lemma-abelian-sheaves-space}. Il suffit donc de montrer que, pour toute suite exacte courte
$$
0 \to \mathcal{F} \to \mathcal{G} \to \mathcal{H} \to 0
$$
avec $\mathcal{F}$ et $\mathcal{G}$ flasques, $\mathcal{H}$
est flasque et la suite reste exacte après prise de sections sur tout ouvert de $X$. En fait, la seconde assertion implique la première. Soit $U \subset X$ un sous-espace ouvert. Soit $s \in \mathcal{H}(U)$. Nous montrerons que nous pouvons relever $s$ en une section de $\mathcal{G}$
sur $U$. Pour ce faire, considérons l'ensemble $T$ des couples $(V, t)$
où $V \subset U$ est ouvert et $t \in \mathcal{G}(V)$ est une section envoyée sur $s|_V$ dans $\mathcal{H}$. Munissons $T$ de l'ordre partiel défini par
$(V, t) \leq (V', t')$ si et seulement si $V \subset V'$ et $t'|_V = t$. Si $(V_\alpha, t_\alpha)$, $\alpha \in A$,
est un sous-ensemble totalement ordonné de $T$, alors $V = \bigcup V_\alpha$
est ouvert et il existe une section unique $t \in \mathcal{G}(V)$
se restreignant à $t_\alpha$ sur $V_\alpha$, grâce à la condition de recollement pour
$\mathcal{G}$. Ainsi, le lemme de Zorn fournit un élément maximal
$(V, t)$ de $T$. Nous allons montrer que $V = U$, ce qui achèvera la démonstration. En effet, choisissons $x \in U$. Nous pouvons trouver un petit voisinage ouvert
$W \subset U$ de $x$ et $t' \in \mathcal{G}(W)$ dont l'image est $s|_W$
dans $\mathcal{H}$. Alors $t'|_{W \cap V} - t|_{W \cap V}$ a pour image zéro dans $\mathcal{H}$ et provient donc d'une section
$r' \in \mathcal{F}(W \cap V)$. Comme $\mathcal{F}$ est flasque, nous trouvons une section $r \in \mathcal{F}(W)$ dont la restriction est $r'$
sur $W \cap V$. En modifiant $t'$ par l'image de $r$, nous pouvons supposer que $t$ et $t'$ se restreignent à la même section sur
$W \cap V$. Grâce à la condition de recollement pour $\mathcal{G}$,
nous pouvons trouver une section $\tilde t$ de $\mathcal{G}$ sur
$W \cup V$ qui se restreint à $t$ et à $t'$. Par maximalité de $(V, t)$, nous voyons que $V \cup W = V$. Ainsi, $x \in V$ et la démonstration est achevée.

\end{proof}

\noindent
Le lemme suivant ne vaut pas pour les préfaisceaux flasques.

\begin{lemma}

\label{lemma-flasque-acyclic-cech}
Soit $(X, \mathcal{O}_X)$ un espace annelé. Soit $\mathcal{F}$ un faisceau de $\mathcal{O}_X$-modules. Soit $\mathcal{U} : U = \bigcup U_i$ un recouvrement ouvert. Si $\mathcal{F}$ est flasque, alors $\check{H}^p(\mathcal{U}, \mathcal{F}) = 0$ pour $p > 0$.
\end{lemma}

\begin{proof}

Les préfaisceaux $\underline{H}^q(\mathcal{F})$ utilisés dans l'énoncé du lemme \ref{lemma-cech-spectral-sequence} sont nuls d'après le lemme \ref{lemma-flasque-acyclic}. Par conséquent, $\check{H}^p(U, \mathcal{F}) = H^p(U, \mathcal{F}) = 0$
d'après le lemme \ref{lemma-flasque-acyclic}, encore une fois.

\end{proof}

\begin{lemma}

\label{lemma-flasque-acyclic-pushforward}
Soit $f : (X, \mathcal{O}_X) \to (Y, \mathcal{O}_Y)$ un morphisme d'espaces annelés. Soit $\mathcal{F}$ un faisceau de $\mathcal{O}_X$-modules. Si $\mathcal{F}$ est flasque, alors $R^pf_*\mathcal{F} = 0$ pour $p > 0$.
\end{lemma}

\begin{proof}
Cela résulte immédiatement des lemmes \ref{lemma-describe-higher-direct-images} et \ref{lemma-flasque-acyclic}.
\end{proof}

\noindent
Le lemme suivant peut être démontré par un argument élémentaire d'induction pour les recouvrements finis ; comparer avec la discussion de la cohomologie de {\v C}ech dans \cite{FOAG}.

\begin{lemma}

\label{lemma-vanishing-ravi}
Soit $X$ un espace topologique. Soit $\mathcal{F}$ un faisceau abélien sur $X$. Soit $\mathcal{U} : U = \bigcup_{i \in I} U_i$ un recouvrement ouvert. Supposons que les applications de restriction $\mathcal{F}(U) \to \mathcal{F}(U')$ soient surjectives pour toute union arbitraire $U'$ d'ouverts de la forme $U_{i_0 \ldots i_p}$. Alors $\check{H}^p(\mathcal{U}, \mathcal{F})$ s'annule pour $p > 0$.
\end{lemma}

\begin{proof}

Soit $Y$ l'ensemble de sous-ensembles non vides de $I$. Nous utiliserons les lettres
$A, B, C, \ldots$ pour désigner les éléments de $Y$, c'est-à-dire les sous-ensembles non vides de $I$. Pour un sous-ensemble fini non vide $J \subset I$, posons
$$
V_J = \{A \in Y \mid J \subset A\}
$$
Cela signifie que $V_{\{i\}} = \{A \in Y \mid i \in A\}$ et
$V_J = \bigcap_{j \in J} V_{\{j\}}$. Alors $V_J \subset V_K$ si et seulement si $J \supset K$. Il existe une topologie unique sur $Y$ telle que la collection des sous-ensembles $V_J$ forme une base de la topologie de $Y$. Tout ouvert est de la forme
$$
V = \bigcup\nolimits_{t \in T} V_{J_t}
$$
pour une famille de sous-ensembles finis $J_t$. Si $J_t \subset J_{t'}$,
alors nous pouvons retirer $J_{t'}$ de la famille sans modifier $V$. Ainsi, nous pouvons supposer qu'il n'y a pas d'inclusions parmi les $J_t$. Dans ce cas, les éléments minimaux de $V$ sont les ensembles $A = J_t$. Par conséquent, nous pouvons retrouver la famille $(J_t)_{t \in T}$ à partir de l'ouvert $V$.

\medskip\noindent

Nous pouvons décrire complètement les recouvrements ouverts de $Y$. Premièrement, puisque les éléments $A \in Y$ sont des sous-ensembles non vides de $I$, nous avons
$$
Y = \bigcup\nolimits_{i \in I} V_{\{i\}}
$$
Pour comprendre les autres recouvrements, soit $V$ comme ci-dessus et soit $V_s \subset Y$
un ouvert correspondant à la famille $(J_{s, t})_{t \in T_s}$. Alors
$$
V = \bigcup\nolimits_{s \in S} V_s
$$
si et seulement si, pour chaque $t \in T$, il existe un $s \in S$ et
$t_s \in T_s$ tel que $J_t = J_{s, t_s}$. Comme la famille
$(J_t)_{t \in T}$ est minimale, l'élément minimal $A = J_t$
doit appartenir à $V_s$ pour un certain $s$, donc $A \in V_{J_{t_s}}$ pour un certain
$t_s \in T_s$. Mais, comme $A$ est également minimal dans $V_s$, nous concluons que $J_{t_s} = J_t$.

\medskip\noindent

Ensuite, nous associons à chaque ouvert de $Y$ un ouvert de $X$. Plus précisément, nous envoyons
$Y$ sur $U$ et adoptons la règle
$$
V_J \mapsto U_J = \bigcap\nolimits_{i \in J} U_i
$$
sur les ouverts $V_J$, et nous l'étendons à des ouverts arbitraires $V$ par la règle
$$
V = \bigcup\nolimits_{t \in T} V_{J_t}
\mapsto
\bigcup\nolimits_{t \in T} U_{J_t}
$$
La classification des recouvrements ouverts de $Y$ ci-dessus montre que cette règle transforme les recouvrements ouverts en recouvrements ouverts. Nous obtenons ainsi un faisceau abélien $\mathcal{G}$ sur $Y$ en posant
$\mathcal{G}(Y) = \mathcal{F}(U)$ et pour
$V = \bigcup\nolimits_{t \in T} V_{J_t}$, en posant
$$
\mathcal{G}(V) = \mathcal{F}\left(\bigcup\nolimits_{t \in T} U_{J_t}\right)
$$
et en utilisant les applications de restriction de $\mathcal{F}$.

\medskip\noindent
Ces préliminaires étant établis, nous pouvons démontrer le lemme comme suit. Nous avons un recouvrement ouvert $\mathcal{V} : Y = \bigcup_{i \in I} V_{\{i\}}$ de $Y$. Par construction, nous avons une égalité $$
\check{C}^\bullet(\mathcal{V}, \mathcal{G}) =
\check{C}^\bullet(\mathcal{U}, \mathcal{F})
$$ de complexes de {\v C}ech. Puisque le faisceau $\mathcal{G}$ est flasque sur $Y$ (d'après notre hypothèse sur $\mathcal{F}$ dans l'énoncé du lemme), l'annulation résulte du lemme \ref{lemma-flasque-acyclic-cech}.
\end{proof}



\section{La suite spectrale de Leray}

\label{section-Leray}

\begin{lemma}

\label{lemma-before-Leray}

Soit $f : X \to Y$ un morphisme d'espaces annelés. Il existe un diagramme commutatif
$$
\xymatrix{
D^{+}(X) \ar[rr]_-{R\Gamma(X, -)} \ar[d]_{Rf_*} & &
D^{+}(\mathcal{O}_X(X)) \ar[d]^{\text{restriction}} \\
D^{+}(Y) \ar[rr]^-{R\Gamma(Y, -)} & &
D^{+}(\mathcal{O}_Y(Y))
}
$$
Plus généralement, pour tout ouvert $V \subset Y$ et $U = f^{-1}(V)$, il existe un diagramme commutatif
$$
\xymatrix{
D^{+}(X) \ar[rr]_-{R\Gamma(U, -)} \ar[d]_{Rf_*} & &
D^{+}(\mathcal{O}_X(U)) \ar[d]^{\text{restriction}} \\
D^{+}(Y) \ar[rr]^-{R\Gamma(V, -)} & &
D^{+}(\mathcal{O}_Y(V))
}
$$
Voir aussi la remarque \ref{remark-elucidate-lemma} pour plus d'explications.

\end{lemma}

\begin{proof}

Soit $\Gamma_{res} : \textit{Mod}(\mathcal{O}_X) \to \text{Mod}_{\mathcal{O}_Y(Y)}$
le foncteur qui associe à un $\mathcal{O}_X$-module $\mathcal{F}$
les sections globales de $\mathcal{F}$, considérées comme un $\mathcal{O}_Y(Y)$-module via l'application $f^\sharp : \mathcal{O}_Y(Y) \to \mathcal{O}_X(X)$. Soit $restriction : \text{Mod}_{\mathcal{O}_X(X)} \to \text{Mod}_{\mathcal{O}_Y(Y)}$
le foncteur de restriction des scalaires induit par
$f^\sharp : \mathcal{O}_Y(Y) \to \mathcal{O}_X(X)$. Notons que $restriction$
est exact, si bien que son foncteur dérivé à droite se calcule en appliquant simplement le foncteur de restriction, voir Catégories dérivées, lemme \ref{derived-lemma-right-derived-exact-functor}. Il est clair que
$$
\Gamma_{res}
=
restriction \circ \Gamma(X, -)
=
\Gamma(Y, -) \circ f_*
$$
Nous affirmons que le lemme \ref{derived-lemma-compose-derived-functors} de Catégories dérivées s'applique aux deux compositions.
Pour la première, cela résulte clairement des remarques ci-dessus. Pour la seconde, cela découle du lemme \ref{lemma-pushforward-injective}, qui implique que les $\mathcal{O}_X$-modules injectifs sont envoyés sur des faisceaux acycliques pour $\Gamma(Y, -)$ dans la catégorie des faisceaux sur $Y$.

\end{proof}

\begin{remark}

\label{remark-elucidate-lemma}
Voici une explication concrète du lemme \ref{lemma-before-Leray}. Étant donnés $f : X \to Y$, $\mathcal{F} \in \textit{Mod}(\mathcal{O}_X)$ et une résolution injective $\mathcal{F} \to \mathcal{I}^\bullet$, on a $$
\begin{matrix}
R\Gamma(X, \mathcal{F}) & \text{est représenté par} &
\Gamma(X, \mathcal{I}^\bullet) \\
Rf_*\mathcal{F} & \text{est représenté par} & f_*\mathcal{I}^\bullet \\
R\Gamma(Y, Rf_*\mathcal{F}) & \text{est représenté par} &
\Gamma(Y, f_*\mathcal{I}^\bullet)
\end{matrix}
$$ où la dernière assertion résulte du lemme d'acyclicité de Leray (Catégories dérivées, lemme \ref{derived-lemma-leray-acyclicity}) et du lemme \ref{lemma-pushforward-injective}. Enfin, combinons cela avec l'observation immédiate que $$
\Gamma(X, \mathcal{I}^\bullet)
=
\Gamma(Y, f_*\mathcal{I}^\bullet).
$$ ce qui donne la commutativité du diagramme du lemme.
\end{remark}

\begin{lemma}

\label{lemma-modules-abelian}

Soit $X$ un espace annelé. Soit $\mathcal{F}$ un $\mathcal{O}_X$-module.

\begin{enumerate}

\item Les groupes de cohomologie $H^i(U, \mathcal{F})$, pour tout ouvert $U \subset X$, calculés en considérant $\mathcal{F}$ comme un $\mathcal{O}_X$-module ou comme un faisceau abélien sont identiques.
\item Soit $f : X \to Y$ un morphisme d'espaces annelés. Les images directes supérieures $R^if_*\mathcal{F}$ de $\mathcal{F}$
calculées en considérant celui-ci comme un $\mathcal{O}_X$-module ou comme un faisceau abélien sont identiques.

\end{enumerate}

Des énoncés analogues valent pour les complexes bornés inférieurement de $\mathcal{O}_X$-modules.

\end{lemma}

\begin{proof}
Considérons le morphisme d'espaces annelés $(X, \mathcal{O}_X) \to (X, \underline{\mathbf{Z}}_X)$ donné par l'identité sur l'espace topologique sous-jacent et par l'unique morphisme de faisceaux d'anneaux $\underline{\mathbf{Z}}_X \to \mathcal{O}_X$. Soit $\mathcal{F}$ un $\mathcal{O}_X$-module. Notons $\mathcal{F}_{ab}$ le même faisceau, vu comme un $\underline{\mathbf{Z}}_X$-module, c'est-à-dire comme un faisceau de groupes abéliens. Soit $\mathcal{F} \to \mathcal{I}^\bullet$ une résolution injective. D'après la remarque \ref{remark-elucidate-lemma}, $\Gamma(X, \mathcal{I}^\bullet)$ calcule à la fois $R\Gamma(X, \mathcal{F})$ et $R\Gamma(X, \mathcal{F}_{ab})$. Ceci prouve (1).

\medskip\noindent

Pour prouver (2), nous utilisons (1) et le lemme \ref{lemma-describe-higher-direct-images}. Le résultat suit immédiatement.

\end{proof}

\begin{lemma}
[suite spectrale de Leray]
\label{lemma-Leray}
Soit $f : X \to Y$ un morphisme d'espaces annelés. Soit $\mathcal{F}^\bullet$ un complexe borné inférieurement de $\mathcal{O}_X$-modules. Il existe une suite spectrale $$
E_2^{p, q} = H^p(Y, R^qf_*(\mathcal{F}^\bullet))
$$ convergeant vers $H^{p + q}(X, \mathcal{F}^\bullet)$.
\end{lemma}

\begin{proof}
Il s'agit simplement de la suite spectrale de Grothendieck des Catégories dérivées, lemme \ref{derived-lemma-grothendieck-spectral-sequence}, issue de la composition des foncteurs $\Gamma_{res} = \Gamma(Y, -) \circ f_*$, où $\Gamma_{res}$ est celui de la démonstration du lemme \ref{lemma-before-Leray}. Pour vérifier les hypothèses des Catégories dérivées, lemme \ref{derived-lemma-grothendieck-spectral-sequence}, voir la démonstration du lemme \ref{lemma-before-Leray} ou la remarque \ref{remark-elucidate-lemma}.
\end{proof}

\begin{remark}

\label{remark-Leray-ss-more-structure}

La suite spectrale de Leray, telle que nous l'avons construite dans le lemme \ref{lemma-Leray},
est une suite spectrale de $\Gamma(Y, \mathcal{O}_Y)$-modules. Il est toutefois assez facile de voir qu'il s'agit en fait d'une suite spectrale de
$\Gamma(X, \mathcal{O}_X)$-modules. Par exemple $f$ donne lieu à un morphisme d'espaces annelés
$f' :  (X, \mathcal{O}_X) \to (Y, f_*\mathcal{O}_X)$. D'après le lemme \ref{lemma-modules-abelian}, les termes $E_r^{p, q}$ de la suite spectrale de Leray associés au $\mathcal{O}_X$-module $\mathcal{F}$
et $f$ sont identiques à ceux associés à $\mathcal{F}$ et $f'$
au moins pour $r \geq 2$. En effet, ils coïncident tous deux avec les termes de la suite spectrale de Leray de $\mathcal{F}$ considéré comme faisceau abélien. Puisque $(f_*\mathcal{O}_X)(Y) = \mathcal{O}_X(X)$, le résultat en découle. La suite spectrale de Leray possède souvent une structure supplémentaire.

\end{remark}

\begin{lemma}

\label{lemma-apply-Leray}
Soit $f : X \to Y$ un morphisme d'espaces annelés. Soit $\mathcal{F}$ un $\mathcal{O}_X$-module.
\begin{enumerate}
\item Si $R^qf_*\mathcal{F} = 0$ pour $q > 0$, alors $H^p(X, \mathcal{F}) = H^p(Y, f_*\mathcal{F})$ pour tous $p$.
\item Si $H^p(Y, R^qf_*\mathcal{F}) = 0$ pour tous $q$ et $p > 0$, alors $H^q(X, \mathcal{F}) = H^0(Y, R^qf_*\mathcal{F})$ pour tous $q$.
\end{enumerate}

\end{lemma}

\begin{proof}

Ce sont deux conditions simples qui forcent la suite spectrale de Leray à dégénérer en $E_2$. On peut aussi démontrer ces faits directement, sans utiliser la suite spectrale ; c'est un bon exercice de cohomologie des faisceaux.

\end{proof}

\begin{lemma}

\label{lemma-higher-direct-images-compose}

\begin{slogan}
Le foncteur dérivé total d'une composition est la composition des foncteurs dérivés totaux.
\end{slogan}
Soient $f : X \to Y$ et $g : Y \to Z$ des morphismes d'espaces annelés. Alors $Rg_* \circ Rf_* = R(g \circ f)_*$ comme foncteurs de $D^{+}(X) \to D^{+}(Z)$.
\end{lemma}

\begin{proof}
Nous allons appliquer Catégories dérivées, lemme \ref{derived-lemma-compose-derived-functors}. Il est clair que $g_* \circ f_* = (g \circ f)_*$ ; voir Faisceaux, lemme \ref{sheaves-lemma-pushforward-composition}. Il reste à montrer que $f_*\mathcal{I}$ est $g_*$-acyclique. Ceci résulte du lemme \ref{lemma-pushforward-injective} et de la description des images directes supérieures $R^ig_*$ dans le lemme \ref{lemma-describe-higher-direct-images}.
\end{proof}

\begin{lemma}
[Suite spectrale de Leray relative]
\label{lemma-relative-Leray}
Soient $f : X \to Y$ et $g : Y \to Z$ des morphismes d'espaces annelés. Soit $\mathcal{F}$ un $\mathcal{O}_X$-module. Il existe une suite spectrale telle que $$
E_2^{p, q} = R^pg_*(R^qf_*\mathcal{F})
$$ convergeant vers $R^{p + q}(g \circ f)_*\mathcal{F}$. Cette suite spectrale est fonctorielle en $\mathcal{F}$, et il existe une version pour les complexes de $\mathcal{O}_X$-modules.
\end{lemma}

\begin{proof}
Il s'agit de la suite spectrale de Grothendieck associée à la composition des foncteurs ; elle résulte du lemme \ref{lemma-higher-direct-images-compose} et des Catégories dérivées, lemme \ref{derived-lemma-grothendieck-spectral-sequence}.
\end{proof}













\section{Fonctorialité de la cohomologie}

\label{section-functoriality}

\begin{lemma}

\label{lemma-functoriality}

Soit $f : X \to Y$ un morphisme d'espaces annelés. Soient $\mathcal{G}^\bullet$, respectivement $\mathcal{F}^\bullet$, des complexes bornés inférieurement de $\mathcal{O}_Y$-modules, respectivement de $\mathcal{O}_X$-modules.
Soit $\varphi : \mathcal{G}^\bullet \to f_*\mathcal{F}^\bullet$ un
morphisme de complexes. Il existe un morphisme canonique
$$
\mathcal{G}^\bullet
\longrightarrow
Rf_*(\mathcal{F}^\bullet)
$$
dans $D^{+}(Y)$. De plus, cette construction est fonctorielle par rapport au triple
$(\mathcal{G}^\bullet, \mathcal{F}^\bullet, \varphi)$.

\end{lemma}

\begin{proof}
Choisissons une résolution injective $\mathcal{F}^\bullet \to \mathcal{I}^\bullet$. Par définition, $Rf_*(\mathcal{F}^\bullet)$ est représenté par $f_*\mathcal{I}^\bullet$ dans $K^{+}(\mathcal{O}_Y)$. La composée $$
\mathcal{G}^\bullet \to f_*\mathcal{F}^\bullet \to f_*\mathcal{I}^\bullet
$$ est un morphisme dans $K^{+}(Y)$ qui se transforme en morphisme du lemme en appliquant le foncteur de localisation $j_Y : K^{+}(Y) \to D^{+}(Y)$.
\end{proof}

\noindent

Soit $f : X \to Y$ un morphisme d'espaces annelés. Soit $\mathcal{G}$ un $\mathcal{O}_Y$-module et soit
$\mathcal{F}$ un $\mathcal{O}_X$-module. Rappelons qu'un
$f$-morphisme $\varphi$ de $\mathcal{G}$ vers $\mathcal{F}$ est un morphisme
$\varphi : \mathcal{G} \to f_*\mathcal{F}$, ou, ce qui revient au même, un morphisme $\varphi : f^*\mathcal{G} \to \mathcal{F}$. Voir Faisceaux, définition \ref{sheaves-definition-f-map}. Un tel $f$-morphisme donne lieu à un morphisme de complexes

\begin{equation}
\label{equation-functorial-derived}
\varphi :
R\Gamma(Y, \mathcal{G})
\longrightarrow
R\Gamma(X, \mathcal{F})
\end{equation}

dans $D^{+}(\mathcal{O}_Y(Y))$. Plus précisément, nous utilisons le morphisme
$\mathcal{G} \to Rf_*\mathcal{F}$ dans $D^{+}(Y)$ du lemme \ref{lemma-functoriality}, puis nous appliquons $R\Gamma(Y, -)$. D'après le lemme \ref{lemma-before-Leray}, nous avons
$R\Gamma(X, \mathcal{F}) = R\Gamma(Y, Rf_*\mathcal{F})$
et nous obtenons la flèche affichée. Nous la décrivons complètement dans la remarque \ref{remark-explain-arrow} ci-dessous. Elle induit en particulier des morphismes en cohomologie

\begin{equation}
\label{equation-functorial}
\varphi : H^i(Y, \mathcal{G}) \longrightarrow H^i(X, \mathcal{F}).
\end{equation}

\begin{remark}

\label{remark-explain-arrow}

Soit $f : X \to Y$ un morphisme d'espaces annelés. Soit $\mathcal{G}$ un $\mathcal{O}_Y$-module. Soit $\mathcal{F}$ un $\mathcal{O}_X$-module. Soit $\varphi$ un $f$-morphisme de $\mathcal{G}$ vers $\mathcal{F}$. Choisissons une résolution $\mathcal{F} \to \mathcal{I}^\bullet$
par un complexe de $\mathcal{O}_X$-modules injectifs. Choisissons des résolutions $\mathcal{G} \to \mathcal{J}^\bullet$ et
$f_*\mathcal{I}^\bullet \to (\mathcal{J}')^\bullet$ par des complexes de $\mathcal{O}_Y$-modules injectifs. D'après Catégories dérivées, lemme \ref{derived-lemma-morphisms-lift},
il existe un morphisme de complexes
$\beta$ tel que le diagramme

\begin{equation}
\label{equation-choice}
\xymatrix{
\mathcal{G} \ar[d] \ar[r] &
f_*\mathcal{F} \ar[r] &
f_*\mathcal{I}^\bullet \ar[d] \\
\mathcal{J}^\bullet \ar[rr]^\beta & &
(\mathcal{J}')^\bullet
}
\end{equation}

En appliquant les foncteurs de sections globales, nous obtenons un diagramme
$$
\xymatrix{
 & & \Gamma(Y, f_*\mathcal{I}^\bullet) \ar[d]_{qis} \ar@{=}[r] &
\Gamma(X, \mathcal{I}^\bullet) \\
\Gamma(Y, \mathcal{J}^\bullet) \ar[rr]^\beta & &
\Gamma(Y, (\mathcal{J}')^\bullet) &
}
$$
Le complexe en bas à gauche représente $R\Gamma(Y, \mathcal{G})$
et le complexe en haut à droite représente $R\Gamma(X, \mathcal{F})$. La flèche verticale est un quasi-isomorphisme d'après le lemme \ref{lemma-before-Leray} ; elle devient inversible après application du foncteur de localisation
$K^{+}(\mathcal{O}_Y(Y)) \to D^{+}(\mathcal{O}_Y(Y))$. La flèche (\ref{equation-functorial-derived}) est la composée de la flèche horizontale avec l'inverse de la flèche verticale.

\end{remark}





\section{Raffinements et cohomologie de {\v C}ech}

\label{section-refinements-cech}

\noindent

Soit $(X, \mathcal{O}_X)$ un espace annelé.
$\mathcal{U} : X = \bigcup_{i \in I} U_i$ et
$\mathcal{V} : X = \bigcup_{j \in J} V_j$ soient des recouvrements ouverts. Supposons que $\mathcal{U}$ soit un raffinement de $\mathcal{V}$. Choisissons une application $c : I \to J$ telle que $U_i \subset V_{c(i)}$
pour tous $i \in I$. Cela induit un morphisme de complexes de {\v C}ech
$$
\gamma :
\check{\mathcal{C}}^\bullet(\mathcal{V}, \mathcal{F})
\longrightarrow
\check{\mathcal{C}}^\bullet(\mathcal{U}, \mathcal{F}),
\quad
(\xi_{j_0 \ldots j_p})
\longmapsto
(\xi_{c(i_0) \ldots c(i_p)}|_{U_{i_0 \ldots i_p}})
$$
fonctoriel par rapport au faisceau de $\mathcal{O}_X$-modules $\mathcal{F}$. Supposons que $c' : I \to J$ soit une seconde application telle que
$U_i \subset V_{c'(i)}$ pour tous $i \in I$. Alors les morphismes correspondants
$\gamma$ et $\gamma'$ sont homotopes.
$\gamma - \gamma' = \text{d} \circ h + h \circ \text{d}$
avec
$h : \check{\mathcal{C}}^{p + 1}(\mathcal{V}, \mathcal{F}) \to
\check{\mathcal{C}}^p(\mathcal{U}, \mathcal{F})$
donné par la formule
$$
h(\alpha)_{i_0 \ldots i_p} =
\sum\nolimits_{a = 0}^{p}
(-1)^a
\alpha_{c(i_0)\ldots c(i_a) c'(i_a) \ldots c'(i_p)}
$$
Nous omettons le calcul qui le vérifie ; voir la discussion qui suit (\ref{equation-transformation}) pour une démonstration dans un cadre plus général. En particulier, le morphisme induit sur les groupes de cohomologie de {\v C}ech est indépendant du choix de $c$. En outre, si
$\mathcal{W} : X = \bigcup_{k \in K} W_k$ est un troisième recouvrement ouvert et si $\mathcal{V}$ est un raffinement de $\mathcal{W}$, alors la composée des morphismes
$$
\check{\mathcal{C}}^\bullet(\mathcal{W}, \mathcal{F})
\longrightarrow
\check{\mathcal{C}}^\bullet(\mathcal{V}, \mathcal{F})
\longrightarrow
\check{\mathcal{C}}^\bullet(\mathcal{U}, \mathcal{F})
$$
associés aux applications $I \to J$ et $J \to K$ est le morphisme associé à la composée $I \to K$. Nous pouvons donc définir les groupes de cohomologie de {\v C}ech
$$
\check{H}^p(X, \mathcal{F}) =
\colim_\mathcal{U} \check{H}^p(\mathcal{U}, \mathcal{F})
$$
où la colimite porte sur tous les recouvrements ouverts de $X$, ordonnés par raffinement.

\medskip\noindent

Les morphismes $\gamma$ définis ci-dessus sont compatibles avec le morphisme vers la cohomologie ; autrement dit, la composée
$$
\check{H}^p(\mathcal{V}, \mathcal{F}) \to
\check{H}^p(\mathcal{U}, \mathcal{F})
\xrightarrow{\text{Lemme \ref{lemma-cech-cohomology}}}
H^p(X, \mathcal{F})
$$
est le morphisme canonique du premier groupe vers la cohomologie du lemme \ref{lemma-cech-cohomology}. Nous le démontrerons ci-dessous dans un cadre légèrement plus général. Il en résulte un morphisme bien défini

\begin{equation}
\label{equation-cech-to-cohomology}
\check{H}^p(X, \mathcal{F}) =
\colim_\mathcal{U} \check{H}^p(\mathcal{U}, \mathcal{F})
\longrightarrow
H^p(X, \mathcal{F})
\end{equation}

de la cohomologie de {\v C}ech vers la cohomologie.

\begin{lemma}

\label{lemma-functoriality-cech}

Soit $f : X \to Y$ un morphisme d'espaces annelés. Soit $\varphi : f^*\mathcal{G} \to \mathcal{F}$ un $f$-morphisme d'un $\mathcal{O}_Y$-module $\mathcal{G}$ vers un
$\mathcal{O}_X$-module $\mathcal{F}$. Soient $\mathcal{U} : X = \bigcup_{i \in I} U_i$ et
$\mathcal{V} : Y = \bigcup_{j \in J} V_j$ des recouvrements ouverts. Supposons que $\mathcal{U}$ soit un raffinement de
$f^{-1}\mathcal{V} : X = \bigcup_{j \in J} f^{-1}(V_j)$. Dans ce cas, il existe un diagramme commutatif
$$
\xymatrix{
\check{\mathcal{C}}^\bullet(\mathcal{U}, \mathcal{F}) \ar[r] &
R\Gamma(X, \mathcal{F}) \\
\check{\mathcal{C}}^\bullet(\mathcal{V}, \mathcal{G}) \ar[r]
\ar[u]^\gamma &
R\Gamma(Y, \mathcal{G}) \ar[u]
}
$$
dans $D^{+}(\mathcal{O}_X(X))$, où les flèches horizontales sont celles du lemme \ref{lemma-cech-cohomology} et la flèche verticale de droite est donnée par (\ref{equation-functorial-derived}). En particulier, nous obtenons des diagrammes commutatifs de groupes de cohomologie
$$
\xymatrix{
\check{H}^p(\mathcal{U}, \mathcal{F}) \ar[r] &
H^p(X, \mathcal{F}) \\
\check{H}^p(\mathcal{V}, \mathcal{G}) \ar[r]
\ar[u]^\gamma &
H^p(Y, \mathcal{G}) \ar[u]
}
$$
où la flèche verticale de droite est celle de (\ref{equation-functorial}).

\end{lemma}

\begin{proof}

Définissons d'abord la flèche verticale de gauche. Choisissons une application
$c : I \to J$ telle que $U_i \subset f^{-1}(V_{c(i)})$ pour tous
$i \in I$. En degré $p$ nous définissons l'application par la règle
$$
\gamma(s)_{i_0 \ldots i_p} = \varphi(s)_{c(i_0) \ldots c(i_p)}
$$
Ceci est bien défini, car $\varphi$ induit des morphismes
$\mathcal{G}(V_{c(i_0) \ldots c(i_p)}) \to \mathcal{F}(U_{i_0 \ldots i_p})$
par hypothèse. Cela définit clairement un morphisme de complexes. Choisissons des résolutions injectives
$\mathcal{F} \to \mathcal{I}^\bullet$ sur $X$ et
$\mathcal{G} \to J^\bullet$ sur $Y$. Comme dans la démonstration du lemme \ref{lemma-cech-cohomology}, introduisons les bicomplexes $A^{\bullet, \bullet}$ et $B^{\bullet, \bullet}$
avec des termes
$$
B^{p, q} = \check{\mathcal{C}}^p(\mathcal{V}, \mathcal{J}^q)
\quad
\text{et}
\quad
A^{p, q} = \check{\mathcal{C}}^p(\mathcal{U}, \mathcal{I}^q).
$$
Comme dans la remarque \ref{remark-explain-arrow} ci-dessus, choisissons également une résolution injective
$f_*\mathcal{I} \to (\mathcal{J}')^\bullet$ sur $Y$ et un morphisme de complexes $\beta : \mathcal{J} \to (\mathcal{J}')^\bullet$
rendant (\ref{equation-choice}) commutatif. Introduisons deux autres bicomplexes, $(B')^{\bullet, \bullet}$ et
$(B'')^{\bullet, \bullet}$ avec
$$
(B')^{p, q} = \check{\mathcal{C}}^p(\mathcal{V}, (\mathcal{J}')^q)
\quad
\text{et}
\quad
(B'')^{p, q} = \check{\mathcal{C}}^p(\mathcal{V}, f_*\mathcal{I}^q).
$$
Il existe un $f$-morphisme de complexes de
$f_*\mathcal{I}^\bullet$ vers $\mathcal{I}^\bullet$. La même formule que ci-dessus définit donc un morphisme de bicomplexes
$$
\gamma : (B'')^{\bullet, \bullet} \longrightarrow A^{\bullet, \bullet}.
$$
Considérons le diagramme de complexes
$$
\xymatrix{
\check{\mathcal{C}}^\bullet(\mathcal{U}, \mathcal{F})
\ar[r] &
\text{Tot}(A^{\bullet, \bullet}) & & &
\Gamma(X, \mathcal{I}^\bullet) \ar[lll]^{qis}
\ar@{=}[ddl]\\
\check{\mathcal{C}}^\bullet(\mathcal{V}, \mathcal{G})
\ar[r] \ar[u]^\gamma &
\text{Tot}(B^{\bullet, \bullet}) \ar[r]^\beta &
\text{Tot}((B')^{\bullet, \bullet}) &
\text{Tot}((B'')^{\bullet, \bullet}) \ar[l] \ar[llu]_{s\gamma} \\
& \Gamma(Y, \mathcal{J}^\bullet) \ar[u]^{qis} \ar[r]^\beta &
\Gamma(Y, (\mathcal{J}')^\bullet) \ar[u] &
\Gamma(Y, f_*\mathcal{I}^\bullet) \ar[u] \ar[l]_{qis}
}
$$
Les deux flèches horizontales de cibles $\text{Tot}(A^{\bullet, \bullet})$ et
$\text{Tot}(B^{\bullet, \bullet})$
sont celles expliquées dans le lemme \ref{lemma-cech-cohomology}. La partie supérieure gauche, un pentagone, est commutative parce que (\ref{equation-choice}) l'est. Les deux carrés inférieurs sont manifestement commutatifs. Les définitions montrent aussi immédiatement que la partie supérieure droite, un carré, est commutative. Le lemme résulte alors des définitions et du fait que les deux chemins extérieurs allant de $\check{\mathcal{C}}^\bullet(\mathcal{V}, \mathcal{G})$
à $\Gamma(X, \mathcal{I}^\bullet)$, par le haut ou par le bas, coïncident après inversion des quasi-isomorphismes rencontrés.

\end{proof}





\section{Cohomologie sur les espaces séparés quasi compacts}

\label{section-cohomology-LC}

\noindent
Sur un tel espace, la cohomologie de {\v C}ech coïncide avec la cohomologie usuelle.

\begin{lemma}

\label{lemma-cech-always}

Soit $X$ un espace topologique et soit $\mathcal{F}$ un faisceau abélien sur cet espace. Alors l'application $\check{H}^1(X, \mathcal{F}) \to H^1(X, \mathcal{F})$ définie dans (\ref{equation-cech-to-cohomology}) est un isomorphisme.

\end{lemma}

\begin{proof}
Soit $\mathcal{U}$ un recouvrement ouvert de $X$. D'après le lemme \ref{lemma-cech-spectral-sequence}, il existe une suite exacte $$
0 \to \check{H}^1(\mathcal{U}, \mathcal{F}) \to H^1(X, \mathcal{F})
\to \check{H}^0(\mathcal{U}, \underline{H}^1(\mathcal{F}))
$$ L'application est donc injective. Pour établir la surjectivité, il suffit de montrer que tout élément de $\check{H}^0(\mathcal{U}, \underline{H}^1(\mathcal{F}))$ s'envoie sur zéro après remplacement de $\mathcal{U}$ par un raffinement. Cela résulte immédiatement des définitions et du fait que $\underline{H}^1(\mathcal{F})$ est un préfaisceau de groupes abéliens dont la faisceautisation est nulle par le caractère local de la cohomologie ; voir le lemme \ref{lemma-kill-cohomology-class-on-covering}.
\end{proof}

\begin{lemma}

\label{lemma-cech-Hausdorff-quasi-compact}
Soit $X$ un espace topologique séparé et quasi-compact. Soit $\mathcal{F}$ un faisceau abélien sur $X$. Alors l'application $\check{H}^n(X, \mathcal{F}) \to H^n(X, \mathcal{F})$ définie dans (\ref{equation-cech-to-cohomology}) est un isomorphisme pour tout $n$.
\end{lemma}

\begin{proof}
Nous savons déjà que $\check{H}^n(X, -) \to H^n(X, -)$ est un isomorphisme de foncteurs pour $n = 0, 1$ ; voir le lemme \ref{lemma-cech-always}. Les foncteurs $H^n(X, -)$ forment un $\delta$-foncteur universel ; voir Catégories dérivées, lemme \ref{derived-lemma-higher-derived-functors}. Si nous montrons que $\check{H}^n(X, -)$ forme un $\delta$-foncteur universel et que $\check{H}^n(X, -) \to H^n(X, -)$ est compatible avec les morphismes de bord, alors l'application sera automatiquement un isomorphisme par unicité des $\delta$-foncteurs universels ; voir Homologie, lemme \ref{homology-lemma-uniqueness-universal-delta-functor}.

\medskip\noindent

Soit $0 \to \mathcal{F} \to \mathcal{G} \to \mathcal{H} \to 0$
une suite exacte courte de faisceaux abéliens sur $X$. Soit $\mathcal{U} : X = \bigcup_{i \in I} U_i$ un recouvrement ouvert. On obtient alors un complexe de complexes
$$
0 \to \check{\mathcal{C}}^\bullet(\mathcal{U}, \mathcal{F}) \to
\check{\mathcal{C}}^\bullet(\mathcal{U}, \mathcal{G}) \to
\check{\mathcal{C}}^\bullet(\mathcal{U}, \mathcal{H}) \to 0
$$
qui n'est en général pas exact à droite. Cette suite définit les morphismes
$$
\check{H}^n(\mathcal{U}, \mathcal{F}) \to
\check{H}^n(\mathcal{U}, \mathcal{G}) \to
\check{H}^n(\mathcal{U}, \mathcal{H})
$$
mais ne suffit pas à définir un morphisme de bord
$\delta : \check{H}^n(\mathcal{U}, \mathcal{H}) \to
\check{H}^{n + 1}(\mathcal{U}, \mathcal{F})$. En effet une telle chose n'existera pas en général. Cependant, compte tenu d'un élément $\overline{h} \in \check{H}^n(\mathcal{U}, \mathcal{H})$
représenté par le cocycle
$h = (h_{i_0 \ldots i_n})$
nous pouvons choisir des recouvrements ouverts
$$
U_{i_0 \ldots i_n} = \bigcup W_{i_0 \ldots i_n, k}
$$
de sorte que $h_{i_0 \ldots i_n}|_{W_{i_0 \ldots i_n, k}}$
se relève en une section de $\mathcal{G}$ sur $W_{i_0 \ldots i_n, k}$. D'après Topologie, lemme \ref{topology-lemma-refine-covering},
(c'est ici que nous utilisons l'hypothèse que $X$ est séparé et quasi-compact), nous pouvons choisir un recouvrement ouvert $\mathcal{V} : X = \bigcup_{j \in J} V_j$
et $\alpha : J \to I$ tels que $V_j \subset U_{\alpha(j)}$
(c'est un raffinement) et tel que pour tous $j_0, \ldots, j_n \in J$
il existe un $k$ tel que
$V_{j_0 \ldots j_n} \subset W_{\alpha(j_0) \ldots \alpha(j_n), k}$. Nous obtenons des morphismes de complexes
$$
\xymatrix{
0 \ar[r] &
\check{\mathcal{C}}^\bullet(\mathcal{U}, \mathcal{F}) \ar[d] \ar[r] &
\check{\mathcal{C}}^\bullet(\mathcal{U}, \mathcal{G}) \ar[d] \ar[r] &
\check{\mathcal{C}}^\bullet(\mathcal{U}, \mathcal{H}) \ar[d] \ar[r] &
0 \\
0 \ar[r] &
\check{\mathcal{C}}^\bullet(\mathcal{V}, \mathcal{F}) \ar[r] &
\check{\mathcal{C}}^\bullet(\mathcal{V}, \mathcal{G}) \ar[r] &
\check{\mathcal{C}}^\bullet(\mathcal{V}, \mathcal{H}) \ar[r] &
0
}
$$
En fait, les flèches verticales sont les morphismes de complexes utilisés pour définir les morphismes de transition entre les groupes de cohomologie de {\v C}ech. Notre choix de raffinement montre que nous pouvons choisir
$$
g_{j_0 \ldots j_n} \in
\mathcal{G}(V_{j_0 \ldots j_n}),\quad
g_{j_0 \ldots j_n} \longmapsto
h_{\alpha(j_0) \ldots \alpha(j_n)}|_{V_{j_0 \ldots j_n}}
$$
La cochaîne $g = (g_{j_0 \ldots j_n})$ n'est pas nécessairement un cocycle, mais son cobord de {\v C}ech $\text{d}(g)$
s'envoie sur zéro dans $\check{\mathcal{C}}^{n + 1}(\mathcal{V}, \mathcal{H})$
(par le diagramme commutatif ci-dessus et parce que $h$ est un cocycle). Ainsi, $\text{d}(g)$ est un cocycle dans
$\check{\mathcal{C}}^\bullet(\mathcal{V}, \mathcal{F})$. Cela nous permet de définir
$$
\delta(\overline{h}) = \text{classe de }\text{d}(g)\text{ dans }
\check{H}^{n + 1}(\mathcal{V}, \mathcal{F})
$$
Étant donné un élément $\xi \in \check{H}^n(X, \mathcal{G})$,
nous choisissons un recouvrement ouvert $\mathcal{U}$ et un élément
$\overline{h} \in \check{H}^n(\mathcal{U}, \mathcal{G})$
qui s'envoie sur $\xi$ dans la colimite définissant la cohomologie de {\v C}ech. Choisissons ensuite $\mathcal{V}$ et $g$ comme ci-dessus, et posons
$\delta(\xi)$ égal à l'image de $\delta(\overline{h})$
dans $\check{H}^n(X, \mathcal{F})$. Il reste à vérifier plusieurs propriétés, toutes immédiates. Il faut notamment montrer que la construction est indépendante du choix de
$\mathcal{U}, \overline{h}, \mathcal{V}, \alpha : J \to I, g$. La classe de $\text{d}(g)$ est indépendante du choix des relèvements
$g_{i_0 \ldots i_n}$, car la différence est un cobord. L'application ne dépend pas de $\alpha$\footnote{Cette vérification est importante : la non-unicité de $\alpha$ est le seul obstacle à prendre la colimite des complexes de {\v C}ech sur tous les recouvrements de $X$ afin d'obtenir une suite exacte courte de complexes calculant la cohomologie de {\v C}ech.}, car deux choix de $\alpha$ donnent des morphismes verticaux homotopes entre les complexes du diagramme ci-dessus ; voir la section \ref{section-refinements-cech}. Pour les autres choix intervenant dans la construction, toute famille finie de recouvrements ouverts de $X$ admet un raffinement ouvert commun. L'additivité se vérifie de la même manière. Enfin, il faut vérifier que les applications
$\check{H}^n(X, -) \to H^n(X, -)$ sont compatibles avec les morphismes de bord. Pour ce faire, nous choisissons des résolutions injectives
$$
\xymatrix{
0 \ar[r] &
\mathcal{F} \ar[r] \ar[d] &
\mathcal{G} \ar[r] \ar[d] &
\mathcal{H} \ar[r] \ar[d] &
0 \\
0 \ar[r] &
\mathcal{I}_1^\bullet \ar[r] &
\mathcal{I}_2^\bullet \ar[r] &
\mathcal{I}_3^\bullet \ar[r] &
0
}
$$
comme dans Catégories dérivées, lemme \ref{derived-lemma-injective-resolution-ses}. Ceci donne un diagramme commutatif
$$
\xymatrix{
0 \ar[r] &
\check{\mathcal{C}}^\bullet(\mathcal{U}, \mathcal{F}) \ar[r] \ar[d] &
\check{\mathcal{C}}^\bullet(\mathcal{U}, \mathcal{F}) \ar[r] \ar[d] &
\check{\mathcal{C}}^\bullet(\mathcal{U}, \mathcal{F}) \ar[r] \ar[d] &
0 \\
0 \ar[r] &
\text{Tot}(\check{\mathcal{C}}^\bullet(\mathcal{U}, \mathcal{I}_1^\bullet))
\ar[r] &
\text{Tot}(\check{\mathcal{C}}^\bullet(\mathcal{U}, \mathcal{I}_2^\bullet))
\ar[r] &
\text{Tot}(\check{\mathcal{C}}^\bullet(\mathcal{U}, \mathcal{I}_3^\bullet))
\ar[r] &
0
}
$$
Ici $\mathcal{U}$ est un recouvrement ouvert comme ci-dessus et les flèches verticales sont celles utilisées pour définir les morphismes
$\check{H}^n(\mathcal{U}, -) \to H^n(X, -)$ ; voir le lemme \ref{lemma-cech-cohomology}. Le complexe inférieur est exact, car la suite de complexes injectifs est scindée exacte terme à terme. Le morphisme de bord en cohomologie est donc calculé par la procédure usuelle appliquée à cette suite exacte inférieure ; voir Homologie, lemme \ref{homology-lemma-long-exact-sequence-cochain}. Il en va de même après passage au raffinement
$\mathcal{V}$, où le morphisme de bord en cohomologie de {\v C}ech a été défini. Les morphismes de bord coïncident donc, puisqu'ils reposent sur la même construction (lorsque le premier est défini sur un élément de la cohomologie de {\v C}ech associé à un recouvrement ouvert donné). Ceci achève la construction de la structure de $\delta$-foncteur sur la cohomologie de {\v C}ech et établit sa compatibilité avec la
structure de $\delta$-foncteur sur la cohomologie usuelle.

\medskip\noindent
Enfin, le lemme \ref{lemma-injective-trivial-cech} montre que la cohomologie supérieure de {\v C}ech est nulle sur les faisceaux injectifs. La cohomologie de {\v C}ech forme donc un $\delta$-foncteur universel, d'après Homologie, lemme \ref{homology-lemma-efface-implies-universal}.
\end{proof}

\begin{lemma}

\label{lemma-cohomology-of-closed}

\begin{reference}
\cite[Expose V bis, 4.1.3]{SGA4}
\end{reference}
Soit $X$ un espace topologique. Soit $Z \subset X$ un sous-ensemble quasi-compact tel que deux points quelconques de $Z$ possèdent des voisinages ouverts disjoints dans $X$. Pour tout faisceau abélien $\mathcal{F}$ sur $X$, le morphisme canonique $$
\colim H^p(U, \mathcal{F})
\longrightarrow
H^p(Z, \mathcal{F}|_Z)
$$ où la colimite porte sur les voisinages ouverts $U$ de $Z$ dans $X$, est un isomorphisme.
\end{lemma}

\begin{proof}

Commençons par le cas $p = 0$. L'injectivité découle de la définition de $\mathcal{F}|_Z$ et vaut en général (pour tout sous-ensemble de tout espace topologique $X$). Supposons ensuite que
$s \in H^0(Z, \mathcal{F}|_Z)$. Nous pouvons trouver des ouverts $U_i \subset X$
tels que $Z \subset \bigcup U_i$ et tels que $s|_{Z \cap U_i}$
vient de $s_i \in \mathcal{F}(U_i)$. Il s'ensuit qu'il existe des ouvertures $W_{ij} \subset U_i \cap U_j$ avec
$W_{ij} \cap Z = U_i \cap U_j \cap Z$ tels que
$s_i|_{W_{ij}} = s_j|_{W_{ij}}$. En appliquant Topologie, lemme
\ref{topology-lemma-lift-covering-of-quasi-compact-hausdorff-subset}
nous trouvons des ouverts $V_i$ de $X$ tels que $V_i \subset U_i$ et $V_i \cap V_j \subset W_{ij}$. Nous voyons donc que
$s_i|_{V_i}$ colle sur une section de $\mathcal{F}$ sur le voisinage ouvert $\bigcup V_i$ de $Z$.

\medskip\noindent
Pour achever la démonstration, il suffit de montrer que, si $\mathcal{I}$ est un faisceau abélien injectif sur $X$, alors $H^p(Z, \mathcal{I}|_Z) = 0$ pour $p > 0$. Cela résulte des suites exactes courtes et du décalage des degrés ; nous omettons les détails. Supposons donc que $\overline{\xi}$ soit un élément de $H^p(Z, \mathcal{I}|_Z)$ pour un certain $p > 0$. D'après le lemme \ref{lemma-cech-Hausdorff-quasi-compact}, l'élément $\overline{\xi}$ provient de $\check{H}^p(\mathcal{V}, \mathcal{I}|_Z)$ pour un recouvrement ouvert $\mathcal{V} : Z = \bigcup V_i$ de $Z$. Écrivons $\overline{\xi}$ comme l'image de la classe d'un cocycle $\xi = (\xi_{i_0 \ldots i_p})$ de $\check{\mathcal{C}}^p(\mathcal{V}, \mathcal{I}|_Z)$.

\medskip\noindent
Soit $\mathcal{I}' \subset \mathcal{I}|_Z$ le sous-préfaisceau défini par la formule $$
\mathcal{I}'(V) =
\{s \in \mathcal{I}|_Z(V) \mid
\exists (U, t),\ U \subset X\text{ ouvert},
\ t \in \mathcal{I}(U),\ V = Z \cap U,\ s = t|_{Z \cap U} \}
$$ Alors $\mathcal{I}|_Z$ est la faisceautisation de $\mathcal{I}'$. Ainsi, pour chaque $(p + 1)$-uplet $i_0 \ldots i_p$, nous pouvons trouver un recouvrement ouvert $V_{i_0 \ldots i_p} = \bigcup W_{i_0 \ldots i_p, k}$ tel que $\xi_{i_0 \ldots i_p}|_{W_{i_0 \ldots i_p, k}}$ soit une section de $\mathcal{I}'$. En appliquant Topologie, lemme \ref{topology-lemma-refine-covering}, nous pouvons, après avoir raffiné $\mathcal{V}$, supposer que chaque $\xi_{i_0 \ldots i_p}$ est une section du préfaisceau $\mathcal{I}'$.

\medskip\noindent

Écrivons $V_i = Z \cap U_i$ pour certains ouverts $U_i \subset X$. Puisque $\mathcal{I}$ est flasque (lemme \ref{lemma-injective-flasque}) et que $\xi_{i_0 \ldots i_p}$ est une section de $\mathcal{I}'$,
pour chaque $(p + 1)$-uplet $i_0 \ldots i_p$, nous pouvons choisir une section $s_{i_0 \ldots i_p} \in \mathcal{I}(U_{i_0 \ldots i_p})$
dont la restriction est $\xi_{i_0 \ldots i_p}$ sur
$V_{i_0 \ldots i_p} = Z \cap U_{i_0 \ldots i_p}$. (On peut éviter ici d'utiliser le fait que les injectifs sont flasques en appliquant une fois de plus Topologie, lemme
\ref{topology-lemma-lift-covering-of-quasi-compact-hausdorff-subset}.) Soit $s = (s_{i_0 \ldots i_p})$ la cochaîne correspondante pour le recouvrement ouvert $U = \bigcup U_i$. Puisque $\text{d}(\xi) = 0$, les sections
$\text{d}(s)_{i_0 \ldots i_{p + 1}}$ se restreignent à zéro sur $Z \cap U_{i_0 \ldots i_{p + 1}}$. Par conséquent, d'après les remarques initiales de la démonstration, il existe des ouverts
$W_{i_0 \ldots i_{p + 1}} \subset U_{i_0 \ldots i_{p + 1}}$
avec $Z \cap W_{i_0 \ldots i_{p + 1}} = Z \cap U_{i_0 \ldots i_{p + 1}}$
tels que $\text{d}(s)_{i_0 \ldots i_{p + 1}}|_{W_{i_0 \ldots i_{p + 1}}} = 0$. D'après Topologie, lemme
\ref{topology-lemma-lift-covering-of-quasi-compact-hausdorff-subset}
nous pouvons trouver $U'_i \subset U_i$ tels que $Z \subset \bigcup U'_i$
et $U'_{i_0 \ldots i_{p + 1}} \subset W_{i_0 \ldots i_{p + 1}}$. Alors $s' = (s'_{i_0 \ldots i_p})$ avec
$s'_{i_0 \ldots i_p} = s_{i_0 \ldots i_p}|_{U'_{i_0 \ldots i_p}}$
est un cocycle à valeurs dans $\mathcal{I}$ pour le recouvrement ouvert
$U' = \bigcup U'_i$ d'un voisinage ouvert de $Z$. Puisque les groupes de cohomologie supérieure de {\v C}ech de $\mathcal{I}$ sont nuls (lemme \ref{lemma-injective-trivial-cech}), nous concluons que $s'$ est un cobord. Il s'ensuit que l'image de
$\xi$ dans le complexe de {\v C}ech associé au recouvrement ouvert
$Z = \bigcup Z \cap U'_i$ est un cobord, ce qui achève la démonstration.

\end{proof}





\section{L'application de changement de base}

\label{section-base-change-map}

\noindent

Nous devons savoir construire le morphisme de changement de base dans certains cas. Comme nous n'avons pas encore étudié l'image inverse dérivée, nous ne considérons ici que le changement de base par un morphisme plat d'espaces annelés. Avant d'énoncer le résultat, examinons l'image inverse plate sur la catégorie dérivée. Soit $g : X \to Y$ un morphisme plat d'espaces annelés. D'après Faisceaux de modules, lemme \ref{modules-lemma-pullback-flat},
le foncteur $g^* : \textit{Mod}(\mathcal{O}_Y) \to
\textit{Mod}(\mathcal{O}_X)$ est exact. Il admet donc un foncteur dérivé
$$
g^* : D^{+}(Y) \to D^{+}(X)
$$
qui se calcule simplement en appliquant l'image inverse à un représentant de l'objet de $D^{+}(Y)$ considéré ; voir Catégories dérivées, lemme \ref{derived-lemma-right-derived-exact-functor}. Ainsi, comme annoncé, nous notons ce foncteur $g^*$ plutôt que
$Lg^*$.

\begin{lemma}

\label{lemma-base-change-map-flat-case}
Considérons le diagramme commutatif suivant d'espaces annelés : $$
\xymatrix{
X' \ar[r]_{g'} \ar[d]_{f'} &
X \ar[d]^f \\
S' \ar[r]^g &
S
}
$$ Soit $\mathcal{F}^\bullet$ un complexe borné inférieurement de $\mathcal{O}_X$-modules. Supposons $g$ et $g'$ plats. Il existe alors un morphisme canonique de changement de base $$
g^*Rf_*\mathcal{F}^\bullet
\longrightarrow
R(f')_*(g')^*\mathcal{F}^\bullet
$$ dans $D^{+}(S')$.
\end{lemma}

\begin{proof}

Choisissons des résolutions injectives $\mathcal{F}^\bullet \to \mathcal{I}^\bullet$
et $(g')^*\mathcal{F}^\bullet \to \mathcal{J}^\bullet$. D'après le lemme \ref{lemma-pushforward-injective-flat},
$(g')_*\mathcal{J}^\bullet$ est un complexe de modules injectifs représentant
$R(g')_*(g')^*\mathcal{F}^\bullet$. D'où par catégories dérivées, lemmes \ref{derived-lemma-morphisms-lift}
et \ref{derived-lemma-morphisms-equal-up-to-homotopy}
le morphisme $\beta$ du diagramme
$$
\xymatrix{
(g')_*(g')^*\mathcal{F}^\bullet \ar[r] &
(g')_*\mathcal{J}^\bullet \\
\mathcal{F}^\bullet \ar[u]^{adjunction} \ar[r] &
\mathcal{I}^\bullet \ar[u]_\beta
}
$$
existe et est unique à homotopie près. En poussant vers $S$, nous obtenons
$$
f_*\beta :
f_*\mathcal{I}^\bullet
\longrightarrow
f_*(g')_*\mathcal{J}^\bullet
=
g_*(f')_*\mathcal{J}^\bullet
$$
Par adjonction entre $g^*$ et $g_*$, nous obtenons un morphisme de complexes
$g^*f_*\mathcal{I}^\bullet \to (f')_*\mathcal{J}^\bullet$. Ce morphisme est unique à homotopie près, puisque le seul choix effectué au cours de la construction est celui du morphisme $\beta$
et tout a été fait au niveau des complexes.

\end{proof}

\begin{remark}

\label{remark-correct-version-base-change-map}
La version « correcte » du morphisme de changement de base est le morphisme $$
Lg^* Rf_* \mathcal{F}^\bullet
\longrightarrow
R(f')_* L(g')^*\mathcal{F}^\bullet.
$$ Sa construction fait intervenir des complexes non bornés ; voir la remarque \ref{remark-base-change}.
\end{remark}





\section{Changement de base propre en topologie}

\label{section-proper-base-change}

\noindent
Dans cette section, nous démontrons une version très générale du théorème de changement de base propre en topologie. Elle affirme que les germes des images directes supérieures $R^pf_*$ peuvent se calculer sur la fibre.

\begin{lemma}

\label{lemma-proper-base-change}
Soit $f : (X, \mathcal{O}_X) \to (Y, \mathcal{O}_Y)$ un morphisme d'espaces annelés. Soit $y \in Y$. Supposons que
\begin{enumerate}
\item $f$ est fermé,
\item $f$ est séparé, et
\item $f^{-1}(y)$ est quasi-compact.
\end{enumerate}
Alors, pour $E$ dans $D^+(\mathcal{O}_X)$, nous avons $(Rf_*E)_y = R\Gamma(f^{-1}(y), E|_{f^{-1}(y)})$ dans $D^+(\mathcal{O}_{Y, y})$.
\end{lemma}

\begin{proof}

Le morphisme de changement de base du lemme \ref{lemma-base-change-map-flat-case}
donne un morphisme canonique $(Rf_*E)_y \to R\Gamma(f^{-1}(y), E|_{f^{-1}(y)})$. Pour montrer que ce morphisme est un isomorphisme, représentons $E$ par un complexe borné inférieurement de modules injectifs $\mathcal{I}^\bullet$. Posons $Z = f^{-1}(\{y\})$. Les hypothèses du lemme \ref{lemma-cohomology-of-closed}
sont satisfaites ; voir Topologie, lemme \ref{topology-lemma-separated}. Les restrictions
$\mathcal{I}^n|_Z$ sont acycliques pour $\Gamma(Z, -)$. Ainsi $R\Gamma(Z, E|_Z)$ est représenté par le complexe $\Gamma(Z, \mathcal{I}^\bullet|_Z)$ ; voir Catégories dérivées, lemme \ref{derived-lemma-leray-acyclicity}. En d'autres termes, nous devons montrer que le morphisme
$$
\colim_V \mathcal{I}^\bullet(f^{-1}(V))
\longrightarrow
\Gamma(Z, \mathcal{I}^\bullet|_Z)
$$
est un isomorphisme. D'après le lemme \ref{lemma-cohomology-of-closed},
nous voyons qu'il suffit de montrer que la collection de voisinages ouverts
$f^{-1}(V)$ de $Z = f^{-1}(\{y\})$
est cofinal dans le système de tous les voisinages ouverts. Si $f^{-1}(\{y\}) \subset U$ est un voisinage ouvert, alors, puisque $f$ est fermé, l'ensemble $V = Y \setminus f(X \setminus U)$ est un voisinage ouvert de $y$ tel que $f^{-1}(V) \subset U$. Cela démontre le lemme.

\end{proof}

\begin{theorem}
[changement de base propre]
\label{theorem-proper-base-change}

\begin{reference}
\cite[Expose V bis, 4.1.1]{SGA4}
\end{reference}
Considérez un carré cartésien d'espaces topologiques $$
\xymatrix{
X' = Y' \times_Y X \ar[d]_{f'} \ar[r]_-{g'} & X \ar[d]^f \\
Y' \ar[r]^g & Y
}
$$ Supposons que $f$ soit propre. Soit $E$ un objet de $D^+(X)$. Alors le morphisme de changement de base $$
g^{-1}Rf_*E \longrightarrow Rf'_*(g')^{-1}E
$$ du lemme \ref{lemma-base-change-map-flat-case} est un isomorphisme dans $D^+(Y')$.
\end{theorem}

\begin{proof}

Soit $y' \in Y'$ un point d'image $y \in Y$. Il suffit de montrer que le morphisme de changement de base induit un isomorphisme sur les germes en $y'$. Puisque $f$ est propre, il en va de même de $f'$ ; les fibres de $f$ et de $f'$ sont quasi-compactes, et $f$ et $f'$ sont fermés. Voir Topologie, théorème \ref{topology-theorem-characterize-proper} et lemme \ref{topology-lemma-base-change-separated}. Nous pouvons donc appliquer deux fois le lemme \ref{lemma-proper-base-change} pour obtenir
$$
(Rf'_*(g')^{-1}E)_{y'} = R\Gamma((f')^{-1}(y'), (g')^{-1}E|_{(f')^{-1}(y')})
$$
et
$$
(Rf_*E)_y = R\Gamma(f^{-1}(y), E|_{f^{-1}(y)})
$$
Le morphisme induit sur les fibres $(f')^{-1}(y') \to f^{-1}(y)$ est un homéomorphisme d'espaces topologiques, et l'image inverse de
$E|_{f^{-1}(y)}$ est $(g')^{-1}E|_{(f')^{-1}(y')}$. Le résultat en découle.

\end{proof}

\begin{lemma}
[changement de base propre pour faisceaux d'ensembles]
\label{lemma-proper-base-change-sheaves-of-sets}
Considérez un carré cartésien d'espaces topologiques $$
\xymatrix{
X' \ar[d]_{f'} \ar[r]_-{g'} & X \ar[d]^f \\
Y' \ar[r]^g & Y
}
$$ Supposons que $f$ soit propre. Alors $g^{-1}f_*\mathcal{F} = f'_*(g')^{-1}\mathcal{F}$ pour tout faisceau d'ensembles $\mathcal{F}$ sur $X$.
\end{lemma}

\begin{proof}
On raisonne exactement comme dans la démonstration du théorème \ref{theorem-proper-base-change} : il suffit d'établir $(f_*\mathcal{F})_y = \Gamma(X_y, \mathcal{F}|_{X_y})$. Comme dans le lemme \ref{lemma-proper-base-change}, on se ramène ensuite au cas $p = 0$ du lemme \ref{lemma-cohomology-of-closed} pour les faisceaux d'ensembles. La première partie de la démonstration du lemme \ref{lemma-cohomology-of-closed} s'applique aux faisceaux d'ensembles et achève la preuve. Nous omettons quelques détails.
\end{proof}



\section{Cohomologie et colimites}

\label{section-limits}

\noindent

Soit $X$ un espace annelé. Soit $(\mathcal{F}_i, \varphi_{ii'})$ un système de faisceaux de $\mathcal{O}_X$-modules indexé par l'ensemble dirigé $I$ ; voir Catégories, section \ref{categories-section-posets-limits}. Pour chaque $i$, le morphisme canonique
$\mathcal{F}_i \to \colim_i \mathcal{F}_i$ induit un morphisme canonique
$$
\colim_i H^p(X, \mathcal{F}_i)
\longrightarrow
H^p(X, \colim_i \mathcal{F}_i)
$$
pour chaque $p \geq 0$. Il existe bien sûr un morphisme analogue pour chaque $U \subset X$. En général, ces morphismes ne sont pas des isomorphismes, même pour $p = 0$. Dans cette section, nous généralisons le résultat de Faisceaux, \ref{sheaves-lemma-directed-colimits-sections}. Voir aussi Faisceaux de modules, lemme \ref{modules-lemma-finite-presentation-quasi-compact-colimit}
(dans le cas spécial $\mathcal{G} = \mathcal{O}_X$).

\begin{lemma}

\label{lemma-quasi-separated-cohomology-colimit}
Soit $X$ un espace annelé. Supposons que l'espace topologique sous-jacent à $X$ possède les propriétés suivantes :
\begin{enumerate}
\item il existe une base de sous-ensembles ouverts quasi compacts, et
\item l'intersection de deux ouverts quasi compacts est quasi-compacte.
\end{enumerate}
Alors, pour tout système dirigé $(\mathcal{F}_i, \varphi_{ii'})$ de faisceaux de $\mathcal{O}_X$-modules et pour tout ouvert quasi compact $U \subset X$, le morphisme canonique $$
\colim_i H^q(U, \mathcal{F}_i)
\longrightarrow
H^q(U, \colim_i \mathcal{F}_i)
$$ est un isomorphisme pour chaque $q \geq 0$.
\end{lemma}

\begin{proof}
Il importe de raisonner simultanément pour tous les ouverts quasi compacts $U \subset X$. Pour $q = 0$ et tout ouvert quasi compact $U \subset X$, le résultat découle de Faisceaux, lemme \ref{sheaves-lemma-directed-colimits-sections}, combiné avec Topologie, lemme \ref{topology-lemma-topology-quasi-separated-scheme}. Supposons le résultat établi pour tout $q \leq q_0$ et démontrons-le pour $q = q_0 + 1$.

\medskip\noindent

Par convention, l'ensemble d'indices $I$ est dirigé, et tout système de $\mathcal{O}_X$-modules $(\mathcal{F}_i, \varphi_{ii'})$
indexé par $I$ est un système dirigé. D'après Injectifs, lemme \ref{injectives-lemma-sheaves-modules-space}, la catégorie des $\mathcal{O}_X$-modules possède des plongements injectifs fonctoriels. Ainsi, pour tout système $(\mathcal{F}_i, \varphi_{ii'})$, il existe un système $(\mathcal{I}_i, \varphi_{ii'})$ tel que chaque $\mathcal{I}_i$ soit un $\mathcal{O}_X$-module injectif et un morphisme de systèmes formé de morphismes injectifs de $\mathcal{O}_X$-modules
$\mathcal{F}_i \to \mathcal{I}_i$. Notons $\mathcal{Q}_i$ le conoyau ; nous obtenons ainsi des suites exactes courtes
$$
0 \to
\mathcal{F}_i \to
\mathcal{I}_i \to
\mathcal{Q}_i \to 0.
$$
Nous affirmons que la séquence
$$
0 \to
\colim_i \mathcal{F}_i \to
\colim_i \mathcal{I}_i \to
\colim_i \mathcal{Q}_i \to 0.
$$
est également une suite exacte courte de $\mathcal{O}_X$-modules ; cela se vérifie sur les germes. Voir \ref{sheaves-section-limits-presheaves}
et \ref{sheaves-section-limits-sheaves}
puisque la colimite dirigée d'une suite exacte courte de groupes abéliens est exacte (voir Algèbre, lemme \ref{algebra-lemma-directed-colimit-exact}), le résultat en découle. Nous affirmons que
$H^q(U, \colim_i \mathcal{I}_i) = 0$ pour tout ouvert quasi compact $U \subset X$ et tout $q \geq 1$. Admettons provisoirement cette assertion et considérons le diagramme
$$
\xymatrix{
\colim_i H^{q_0}(U, \mathcal{I}_i) \ar[d] \ar[r] &
\colim_i H^{q_0}(U, \mathcal{Q}_i) \ar[d] \ar[r] &
\colim_i H^{q_0 + 1}(U, \mathcal{F}_i) \ar[d] \ar[r] &
0 \ar[d] \\
H^{q_0}(U, \colim_i \mathcal{I}_i) \ar[r] &
H^{q_0}(U, \colim_i \mathcal{Q}_i) \ar[r] &
H^{q_0 + 1}(U, \colim_i \mathcal{F}_i) \ar[r] &
0
}
$$
Le zéro dans le coin inférieur droit provient de l'assertion, et celui du coin supérieur droit du fait que les faisceaux
$\mathcal{I}_i$ sont injectifs. La ligne supérieure est exacte par Algèbre, lemme \ref{algebra-lemma-directed-colimit-exact}. Le lemme du serpent donne alors le résultat pour $q = q_0 + 1$.

\medskip\noindent

Il reste à démontrer l'assertion. Nous appliquons le lemme \ref{lemma-cech-vanish-basis}. Soit $\mathcal{B}$ l'ensemble de tous les ouverts quasi compacts de $X$ ; par hypothèse, c'est une base de la topologie de $X$. Soit $\text{Cov}$ l'ensemble des recouvrements ouverts finis
$\mathcal{U} : U = \bigcup_{j = 1, \ldots, m} U_j$ tels que $U$ et chaque $U_j$ soient des ouverts quasi compacts de $X$. D'après le cas $q = 0$,
pour $\mathcal{U} \in \text{Cov}$, nous avons
$$
\check{\mathcal{C}}^\bullet(\mathcal{U}, \colim_i \mathcal{I}_i)
=
\colim_i \check{\mathcal{C}}^\bullet(\mathcal{U}, \mathcal{I}_i)
$$
parce que toutes les intersections multiples $U_{j_0 \ldots j_p}$
sont quasi-compactes. D'après le lemme \ref{lemma-injective-trivial-cech},
chacun des complexes de la colimite de complexes de {\v C}ech est acyclique en degré $\geq 1$. Par Algèbre, lemme \ref{algebra-lemma-directed-colimit-exact},
le complexe de {\v C}ech
$\check{\mathcal{C}}^\bullet(\mathcal{U}, \colim_i \mathcal{I}_i)$
est acyclique en degrés $\geq 1$. En d'autres termes, nous voyons que
$\check{H}^p(\mathcal{U},  \colim_i \mathcal{I}_i) = 0$
pour tout $p \geq 1$. Les hypothèses du lemme \ref{lemma-cech-vanish-basis} sont donc satisfaites, ce qui démontre l'assertion.

\end{proof}

\begin{lemma}

\label{lemma-higher-direct-image-colimit}
Soit $f : X \to Y$ une application continue d'espaces topologiques. Soit $(\mathcal{F}_i, \varphi_{ii'})$ un système de faisceaux abéliens sur $X$. Posons $\mathcal{F} = \colim \mathcal{F}_i$. Soit $p \geq 0$ un entier. Supposons que les ouverts $V \subset Y$ tels que $H^p(f^{-1}(V), \mathcal{F}) = \colim H^p(f^{-1}(V), \mathcal{F}_i)$ forment une base de la topologie de $Y$. Alors $R^pf_*\mathcal{F} = \colim R^pf_*\mathcal{F}_i$.
\end{lemma}

\begin{proof}

Rappelons que $R^pf_*\mathcal{F}$ est la faisceautisation du préfaisceau
$\mathcal{G}$ qui associe à $V$ le groupe $H^p(f^{-1}(V), \mathcal{F})$ ; voir le lemme \ref{lemma-describe-higher-direct-images}. De même, 
$R^pf_*\mathcal{F}_i$ est la faisceautisation du préfaisceau
$\mathcal{G}_i$ qui associe à $V$ le groupe $H^p(f^{-1}(V), \mathcal{F}_i)$. Rappelons que la faisceautisation est adjointe à gauche au foncteur d'inclusion des faisceaux dans les préfaisceaux ; voir Faisceaux, section
\ref{sheaves-section-sheafification}. La faisceautisation commute donc aux colimites ; voir Catégories, lemme \ref{categories-lemma-adjoint-exact}. Il suffit ainsi de montrer que le morphisme de préfaisceaux, la colimite étant prise dans la catégorie des préfaisceaux,
$$
\colim \mathcal{G}_i \longrightarrow \mathcal{G}
$$
induit un isomorphisme sur les faisceautisés. Il suffit pour cela de montrer que les préfaisceaux $\mathcal{G}$ et $\colim \mathcal{G}_i$
coïncident sur une base de la topologie de $Y$. En effet, les germes de leurs faisceautisés, calculables directement à partir des valeurs des préfaisceaux sur les éléments de la base, coïncident alors. L'égalité requise est précisément l'hypothèse du lemme.

\end{proof}

\noindent

Formulons maintenant l'analogue, pour la cohomologie, de Faisceaux, lemme \ref{sheaves-lemma-descend-opens}.
Soit $X$ un espace spectral, écrit comme limite cofiltrée d'espaces spectraux $X_i$ pour un diagramme dont les morphismes de transition sont spectraux, comme dans Topologie, lemme \ref{topology-lemma-directed-inverse-limit-spectral-spaces}. Supposons donnés

\begin{enumerate}

\item un faisceau abélien $\mathcal{F}_i$ sur $X_i$ pour tout
$i \in \Ob(\mathcal{I})$,
\item pour $a : j \to i$, un $f_a$-morphisme
$\varphi_a : \mathcal{F}_i \to \mathcal{F}_j$ de faisceaux abéliens (voir Faisceaux, définition \ref{sheaves-definition-f-map})

\end{enumerate}

tels que $\varphi_c = \varphi_b \circ \varphi_a$
dès que $c = a \circ b$. Posons $\mathcal{F} = \colim p_i^{-1}\mathcal{F}_i$
sur $X$.

\begin{lemma}

\label{lemma-colimit}
Dans la situation ci-dessus, soit $i \in \Ob(\mathcal{I})$ et soit $U_i \subset X_i$ un ouvert quasi compact. Alors $$
\colim_{a : j \to i} H^p(f_a^{-1}(U_i), \mathcal{F}_j) =
H^p(p_i^{-1}(U_i), \mathcal{F})
$$ pour tous $p \geq 0$. En particulier, nous avons $H^p(X, \mathcal{F}) = \colim H^p(X_i, \mathcal{F}_i)$.
\end{lemma}

\begin{proof}
Le cas $p = 0$ est Faisceaux, lemme \ref{sheaves-lemma-descend-opens}.

\medskip\noindent

Dans ce paragraphe, montrons que l'on peut trouver un morphisme de systèmes
$(\gamma_i) : (\mathcal{F}_i, \varphi_a) \to (\mathcal{G}_i, \psi_a)$
où $\mathcal{G}_i$ est un faisceau abélien injectif et $\gamma_i$ est injectif. Pour chaque $i$, choisissons une injection $\mathcal{F}_i \to \mathcal{I}_i$,
où $\mathcal{I}_i$ est un faisceau abélien injectif sur $X_i$. Nous pouvons alors considérer la famille de morphismes
$$
\gamma_i :
\mathcal{F}_i
\longrightarrow
\prod\nolimits_{b : k \to i} f_{b, *}\mathcal{I}_k = \mathcal{G}_i
$$
où les morphismes composantes sont adjoints aux morphismes
$f_b^{-1}\mathcal{F}_i \to \mathcal{F}_k \to \mathcal{I}_k$. Pour $a : j \to i$ dans $\mathcal{I}$, il existe un morphisme canonique
$$
\psi_a : f_a^{-1}\mathcal{G}_i \to \mathcal{G}_j
$$
dont les composantes sont les morphismes canoniques
$f_b^{-1}f_{a \circ b, *}\mathcal{I}_k \to f_{b, *}\mathcal{I}_k$
pour $b : k \to j$. Nous obtenons ainsi une injection
$\{\gamma_i\} : \{\mathcal{F}_i, \varphi_a) \to (\mathcal{G}_i, \psi_a)$
de systèmes de faisceaux abéliens. Le faisceau de groupes abéliens $\mathcal{G}_i$ est injectif sur $X_i$ ; voir le lemme \ref{lemma-pushforward-injective-flat} et Homologie, lemme \ref{homology-lemma-product-injectives}. Ceci termine la construction.

\medskip\noindent
En raisonnant exactement comme dans la démonstration du lemme \ref{lemma-quasi-separated-cohomology-colimit}, il suffit de montrer que $H^p(X, \colim f_i^{-1}\mathcal{G}_i) = 0$ pour $p > 0$.

\medskip\noindent

Posons $\mathcal{G} = \colim f_i^{-1}\mathcal{G}_i$. Pour établir l'annulation de la cohomologie de $\mathcal{G}$ sur tout ouvert quasi compact de $X$, il suffit de montrer que la cohomologie de {\v C}ech de
$\mathcal{G}$, pour tout recouvrement ouvert $\mathcal{U}$ d'un ouvert quasi compact de
$X$ par un nombre fini d'ouverts quasi compacts, est nulle ; voir le lemme \ref{lemma-cech-vanish-basis}. Un tel recouvrement est l'image inverse par $p_i$ d'un recouvrement $\mathcal{U}_i$
de l'espace $X_i$ pour un certain $i$, d'après Topologie, lemme \ref{topology-lemma-descend-opens}. Nous avons
$$
\check{\mathcal{C}}^\bullet(\mathcal{U}, \mathcal{G}) =
\colim_{a : j \to i}
\check{\mathcal{C}}^\bullet(f_a^{-1}(\mathcal{U}_i), \mathcal{G}_j)
$$
d'après le cas $p = 0$. Le membre de droite est une colimite filtrante de complexes, chacun acyclique en degrés positifs d'après le lemme \ref{lemma-injective-trivial-cech}. On conclut alors par Algèbre, lemme \ref{algebra-lemma-directed-colimit-exact}.

\end{proof}
\section{Annulation sur les espaces topologiques noethériens}

\label{section-vanishing-Noetherian}

\noindent
L'objectif est de prouver un théorème de Grothendieck, à savoir la proposition \ref{proposition-vanishing-Noetherian}. Voir \cite{Tohoku}.

\begin{lemma}

\label{lemma-cohomology-and-closed-immersions}
Soit $i : Z \to X$ une immersion fermée d'espaces topologiques. Pour tout faisceau abélien $\mathcal{F}$ sur $Z$, on a $H^p(Z, \mathcal{F}) = H^p(X, i_*\mathcal{F})$.
\end{lemma}

\begin{proof}
Cela résulte de l'exactitude de $i_*$ (voir Modules, lemme \ref{modules-lemma-i-star-exact}); par conséquent $R^pi_* = 0$ en tant que foncteur (Catégories dérivées, lemme \ref{derived-lemma-right-derived-exact-functor}). Nous pouvons donc appliquer le lemme \ref{lemma-apply-Leray}.
\end{proof}

\begin{lemma}

\label{lemma-irreducible-constant-cohomology-zero}
Soit $X$ un espace topologique irréductible. Alors $H^p(X, \underline{A}) = 0$ pour tout $p > 0$ et tout groupe abélien $A$.
\end{lemma}

\begin{proof}

Rappelons que $\underline{A}$ est le faisceau constant défini dans Faisceaux, définition \ref{sheaves-definition-constant-sheaf}. Comme $X$ est irréductible, tout ouvert non vide $U$ est irréductible et, a fortiori, connexe. Si $U \subset X$
est un ouvert non vide, alors $\underline{A}(U) = A$. De plus, $\underline{A}(\emptyset) = 0$. Ainsi $\underline{A}$
est un faisceau abélien flasque sur $X$. L'annulation résulte du lemme \ref{lemma-flasque-acyclic}.

\end{proof}

\begin{lemma}

\label{lemma-subsheaf-of-constant-sheaf}

\begin{reference}

\cite[Page 168]{Tohoku}.

\end{reference}

Soit $X$ un espace topologique tel que l'intersection de deux ouverts quasi-compacts est quasi-compact.
$\mathcal{F} \subset \underline{\mathbf{Z}}$
un sous-faisceau engendré par un nombre fini de sections sur des ouverts quasi-compacts.
$$
(0) = \mathcal{F}_0 \subset \mathcal{F}_1 \subset \ldots \subset
\mathcal{F}_n = \mathcal{F}
$$
par des sous-faisceaux abéliens tels que, pour chaque $0 < i \leq n$,
il existe une suite exacte courte
$$
0 \to j'_!\underline{\mathbf{Z}}_V \to j_!\underline{\mathbf{Z}}_U \to
\mathcal{F}_i/\mathcal{F}_{i - 1} \to 0
$$
avec $j : U \to X$ et $j' : V \to X$ l'inclusion des ouverts quasi-compacts dans $X$.

\end{lemma}

\begin{proof}

Supposons que $\mathcal{F}$ soit engendré par les sections $s_1, \ldots, s_t$ sur les ouverts quasi-compacts $U_1, \ldots, U_t$. Comme $U_i$ est quasi-compact et que
$s_i$ est une fonction localement constante à valeurs dans $\mathbf{Z}$, nous pouvons, quitte à remplacer $U_i$ par les morceaux d'une décomposition finie en parties ouvertes et fermées, supposer que $s_i$ est constante, disons $s_i = n_i$ avec $n_i \in \mathbf{Z}$. Nous pouvons supprimer
$(U_i, n_i)$ de la liste lorsque $n_i = 0$. En changeant les signes si nécessaire, nous pouvons aussi supposer $n_i > 0$. Ensuite, pour tout sous-ensemble $I \subset \{1, \ldots, t\}$,
nous pouvons ajouter à la liste $\bigcap_{i \in I} U_i$ et $\gcd(n_i, i \in I)$. Après cette opération, la liste $(U_1, n_1), \ldots, (U_t, n_t)$
possède la propriété suivante : pour $x \in X$, posons $I_x = \{i \in \{1, \ldots, t\} \mid x \in U_i\}$. Alors $\gcd(n_i, i \in I_x)$ est égal à l'un des $n_i$, pour un certain $i \in I_x$.

\medskip\noindent
Pour filtration, prenons $\mathcal{F}_0 = (0)$ et $\mathcal{F}_n$ le sous-faisceau engendré par les sections $n_i$ sur $U_i$ pour les $i$ tels que $n_i \leq n$. Il est clair que $\mathcal{F}_n = \mathcal{F}$ pour $n \gg 0$. Le quotient $\mathcal{F}_n/\mathcal{F}_{n - 1}$ est engendré par la section $n$ sur $U = \bigcup_{n_i \leq n} U_i$, et le noyau de $j_!\underline{\mathbf{Z}}_U \to \mathcal{F}_n/\mathcal{F}_{n - 1}$ est engendré par la section $n$ sur $V = \bigcup_{n_i \leq n - 1} U_i$. On obtient ainsi une suite exacte courte comme dans l'énoncé du lemme.
\end{proof}

\begin{lemma}

\label{lemma-vanishing-generated-one-section}

\begin{reference}

Il s'agit d'un cas particulier: \cite[proposition 3.6.1]{Tohoku}.

\end{reference}

Soit $X$ un espace topologique et soit $d \geq 0$ un entier. Supposons que

\begin{enumerate}

\item  $X$ est quasi-compact,
\item les ouverts quasi-compacts constituent une base pour $X$,
\item l'intersection de deux ouverts quasi-compacts est quasi-compact.
\item  $H^p(X, j_!\underline{\mathbf{Z}}_U) = 0$ pour tous $p > d$
et tout ouvert quasi-compact $j : U \to X$.

\end{enumerate}

Alors $H^p(X, \mathcal{F}) = 0$ pour tout $p > d$
et tout faisceau abélien $\mathcal{F}$ sur $X$.

\end{lemma}

\begin{proof}

Posons $S = \coprod_{U \subset X} \mathcal{F}(U)$, où $U$ parcourt les ouverts quasi-compacts de $X$. Pour tout sous-ensemble fini $A = \{s_1, \ldots, s_n\} \subset S$, soit $\mathcal{F}_A$ le sous-faisceau de $\mathcal{F}$ engendré par les $s_i$ (voir Modules, définition \ref{modules-definition-generated-by-local-sections}). Si $A \subset A'$, alors $\mathcal{F}_A \subset \mathcal{F}_{A'}$. Ainsi, les $\{\mathcal{F}_A\}$ forment un système indexé par l'ensemble ordonné filtrant des sous-ensembles finis de $S$. D'après Modules, lemme \ref{modules-lemma-generated-by-local-sections-stalk},
il est clair que
$$
\colim_A \mathcal{F}_A = \mathcal{F}
$$
comme on le voit sur les germes. D'après le lemme \ref{lemma-quasi-separated-cohomology-colimit}, on a
$$
H^p(X, \mathcal{F}) =
\colim_A H^p(X, \mathcal{F}_A)
$$
Il suffit donc de prouver l'annulation pour les faisceaux abéliens
$\mathcal{F}_A$. Autrement dit, il suffit de traiter le cas où $\mathcal{F}$ est engendré par un nombre fini de sections locales sur des ouverts quasi-compacts de $X$.

\medskip\noindent
Supposons que $\mathcal{F}$ soit engendré par les sections locales $s_1, \ldots, s_n$. Soit $\mathcal{F}' \subset \mathcal{F}$ le sous-faisceau engendré par $s_1, \ldots, s_{n - 1}$. Nous avons alors une suite exacte courte $$
0 \to \mathcal{F}' \to \mathcal{F} \to \mathcal{F}/\mathcal{F}' \to 0
$$ La suite exacte longue de cohomologie montre qu'il suffit d'établir l'annulation pour $\mathcal{F}'$ et $\mathcal{F}/\mathcal{F}'$, qui sont engendrés par moins de $n$ sections locales. Il suffit donc de traiter les faisceaux engendrés par une seule section locale au plus. Ce sont précisément les quotients des faisceaux $j_!\underline{\mathbf{Z}}_U$, où $U$ est un ouvert quasi-compact de $X$.

\medskip\noindent

Supposons maintenant donnée une suite exacte courte
$$
0 \to \mathcal{K} \to j_!\underline{\mathbf{Z}}_U \to \mathcal{F} \to 0
$$
où $U$ est un ouvert quasi-compact de $X$. Il suffit de montrer que $H^q(X, \mathcal{K})$ est nul pour $q \geq d + 1$. Comme ci-dessus, écrivons $\mathcal{K}$ comme colimite filtrante de sous-faisceaux $\mathcal{K}'$ engendrés par un nombre fini de sections sur des ouverts quasi-compacts. Alors $\mathcal{F}$ est la colimite filtrante des faisceaux $j_!\underline{\mathbf{Z}}_U/\mathcal{K}'$. Nous nous ramenons ainsi au cas où $\mathcal{K}$ est engendré par un nombre fini de sections sur des ouverts quasi-compacts. Remarquons que $\mathcal{K}$
est un sous-faisceau de $\underline{\mathbf{Z}}_X$. Le lemme \ref{lemma-subsheaf-of-constant-sheaf} fournit donc une filtration finie de $\mathcal{K}$ dont chaque quotient successif $\mathcal{Q}$ s'insère dans une suite exacte courte
$$
0 \to j''_!\underline{\mathbf{Z}}_W \to
j'_!\underline{\mathbf{Z}}_V \to \mathcal{Q} \to 0
$$
avec $j'' : W \to X$ et $j' : V \to X$ les inclusions des ouverts quasi-compacts. L'annulation de $H^p(X, \mathcal{Q})$ pour $p > d$ découle alors de l'hypothèse du lemme appliquée aux groupes de cohomologie de $j''_!\underline{\mathbf{Z}}_W$ et $j'_!\underline{\mathbf{Z}}_V$. En revenant à $\mathcal{K}$, un raisonnement par récurrence utilisant la suite exacte longue de cohomologie donne l'annulation cherchée.

\end{proof}

\begin{example}

\label{example-datta}

Soit $X = \mathbf{N}$ doté de la topologie dont les ouverts sont
$\emptyset$, $X$, $U_n = \{i \mid i \leq n\}$ pour $n \geq 1$. Un faisceau abélien $\mathcal{F}$ sur $X$ équivaut à la donnée d'un système projectif de groupes abéliens $A_n = \mathcal{F}(U_n)$ et
$\Gamma(X, \mathcal{F}) = \lim A_n$. Comme le foncteur de limite projective n'est pas exact sur la catégorie des systèmes projectifs, il existe un faisceau abélien dont $H^1$ est non nul. Enfin, le lecteur peut vérifier que $H^p(X, j_!\mathbf{Z}_U) = 0$,
$p \geq 1$ si $j : U = U_n \to X$ est l'inclusion. Ainsi, nous voyons que $X$ est un exemple d'espace satisfaisant aux conditions (2), (3) et (4) du lemme \ref{lemma-vanishing-generated-one-section} pour $d = 0$
mais non à sa conclusion.

\end{example}

\begin{lemma}

\label{lemma-subsheaf-irreducible}
Soit $X$ un espace topologique irréductible. Soit $\mathcal{H} \subset \underline{\mathbf{Z}}$ un sous-faisceau abélien du faisceau constant. Il existe alors un ouvert non vide $U \subset X$ tel que $\mathcal{H}|_U = \underline{d\mathbf{Z}}_U$ pour un certain $d \in \mathbf{Z}$.
\end{lemma}

\begin{proof}
Rappelons que $\underline{\mathbf{Z}}(V) = \mathbf{Z}$ pour tout ouvert non vide $V$ de $X$ (voir la preuve du lemme \ref{lemma-irreducible-constant-cohomology-zero}). Si $\mathcal{H} = 0$, le lemme vaut avec $d = 0$. Si $\mathcal{H} \not = 0$, il existe un ouvert non vide $U \subset X$ tel que $\mathcal{H}(U) \not = 0$. Écrivons $\mathcal{H}(U) = n\mathbf{Z}$ pour un certain $n \geq 1$. On a donc $\underline{n\mathbf{Z}}_U
\subset \mathcal{H}|_U \subset
\underline{\mathbf{Z}}_U$. Si la première inclusion est stricte, nous pouvons trouver un $U' \subset U$ non vide et un entier $1 \leq n' < n$ tel que $\underline{n'\mathbf{Z}}_{U'}
\subset \mathcal{H}|_{U'} \subset
\underline{\mathbf{Z}}_{U'}$. Ce processus s'arrête après un nombre fini d'étapes, ce qui prouve le lemme.
\end{proof}

\begin{proposition}
[Grothendieck]
\label{proposition-vanishing-Noetherian}

\begin{reference}
\cite[théorème 3.6.5]{Tohoku}.
\end{reference}
Soit $X$ un espace topologique noethérien. Si $\dim(X) \leq d$, alors $H^p(X, \mathcal{F}) = 0$ pour tous les $p > d$ et tout faisceau abélien $\mathcal{F}$ sur $X$.
\end{proposition}

\begin{proof}
Nous prouvons ce lemme par induction sur $d$. Fixez donc $d$ et supposez que le lemme tient pour tous les espaces topologiques noethériens de dimension $< d$.

\medskip\noindent
Soit $\mathcal{F}$ un faisceau abélien sur $X$. Supposons que $U \subset X$ soit un ouvert et notons $Z \subset X$ son complémentaire fermé. Notons $j : U \to X$ et $i : Z \to X$ les applications d'inclusion. On a alors une suite exacte courte $$
0 \to j_{!}j^*\mathcal{F} \to \mathcal{F} \to i_*i^*\mathcal{F} \to 0
$$ voir Modules, lemme \ref{modules-lemma-canonical-exact-sequence}. Notez que $j_!j^*\mathcal{F}$ est pris en charge sur la fermeture topologique $Z'$ de $U$, c'est-à-dire qu'il est de la forme $i'_*\mathcal{F}'$ pour certains faisceau abélien $\mathcal{F}'$ sur $Z'$, où $i' : Z' \to X$ est l'inclusion.

\medskip\noindent

Nous pouvons utiliser cela pour nous ramener au cas où $X$ est irréductible. En effet, d'après Topologie, lemme \ref{topology-lemma-Noetherian},

$X$ possède un nombre fini de composantes irréductibles. Si $X$ en possède plus d'une, soit $Z \subset X$ une composante irréductible de $X$
et posons $U = X \setminus Z$. D'après ce qui précède et la suite exacte longue de cohomologie, il suffit de prouver l'annulation de
$H^p(X, i_*i^*\mathcal{F})$ et $H^p(X, i'_*\mathcal{F}')$ pour $p > d$. D'après le lemme \ref{lemma-cohomology-and-closed-immersions} il suffit de prouver
$H^p(Z, i^*\mathcal{F})$ et de $H^p(Z', \mathcal{F}')$ pour $p > d$. Comme $Z'$ et $Z$ ont moins de composantes irréductibles, nous nous ramenons bien au cas où $X$ est irréductible.

\medskip\noindent
Si $d = 0$ et si $X$ est irréductible, alors $X$ est le seul ouvert non vide de $X$. Par conséquent, tout faisceau est constant et les groupes de cohomologie supérieurs s'annulent (par exemple d'après le lemme \ref{lemma-irreducible-constant-cohomology-zero}).

\medskip\noindent
Supposons que $X$ soit irréductible de dimension $d > 0$. D'après le lemme \ref{lemma-vanishing-generated-one-section}, nous nous ramenons au cas où $\mathcal{F} = j_!\underline{\mathbf{Z}}_U$ pour un ouvert $U \subset X$. Dans ce cas, considérons la suite exacte courte $$
0 \to j_!(\underline{\mathbf{Z}}_U) \to
\underline{\mathbf{Z}}_X \to i_*\underline{\mathbf{Z}}_Z \to 0
$$ où $Z = X \setminus U$. D'après le lemme \ref{lemma-irreducible-constant-cohomology-zero}, on a l'annulation de $H^p(X, \underline{\mathbf{Z}}_X)$ pour tout $p \geq 1$. Par récurrence, $H^p(X, i_*\underline{\mathbf{Z}}_Z) = H^p(Z, \underline{\mathbf{Z}}_Z) = 0$ pour $p \geq d$. Le résultat découle donc de la suite exacte longue de cohomologie.
\end{proof}





\section{Cohomologie à l'appui dans une partie fermée}

\label{section-cohomology-support}

\noindent
Cette section traite juste du minimum -- la discussion se poursuivra dans la section \ref{section-cohomology-support-bis}.

\medskip\noindent
Soit $X$ un espace topologique, soit $Z \subset X$ une partie fermée et soit $\mathcal{F}$ un faisceau abélien sur $X$. Posons $$
\Gamma_Z(X, \mathcal{F}) =
\{s \in \mathcal{F}(X) \mid \text{Supp}(s) \subset Z\}
$$ le sous-ensemble des sections dont le support est contenu dans $Z$. Le support d'une section est défini dans Modules, définition \ref{modules-definition-support}. Modules, lemme \ref{modules-lemma-support-section-closed} implique que $\Gamma_Z(X, \mathcal{F})$ est un sous-groupe de $\Gamma(X, \mathcal{F})$. Le même lemme garantit que l'association $\mathcal{F} \mapsto \Gamma_Z(X, \mathcal{F})$ est fonctorielle en $\mathcal{F}$. Ce foncteur est exact à gauche, mais n'est pas exact en général.

\medskip\noindent

Comme la catégorie des faisceaux abéliens possède assez d'injectifs (Injectifs, lemme \ref{injectives-lemma-abelian-sheaves-space}), on obtient un foncteur dérivé à droite
$$
R\Gamma_Z(X, -) : D^+(X) \longrightarrow D^+(\textit{Ab})
$$
d'après Catégories dérivées, lemme \ref{derived-lemma-enough-injectives-right-derived}. La valeur de $R\Gamma_Z(X, -)$ sur un objet $K$ se calcule en représentant
$K$ par un complexe borné inférieurement $\mathcal{I}^\bullet$ de faisceaux abéliens injectifs, puis en prenant $\Gamma_Z(X, \mathcal{I}^\bullet)$; voir Catégories dérivées, lemme \ref{derived-lemma-injective-acyclic}. Les groupes de cohomologie d'un faisceau abélien $\mathcal{F}$
à support dans $Z$ sont définis par
$H^q_Z(X, \mathcal{F}) = R^q\Gamma_Z(X, \mathcal{F})$.

\medskip\noindent
Soit $\mathcal{I}$ un faisceau abélien injectif sur $X$, et posons $U = X \setminus Z$. L'application de restriction $\mathcal{I}(X) \to \mathcal{I}(U)$ est surjective (lemme \ref{lemma-injective-restriction-surjective}) et a pour noyau $\Gamma_Z(X, \mathcal{I})$. Il s'ensuit immédiatement que, pour $K \in D^+(X)$, on a un triangle distingué $$
R\Gamma_Z(X, K) \to R\Gamma(X, K) \to R\Gamma(U, K) \to R\Gamma_Z(X, K)[1]
$$ dans $D^+(\textit{Ab})$. Par conséquent, nous obtenons une suite exacte longue de cohomologie $$
\ldots \to H^i_Z(X, K) \to H^i(X, K) \to H^i(U, K) \to
H^{i + 1}_Z(X, K) \to \ldots
$$ pour tout $K$ dans $D^+(X)$.

\medskip\noindent

Pour un faisceau abélien $\mathcal{F}$ sur $X$, nous pouvons considérer le
{\it sous-faisceau des sections à support dans $Z$}, dénoté
$\mathcal{H}_Z(\mathcal{F})$, défini par la règle
$$
\mathcal{H}_Z(\mathcal{F})(U) =
\{s \in \mathcal{F}(U) \mid \text{Supp}(s) \subset U \cap Z\} =
\Gamma_{Z \cap U}(U, \mathcal{F}|_U)
$$
En utilisant l'équivalence de Modules, lemme \ref{modules-lemma-i-star-exact},
nous pouvons considérer $\mathcal{H}_Z(\mathcal{F})$ comme un faisceau abélien sur $Z$; voir Modules, remarque \ref{modules-remark-sections-support-in-closed}. Ainsi, nous obtenons un foncteur
$$
\textit{Ab}(X) \longrightarrow \textit{Ab}(Z),\quad
\mathcal{F} \longmapsto
\mathcal{H}_Z(\mathcal{F})\text{ considéré comme faisceau sur }Z
$$
Ce foncteur est exact à gauche, mais n'est pas exact en général. Comme ci-dessus, nous obtenons un foncteur dérivé à droite
$$
R\mathcal{H}_Z : D^+(X) \longrightarrow D^+(Z)
$$
et nous posons
$\mathcal{H}^q_Z(\mathcal{F}) = R^q\mathcal{H}_Z(\mathcal{F})$, de sorte que
$\mathcal{H}^0_Z(\mathcal{F}) = \mathcal{H}_Z(\mathcal{F})$.

\medskip\noindent
On a $\Gamma_Z(X, \mathcal{F}) = \Gamma(Z, \mathcal{H}_Z(\mathcal{F}))$ pour tout faisceau abélien $\mathcal{F}$. D'après le lemme \ref{lemma-sections-with-support-acyclic}, le foncteur $\mathcal{H}_Z$ transforme les faisceaux abéliens injectifs en faisceaux acycliques pour $\Gamma(Z, -)$. Le lemme de Catégories dérivées \ref{derived-lemma-grothendieck-spectral-sequence} fournit donc une suite spectrale de Grothendieck convergente $$
E_2^{p, q} = H^p(Z, \mathcal{H}^q_Z(K)) \Rightarrow H^{p + q}_Z(X, K)
$$ fonctorielle en $K$ dans $D^+(X)$.

\begin{lemma}

\label{lemma-sections-with-support-acyclic}
Soit $i : Z \to X$ l'inclusion d'une partie fermée, et soit $\mathcal{I}$ un faisceau abélien injectif sur $X$. Alors $\mathcal{H}_Z(\mathcal{I})$ est un faisceau abélien injectif sur $Z$.
\end{lemma}

\begin{proof}
Cela résulte du lemme d'Homologie \ref{homology-lemma-adjoint-preserve-injectives}, car $\mathcal{H}_Z(-)$ est adjoint à droite du foncteur exact $i_*$. Voir Modules, lemmes \ref{modules-lemma-i-star-exact} et \ref{modules-lemma-i-star-right-adjoint}.
\end{proof}



\section{Cohomologie sur les espaces spectraux}

\label{section-spectral}

\noindent
Un résultat clé sur la cohomologie des espaces spectraux est le lemme \ref{lemma-colimit}, qui affirme, en termes informels, que la cohomologie commute aux limites cofiltrantes dans la catégorie des espaces spectraux définie dans Topologie, définition \ref{topology-definition-spectral-space}. On peut l'appliquer pour obtenir les analogues suivants des lemmes \ref{lemma-cohomology-of-closed} et \ref{lemma-proper-base-change}.

\begin{lemma}

\label{lemma-cohomology-of-neighbourhoods-of-closed}
Soit $X$ un espace spectral, soit $\mathcal{F}$ un faisceau abélien sur $X$, soit $E \subset X$ un sous-ensemble quasi-compact et soit $W \subset X$ l'ensemble des points de $X$ qui se spécialisent en un point de $E$.
\begin{enumerate}
\item $H^p(W, \mathcal{F}|_W) = \colim H^p(U, \mathcal{F})$ où la colimite est sur les voisinages quasi-compacts ouverts de $E$,
\item $H^p(W \setminus E, \mathcal{F}|_{W \setminus E}) =
\colim H^p(U \setminus E, \mathcal{F}|_{U \setminus E})$ si $E$ est un sous-ensemble constructible.
\end{enumerate}

\end{lemma}

\begin{proof}
D'après Topologie, lemme \ref{topology-lemma-make-spectral-space}, on a $W = \lim U$, où la limite porte sur les ouverts quasi-compacts contenant $E$. Chaque $U$ est un espace spectral d'après Topologie, lemme \ref{topology-lemma-spectral-sub}. Le lemme \ref{lemma-colimit} donne donc (1). La même démonstration vaut pour (2), en utilisant Topologie, lemme \ref{topology-lemma-make-spectral-space-minus}.
\end{proof}

\begin{lemma}

\label{lemma-proper-base-change-spectral}
Soit $f : X \to Y$ une application spectrale entre espaces spectraux. Soit $y \in Y$, et soit $E \subset Y$ l'ensemble des points qui se spécialisent en $y$. Pour tout faisceau abélien $\mathcal{F}$ sur $X$, on a $(R^pf_*\mathcal{F})_y = H^p(f^{-1}(E), \mathcal{F}|_{f^{-1}(E)})$.
\end{lemma}

\begin{proof}

On a $E = \bigcap V$, où $V$ parcourt les voisinages ouverts quasi-compacts de $y$ dans $Y$. Ainsi $f^{-1}(E) = \bigcap f^{-1}(V)$, donc $f^{-1}(E) = \lim f^{-1}(V)$ comme espace topologique. Puisque $f$ est spectrale, chaque $f^{-1}(V)$ est encore un espace spectral (Topologie, lemme \ref{topology-lemma-spectral-sub}). On en déduit que $f^{-1}(E)$ est un espace spectral et que
$$
H^p(f^{-1}(E), \mathcal{F}|_{f^{-1}(E)}) =
\colim H^p(f^{-1}(V), \mathcal{F})
$$
par lemme \ref{lemma-colimit}. D'autre part, le germe de
$R^pf_*\mathcal{F}$ en $y$ est donné par la colimite du membre de droite.

\end{proof}

\begin{lemma}

\label{lemma-vanishing-for-profinite}
Soit $X$ un espace topologique profini. Alors $H^q(X, \mathcal{F}) = 0$ pour tout $q > 0$ et tout faisceau abélien $\mathcal{F}$.
\end{lemma}

\begin{proof}
Tout recouvrement ouvert de $X$ admet un raffinement qui est une décomposition finie en parties ouvertes disjointes; voir Topologie, lemme \ref{topology-lemma-profinite-refine-open-covering}. Ainsi, si $\mathcal{F} \to \mathcal{G}$ est une surjection de faisceaux abéliens sur $X$, l'application $\mathcal{F}(X) \to \mathcal{G}(X)$ est surjective. Autrement dit, le foncteur des sections globales est exact. Ses foncteurs dérivés supérieurs sont donc nuls; voir Catégories dérivées, lemme \ref{derived-lemma-right-derived-exact-functor}.
\end{proof}

\noindent
Le résultat suivant sur l'annulation cohomologique améliore le résultat de Grothendieck (proposition \ref{proposition-vanishing-Noetherian}) et se trouve dans \cite{Scheiderer}.

\begin{proposition}

\label{proposition-cohomological-dimension-spectral}

\begin{reference}

La partie 1 est le théorème principal de \cite{Scheiderer}.

\end{reference}

Soit $X$ un espace spectral de dimension de Krull $d$. Soit $\mathcal{F}$ un faisceau abélien sur $X$.

\begin{enumerate}

\item  $H^q(X, \mathcal{F}) = 0$ pour $q > d$,
\item  $H^d(X, \mathcal{F}) \to H^d(U, \mathcal{F})$ est surjectif pour chaque ouvert quasi-compact $U \subset X$,
\item  $H^q_Z(X, \mathcal{F}) = 0$ pour $q > d$ et toute partie fermée constructible $Z \subset X$.

\end{enumerate}

\end{proposition}

\begin{proof}
Nous prouvons ce résultat par induction sur $d$.

\medskip\noindent

Si $d = 0$, alors $X$ est un espace profini; voir Topologie, lemme \ref{topology-lemma-characterize-profinite-spectral}. Le point (1) résulte donc du lemme \ref{lemma-vanishing-for-profinite}. Si $U \subset X$ est un ouvert quasi-compact, alors $U$ est aussi fermé, puisqu'il est quasi-compact dans un espace séparé. Ainsi $X = U \amalg (X \setminus U)$, ce qui prouve (2). Pour $Z$ comme dans (3), considérons la suite exacte longue
$$
H^{q - 1}(X, \mathcal{F}) \to
H^{q - 1}(X \setminus Z, \mathcal{F}) \to
H^q_Z(X, \mathcal{F}) \to H^q(X, \mathcal{F})
$$
Comme $X$ et $U = X \setminus Z$ sont profinis (en effet, $U$ est quasi-compact puisque $Z$ est constructible), les points (2) et (1) donnent l'annulation cherchée des groupes de cohomologie à support dans $Z$.

\medskip\noindent

Étape d'induction. Supposons $d \geq 1$ et la proposition vraie pour tout espace spectral de dimension $< d$. Prouvons d'abord (2) pour $X$. Soit $U$ un ouvert quasi-compact et soit $\xi \in H^d(U, \mathcal{F})$. Posons $Z = X \setminus U$, et soit $W \subset X$ l'ensemble des points qui se spécialisent dans $Z$. D'après le lemme \ref{lemma-cohomology-of-neighbourhoods-of-closed}, on a
$$
H^d(W \setminus Z, \mathcal{F}|_{W \setminus Z}) =
\colim_{Z \subset V} H^d(V \setminus Z, \mathcal{F})
$$
où la colimite porte sur les voisinages ouverts quasi-compacts $V$ de $Z$ dans $X$.
D'après Topologie, lemme \ref{topology-lemma-make-spectral-space},
$W \setminus Z$ est un espace spectral. Comme tout point de $W$ se spécialise en un point de $Z$,
$W \setminus Z$ est de dimension de Krull $< d$. Par l'hypothèse d'induction, l'image de $\xi$ dans
$H^d(W \setminus Z, \mathcal{F}|_{W \setminus Z})$ est nulle. Par la formule affichée, il existe un ouvert quasi-compact $Z \subset V \subset X$
tel que $\xi|_{V \setminus Z} = 0$. Comme $V \setminus Z = V \cap U$, la suite de Mayer–Vietoris (lemme \ref{lemma-mayer-vietoris}) pour le recouvrement $X = U \cup V$
fournit un élément $\tilde \xi \in H^d(X, \mathcal{F})$ qui se restreint à $\xi$ sur $U$ et à zéro sur $V$. Le point (2) est donc établi.

\medskip\noindent

Preuve de (1), en supposant (2). Choisissons une résolution injective
$\mathcal{F} \to \mathcal{I}^\bullet$. Posons
$$
\mathcal{G} = \Im(\mathcal{I}^{d - 1} \to \mathcal{I}^d) =
\Ker(\mathcal{I}^d \to \mathcal{I}^{d + 1})
$$
Pour tout ouvert quasi-compact $U \subset X$, nous avons le morphisme suivant de suites exactes :
$$
\xymatrix{
\mathcal{I}^{d - 1}(X) \ar[r] \ar[d] &
\mathcal{G}(X) \ar[r] \ar[d] &
H^d(X, \mathcal{F}) \ar[d] \ar[r] & 0 \\
\mathcal{I}^{d - 1}(U) \ar[r] &
\mathcal{G}(U) \ar[r] &
H^d(U, \mathcal{F}) \ar[r] & 0
}
$$
Le faisceau $\mathcal{I}^{d - 1}$ est flasque, par le lemme \ref{lemma-injective-flasque} et puisque $d \geq 1$. Le point (2) montre que la flèche verticale de droite est surjective. Une chasse au diagramme montre alors que l'application
$\mathcal{G}(X) \to \mathcal{G}(U)$ est surjective. D'après le lemme \ref{lemma-vanishing-ravi}, on en déduit que
$\check{H}^q(\mathcal{U}, \mathcal{G}) = 0$ pour $q > 0$ et tout recouvrement fini $\mathcal{U} : U = U_1 \cup \ldots \cup U_n$
d'un ouvert quasi-compact par des ouverts quasi-compacts. D'après le lemme \ref{lemma-cech-vanish-basis}, on a $H^q(U, \mathcal{G}) = 0$
pour tout $q > 0$ et tout ouvert quasi-compact $U$ de $X$. Le lemme d'acyclicité de Leray (Catégories dérivées, lemme \ref{derived-lemma-leray-acyclicity}) donne alors
$$
H^q(X, \mathcal{F}) =
H^q\left(
\Gamma(X, \mathcal{I}^0) \to \ldots \to
\Gamma(X, \mathcal{I}^{d - 1}) \to \Gamma(X, \mathcal{G})
\right)
$$
En particulier, ce groupe de cohomologie est nul si $q > d$.

\medskip\noindent

Démonstration de (3). Pour $Z$ comme dans (3), considérons la suite exacte longue
$$
H^{q - 1}(X, \mathcal{F}) \to
H^{q - 1}(X \setminus Z, \mathcal{F}) \to
H^q_Z(X, \mathcal{F}) \to H^q(X, \mathcal{F})
$$
Comme $X$ et $U = X \setminus Z$ sont des espaces spectraux (Topologie, lemme \ref{topology-lemma-spectral-sub}) de dimension $\leq d$,
les points (2) et (1) donnent l'annulation cherchée.

\end{proof}













\section{Le complexe de {\v C}
ech modifié}
\label{section-alternating-cech}

\noindent
Cette section compare le complexe {\v C}ech avec le complexe {\v C}ech alternatif et certains complexes connexes.

\medskip\noindent
Soit $X$ un espace topologique et soit $\mathcal{U} : U = \bigcup_{i \in I} U_i$ un recouvrement ouvert. Pour $p \geq 0$, posons $$
\check{\mathcal{C}}_{alt}^p(\mathcal{U}, \mathcal{F})
=
\left\{
\begin{matrix}
s \in  \check{\mathcal{C}}^p(\mathcal{U}, \mathcal{F})
\text{ tel que }
s_{i_0 \ldots i_p} = 0 \text{ si } i_n = i_m \text{ pour certains } n \not = m\\
\text{ et }
s_{i_0\ldots i_n \ldots i_m \ldots i_p}
=
-s_{i_0\ldots i_m \ldots i_n \ldots i_p}
\text{ dans tous les cas.}
\end{matrix}
\right\}
$$ Nous omettons de vérifier que le différentiel $d$ de l'équation (\ref{equation-d-cech}) envoie $\check{\mathcal{C}}^p_{alt}(\mathcal{U}, \mathcal{F})$ dans $\check{\mathcal{C}}^{p + 1}_{alt}(\mathcal{U}, \mathcal{F})$.

\begin{definition}

\label{definition-alternating-cech-complex}
Soit $X$ un espace topologique, soit $\mathcal{U} : U = \bigcup_{i \in I} U_i$ un recouvrement ouvert, et soit $\mathcal{F}$ un préfaisceau abélien sur $X$. Le complexe $\check{\mathcal{C}}_{alt}^\bullet(\mathcal{U}, \mathcal{F})$ est le {\it complexe de {\v C}ech alterné} associé à $\mathcal{F}$ et au recouvrement $\mathcal{U}$.
\end{definition}

\noindent
Il y a donc un morphisme canonique des complexes $$
\check{\mathcal{C}}_{alt}^\bullet(\mathcal{U}, \mathcal{F})
\longrightarrow
\check{\mathcal{C}}^\bullet(\mathcal{U}, \mathcal{F})
$$ à savoir l'inclusion du complexe {\v C}ech alternatif dans le complexe {\v C}ech habituel.

\medskip\noindent
Supposons le recouvrement $\mathcal{U} : U = \bigcup_{i \in I} U_i$ muni d'un ordre total $<$ sur $I$. Posons alors $$
\check{\mathcal{C}}_{ord}^p(\mathcal{U}, \mathcal{F})
=
\prod\nolimits_{(i_0, \ldots, i_p) \in I^{p + 1}, i_0 < \ldots < i_p}
\mathcal{F}(U_{i_0\ldots i_p}).
$$ C'est un groupe abélien. Pour $s \in \check{\mathcal{C}}_{ord}^p(\mathcal{U}, \mathcal{F})$, nous notons $s_{i_0\ldots i_p}$ sa valeur dans $\mathcal{F}(U_{i_0\ldots i_p})$. Nous définissons $$
d : \check{\mathcal{C}}_{ord}^p(\mathcal{U}, \mathcal{F})
\longrightarrow
\check{\mathcal{C}}_{ord}^{p + 1}(\mathcal{U}, \mathcal{F})
$$ par la formule $$
d(s)_{i_0\ldots i_{p + 1}}
=
\sum\nolimits_{j = 0}^{p + 1}
(-1)^j
s_{i_0\ldots \hat i_j \ldots i_{p + 1}}|_{U_{i_0\ldots i_{p + 1}}}
$$ pour n'importe quelle $i_0 < \ldots < i_{p + 1}$. Notez que cette formule est identique à l'équation (\ref{equation-d-cech}). Il est simple de voir que $d \circ d = 0$. En d'autres termes, $\check{\mathcal{C}}_{ord}^\bullet(\mathcal{U}, \mathcal{F})$ est un complexe.

\begin{definition}

\label{definition-ordered-cech-complex}
Soit $X$ un espace topologique, soit $\mathcal{U} : U = \bigcup_{i \in I} U_i$ un recouvrement ouvert, et supposons $I$ muni d'un ordre total. Soit $\mathcal{F}$ un préfaisceau abélien sur $X$. Le complexe $\check{\mathcal{C}}_{ord}^\bullet(\mathcal{U}, \mathcal{F})$ est le {\it complexe de {\v C}ech ordonné} associé à $\mathcal{F}$, au recouvrement $\mathcal{U}$ et à l'ordre total donné sur $I$.
\end{definition}

\noindent
Ce complexe est parfois appelé le complexe alterné de {\v C}ech. La raison est qu'il existe une application de comparaison évidente entre le complexe ordonné de {\v C}ech et le complexe alterné de {\v C}ech. Plus précisément, considérons l'application $$
c :
\check{\mathcal{C}}_{ord}^\bullet(\mathcal{U}, \mathcal{F})
\longrightarrow
\check{\mathcal{C}}^\bullet(\mathcal{U}, \mathcal{F})
$$ donnée par la règle $$
c(s)_{i_0\ldots i_p} =
\left\{
\begin{matrix}
0 &
\text{si} &
i_n = i_m \text{ pour certains } n \not = m\\
\text{sgn}(\sigma) s_{i_{\sigma(0)}\ldots i_{\sigma(p)}} &
\text{si} &
i_{\sigma(0)} < i_{\sigma(1)} < \ldots < i_{\sigma(p)}
\end{matrix}
\right.
$$ Ici, $\sigma$ désigne une permutation de $\{0, \ldots, p\}$ et $\text{sgn}(\sigma)$ désigne son signe. Les complexes alterné et ordonné de {\v C}ech sont souvent identifiés dans la littérature au moyen de l'application $c$. Nous avons en effet le lemme élémentaire suivant.

\begin{lemma}

\label{lemma-ordered-alternating}
Soit $X$ un espace topologique, soit $\mathcal{U} : U = \bigcup_{i \in I} U_i$ un recouvrement ouvert, et supposons $I$ muni d'un ordre total. L'application $c$ est un morphisme de complexes. Elle induit même un isomorphisme $$
c : \check{\mathcal{C}}_{ord}^\bullet(\mathcal{U}, \mathcal{F})
\to \check{\mathcal{C}}_{alt}^\bullet(\mathcal{U}, \mathcal{F})
$$ de complexes.
\end{lemma}

\begin{proof}
Démonstration omise.
\end{proof}

\noindent
Il y a aussi une application $$
\pi :
\check{\mathcal{C}}^\bullet(\mathcal{U}, \mathcal{F})
\longrightarrow
\check{\mathcal{C}}_{ord}^\bullet(\mathcal{U}, \mathcal{F})
$$ qui est décrite par la règle $$
\pi(s)_{i_0\ldots i_p} = s_{i_0\ldots i_p}
$$ chaque fois que $i_0 < i_1 < \ldots < i_p$.

\begin{lemma}

\label{lemma-project-to-ordered}
Soit $X$ un espace topologique, soit $\mathcal{U} : U = \bigcup_{i \in I} U_i$ un recouvrement ouvert, et supposons $I$ muni d'un ordre total. La projection $\pi : \check{\mathcal{C}}^\bullet(\mathcal{U}, \mathcal{F})
\to \check{\mathcal{C}}_{ord}^\bullet(\mathcal{U}, \mathcal{F})$ est un morphisme de complexes. Il induit un isomorphisme $$
\pi : \check{\mathcal{C}}_{alt}^\bullet(\mathcal{U}, \mathcal{F})
\to \check{\mathcal{C}}_{ord}^\bullet(\mathcal{U}, \mathcal{F})
$$ de complexes qui est un inverse gauche au morphisme $c$.
\end{lemma}

\begin{proof}
Démonstration omise.
\end{proof}

\begin{remark}

\label{remark-compared-ordered-complexes}
Cela signifie que si nous avons deux commandes totales $<_1$ et $<_2$ sur le jeu d'index $I$, alors nous obtenons un isomorphisme des complexes $\tau = \pi_2 \circ c_1 :
\check{\mathcal{C}}_{ord\text{-}1}(\mathcal{U}, \mathcal{F}) \to
\check{\mathcal{C}}_{ord\text{-}2}(\mathcal{U}, \mathcal{F})$. Il est clair que $$
\tau(s)_{i_0 \ldots i_p} =
\text{sign}(\sigma) s_{i_{\sigma(0)} \ldots i_{\sigma(p)}}
$$ où $i_0 <_1 i_1 <_1 \ldots <_1 i_p$ et $i_{\sigma(0)} <_2 i_{\sigma(1)} <_2 \ldots <_2 i_{\sigma(p)}$. C'est en ce sens que le complexe de {\v C}ech ordonné ne dépend pas de l'ordre total choisi.
\end{remark}

\begin{lemma}

\label{lemma-alternating-usual}
Soit $X$ un espace topologique, soit $\mathcal{U} : U = \bigcup_{i \in I} U_i$ un recouvrement ouvert, et supposons $I$ muni d'un ordre total. L'application $c \circ \pi$ est homotope à l'identité de $\check{\mathcal{C}}^\bullet(\mathcal{U}, \mathcal{F})$. En particulier, l'inclusion $\check{\mathcal{C}}_{alt}^\bullet(\mathcal{U}, \mathcal{F}) \to
\check{\mathcal{C}}^\bullet(\mathcal{U}, \mathcal{F})$ est une équivalence d'homotopie.
\end{lemma}

\begin{proof}
Pour tout multi-index $(i_0, \ldots, i_p) \in I^{p + 1}$ il existe une permutation unique $\sigma : \{0, \ldots, p\} \to \{0, \ldots, p\}$ telle que $$
i_{\sigma(0)} \leq i_{\sigma(1)} \leq \ldots \leq i_{\sigma(p)}
\quad
\text{et}
\quad
\sigma(j) < \sigma(j + 1)
\quad
\text{si}
\quad
i_{\sigma(j)} = i_{\sigma(j + 1)}.
$$ Nous dénotons cette permutation $\sigma = \sigma^{i_0 \ldots i_p}$.

\medskip\noindent
Pour toute permutation $\sigma : \{0, \ldots, p\} \to \{0, \ldots, p\}$ et toute $a$, $0 \leq a \leq p$ nous dénotons $\sigma_a$ la permutation unique de $\{0, \ldots, p\}$ telle que $\sigma_a(j) = \sigma(j)$ pour $0 \leq j < a$ et telle que $\sigma_a(a) < \sigma_a(a + 1) < \ldots < \sigma_a(p)$. Donc si $p = 3$ et $\sigma$, $\tau$ sont donnés par $$
\begin{matrix}
\text{id} & 0 & 1 & 2 & 3 \\
\sigma & 3 & 2 & 1 & 0
\end{matrix}
\quad \text{et} \quad
\begin{matrix}
\text{id} & 0 & 1 & 2 & 3 \\
\tau & 3 & 0 & 2 & 1
\end{matrix}
$$ alors nous avons $$
\begin{matrix}
\text{id} & 0 & 1 & 2 & 3 \\
\sigma_0 & 0 & 1 & 2 & 3 \\
\sigma_1 & 3 & 0 & 1 & 2 \\
\sigma_2 & 3 & 2 & 0 & 1 \\
\sigma_3 & 3 & 2 & 1 & 0 \\
\end{matrix}
\quad \text{et} \quad
\begin{matrix}
\text{id} & 0 & 1 & 2 & 3 \\
\tau_0 & 0 & 1 & 2 & 3 \\
\tau_1 & 3 & 0 & 1 & 2 \\
\tau_2 & 3 & 0 & 1 & 2 \\
\tau_3 & 3 & 0 & 2 & 1 \\
\end{matrix}
$$ Il est clair que toujours $\sigma_0 = \text{id}$ et $\sigma_p = \sigma$.

\medskip\noindent

Après avoir introduit cette notation, nous définissons pour
$s \in \check{\mathcal{C}}^{p + 1}(\mathcal{U}, \mathcal{F})$
l'élément $h(s) \in \check{\mathcal{C}}^p(\mathcal{U}, \mathcal{F})$
comme l'élément dont les composantes sont

\begin{equation}
\label{equation-first-homotopy}
h(s)_{i_0\ldots i_p} =
\sum\nolimits_{0 \leq a \leq p}
(-1)^a \text{sign}(\sigma_a)
s_{i_{\sigma(0)} \ldots i_{\sigma(a)} i_{\sigma_a(a)} \ldots i_{\sigma_a(p)}}
\end{equation}

où $\sigma = \sigma^{i_0 \ldots i_p}$. L'indice
$i_{\sigma(a)}$ survient deux fois dans
$i_{\sigma(0)} \ldots i_{\sigma(a)} i_{\sigma_a(a)} \ldots i_{\sigma_a(p)}$
une fois dans le premier groupe de $a + 1$ indices et une fois dans le second groupe de $p - a + 1$ indices, puisque $\sigma_a(j) = \sigma(a)$ pour un certain
$j \geq a$ par définition de $\sigma_a$. La somme a donc un sens, car chacun des éléments
$s_{i_{\sigma(0)} \ldots i_{\sigma(a)} i_{\sigma_a(a)} \ldots i_{\sigma_a(p)}}$
est défini sur l'ouvert $U_{i_0 \ldots i_p}$. Notons également que, pour $a = 0$, nous obtenons $s_{i_0 \ldots i_p}$ et que, pour $a = p$, nous obtenons
$(-1)^p \text{sign}(\sigma) s_{i_{\sigma(0)} \ldots i_{\sigma(p)}}$.

\medskip\noindent

Nous affirmons que :
$$
(dh + hd)(s)_{i_0 \ldots i_p} =
s_{i_0 \ldots i_p} -
\text{sign}(\sigma) s_{i_{\sigma(0)} \ldots i_{\sigma(p)}}
$$
où $\sigma = \sigma^{i_0 \ldots i_p}$. Nous omettons la vérification de cette revendication. (Il y a un script PARI/gp appelé first-homotopy.gp dans les scripts de sous-répertoire stacks-project qui peuvent être utilisés pour vérifier de nombreuses instances finies de cette revendication. Nous avons écrit ce script pour nous assurer que les signes sont corrects.)
$$
\kappa :
\check{\mathcal{C}}^\bullet(\mathcal{U}, \mathcal{F})
\longrightarrow
\check{\mathcal{C}}^\bullet(\mathcal{U}, \mathcal{F})
$$
pour l'opérateur donné par la règle
$$
\kappa(s)_{i_0 \ldots i_p} =
\text{sign}(\sigma^{i_0 \ldots i_p}) s_{i_{\sigma(0)} \ldots i_{\sigma(p)}}.
$$
L'allégation ci-dessus implique que: $\kappa$ est un morphisme de complexes et que
$\kappa$ est homotope à l'identité du complexe de {\v C}ech. Cela n'implique pas immédiatement le lemme, car l'image de l'opérateur $\kappa$ n'est pas le sous-complexe alterné. Plus précisément, l'image de $\kappa$ est le complexe « semi-alterné »
$\check{\mathcal{C}}_{semi\text{-}alt}^p(\mathcal{U}, \mathcal{F})$
où $s$ est une $p$-cochaîne de ce complexe si et seulement si
$$
s_{i_0 \ldots i_p} = \text{sign}(\sigma) s_{i_{\sigma(0)} \ldots i_{\sigma(p)}}
$$
pour tout $(i_0, \ldots, i_p) \in I^{p + 1}$ avec
$\sigma = \sigma^{i_0 \ldots i_p}$. Nous introduisons une autre variante du complexe de {\v C}ech, à savoir le complexe de {\v C}ech semi-ordonné défini par
$$
\check{\mathcal{C}}_{semi\text{-}ord}^p(\mathcal{U}, \mathcal{F})
=
\prod\nolimits_{i_0 \leq i_1 \leq \ldots \leq i_p}
\mathcal{F}(U_{i_0 \ldots i_p})
$$
Il est facile de voir que l'équation (\ref{equation-d-cech}) définit encore une différentielle, donc un complexe. Il est également clair, comme dans le lemme \ref{lemma-project-to-ordered}, que l'application de projection
$$
\check{\mathcal{C}}_{semi\text{-}alt}^\bullet(\mathcal{U}, \mathcal{F})
\longrightarrow
\check{\mathcal{C}}_{semi\text{-}ord}^\bullet(\mathcal{U}, \mathcal{F})
$$
est un isomorphisme des complexes.

\medskip\noindent

Par conséquent, le lemme suit si nous pouvons montrer que l'application d'inclusion évidente
$$
\check{\mathcal{C}}_{ord}^p(\mathcal{U}, \mathcal{F})
\longrightarrow
\check{\mathcal{C}}_{semi\text{-}ord}^p(\mathcal{U}, \mathcal{F})
$$
est une équivalence d'homotopie. Pour cela nous utilisons l'homotopie

\begin{equation}
\label{equation-second-homotopy}
h(s)_{i_0 \ldots i_p} =
\left\{
\begin{matrix}
0 & \text{si} & i_0 < i_1 < \ldots < i_p \\
(-1)^a s_{i_0 \ldots i_{a - 1} i_a i_a i_{a + 1} \ldots i_p}
& \text{si} & i_0 < i_1 < \ldots < i_{a - 1} < i_a = i_{a + 1}
\end{matrix}
\right.
\end{equation}

Nous affirmons que :
$$
(dh + hd)(s)_{i_0 \ldots i_p} =
\left\{
\begin{matrix}
0 & \text{si} & i_0 < i_1 < \ldots < i_p \\
s_{i_0 \ldots i_p}
& \text{sinon} &
\end{matrix}
\right.
$$
Nous omettons la vérification. (Il y a un script PARI/gp appelé second-homotopy.gp dans les scripts de sous-répertoire stacks-project qui peuvent être utilisés pour vérifier de nombreuses instances de cette revendication. Nous avons écrit ce script pour nous assurer que les signes sont corrects.) La revendication montre clairement que la composition
$$
\check{\mathcal{C}}_{semi\text{-}ord}^\bullet(\mathcal{U}, \mathcal{F})
\longrightarrow
\check{\mathcal{C}}_{ord}^\bullet(\mathcal{U}, \mathcal{F})
\longrightarrow
\check{\mathcal{C}}_{semi\text{-}ord}^\bullet(\mathcal{U}, \mathcal{F})
$$
de la projection avec l'inclusion naturelle est homotopique à l'application d'identité comme souhaité.

\end{proof}

\begin{lemma}

\label{lemma-alternating-cech-trivial}

Soit $X$ un espace topologique, soit $\mathcal{F}$ un préfaisceau abélien sur $X$ et soit $\mathcal{U} : U = \bigcup_{i \in I} U_i$ un recouvrement ouvert. Si
$U_i = U$ pour un certain $i \in I$, alors le complexe de {\v C}ech alterné augmenté
$$
\mathcal{F}(U) \to \check{\mathcal{C}}_{alt}^\bullet(\mathcal{U}, \mathcal{F})
$$
obtenu en plaçant $\mathcal{F}(U)$ en degré $-1$, avec pour différentielle l'application canonique de $\mathcal{F}(U)$ dans
$\check{\mathcal{C}}^0(\mathcal{U}, \mathcal{F})$
est homotopiquement équivalent à $0$. De même, pour tout ordre total sur $I$,
le complexe de {\v C}ech ordonné augmenté
$$
\mathcal{F}(U) \to
\check{\mathcal{C}}_{ord}^\bullet(\mathcal{U}, \mathcal{F})
$$
est homotopiquement équivalent à $0$.

\end{lemma}

\begin{proof}
[Première preuve] Combiner les lemmes \ref{lemma-cech-trivial} et \ref{lemma-alternating-usual}.
\end{proof}

\begin{proof}
[Deuxième preuve] Puisque les complexes de {\v C}ech alterné et ordonné sont isomorphes, il suffit de traiter le complexe ordonné. Nous emploierons la notation usuelle : une cochaîne $s$ de degré $p$
du complexe de {\v C}ech ordonné augmenté est de la forme
$s = (s_{i_0 \ldots i_p})$ où $s_{i_0 \ldots i_p}$ est dans
$\mathcal{F}(U_{i_0 \ldots i_p})$ et $i_0 < \ldots < i_p$. Avec cette notation nous avons
$$
d(x)_{i_0 \ldots i_{p + 1}} =
\sum\nolimits_j (-1)^j x_{i_0 \ldots \hat i_j \ldots i_p}
$$
Fixons un indice $i \in I$ tel que $U = U_i$. Nous utilisons comme homotopie les applications
$$
h : \text{cochaînes de degré }p + 1 \to \text{cochaînes de degré }p
$$
données par la règle
$$
h(s)_{i_0 \ldots i_p} = 0 \text{ si } i \in \{i_0, \ldots, i_p\}
\text{ et }
h(s)_{i_0 \ldots i_p} = 
(-1)^j s_{i_0 \ldots i_j i i_{j + 1} \ldots i_p} \text{ sinon}
$$
Ici $j$ est l'unique indice tel que $i_j < i < i_{j + 1}$ dans le second cas; puisque $U = U_i$, nous avons aussi l'égalité
$$
\mathcal{F}(U_{i_0 \ldots i_p}) =
\mathcal{F}(U_{i_0 \ldots i_j i i_{j + 1} \ldots i_p})
$$
qui permet de considérer
$(-1)^j s_{i_0 \ldots i_j i i_{j + 1} \ldots i_p}$
comme un élément de $\mathcal{F}(U_{i_0 \ldots i_p})$. Nous allons montrer par un calcul que $d h + h d$ est l'opposé de l'identité, ce qui achèvera la preuve. Fixons une cochaîne $s$ de degré $p$ et des éléments
$i_0 < \ldots < i_p$ de $I$.

\medskip\noindent
Cas I : $i \in \{i_0, \ldots, i_p\}$. Écrivons $i = i_t$. Alors $h(d(s))_{i_0 \ldots i_p} = 0$. D'autre part, on a $$
d(h(s))_{i_0 \ldots i_p} =
\sum (-1)^j h(s)_{i_0 \ldots \hat i_j \ldots i_p} =
(-1)^t h(s)_{i_0 \ldots \hat i \ldots i_p} =
(-1)^t (-1)^{t - 1} s_{i_0 \ldots i_p}
$$ Ainsi $(dh + hd)(s)_{i_0 \ldots i_p} = -s_{i_0 \ldots i_p}$ comme désiré.

\medskip\noindent

Cas II : $i \not \in \{i_0, \ldots, i_p\}$. Soit $j$ tel que
$i_j < i < i_{j + 1}$. Alors nous voyons que

\begin{align*}
h(d(s))_{i_0 \ldots i_p}
& =
(-1)^j d(s)_{i_0 \ldots i_j i i_{j + 1} \ldots i_p} \\
& =
\sum\nolimits_{j' \leq j} (-1)^{j + j'}
s_{i_0 \ldots \hat i_{j'} \ldots i_j i i_{j + 1} \ldots i_p} -
s_{i_0 \ldots i_p} \\
&
+ \sum\nolimits_{j' > j} (-1)^{j + j' + 1}
s_{i_0 \ldots i_j i i_{j + 1} \ldots \hat i_{j'} \ldots i_p}
\end{align*}

D'autre part, nous avons

\begin{align*}
d(h(s))_{i_0 \ldots i_p}
& =
\sum\nolimits_{j'} (-1)^{j'} h(s)_{i_0 \ldots \hat i_{j'} \ldots i_p} \\
& =
\sum\nolimits_{j' \leq j} (-1)^{j' + j - 1}
s_{i_0 \ldots \hat i_{j'} \ldots i_j i i_{j + 1} \ldots i_p} \\
& +
\sum\nolimits_{j' > j} (-1)^{j' + j}
s_{i_0 \ldots i_j i i_{j + 1} \ldots \hat i_{j'} \ldots i_p}
\end{align*}

En les ajoutant, nous obtenons
$(dh + hd)(s)_{i_0 \ldots i_p} = - s_{i_0 \ldots i_p}$
comme désiré.

\end{proof}





\section{Un autre point de vue sur le complexe de {\v C}
ech}
\label{section-locally-finite-cech}

\noindent
Dans cette section, nous discutons d'une autre façon d'établir la relation entre le complexe {\v C}ech et la cohomologie.

\begin{lemma}

\label{lemma-covering-resolution}

Soit $X$ un espace annelé. Soit $\mathcal{U} : X = \bigcup_{i \in I} U_i$
un recouvrement ouvert de $X$. Soit $\mathcal{F}$ un $\mathcal{O}_X$-module. Notons $\mathcal{F}_{i_0 \ldots i_p}$ la restriction de
$\mathcal{F}$ à $U_{i_0 \ldots i_p}$. Il existe un complexe
${\mathfrak C}^\bullet(\mathcal{U}, \mathcal{F})$
de $\mathcal{O}_X$-modules avec
$$
{\mathfrak C}^p(\mathcal{U}, \mathcal{F}) =
\prod\nolimits_{i_0 \ldots i_p}
(j_{i_0 \ldots i_p})_* \mathcal{F}_{i_0 \ldots i_p}
$$
et différentiel
$d : {\mathfrak C}^p(\mathcal{U}, \mathcal{F})
\to {\mathfrak C}^{p + 1}(\mathcal{U}, \mathcal{F})$
comme dans l'équation (\ref{equation-d-cech}). En outre, il existe une application canonique
$$
\mathcal{F} \to {\mathfrak C}^\bullet(\mathcal{U}, \mathcal{F})
$$
qui est un quasi-isomorphisme, c'est-à-dire,
${\mathfrak C}^\bullet(\mathcal{U}, \mathcal{F})$
est une résolution de $\mathcal{F}$.

\end{lemma}

\begin{proof}

Nous vérifions
$$
0 \to \mathcal{F} \to \mathfrak{C}^0(\mathcal{U}, \mathcal{F}) \to
\mathfrak{C}^1(\mathcal{U}, \mathcal{F}) \to  \ldots
$$
est exact sur les germes. $x \in X$ et choisir $i_{\text{fix}} \in I$
de telle manière que $x \in U_{i_{\text{fix}}}$. Ensuite, définir 
$$
h : \mathfrak{C}^p(\mathcal{U}, \mathcal{F})_x
\to \mathfrak{C}^{p - 1}(\mathcal{U}, \mathcal{F})_x
$$
comme suit: $s \in \mathfrak{C}^p(\mathcal{U}, \mathcal{F})_x$, prendre un représentant
$$
\widetilde{s} \in
\mathfrak{C}^p(\mathcal{U}, \mathcal{F})(V) =
\prod\nolimits_{i_0 \ldots i_p}
\mathcal{F}(V \cap U_{i_0} \cap \ldots \cap U_{i_p})
$$
défini sur un voisinage $V$ de $x$, et posons
$$
h(s)_{i_0 \ldots i_{p - 1}} =
\widetilde{s}_{i_{\text{fix}} i_0 \ldots i_{p - 1}, x}.
$$
Par la même formule (pour $p = 0$) nous obtenons une application
$\mathfrak{C}^{0}(\mathcal{U},\mathcal{F})_x \to \mathcal{F}_x$. Nous calculons formellement comme suit :

\begin{align*}
(dh + hd)(s)_{i_0 \ldots i_p}
& =
\sum\nolimits_{j = 0}^p
(-1)^j
h(s)_{i_0 \ldots \hat i_j \ldots i_p}
+
d(s)_{i_{\text{fix}} i_0 \ldots i_p}\\
& =
\sum\nolimits_{j = 0}^p
(-1)^j
s_{i_{\text{fix}} i_0 \ldots \hat i_j \ldots i_p}
+
s_{i_0 \ldots i_p}
+
\sum\nolimits_{j = 0}^p
(-1)^{j + 1}
s_{i_{\text{fix}} i_0 \ldots \hat i_j \ldots i_p} \\
& =
s_{i_0 \ldots i_p}
\end{align*}

Ceci montre $h$ est une homotopie de l'application d'identité du complexe étendu
$$
0 \to \mathcal{F}_x \to \mathfrak{C}^0(\mathcal{U}, \mathcal{F})_x
\to \mathfrak{C}^1(\mathcal{U}, \mathcal{F})_x \to \ldots
$$
à zéro et nous concluons.

\end{proof}

\noindent

Ce lemme permet de redémontrer facilement la suite spectrale reliant la cohomologie de {\v C}ech à la cohomologie du lemme \ref{lemma-cech-spectral-sequence}. Soient $X$, $\mathcal{U}$ et $\mathcal{F}$ comme dans le lemme \ref{lemma-covering-resolution},
et soit $\mathcal{F} \to \mathcal{I}^\bullet$ une résolution injective. Considérons alors le bicomplexe
$$
A^{\bullet, \bullet} =
\Gamma(X, {\mathfrak C}^\bullet(\mathcal{U}, \mathcal{I}^\bullet)).
$$
Par construction nous avons
$$
A^{p, q} = \prod\nolimits_{i_0 \ldots i_p} \mathcal{I}^q(U_{i_0 \ldots i_p})
$$
Considérons les deux suites spectrales associées à ce bicomplexe dans Homologie, section \ref{homology-section-double-complex}; voir notamment Homologie, lemme \ref{homology-lemma-ss-double-complex}. Pour la suite spectrale $({}'E_r, {}'d_r)_{r \geq 0}$, nous obtenons
${}'E_2^{p, q} = \check{H}^p(\mathcal{U}, \underline{H}^q(\mathcal{F}))$
car la formation des produits est exacte (Homologie, lemme \ref{homology-lemma-product-abelian-groups-exact}). Pour la suite spectrale $({}''E_r, {}''d_r)_{r \geq 0}$, nous obtenons
${}''E_2^{p, q} = 0$ si $p > 0$ et ${}''E_2^{0, q} = H^q(X, \mathcal{F})$. En effet, pour $q$ fixé, le complexe de faisceaux
${\mathfrak C}^\bullet(\mathcal{U}, \mathcal{I}^q)$
est une résolution (lemme \ref{lemma-covering-resolution}) du faisceau injectif $\mathcal{I}^q$
par des faisceaux injectifs (d'après les lemmes \ref{lemma-cohomology-of-open} et
\ref{lemma-pushforward-injective-flat}
et Homologie, lemme \ref{homology-lemma-product-injectives}). Par conséquent, la cohomologie de
$\Gamma(X, {\mathfrak C}^\bullet(\mathcal{U}, \mathcal{I}^q))$
est nulle en degrés positifs et égale à $\Gamma(X, \mathcal{I}^q)$
en degré $0$. En prenant la cohomologie pour la différentielle suivante, on obtient l'assertion sur la suite spectrale $({}''E_r, {}''d_r)_{r \geq 0}$. Le résultat s'ensuit, puisque les deux suites spectrales convergent vers la cohomologie du complexe total associé à $A^{\bullet, \bullet}$.

\begin{definition}

\label{definition-covering-locally-finite}

Soit $X$ un espace topologique. Un recouvrement ouvert $X = \bigcup_{i \in I} U_i$ est dit
{\it localement fini} si, pour tout $x \in X$, il existe un voisinage ouvert
$W$ de $x$ tel que $\{i \in I \mid W \cap U_i \not = \emptyset\}$ soit fini.

\end{definition}

\begin{remark}

\label{remark-locally-finite-sections}

Soit $X = \bigcup_{i \in I} U_i$ un recouvrement ouvert localement fini. Notons $j_i : U_i \to X$ l'application d'inclusion. Supposons que, pour chaque $i$,
on se donne un faisceau abélien $\mathcal{F}_i$ sur $U_i$. Considérons le faisceau abélien $\mathcal{G} = \bigoplus_{i \in I} (j_i)_*\mathcal{F}_i$. Alors, pour tout ouvert $V \subset X$, on a
$$
\Gamma(V, \mathcal{G}) = \prod\nolimits_{i \in I} \mathcal{F}_i(V \cap U_i).
$$
En d'autres termes, nous avons
$$
\bigoplus\nolimits_{i \in I} (j_i)_*\mathcal{F}_i =
\prod\nolimits_{i \in I} (j_i)_*\mathcal{F}_i
$$
Cela peut sembler étrange jusqu'à ce que l'on se rappelle que la somme directe d'une famille de faisceaux est le faisceau associé au préfaisceau somme directe. Voir la discussion dans Modules, section \ref{modules-section-kernels}. Nous concluons donc que, dans ce cas, les termes du complexe du lemme \ref{lemma-covering-resolution} sont
$$
{\mathfrak C}^p(\mathcal{U}, \mathcal{F}) =
\bigoplus\nolimits_{i_0 \ldots i_p}
(j_{i_0 \ldots i_p})_* \mathcal{F}_{i_0 \ldots i_p}
$$
Ce qui est parfois utile.

\end{remark}











\section{Cohomologie de {\v C}
ech des complexes}
\label{section-cech-cohomology-of-complexes}

\noindent
En général, pour des faisceaux de groupes abéliens ${\mathcal F}$ et ${\mathcal G}$ sur $X$, il existe une application produit cup $$
H^i(X, {\mathcal F}) \times H^j(X, {\mathcal G})
\longrightarrow
H^{i + j}(X, {\mathcal F} \otimes_{\mathbf Z} {\mathcal G}).
$$ Dans cette section, nous la définissons au moyen de cocycles de {\v C}ech, par une formule explicite pour le produit cup. Si l'on craint que la cohomologie ne coïncide pas avec la cohomologie de {\v C}ech, on peut utiliser des hyperrecouvrements tout en conservant la notation des cocycles. Cette méthode a en outre l'avantage de définir le produit cup en hypercohomologie sur tout topos (insérer ici une référence ultérieure).

\medskip\noindent

Soit ${\mathcal F}^\bullet$ un complexe borné inférieurement de préfaisceaux de groupes abéliens sur $X$. On peut souvent calculer $H^n(X, {\mathcal F}^\bullet)$
au moyen de cocycles de {\v C}ech. Soit donc
${\mathcal U} : X = \bigcup_{i \in I} U_i$
un recouvrement ouvert de $X$. Comme le complexe de {\v C}ech
$\check{\mathcal{C}}^\bullet(\mathcal{U}, \mathcal{F})$
(définition \ref{definition-cech-complex}) est fonctoriel en le préfaisceau $\mathcal{F}$, on obtient un bicomplexe
$\check{\mathcal{C}}^\bullet(\mathcal{U}, \mathcal{F}^\bullet)$. Le complexe associé total à
$\check{\mathcal{C}}^\bullet({\mathcal U}, {\mathcal F}^\bullet)$
est le complexe dont le terme de degré $n$ est
$$
\text{Tot}^n(\check{\mathcal{C}}^\bullet({\mathcal U}, {\mathcal F}^\bullet))
=
\bigoplus\nolimits_{p + q = n}
\prod\nolimits_{i_0\ldots i_p} {\mathcal F}^q(U_{i_0\ldots i_p})
$$
Voir Homologie, définition \ref{homology-definition-associated-simple-complex}. Un élément typique de $\text{Tot}^n$ sera noté
$\alpha = \{\alpha_{i_0\ldots i_p}\}$ où
$\alpha_{i_0 \ldots i_p} \in \mathcal{F}^q(U_{i_0\ldots i_p})$. Autrement dit, le degré en $\mathcal{F}$ de $\alpha_{i_0\ldots i_p}$ est
$q = n - p$. Cette notation suppose que l'on garde à l'esprit le degré de $\alpha$;
nous indiquons cette situation par la formule
$\deg_{\mathcal F}(\alpha_{i_0\ldots i_p}) = q$. Selon nos conventions dans Homologie, définition \ref{homology-definition-associated-simple-complex}
la différentielle d'un élément $\alpha$ de degré $n$ est donnée par
$$
d(\alpha)_{i_0\ldots i_{p + 1}}
=
\sum\nolimits_{j = 0}^{p + 1}
(-1)^j \alpha_{i_0 \ldots \hat i_j \ldots i_{p + 1}} + 
(-1)^{p + 1}d_{{\mathcal F}}(\alpha_{i_0 \ldots i_{p + 1}})
$$
où $d_\mathcal{F}$ désigne le différentiel sur le complexe
$\mathcal{F}^\bullet$. L'expression $\alpha_{i_0 \ldots \hat i_j \ldots i_{p + 1}}$ désigne la restriction de $\alpha_{i_0 \ldots \hat i_j \ldots i_{p + 1}}
\in {\mathcal F}(U_{i_0\ldots\hat i_j\ldots i_{p + 1}})$ à
$U_{i_0 \ldots i_{p + 1}}$.

\medskip\noindent

La construction de
$\text{Tot}(\check{\mathcal{C}}^\bullet({\mathcal U}, {\mathcal F}^\bullet))$
est fonctorielle en ${\mathcal F}^\bullet$. Il existe également une transformation fonctorielle

\begin{equation}
\label{equation-global-sections-to-cech}
\Gamma(X, {\mathcal F}^\bullet)
\longrightarrow
\text{Tot}(\check{\mathcal{C}}^\bullet({\mathcal U}, {\mathcal F}^\bullet))
\end{equation}

de complexes définie par la règle suivante :
$s\in \Gamma(X, {\mathcal F}^n)$
est envoyé sur l'élément $\alpha = \{\alpha_{i_0\ldots i_p}\}$
avec $\alpha_{i_0} = s|_{U_{i_0}}$ et $\alpha_{i_0\ldots i_p} = 0$
pour $p > 0$.

\medskip\noindent

Raffinements. Soit ${\mathcal V} = \{ V_j \}_{j\in J}$ un raffinement de ${\mathcal U}$. Cela signifie qu'il existe une application $t: J \to I$
telle que $V_j \subset U_{t(j)}$ pour tout $j\in J$. Elle induit une transformation fonctorielle

\begin{equation}
\label{equation-transformation}
T_t :
\text{Tot}(\check{\mathcal{C}}^\bullet({\mathcal U}, {\mathcal F}^\bullet))
\longrightarrow
\text{Tot}(\check{\mathcal{C}}^\bullet({\mathcal V}, {\mathcal F}^\bullet)).
\end{equation}

défini par la règle
$$
T_t(\alpha)_{j_0\ldots j_p}
=
\alpha_{t(j_0)\ldots t(j_p)}|_{V_{j_0\ldots j_p}}.
$$
Pour deux applications $t, t' : J \to I$ comme ci-dessus, les applications
$T_t$ et $T_{t'}$ ainsi construites sont homotopes. Une homotopie est donnée par
$$
h(\alpha)_{j_0\ldots j_p}
=
\sum\nolimits_{a = 0}^{p}
(-1)^a
\alpha_{t(j_0)\ldots t(j_a) t'(j_a) \ldots t'(j_p)}
$$
pour un élément $\alpha$ de degré $n$. Cela résulte du calcul suivant, où
$\alpha$ est encore un élément de degré $n$ (ainsi $d(\alpha)$
est de degré $n + 1$ et $h(\alpha)$ de degré $n - 1$) :

\begin{align*}
(
d(h(\alpha)) + h(d(\alpha))
)_{j_0\ldots j_p}
= &
\sum\nolimits_{k = 0}^p
(-1)^k
h(\alpha)_{j_0 \ldots \hat j_k \ldots j_p}
+ \\
&
(-1)^p
d_{\mathcal F}(h(\alpha)_{j_0 \ldots j_p})
+ \\
&
\sum\nolimits_{a = 0}^p
(-1)^a
d(\alpha)_{t(j_0) \ldots t(j_a) t'(j_a) \ldots t'(j_p)}
\\
= &
\sum\nolimits_{k = 0}^p
\sum\nolimits_{a = 0}^{k - 1}
(-1)^{k + a}
\alpha_{t(j_0)\ldots t(j_a)t'(j_a)\ldots \hat{t'(j_k)}\ldots t'(j_p)}
+ \\
&
\sum\nolimits_{k = 0}^p
\sum\nolimits_{a = k + 1}^p
(-1)^{k + a - 1}
\alpha_{t(j_0)\ldots \hat{t(j_k)}\ldots t(j_a)t'(j_a)\ldots t'(j_p)}
+ \\
&
\sum\nolimits_{a = 0}^p
(-1)^{p + a}
d_{\mathcal F}(\alpha_{t(j_0)\ldots t(j_a) t'(j_a) \ldots t'(j_p)})
+ \\
&
\sum\nolimits_{a = 0}^p
\sum\nolimits_{k = 0}^a
(-1)^{a + k}
\alpha_{t(j_0)\ldots\hat{t(j_k)}\ldots t(j_a)t'(j_a)\ldots t'(j_p)}
+ \\
&
\sum\nolimits_{a = 0}^p
\sum\nolimits_{k = a}^p
(-1)^{a + k + 1}
\alpha_{t(j_0) \ldots t(j_a) t'(j_a) \ldots \hat{t'(j_k)} \ldots t'(j_p)}
+ \\
&
\sum\nolimits_{a = 0}^p
(-1)^{a + p + 1}
d_{\mathcal F}(\alpha_{t(j_0)\ldots t(j_a) t'(j_a) \ldots t'(j_p)})
\\
= &
\alpha_{t'(j_0)\ldots t'(j_p)} +
(-1)^{2p + 1}\alpha_{t(j_0)\ldots t(j_p)}
\\
= &
T_{t'}(\alpha)_{j_0\ldots j_p} - T_t(\alpha)_{j_0\ldots j_p}
\end{align*}

Nous laissons le lecteur vérifier les annulations. (Notez que les termes ayant les deux $k$ et $a$ dans les 1er, 2e et 4e, 5e summands annuler, sauf ceux où $a = k$ qui ne se produisent que dans les 4ème et 5ème et qui s'annulent l'un contre l'autre, sauf pour les deux termes souhaités.) Il s'ensuit que l'application induite
$$
H^n(T_t) :
H^n(
\text{Tot}(\check{\mathcal{C}}^\bullet({\mathcal U}, {\mathcal F}^\bullet))
)
\to
H^n(
\text{Tot}(\check{\mathcal{C}}^\bullet({\mathcal V}, {\mathcal F}^\bullet))
)
$$
est indépendant du choix de $t$. Nous définissons
{\it hypercohomologie de {\v C}ech} comme la colimite des groupes de cohomologie de {\v C}ech sur tous les raffinements, relativement aux applications $H^\bullet(T_t)$.

\medskip\noindent
En passant à la colimite sur tous les recouvrements ouverts de $X$, le lemme suivant fournit une application de l'hypercohomologie de {\v C}ech vers la cohomologie. Celle-ci est souvent un isomorphisme et l'est toujours si l'on utilise des hyperrecouvrements.

\begin{lemma}

\label{lemma-cech-complex-complex}
Soit $(X, \mathcal{O}_X)$ un espace annelé, soit $\mathcal{U} : X = \bigcup_{i \in I} U_i$ un recouvrement ouvert. Pour tout complexe borné inférieurement $\mathcal{F}^\bullet$ de $\mathcal{O}_X$-modules, il existe une application canonique $$
\text{Tot}(\check{\mathcal{C}}^\bullet(\mathcal{U}, \mathcal{F}^\bullet))
\longrightarrow
R\Gamma(X, \mathcal{F}^\bullet)
$$ fonctorielle en $\mathcal{F}^\bullet$ et compatible avec (\ref{equation-global-sections-to-cech}) et (\ref{equation-transformation}). Il existe une suite spectrale $(E_r, d_r)_{r \geq 0}$ telle que $$
E_2^{p, q} =
H^p(\text{Tot}(\check{\mathcal{C}}^\bullet(\mathcal{U},
\underline{H}^q(\mathcal{F}^\bullet)))
$$ convergeant vers $H^{p + q}(X, \mathcal{F}^\bullet)$.
\end{lemma}

\begin{proof}

Soit ${\mathcal I}^\bullet$ un complexe borné inférieurement de faisceaux injectifs. L'application (\ref{equation-global-sections-to-cech}) pour
$\mathcal{I}^\bullet$ est une application
$\Gamma(X, {\mathcal I}^\bullet) \to
\text{Tot}(\check{\mathcal{C}}^\bullet({\mathcal U}, {\mathcal I}^\bullet))$. C'est un quasi-isomorphisme de complexes de groupes abéliens, comme il résulte d'Homologie, lemme \ref{homology-lemma-double-complex-gives-resolution},
appliqué au complexe double
$\check{\mathcal{C}}^\bullet({\mathcal U}, {\mathcal I}^\bullet)$, en utilisant le lemme \ref{lemma-injective-trivial-cech}. Choisissons un quasi-isomorphisme ${\mathcal F}^\bullet \to {\mathcal I}^\bullet$ de ${\mathcal F}^\bullet$ vers un complexe borné inférieurement de faisceaux injectifs. Le complexe représentant $R\Gamma(X, {\mathcal F}^\bullet)$ est
$\Gamma(X, {\mathcal I}^\bullet)$; on obtient l'application du lemme au moyen de
$$
\text{Tot}(\check{\mathcal{C}}^\bullet({\mathcal U}, {\mathcal F}^\bullet))
\longrightarrow
\text{Tot}(\check{\mathcal{C}}^\bullet({\mathcal U}, {\mathcal I}^\bullet)).
$$
Nous omettons la vérification de la fonctorialité et des compatibilités. Pour construire la suite spectrale du lemme, choisissons une résolution de Cartan–Eilenberg $\mathcal{F}^\bullet \to \mathcal{I}^{\bullet, \bullet}$; voir Catégories dérivées, lemme \ref{derived-lemma-cartan-eilenberg}. Alors $\mathcal{F}^\bullet \to \text{Tot}(\mathcal{I}^{\bullet, \bullet})$
est une résolution injective, et donc
$$
\text{Tot}(\check{\mathcal{C}}^\bullet({\mathcal U},
\text{Tot}({\mathcal I}^{\bullet, \bullet})))
$$
calcule $R\Gamma(X, \mathcal{F}^\bullet)$, comme nous l'avons vu plus haut. D'après Homologie, remarque \ref{homology-remark-triple-complex},
nous pouvons le considérer comme le complexe total associé au triple complexe
$\check{\mathcal{C}}^\bullet({\mathcal U}, {\mathcal I}^{\bullet, \bullet})$
et, en utilisant la même remarque, comme le complexe total associé au bicomplexe $A^{\bullet, \bullet}$ dont les termes sont
$$
A^{n, m} =
\bigoplus\nolimits_{p + q = n}
\check{\mathcal{C}}^p({\mathcal U}, \mathcal{I}^{q, m})
$$
Comme $\mathcal{I}^{q, \bullet}$ est une résolution injective de
$\mathcal{F}^q$, nous pouvons appliquer la première suite spectrale associée à
$A^{\bullet, \bullet}$
(Homologie, lemme \ref{homology-lemma-ss-double-complex}) pour obtenir une suite spectrale
$$
E_1^{n, m} =
\bigoplus\nolimits_{p + q = n}
\check{\mathcal{C}}^p(\mathcal{U}, \underline{H}^m(\mathcal{F}^q))
$$
qui est le terme de degré $n$ du complexe
$\text{Tot}(\check{\mathcal{C}}^\bullet(\mathcal{U},
\underline{H}^m(\mathcal{F}^\bullet))$. On obtient donc
les termes $E_2$ décrits dans le lemme. La convergence résulte d'Homologie, lemme \ref{homology-lemma-first-quadrant-ss}.

\end{proof}

\begin{lemma}

\label{lemma-cech-complex-complex-computes}
Soit $(X, \mathcal{O}_X)$ un espace annelé. Soit $\mathcal{U} : X = \bigcup_{i \in I} U_i$ un recouvrement ouvert. Soit $\mathcal{F}^\bullet$ un complexe borné inférieurement de $\mathcal{O}_X$-modules. Si $H^i(U_{i_0 \ldots i_p}, \mathcal{F}^q) = 0$ pour tout $i > 0$ et tous $p, i_0, \ldots, i_p, q$, alors l'application $
\text{Tot}(\check{\mathcal{C}}^\bullet(\mathcal{U}, \mathcal{F}^\bullet))
\to
R\Gamma(X, \mathcal{F}^\bullet)
$ du lemme \ref{lemma-cech-complex-complex} est un isomorphisme.
\end{lemma}

\begin{proof}
Cela résulte immédiatement de la suite spectrale du lemme \ref{lemma-cech-complex-complex}.
\end{proof}

\begin{remark}

\label{remark-shift-complex-cech-complex}

Soit $(X, \mathcal{O}_X)$ un espace annelé.
Soit $\mathcal{U} : X = \bigcup_{i \in I} U_i$ un recouvrement ouvert, soit $\mathcal{F}^\bullet$ un complexe borné inférieurement de $\mathcal{O}_X$-modules et soit $b$ un entier. Nous affirmons qu'il existe un diagramme commutatif
$$
\xymatrix{
\text{Tot}(\check{\mathcal{C}}^\bullet(\mathcal{U}, \mathcal{F}^\bullet))[b]
\ar[r] \ar[d]_\gamma &
R\Gamma(X, \mathcal{F}^\bullet)[b] \ar[d] \\
\text{Tot}(\check{\mathcal{C}}^\bullet(\mathcal{U}, \mathcal{F}^\bullet[b]))
\ar[r] &
R\Gamma(X, \mathcal{F}^\bullet[b])
}
$$
dans la catégorie dérivée, où $\gamma$ est l'application de complexes construite dans Homologie, remarque \ref{homology-remark-shift-double-complex}. Cela a un sens puisque le bicomplexe
$\check{\mathcal{C}}^\bullet(\mathcal{U}, \mathcal{F}^\bullet[b])$
est manifestement le même que le bicomplexe
$\check{\mathcal{C}}^\bullet(\mathcal{U}, \mathcal{F}^\bullet)[0, b]$
introduit dans Homologie, remarque \ref{homology-remark-shift-double-complex}. Pour vérifier que le diagramme commute, choisissons une résolution injective
$\mathcal{F}^\bullet \to \mathcal{I}^\bullet$ comme dans la preuve du lemme \ref{lemma-cech-complex-complex}. Une chasse au diagramme montre qu'il suffit de vérifier la commutativité après avoir remplacé $\mathcal{F}^\bullet$
par $\mathcal{I}^\bullet$. Ensuite, nous considérons le diagramme étendu
$$
\xymatrix{
\Gamma(X, \mathcal{I}^\bullet)[b] \ar[r] \ar[d] &
\text{Tot}(\check{\mathcal{C}}^\bullet(\mathcal{U}, \mathcal{I}^\bullet))[b]
\ar[r] \ar[d]_\gamma &
R\Gamma(X, \mathcal{I}^\bullet)[b] \ar[d] \\
\Gamma(X, \mathcal{I}^\bullet[b]) \ar[r] &
\text{Tot}(\check{\mathcal{C}}^\bullet(\mathcal{U}, \mathcal{I}^\bullet[b]))
\ar[r] &
R\Gamma(X, \mathcal{I}^\bullet[b])
}
$$
où les flèches horizontales de gauche sont celles de (\ref{equation-global-sections-to-cech}). Puisque les flèches horizontales sont ici des isomorphismes dans la catégorie dérivée (voir la preuve du lemme \ref{lemma-cech-complex-complex}), il suffit de montrer que le carré de gauche commute. C'est le cas parce que l'application $\gamma$ utilise le signe $1$ sur les facteurs directs
$\check{\mathcal{C}}^0(\mathcal{U}, \mathcal{I}^{q + b})$, voir la formule dans Homologie, remarque \ref{homology-remark-shift-double-complex}.

\end{remark}

\noindent

Soit $X$ un espace topologique, soit $\mathcal{U} : X = \bigcup_{i \in I} U_i$
un recouvrement ouvert et soit $\mathcal{F}^\bullet$ un complexe borné inférieurement de préfaisceaux de groupes abéliens.
$\tau :
\text{Tot}(\check{\mathcal{C}}^\bullet({\mathcal U}, {\mathcal F}^\bullet))
\to
\text{Tot}(\check{\mathcal{C}}^\bullet({\mathcal U}, {\mathcal F}^\bullet))$
défini par
$$
\tau(\alpha)_{i_0 \ldots i_p} = (-1)^{p(p + 1)/2} \alpha_{i_p \ldots i_0}.
$$
Alors, pour tout élément $\alpha$ de degré $n$, on a

\begin{align*}
& d(\tau(\alpha))_{i_0 \ldots i_{p + 1}} \\
& =
\sum\nolimits_{j = 0}^{p + 1}
(-1)^j
\tau(\alpha)_{i_0 \ldots \hat i_j \ldots i_{p + 1}}
+
(-1)^{p + 1}
d_{\mathcal F}(\tau(\alpha)_{i_0 \ldots i_{p + 1}})
\\
& =
\sum\nolimits_{j = 0}^{p + 1}
(-1)^{j + \frac{p(p + 1)}{2}}
\alpha_{i_{p + 1} \ldots \hat i_j \ldots i_0}
+
(-1)^{p + 1 + \frac{(p + 1)(p + 2)}{2}}
d_{\mathcal F}(\alpha_{i_{p + 1} \ldots i_0})
\end{align*}

D'autre part, nous avons

\begin{align*}
& \tau(d(\alpha))_{i_0\ldots i_{p + 1}} \\
& =
(-1)^{\frac{(p + 1)(p + 2)}{2}} d(\alpha)_{i_{p + 1} \ldots i_0}
\\
& =
(-1)^{\frac{(p + 1)(p + 2)}{2}}
\left(
\sum\nolimits_{j = 0}^{p + 1}
(-1)^j
\alpha_{i_{p + 1}\ldots \hat i_{p + 1 - j} \ldots i_0}
+
(-1)^{p + 1}
d_{\mathcal F}(\alpha_{i_{p + 1}\ldots i_0})
\right)
\end{align*}

Ainsi, nous concluons que $d(\tau(\alpha)) = \tau(d(\alpha))$
parce que $p(p + 1)/2 \equiv (p + 1)(p + 2)/2 + p + 1 \bmod 2$. En d'autres termes
$\tau$ est un endomorphisme du complexe
$\text{Tot}(\check{\mathcal{C}}^\bullet({\mathcal U}, {\mathcal F}^\bullet))$. Notez que le schéma
$$
\begin{matrix}
\Gamma(X, {\mathcal F}^\bullet) &
\longrightarrow &
\text{Tot}(\check{\mathcal{C}}^\bullet({\mathcal U}, {\mathcal F}^\bullet)) \\
\downarrow \text{id} & & \downarrow \tau \\
\Gamma(X, {\mathcal F}^\bullet) &
\longrightarrow &
\text{Tot}(\check{\mathcal{C}}^\bullet({\mathcal U}, {\mathcal F}^\bullet))
\end{matrix}
$$
commute. En outre, $\tau$ est manifestement compatible avec les raffinements. Cela suggère que $\tau$ agit comme l'identité sur la cohomologie de {\v C}ech, c'est-à-dire dans la colimite, dès que l'hypercohomologie de {\v C}ech coïncide avec l'hypercohomologie; c'est toujours le cas si l'on utilise des hyperrecouvrements. En fait, $\tau$ est homotope à l'identité sur le complexe total de {\v C}ech
$\text{Tot}(\check{\mathcal{C}}^\bullet({\mathcal U}, {\mathcal F}^\bullet))$. Pour prouver cela, nous utilisons comme homotopie
$$
h(\alpha)_{i_0 \ldots i_p}
=
\sum\nolimits_{a = 0}^p
\epsilon_p(a)
\alpha_{i_0 \ldots i_a i_p \ldots i_a}
\quad\text{avec}\quad
\epsilon_p(a) = (-1)^{\frac{(p - a)(p - a - 1)}{2} + p}
$$
pour $\alpha$ du degré $n$. Comme d'habitude nous omettons d'écrire
$|_{U_{i_0 \ldots i_p}}$. Cela fonctionne à cause du calcul suivant, encore avec
$\alpha$ un élément de degré $n$:

\begin{align*}
(d(h(\alpha)) + h(d(\alpha)))_{i_0 \ldots i_p}
= &
\sum\nolimits_{k = 0}^p
(-1)^k
h(\alpha)_{i_0 \ldots \hat i_k \ldots i_p}
+ \\
&
(-1)^p
d_{\mathcal F}(h(\alpha)_{i_0 \ldots i_p})
+ \\
&
\sum\nolimits_{a = 0}^p
\epsilon_p(a)
d(\alpha)_{i_0 \ldots i_a i_p \ldots i_a}
\\
= &
\sum\nolimits_{k = 0}^p
\sum\nolimits_{a = 0}^{k - 1}
(-1)^k \epsilon_{p - 1}(a)
\alpha_{i_0 \ldots i_a i_p \ldots \hat{i_k} \ldots i_a}
+ \\
&
\sum\nolimits_{k = 0}^p
\sum\nolimits_{a = k + 1}^p
(-1)^k \epsilon_{p - 1}(a - 1)
\alpha_{i_0 \ldots \hat{i_k} \ldots i_a i_p \ldots i_a}
+ \\
&
\sum\nolimits_{a = 0}^p
(-1)^p \epsilon_p(a)
d_{\mathcal F}(\alpha_{i_0 \ldots i_a i_p \ldots i_a})
+ \\
&
\sum\nolimits_{a = 0}^p
\sum\nolimits_{k = 0}^a
\epsilon_p(a) (-1)^k
\alpha_{i_0 \ldots \hat{i_k} \ldots i_a i_p \ldots i_a}
+ \\
&
\sum\nolimits_{a = 0}^p
\sum\nolimits_{k = a}^p
\epsilon_p(a) (-1)^{p + a + 1 - k}
\alpha_{i_0 \ldots i_a i_p \ldots \hat{i_k} \ldots i_a}
+ \\
&
\sum\nolimits_{a = 0}^p
\epsilon_p(a) (-1)^{p + 1}
d_{\mathcal F}(\alpha_{i_0 \ldots i_a i_p \ldots i_a})
\\
= &
\epsilon_p(0) \alpha_{i_p \ldots i_0} +
\epsilon_p(p) (-1)^{p + 1} \alpha_{i_0 \ldots i_p} \\
= &
(-1)^{\frac{p(p + 1)}{2}}\alpha_{i_p \ldots i_0}
- \alpha_{i_0 \ldots i_p}
\end{align*}

Les annulations suivent parce que
$$
(-1)^k \epsilon_{p - 1}(a) + \epsilon_p(a)(-1)^{p + a + 1 - k} = 0
\quad\text{et}\quad
(-1)^k\epsilon_{p - 1}(a - 1) + \epsilon_p(a) (-1)^k = 0
$$
Nous laissons le lecteur vérifier les annulations.

\medskip\noindent

Supposons que nous ayons deux complexes bornés inférieurement de faisceaux abéliens
${\mathcal F}^\bullet$ et ${\mathcal G}^\bullet$. Nous définissons le complexe
$\text{Tot}({\mathcal F}^\bullet\otimes_{\mathbf Z} {\mathcal G}^\bullet)$
être complexe avec les termes
$\bigoplus_{p + q = n} {\mathcal F}^p \otimes {\mathcal G}^q$
et différentiel selon la règle

\begin{equation}
\label{equation-differential-tensor-product-complexes}
d(\alpha \otimes \beta) =
d(\alpha)\otimes \beta + (-1)^{\deg(\alpha)} \alpha \otimes d(\beta)
\end{equation}

lorsque $\alpha$ et $\beta$ homogène, voir Homologie, définition \ref{homology-definition-associated-simple-complex}.

\medskip\noindent
Supposons que $M^\bullet$ et $N^\bullet$ soient deux complexes bornés inférieurement de groupes abéliens. Si $m$, respectivement $n$, est un cocycle de $M^\bullet$, respectivement de $N^\bullet$, il est immédiat que $m \otimes n$ est un cocycle de $\text{Tot}(M^\bullet\otimes N^\bullet)$. On obtient ainsi un produit cup $$
H^i(M^\bullet) \times H^j(N^\bullet)
\longrightarrow
H^{i + j}(Tot(M^\bullet\otimes N^\bullet)).
$$ Ceci est discuté également dans Compléments sur l'algèbre, section \ref{more-algebra-section-products-tor}.

\medskip\noindent

Ainsi, la construction du produit de la coupe dans l'hypercohomologie des complexes repose sur une construction d'une application des complexes

\begin{equation}
\label{equation-needs-signs}
\text{Tot}\left(
\text{Tot}(\check{\mathcal{C}}^\bullet({\mathcal U}, {\mathcal F}^\bullet))
\otimes_{\mathbf Z}
\text{Tot}(\check{\mathcal{C}}^\bullet({\mathcal U}, {\mathcal G}^\bullet))
\right)
\longrightarrow
\text{Tot}(
\check{\mathcal{C}}^\bullet({\mathcal U},
\text{Tot}({\mathcal F}^\bullet\otimes {\mathcal G}^\bullet)
))
\end{equation}

Cette application est dénotée $\cup$ et est donné par la règle
$$
(\alpha \cup \beta)_{i_0 \ldots i_p}
=
\sum\nolimits_{r = 0}^p
\epsilon(n, m, p, r)
\alpha_{i_0 \ldots i_r} \otimes \beta_{i_r \ldots i_p}.
$$
où $\alpha$ est de degré $n$ et $\beta$ de degré $m$
et avec
$$
\epsilon(n, m, p, r) = (-1)^{(p + r)n + rp + r}.
$$
Notez que $\epsilon(n, m, p, n) = 1$. Par conséquent, si
$\mathcal{F}^\bullet = \mathcal{F}[0]$ est le complexe constitué d'un seul faisceau abélien $\mathcal{F}$ placé dans le degré $0$, puis il n'y a pas de signes dans la formule pour $\cup$ (comme dans ce cas
$\alpha_{i_0 \ldots i_r} = 0$ sauf $r = n$). Pour une explication de pourquoi il doit y avoir un signe et comment le calculer voir \cite[Exposé XVII]{SGA4} par Deligne. Pour vérifier (\ref{equation-needs-signs}) est une application des complexes que nous devons montrer que
$$
d(\alpha \cup \beta) =
d(\alpha) \cup \beta +
(-1)^{\deg(\alpha)} \alpha \cup d(\beta)
$$
par la définition du différentiel
$\text{Tot}(
\text{Tot}(\check{\mathcal{C}}^\bullet({\mathcal U}, {\mathcal F}^\bullet))
\otimes_{\mathbf Z}
\text{Tot}(\check{\mathcal{C}}^\bullet({\mathcal U}, {\mathcal G}^\bullet))
)$
tel qu'indiqué dans Homologie, définition \ref{homology-definition-associated-simple-complex}. Nous calculons d'abord

\begin{align*}
d(\alpha \cup \beta)_{i_0 \ldots i_{p + 1}}
= &
\sum\nolimits_{j = 0}^{p + 1}
(-1)^j
(\alpha \cup \beta)_{i_0 \ldots \hat i_j \ldots i_{p + 1}}
+
(-1)^{p + 1}
d_{{\mathcal F} \otimes {\mathcal G}}
((\alpha \cup \beta)_{i_0 \ldots i_{p + 1}})
\\
= &
\sum\nolimits_{j = 0}^{p + 1}
\sum\nolimits_{r = 0}^{j - 1}
(-1)^j \epsilon(n, m, p, r)
\alpha_{i_0 \ldots i_r} \otimes \beta_{i_r \ldots \hat i_j \ldots i_{p + 1}}
+ \\
&
\sum\nolimits_{j = 0}^{p + 1}
\sum\nolimits_{r = j + 1}^{p + 1}
(-1)^j \epsilon(n, m, p, r - 1)
\alpha_{i_0 \ldots \hat i_j \ldots i_r} \otimes \beta_{i_r \ldots i_{p + 1}}
+ \\
&
\sum\nolimits_{r = 0}^{p + 1}
(-1)^{p + 1} \epsilon(n, m, p + 1, r)
d_{{\mathcal F} \otimes {\mathcal G}}
(\alpha_{i_0 \ldots i_r} \otimes \beta_{i_r \ldots i_{p + 1}})
\end{align*}

et noter que les sommes dans le dernier terme égale
$$
(-1)^{p + 1} \epsilon(n, m, p + 1, r)
\left(
d_{\mathcal F}(\alpha_{i_0 \ldots i_r}) \otimes 
\beta_{i_r \ldots i_{p + 1}} +
(-1)^{n - r}
\alpha_{i_0 \ldots i_r} \otimes d_{\mathcal G}(\beta_{i_r \ldots i_{p + 1}})
\right).
$$
parce que $\deg_\mathcal{F}(\alpha_{i_0 \ldots i_r}) = n - r$. D'un autre côté

\begin{align*}
(d(\alpha) \cup \beta)_{i_0\ldots i_{p + 1}}
= &
\sum\nolimits_{r = 0}^{p + 1}
\epsilon(n + 1, m, p + 1, r)
d(\alpha)_{i_0\ldots i_r} \otimes \beta_{i_r\ldots i_{p + 1}}
\\
= &
\sum\nolimits_{r = 0}^{p + 1}
\sum\nolimits_{j = 0}^{r}
\epsilon(n + 1, m, p + 1, r) (-1)^j
\alpha_{i_0\ldots\hat{i_j}\ldots i_r} \otimes \beta_{i_r\ldots i_{p + 1}}
+ \\
&
\sum\nolimits_{r = 0}^{p + 1}
\epsilon(n + 1, m, p + 1, r) (-1)^r
d_{\mathcal F}(\alpha_{i_0 \ldots i_r}) \otimes \beta_{i_r\ldots i_{p + 1}}
\end{align*}

et

\begin{align*}
(\alpha \cup d(\beta))_{i_0\ldots i_{p + 1}}
= &
\sum\nolimits_{r = 0}^{p + 1}
\epsilon(n, m + 1, p + 1, r)
\alpha_{i_0 \ldots i_r} \otimes d(\beta)_{i_r \ldots i_{p + 1}}
\\
= &
\sum\nolimits_{r = 0}^{p + 1}
\sum\nolimits_{j = r}^{p + 1}
\epsilon(n, m + 1, p + 1, r) (-1)^{j - r}
\alpha_{i_0 \ldots i_r} \otimes \beta_{i_r \ldots \hat{i_j}\ldots i_{p + 1}}
+ \\
&
\sum\nolimits_{r = 0}^{p + 1}
\epsilon(n, m + 1, p + 1, r) (-1)^{p + 1 - r}
\alpha_{i_0 \ldots i_r} \otimes d_{\mathcal G}(\beta_{i_r \ldots i_{p + 1}})
\end{align*}

L'égalité souhaitée tient si nous avons

\begin{align*}
(-1)^{p + 1} \epsilon(n, m, p + 1, r)
& =
\epsilon(n + 1, m, p + 1, r) (-1)^r \\
(-1)^{p + 1} \epsilon(n, m, p + 1, r) (-1)^{n - r}
& =
(-1)^n \epsilon(n, m + 1, p + 1, r) (-1)^{p + 1 - r} \\
\epsilon(n + 1, m, p + 1, r) (-1)^r
& =
(-1)^{1 + n} \epsilon(n, m + 1, p + 1, r - 1) \\
(-1)^j \epsilon(n, m, p, r)
& =
(-1)^n \epsilon(n, m + 1, p + 1, r) (-1)^{j - r} \\
(-1)^j \epsilon(n, m, p, r - 1)
& =
\epsilon(n + 1, m, p + 1, r) (-1)^j
\end{align*}

(La troisième égalité est nécessaire pour obtenir les $r = j$
à partir $d(\alpha) \cup \beta$ et $(-1)^n \alpha \cup d(\beta)$
Nous laissons les vérifications au lecteur. (Alternativement, vérifiez les script signs.gp dans le sous-répertoire scripts du projet Champs.)

\medskip\noindent

Associativité du produit de la tasse. Supposons que ${\mathcal F}^\bullet$,
${\mathcal G}^\bullet$ et ${\mathcal H}^\bullet$ sont nés complexes d'intérieur de groupes abéliens sur $X$. La carte évidente (sans intervention des signes) est un isomorphisme des complexes
$$
\text{Tot}(
\text{Tot}({\mathcal F}^\bullet \otimes_{\mathbf Z} {\mathcal G}^\bullet)
\otimes_{\mathbf Z} {\mathcal H}^\bullet
)
\longrightarrow
\text{Tot}(
{\mathcal F}^\bullet \otimes_{\mathbf Z}
\text{Tot}({\mathcal G}^\bullet \otimes_{\mathbf Z} {\mathcal H}^\bullet)
).
$$
Une autre façon de dire ceci est que le triple complexe
${\mathcal F}^\bullet \otimes_{\mathbf Z} {\mathcal G}^\bullet
\otimes_{\mathbf Z} {\mathcal H}^\bullet$ donne lieu à un complexe bien défini total avec différentiel satisfaisant
$$
d(\alpha \otimes \beta \otimes \gamma) =
d(\alpha) \otimes \beta \otimes \gamma +
(-1)^{\deg(\alpha)} \alpha \otimes d(\beta) \otimes \gamma +
(-1)^{\deg(\alpha) + \deg(\beta)} \alpha \otimes \beta \otimes d(\gamma)
$$
pour les éléments homogènes. En utilisant cette application, il est facile de vérifier que
$$
(\alpha \cup \beta) \cup \gamma = \alpha \cup ( \beta \cup \gamma)
$$
à savoir, si $\alpha$ est de degré $a$, $\beta$ de degré $b$ et
$\gamma$ de degré $c$, alors

\begin{align*}
((\alpha \cup \beta) \cup \gamma)_{i_0 \ldots i_p}
= &
\sum\nolimits_{r = 0}^p
\epsilon(a + b, c, p, r)
(\alpha \cup \beta)_{i_0 \ldots i_r} \otimes \gamma_{i_r \ldots i_p}
\\
= &
\sum\nolimits_{r = 0}^p
\sum\nolimits_{s = 0}^r
\epsilon(a + b, c, p, r) \epsilon(a, b, r, s)
\alpha_{i_0 \ldots i_s} \otimes
\beta_{i_s \ldots i_r} \otimes
\gamma_{i_r \ldots i_p}
\end{align*}

et

\begin{align*}
(\alpha \cup (\beta \cup \gamma)_{i_0\ldots i_p}
= &
\sum\nolimits_{s = 0}^p
\epsilon(a, b + c, p, s)
\alpha_{i_0 \ldots i_s} \otimes (\beta \cup \gamma)_{i_s \ldots i_p}
\\
= &
\sum\nolimits_{s = 0}^p
\sum\nolimits_{r = s}^p
\epsilon(a, b + c, p, s) \epsilon(b, c, p - s, r - s)
\alpha_{i_0 \ldots i_s} \otimes \beta_{i_s \ldots i_r} \otimes
\gamma_{i_r \ldots i_p}
\end{align*}

et un mod trivial $2$ (Alternativement, vérifiez les script signs.gp dans le sous-répertoire scripts du projet Champs.)

\medskip\noindent

Enfin, nous indiquons pourquoi le produit de la coupe conserve une structure commutative gradée, au moins au niveau cohomologique.
$\tau$ présenté ci-dessus. Soit ${\mathcal F}^\bullet$ un complexe borné inférieurement de groupes abéliens, muni d'une multiplication commutative graduée
$$
\wedge^\bullet :
\text{Tot}({\mathcal F}^\bullet\otimes {\mathcal F}^\bullet)
\longrightarrow
{\mathcal F}^\bullet.
$$
Cela signifie ce qui suit : $s$ une section locale de
${\mathcal F}^a$, $t$ une section locale de ${\mathcal F}^b$
Nous avons $s \wedge t$ une section locale de ${\mathcal F}^{a + b}$. Commutatif graduel signifie que nous avons
$s \wedge t = (-1)^{ab} t \wedge s$Depuis... $\wedge$ est une application des complexes que nous avons
$d(s\wedge t) = d(s) \wedge t + (-1)^a s \wedge d(t)$. La composition
$$
\text{Tot}(
\text{Tot}(\check{\mathcal{C}}^\bullet({\mathcal U}, {\mathcal F}^\bullet))
\otimes
\text{Tot}(\check{\mathcal{C}}^\bullet({\mathcal U}, {\mathcal F}^\bullet))
)
\to
\text{Tot}(
\check{\mathcal{C}}^\bullet({\mathcal U},
\text{Tot}({\mathcal F}^\bullet\otimes_{\mathbf Z}{\mathcal F}^\bullet))
)
\to
\text{Tot}(\check{\mathcal{C}}^\bullet({\mathcal U}, {\mathcal F}^\bullet))
$$
induit une tasse produit sur la cohomologie
$$
H^n(
\text{Tot}(\check{\mathcal{C}}^\bullet({\mathcal U}, {\mathcal F}^\bullet))
)
\times
H^m(
\text{Tot}(\check{\mathcal{C}}^\bullet({\mathcal U}, {\mathcal F}^\bullet))
)
\longrightarrow
H^{n + m}(
\text{Tot}(\check{\mathcal{C}}^\bullet({\mathcal U}, {\mathcal F}^\bullet))
)
$$
et donc, après passage à la limite inductive, également un produit sur la cohomologie de {\v C}ech et, par conséquent (en utilisant des hyperrecouvrements si nécessaire), un produit sur la cohomologie de ${\mathcal F}^\bullet$. Nous affirmons que ce produit (sur la cohomologie) est lui aussi commutatif gradué. Pour le prouver, considérons d'abord un élément $\alpha$ de degré $n$ dans
$\text{Tot}(\check{\mathcal{C}}^\bullet({\mathcal U}, {\mathcal F}^\bullet))$
et un élément $\beta$ de degré $m$ dans
$\text{Tot}(\check{\mathcal{C}}^\bullet({\mathcal U}, {\mathcal F}^\bullet))$
et calculons

\begin{align*}
\wedge^\bullet(\alpha \cup \beta)_{i_0 \ldots i_p}
= &
\sum\nolimits_{r = 0}^p
\epsilon(n, m, p, r)
\alpha_{i_0 \ldots i_r} \wedge \beta_{i_r \ldots i_p} \\
= &
\sum\nolimits_{r = 0}^p
\epsilon(n, m, p, r)
(-1)^{\deg(\alpha_{i_0 \ldots i_r})\deg(\beta_{i_r \ldots i_p})}
\beta_{i_r \ldots i_p} \wedge \alpha_{i_0 \ldots i_r}
\end{align*}

parce que $\wedge$ est classé comme commutatif. D'autre part, nous avons

\begin{align*}
\tau(\wedge^\bullet(\tau(\beta) \cup \tau(\alpha)))_{i_0 \ldots i_p}
= &
\chi(p)
\sum\nolimits_{r = 0}^p
\epsilon(m, n, p, r)
\tau(\beta)_{i_p \ldots i_{p - r}} \wedge \tau(\alpha)_{i_{p - r} \ldots i_0}
\\
= &
\chi(p)
\sum\nolimits_{r = 0}^p
\epsilon(m, n, p, r) \chi(r) \chi(p - r)
\beta_{i_{p - r} \ldots i_p} \wedge \alpha_{i_0 \ldots i_{p - r}}
\\
= &
\chi(p)
\sum\nolimits_{r = 0}^p
\epsilon(m, n, p, p - r) \chi(r) \chi(p - r)
\beta_{i_r \ldots i_p} \wedge \alpha_{i_0 \ldots i_r}
\end{align*}

où $\chi(t) = (-1)^{\frac{t(t + 1)}{2}}$. Puisque nous avons prouvé plus tôt que
$\tau$ agit comme l'identité sur la cohomologie, nous devons vérifier que
$$
\epsilon(n, m, p, r)
(-1)^{(n - r)(m - (p - r))}
=
(-1)^{nm} \chi(p)\epsilon(m, n, p, p - r) \chi(r) \chi(p - r)
$$
Un calcul immédiat modulo $2$ le vérifie. On peut aussi consulter le script signs.gp dans le sous-répertoire scripts du projet Champs.

\medskip\noindent
Enfin, étudions la compatibilité du produit cup avec les applications de bord. Supposons que $$
0
\to
{\mathcal F}_1^\bullet
\to
{\mathcal F}_2^\bullet
\to
{\mathcal F}_3^\bullet
\to
0
\quad\text{et}\quad
0
\leftarrow
{\mathcal G}_1^\bullet
\leftarrow
{\mathcal G}_2^\bullet
\leftarrow
{\mathcal G}_3^\bullet
\leftarrow
0
$$ soient des suites exactes courtes de complexes bornés inférieurement de faisceaux abéliens sur $X$. Soit ${\mathcal H}^\bullet$ un autre complexe borné inférieurement de faisceaux abéliens, et supposons que nous ayons des morphismes de complexes $$
\gamma_i :
\text{Tot}({\mathcal F}_i^\bullet \otimes_{\mathbf Z} {\mathcal G}_i^\bullet)
\longrightarrow
{\mathcal H}^\bullet
$$ compatibles avec les morphismes entre ces complexes, c'est-à-dire tels que les diagrammes $$
\xymatrix{
\text{Tot}({\mathcal F}_1^\bullet \otimes_{\mathbf Z} {\mathcal G}_1^\bullet)
\ar[d]_{\gamma_1}
&
\text{Tot}({\mathcal F}_1^\bullet \otimes_{\mathbf Z} {\mathcal G}_2^\bullet)
\ar[l] \ar[d]
\\
\mathcal{H}^\bullet &
\text{Tot}({\mathcal F}_2^\bullet \otimes_{\mathbf Z} {\mathcal G}_2^\bullet)
\ar[l]_-{\gamma_2}
}
$$ et $$
\xymatrix{
\text{Tot}({\mathcal F}_2^\bullet \otimes_{\mathbf Z} {\mathcal G}_2^\bullet)
\ar[d]_{\gamma_2}
&
\text{Tot}({\mathcal F}_2^\bullet \otimes_{\mathbf Z} {\mathcal G}_3^\bullet)
\ar[l] \ar[d]
\\
\mathcal{H}^\bullet &
\text{Tot}({\mathcal F}_3^\bullet \otimes_{\mathbf Z} {\mathcal G}_3^\bullet)
\ar[l]_-{\gamma_3}
}
$$ sont commutatifs.

\begin{lemma}

\label{lemma-compute-sign-cup-product-boundaries}
Dans la situation ci-dessus, supposons que la cohomologie de {\v C}ech coïncide avec la cohomologie pour les faisceaux $\mathcal{F}_i^p$ et $\mathcal{G}_j^q$. Soient $a_3 \in H^n(X, \mathcal{F}_3^\bullet)$ et $b_1 \in H^m(X, \mathcal{G}_1^\bullet)$. Alors $$
\gamma_1( \partial a_3 \cup b_1) =
(-1)^{n + 1} \gamma_3( a_3 \cup \partial b_1)
$$ dans $H^{n + m + 1}(X, \mathcal{H}^\bullet)$, où $\partial$ désigne l'application de bord en cohomologie associée aux suites exactes courtes de complexes ci-dessus.
\end{lemma}

\begin{proof}

Nous utiliserons les conventions et notations suivantes. Nous considérons
${\mathcal F}_1^p$ comme un sous-faisceau de ${\mathcal F}_2^p$ et
${\mathcal G}_3^q$ comme un sous-faisceau de ${\mathcal G}_2^q$. Ainsi, si $s$ est une section locale de ${\mathcal F}_1^p$, nous notons encore $s$ la section correspondante de ${\mathcal F}_2^p$. De même pour les sections locales de ${\mathcal G}_3^q$. En outre, si $s$ est une section locale de ${\mathcal F}_2^p$, nous notons
$\bar s$ son image dans ${\mathcal F}_3^p$. On adopte la même convention pour l'application ${\mathcal G}_2^q \to {\mathcal G}^q_1$. En particulier, si
$s$ est une section locale de ${\mathcal F}_2^p$ et $\bar s = 0$
alors $s$ est une section locale de ${\mathcal F}_1^p$. La commutativité des diagrammes ci-dessus implique, pour les sections locales $s$ de
${\mathcal F}_2^p$ et $t$ de ${\mathcal G}_3^q$ que
$\gamma_2(s \otimes t) = \gamma_3(\bar s \otimes t)$ en tant que sections
${\mathcal H}^{p + q}$.

\medskip\noindent

Soit ${\mathcal U} : X =  \bigcup_{i \in I} U_i$ un recouvrement ouvert de $X$. Supposons que $\alpha_3$, respectivement $\beta_1$, soit un cocycle de degré $n$, respectivement $m$, de
$\text{Tot}(
\check{\mathcal{C}}^\bullet({\mathcal U}, {\mathcal F}_3^\bullet))$, respectivement de $\text{Tot}(
\check{\mathcal{C}}^\bullet({\mathcal U}, {\mathcal G}_1^\bullet))$
représentant $a_3$, respectivement $b_1$. Après avoir raffiné $\mathcal{U}$ si nécessaire, nous pouvons trouver des cochaînes $\alpha_2$, respectivement $\beta_2$, de degré $n$, respectivement $m$, dans
$\text{Tot}(
\check{\mathcal{C}}^\bullet({\mathcal U}, {\mathcal F}_2^\bullet))$, respectivement dans $\text{Tot}(
\check{\mathcal{C}}^\bullet({\mathcal U}, {\mathcal G}_2^\bullet))$
qui s'envoient sur $\alpha_3$, respectivement $\beta_1$. On a alors
$$
\overline{d(\alpha_2)} = d(\bar \alpha_2) = 0
\quad\text{et}\quad
\overline{d(\beta_2)} = d(\bar \beta_2) = 0.
$$
Cela signifie que $\alpha_1 = d(\alpha_2)$ est un cocycle de degré $n + 1$ dans
$\text{Tot}(\check{\mathcal{C}}^\bullet({\mathcal U}, {\mathcal F}_1^\bullet))$
représentant $\partial a_3$. De même, $\beta_3 = d(\beta_2)$ est un cocycle de degré $m + 1$ dans
$\text{Tot}(\check{\mathcal{C}}^\bullet({\mathcal U}, {\mathcal G}_3^\bullet))$
représentant $\partial b_1$. Ainsi nous pouvons calculer

\begin{align*}
d(\gamma_2(\alpha_2 \cup \beta_2))
& =
\gamma_2(d(\alpha_2 \cup \beta_2))
\\
& =
\gamma_2(d(\alpha_2) \cup \beta_2 + (-1)^n \alpha_2 \cup d(\beta_2) )
\\
& =
\gamma_2( \alpha_1 \cup \beta_2)  + (-1)^n \gamma_2( \alpha_2 \cup \beta_3)
\\
& =
\gamma_1(\alpha_1 \cup \beta_1) + (-1)^n \gamma_3(\alpha_3 \cup \beta_3)
\end{align*}

Donc cela nous dit même que le signe est $(-1)^{n + 1}$ comme indiqué dans le lemme\footnote{Le signe dépend de la convention pour les signes dans la suite exacte longue en cohomologie associée à un triangle dans $D(X)$. Les conventions dans le projet Champs sont : a) les triangles distincts correspondent à des séquences exactes par terme et b) les cartes limites dans la suite exacte longue sont données par les cartes dans la lemme du serpent sans intervention de signes. Voir Catégories dérivées, section \ref{derived-section-homotopy-triangulated}.}.

\end{proof}

\begin{lemma}

\label{lemma-boundary-derivation-over-cup-product}

Soit $X$ un espace topologique et soit $\mathcal{O}' \to \mathcal{O}$ une surjection de faisceaux d'anneaux dont le noyau $\mathcal{I} \subset \mathcal{O}'$
est de carré nul. Alors $M = H^1(X, \mathcal{I})$ est un
$R = H^0(X, \mathcal{O})$-module et l'application de bord
$\partial : R \to M$ associée à la suite exacte courte
$$
0 \to \mathcal{I} \to \mathcal{O}' \to \mathcal{O} \to 0
$$
est une dérivation (Algèbre, définition \ref{algebra-definition-derivation}).

\end{lemma}

\begin{proof}

L'application $\mathcal{O}' \to \SheafHom(\mathcal{I}, \mathcal{I})$
se factorise par $\mathcal{O}$ puisque $\mathcal{I} \cdot \mathcal{I} = 0$
par hypothèse. Ainsi, $\mathcal{I}$ est un faisceau de $\mathcal{O}$-modules, ce qui définit la structure de $R$-module sur $M$. L'application de bord est additive; il suffit donc de prouver la règle de Leibniz. Soit $f \in R$. Choisissons un recouvrement ouvert
$\mathcal{U} : X = \bigcup U_i$ tel qu'il existe des éléments
$f_i \in \mathcal{O}'(U_i)$ relevant $f|_{U_i} \in \mathcal{O}(U_i)$. Remarquons que $f_i - f_j$ appartient à $\mathcal{I}(U_i \cap U_j)$. Alors $\partial(f)$ correspond à la classe de cohomologie de {\v C}ech du $1$-cocycle $\alpha$ défini par $\alpha_{i_0i_1} = f_{i_0} - f_{i_1}$. (D'après le lemme \ref{lemma-cech-h1}, le premier groupe de cohomologie de {\v C}ech relatif à $\mathcal{U}$ est un sous-module de $M$.) Soit ensuite $g \in R$ un second élément et supposons, après avoir éventuellement raffiné le recouvrement ouvert, que $g_i \in \mathcal{O}'(U_i)$ relève
$g|_{U_i} \in \mathcal{O}(U_i)$. Alors nous voyons que
$\partial(g)$ est donné par le cocycle $\beta$ avec
$\beta_{i_0i_1} = g_{i_0} - g_{i_1}$. Comme $f_ig_i \in \mathcal{O}'(U_i)$
relève $fg|_{U_i}$, on voit que
$\partial(fg)$ est donné par le cocycle $\gamma$ avec
$$
\gamma_{i_0i_1} = f_{i_0}g_{i_0} - f_{i_1}g_{i_1} =
(f_{i_0} - f_{i_1})g_{i_0} + f_{i_1}(g_{i_0} - g_{i_1}) =
\alpha_{i_0i_1}g + f\beta_{i_0i_1}
$$
par notre définition du $\mathcal{O}$-module structure sur $\mathcal{I}$. Cela prouve la règle de Leibniz et la preuve est complète.

\end{proof}









\section{Résolutions plates}

\label{section-flat}

\noindent
Une référence pour le contenu de cette section est \cite{Spaltenstein}. Soit $(X, \mathcal{O}_X)$ un espace annelé. D'après Modules, lemme \ref{modules-lemma-module-quotient-flat}, tout $\mathcal{O}_X$-module est quotient d'un $\mathcal{O}_X$-module plat. D'après Catégories dérivées, lemme \ref{derived-lemma-subcategory-left-resolution}, tout complexe borné supérieurement de $\mathcal{O}_X$-modules admet une résolution à gauche par un complexe borné supérieurement de $\mathcal{O}_X$-modules plats. Pour les complexes non bornés, toutefois, les résolutions plates ne suffisent pas.

\begin{lemma}

\label{lemma-derived-tor-exact}
Soit $(X, \mathcal{O}_X)$ un espace annelé. Soit $\mathcal{G}^\bullet$ un complexe de $\mathcal{O}_X$-modules. Les foncteurs $$
K(\textit{Mod}(\mathcal{O}_X))
\longrightarrow
K(\textit{Mod}(\mathcal{O}_X)),
\quad
\mathcal{F}^\bullet \longmapsto
\text{Tot}(\mathcal{G}^\bullet \otimes_{\mathcal{O}_X} \mathcal{F}^\bullet)
$$ et $$
K(\textit{Mod}(\mathcal{O}_X))
\longrightarrow
K(\textit{Mod}(\mathcal{O}_X)),
\quad
\mathcal{F}^\bullet \longmapsto
\text{Tot}(\mathcal{F}^\bullet \otimes_{\mathcal{O}_X} \mathcal{G}^\bullet)
$$ sont des foncteurs exacts de catégories triangulées.
\end{lemma}

\begin{proof}
Cela résulte de Catégories dérivées, remarque \ref{derived-remark-double-complex-as-tensor-product-of}.
\end{proof}

\begin{definition}

\label{definition-K-flat}

Soit $(X, \mathcal{O}_X)$ un espace annelé. Un complexe $\mathcal{K}^\bullet$ de $\mathcal{O}_X$-modules est dit {\it K-plat} si, pour tout complexe acyclique $\mathcal{F}^\bullet$
de $\mathcal{O}_X$-modules, le complexe
$$
\text{Tot}(\mathcal{F}^\bullet \otimes_{\mathcal{O}_X} \mathcal{K}^\bullet)
$$
est acyclique.

\end{definition}

\begin{lemma}

\label{lemma-K-flat-quasi-isomorphism}

Soit $(X, \mathcal{O}_X)$ un espace annelé. Soit $\mathcal{K}^\bullet$ un complexe K-plat. Alors le foncteur
$$
K(\textit{Mod}(\mathcal{O}_X))
\longrightarrow
K(\textit{Mod}(\mathcal{O}_X)), \quad
\mathcal{F}^\bullet
\longmapsto
\text{Tot}(\mathcal{F}^\bullet \otimes_{\mathcal{O}_X} \mathcal{K}^\bullet)
$$
transforme les quasi-isomorphismes en quasi-isomorphismes.

\end{lemma}

\begin{proof}

Cela résulte du lemme \ref{lemma-derived-tor-exact}
et le fait que les quasi-isomorphismes sont caractérisés par des cônes acycliques.

\end{proof}

\begin{lemma}

\label{lemma-check-K-flat-stalks}

Soit $(X, \mathcal{O}_X)$ un espace annelé. Soit $\mathcal{K}^\bullet$
un complexe de $\mathcal{O}_X$-modules. Le complexe $\mathcal{K}^\bullet$
est K-plat si et seulement si, pour tout $x \in X$, le complexe
$\mathcal{K}_x^\bullet$ de $\mathcal{O}_{X, x}$-modules est K-plat (Compléments sur l'algèbre, définition \ref{more-algebra-definition-K-flat}).

\end{lemma}

\begin{proof}

Si $\mathcal{K}_x^\bullet$ est K-plat pour tout $x \in X$, alors $\mathcal{K}^\bullet$ est K-plat, car $\otimes$ et les sommes directes commutent à la prise des germes et l'exactitude se vérifie sur les germes; voir Modules, lemme \ref{modules-lemma-abelian}. Réciproquement, supposons $\mathcal{K}^\bullet$ K-plat. Fixons $x \in X$
et soit
$M^\bullet$ un complexe acyclique de $\mathcal{O}_{X, x}$-modules. Alors $i_{x, *}M^\bullet$ est un complexe acyclique de $\mathcal{O}_X$-modules, donc $\text{Tot}(i_{x, *}M^\bullet \otimes_{\mathcal{O}_X} \mathcal{K}^\bullet)$
est acyclique. En prenant le germe en $x$, on voit que
$\text{Tot}(M^\bullet \otimes_{\mathcal{O}_{X, x}} \mathcal{K}_x^\bullet)$
est acyclique.

\end{proof}

\begin{lemma}

\label{lemma-tensor-product-K-flat}

Soit $(X, \mathcal{O}_X)$ un espace annelé. Si $\mathcal{K}^\bullet$ et $\mathcal{L}^\bullet$ sont des complexes K-plats de $\mathcal{O}_X$-modules, alors
$\text{Tot}(\mathcal{K}^\bullet \otimes_{\mathcal{O}_X} \mathcal{L}^\bullet)$
est un complexe K-plat de $\mathcal{O}_X$-modules.

\end{lemma}

\begin{proof}

Cela résulte de l'isomorphisme
$$
\text{Tot}(\mathcal{M}^\bullet \otimes_{\mathcal{O}_X}
\text{Tot}(\mathcal{K}^\bullet \otimes_{\mathcal{O}_X} \mathcal{L}^\bullet))
=
\text{Tot}(\text{Tot}(\mathcal{M}^\bullet \otimes_{\mathcal{O}_X}
\mathcal{K}^\bullet) \otimes_{\mathcal{O}_X} \mathcal{L}^\bullet)
$$
et la définition.

\end{proof}

\begin{lemma}

\label{lemma-K-flat-two-out-of-three}
Soit $(X, \mathcal{O}_X)$ un espace annelé. Soit $(\mathcal{K}_1^\bullet, \mathcal{K}_2^\bullet, \mathcal{K}_3^\bullet)$ un triangle distingué dans $K(\textit{Mod}(\mathcal{O}_X))$. Si deux des trois complexes $\mathcal{K}_i^\bullet$ sont K-plats, alors le troisième l'est aussi.
\end{lemma}

\begin{proof}
Cela résulte du lemme \ref{lemma-derived-tor-exact} et du fait que, dans un triangle distingué de $K(\textit{Mod}(\mathcal{O}_X))$, si deux des trois objets sont acycliques, alors le troisième l'est aussi.
\end{proof}

\begin{lemma}

\label{lemma-K-flat-two-out-of-three-ses}
Soit $(X, \mathcal{O}_X)$ un espace annelé. Soit $0 \to \mathcal{K}_1^\bullet \to \mathcal{K}_2^\bullet \to
\mathcal{K}_3^\bullet \to 0$ une suite exacte courte de complexes, dont les termes de $\mathcal{K}_3^\bullet$ sont des $\mathcal{O}_X$-modules plats. Si deux des trois complexes $\mathcal{K}_i^\bullet$ sont K-plats, alors le troisième l'est aussi.
\end{lemma}

\begin{proof}
D'après Modules, lemme \ref{modules-lemma-flat-tor-zero}, pour tout complexe $\mathcal{L}^\bullet$, nous obtenons une suite exacte courte $$
0 \to
\text{Tot}(\mathcal{L}^\bullet \otimes_{\mathcal{O}_X} \mathcal{K}_1^\bullet)
\to
\text{Tot}(\mathcal{L}^\bullet \otimes_{\mathcal{O}_X} \mathcal{K}_1^\bullet)
\to
\text{Tot}(\mathcal{L}^\bullet \otimes_{\mathcal{O}_X} \mathcal{K}_1^\bullet)
\to 0
$$ de complexes. Le lemme résulte donc de la suite exacte longue des faisceaux de cohomologie et de la définition des complexes K-plats.
\end{proof}

\begin{lemma}

\label{lemma-pullback-K-flat}

Soit $f : (X, \mathcal{O}_X) \to (Y, \mathcal{O}_Y)$ un morphisme d'espaces annelés. L'image inverse d'un complexe K-plat de $\mathcal{O}_Y$-modules est un complexe K-plat de $\mathcal{O}_X$-modules.

\end{lemma}

\begin{proof}
Nous pouvons vérifier ceci sur les germes, voir lemme \ref{lemma-check-K-flat-stalks}. C'est ainsi que ceci résulte de Faisceaux, lemme \ref{sheaves-lemma-stalk-pullback-modules} et Compléments sur l'algèbre, lemme \ref{more-algebra-lemma-base-change-K-flat}.
\end{proof}

\begin{lemma}

\label{lemma-bounded-flat-K-flat}
Soit $(X, \mathcal{O}_X)$ un espace annelé. Tout complexe borné supérieurement de $\mathcal{O}_X$-modules plats est K-plat.
\end{lemma}

\begin{proof}
Nous pouvons vérifier ceci sur les germes, voir lemme \ref{lemma-check-K-flat-stalks}. Ainsi ce lemme suit des Modules, lemme \ref{modules-lemma-flat-stalks-flat} et Compléments sur l'algèbre, lemme \ref{more-algebra-lemma-derived-tor-quasi-isomorphism}.
\end{proof}

\noindent
Dans le lemme suivant, la colimite d'un système de complexes désigne la colimite terme à terme.

\begin{lemma}

\label{lemma-colimit-K-flat}

Soit $(X, \mathcal{O}_X)$ un espace annelé. Soit $\mathcal{K}_1^\bullet \to \mathcal{K}_2^\bullet \to \ldots$
un système de complexes K-plats. Alors $\colim_i \mathcal{K}_i^\bullet$ est K-plat.

\end{lemma}

\begin{proof}

Puisque nous prenons les colimites terme à terme, il est clair que
$$
\colim_i \text{Tot}(
\mathcal{F}^\bullet \otimes_{\mathcal{O}_X} \mathcal{K}_i^\bullet)
=
\text{Tot}(\mathcal{F}^\bullet \otimes_{\mathcal{O}_X}
\colim_i \mathcal{K}_i^\bullet)
$$
Le lemme découle donc de l'exactitude des colimites filtrantes.

\end{proof}

\begin{lemma}

\label{lemma-resolution-by-direct-sums-extensions-by-zero}

Soit $(X, \mathcal{O}_X)$ un espace annelé. Pour tout complexe $\mathcal{G}^\bullet$ de $\mathcal{O}_X$-modules il existe un diagramme commutatif des complexes de $\mathcal{O}_X$-modules
$$
\xymatrix{
\mathcal{K}_1^\bullet \ar[d] \ar[r] &
\mathcal{K}_2^\bullet \ar[d] \ar[r] & \ldots \\
\tau_{\leq 1}\mathcal{G}^\bullet \ar[r] &
\tau_{\leq 2}\mathcal{G}^\bullet \ar[r] & \ldots
}
$$
avec les propriétés suivantes : (1) les flèches verticales sont des quasi-isomorphismes et sont surjectives terme à terme, (2) chaque $\mathcal{K}_n^\bullet$ est un complexe borné supérieurement dont les termes sont des sommes directes de $\mathcal{O}_X$-modules de la forme
$j_{U!}\mathcal{O}_U$, et (3) les applications $\mathcal{K}_n^\bullet \to \mathcal{K}_{n + 1}^\bullet$ sont des injections scindées terme à terme dont les conoyaux sont des sommes directes de
$\mathcal{O}_X$-modules de la forme $j_{U!}\mathcal{O}_U$. En outre, l'application
$\colim \mathcal{K}_n^\bullet \to \mathcal{G}^\bullet$ est un quasi-isomorphisme.

\end{lemma}

\begin{proof}
L'existence du diagramme et les propriétés (1), (2), (3) résultent immédiatement de Modules, lemme \ref{modules-lemma-module-quotient-flat}, et de Catégories dérivées, lemme \ref{derived-lemma-special-direct-system}. L'application induite $\colim \mathcal{K}_n^\bullet \to \mathcal{G}^\bullet$ est un quasi-isomorphisme parce que les colimites filtrantes sont exactes.
\end{proof}

\begin{lemma}

\label{lemma-K-flat-resolution}

Soit $(X, \mathcal{O}_X)$ un espace annelé. Pour tout complexe $\mathcal{G}^\bullet$, il existe un complexe $K$-plat
$\mathcal{K}^\bullet$ dont les termes sont des $\mathcal{O}_X$-modules plats, ainsi qu'un quasi-isomorphisme $\mathcal{K}^\bullet \to \mathcal{G}^\bullet$
surjectif terme à terme.

\end{lemma}

\begin{proof}

Choisissez un diagramme comme dans lemme \ref{lemma-resolution-by-direct-sums-extensions-by-zero}. Chaque complexe $\mathcal{K}_n^\bullet$ est un complexe borné supérieurement de modules plats, voir Modules, lemme \ref{modules-lemma-j-shriek-flat}. D'où $\mathcal{K}_n^\bullet$ est K-plat par lemme \ref{lemma-bounded-flat-K-flat}. Ainsi $\colim \mathcal{K}_n^\bullet$ est K-plat par lemme \ref{lemma-colimit-K-flat}. La carte induite
$\colim \mathcal{K}_n^\bullet \to \mathcal{G}^\bullet$
est un quasi-isomorphisme et surjectif par construction. Propriété (3) de lemme \ref{lemma-resolution-by-direct-sums-extensions-by-zero}
montre que $\colim \mathcal{K}_n^m$ est une somme directe de modules plats et donc plats qui prouve l'affirmation finale.

\end{proof}

\begin{lemma}

\label{lemma-derived-tor-quasi-isomorphism-other-side}

Soit $(X, \mathcal{O}_X)$ un espace annelé.
Soit $\alpha : \mathcal{P}^\bullet \to \mathcal{Q}^\bullet$ un quasi-isomorphisme de complexes K-plats de $\mathcal{O}_X$-modules. Pour tout complexe $\mathcal{F}^\bullet$ de $\mathcal{O}_X$-modules, l'application induite
$$
\text{Tot}(\text{id}_{\mathcal{F}^\bullet} \otimes \alpha) :
\text{Tot}(\mathcal{F}^\bullet \otimes_{\mathcal{O}_X} \mathcal{P}^\bullet)
\longrightarrow
\text{Tot}(\mathcal{F}^\bullet \otimes_{\mathcal{O}_X} \mathcal{Q}^\bullet)
$$
est un quasi-isomorphisme.

\end{lemma}

\begin{proof}

Choisissez un quasi-isomorphisme $\mathcal{K}^\bullet \to \mathcal{F}^\bullet$
où $\mathcal{K}^\bullet$ est un complexe K-plat; voir le lemme \ref{lemma-K-flat-resolution}. Considérons le diagramme commutatif
$$
\xymatrix{
\text{Tot}(\mathcal{K}^\bullet
\otimes_{\mathcal{O}_X} \mathcal{P}^\bullet) \ar[r] \ar[d] &
\text{Tot}(\mathcal{K}^\bullet
\otimes_{\mathcal{O}_X} \mathcal{Q}^\bullet) \ar[d] \\
\text{Tot}(\mathcal{F}^\bullet
\otimes_{\mathcal{O}_X} \mathcal{P}^\bullet) \ar[r] &
\text{Tot}(\mathcal{F}^\bullet
\otimes_{\mathcal{O}_X} \mathcal{Q}^\bullet)
}
$$
Le résultat découle du lemme \ref{lemma-K-flat-quasi-isomorphism} :
les flèches verticales et la flèche horizontale supérieure sont des quasi-isomorphismes.

\end{proof}

\noindent
Soit $(X, \mathcal{O}_X)$ un espace annelé et soit $\mathcal{F}^\bullet$ un objet de $D(\mathcal{O}_X)$. Choisissons une résolution K-plate $\mathcal{K}^\bullet \to \mathcal{F}^\bullet$; voir le lemme \ref{lemma-K-flat-resolution}. Le lemme \ref{lemma-derived-tor-exact} fournit un foncteur exact entre catégories triangulées $$
K(\mathcal{O}_X)
\longrightarrow
K(\mathcal{O}_X),
\quad
\mathcal{G}^\bullet
\longmapsto
\text{Tot}(\mathcal{G}^\bullet \otimes_{\mathcal{O}_X} \mathcal{K}^\bullet)
$$ D'après le lemme \ref{lemma-K-flat-quasi-isomorphism}, ce foncteur induit un foncteur $D(\mathcal{O}_X) \to D(\mathcal{O}_X)$, puisque $D(\mathcal{O}_X)$ est la localisation de $K(\mathcal{O}_X)$ relativement aux quasi-isomorphismes. Comme la catégorie des résolutions $K$-plates de $\mathcal{F}^\bullet$ est cofiltrante et compte tenu du lemme \ref{lemma-derived-tor-quasi-isomorphism-other-side}, le foncteur obtenu, à isomorphisme près, ne dépend pas du choix de $\mathcal{K}^\bullet$.

\begin{definition}

\label{definition-derived-tor}
Soit $(X, \mathcal{O}_X)$ un espace annelé. Soit $\mathcal{F}^\bullet$ un objet de $D(\mathcal{O}_X)$. Le {\it produit tensoriel dérivé} $$
- \otimes_{\mathcal{O}_X}^{\mathbf{L}} \mathcal{F}^\bullet :
D(\mathcal{O}_X)
\longrightarrow
D(\mathcal{O}_X)
$$ est le foncteur exact des catégories triangulées décrites ci-dessus.
\end{definition}

\noindent
Il ressort clairement de nos constructions explicites qu'il y a un isomorphisme canonique $$
\mathcal{F}^\bullet \otimes_{\mathcal{O}_X}^{\mathbf{L}} \mathcal{G}^\bullet
\cong
\mathcal{G}^\bullet \otimes_{\mathcal{O}_X}^{\mathbf{L}} \mathcal{F}^\bullet
$$ pour $\mathcal{G}^\bullet$ et $\mathcal{F}^\bullet$ dans $D(\mathcal{O}_X)$. Ici nous utilisons les règles des signes comme indiqué dans Compléments sur l'algèbre, section \ref{more-algebra-section-sign-rules}. Par conséquent, lorsque nous écrivons $\mathcal{F}^\bullet \otimes_{\mathcal{O}_X}^{\mathbf{L}} \mathcal{G}^\bullet$ nous serons généralement agnostiques sur quelle variable nous utilisons pour définir le produit tensoriel dérivé avec.

\begin{definition}

\label{definition-tor}
Soit $(X, \mathcal{O}_X)$ un espace annelé. Soit $\mathcal{F}$, $\mathcal{G}$ $\mathcal{O}_X$-modules. Les {\it Tor} de $\mathcal{F}$ et $\mathcal{G}$ sont définis par la formule $$
\text{Tor}_p^{\mathcal{O}_X}(\mathcal{F}, \mathcal{G}) =
H^{-p}(\mathcal{F} \otimes_{\mathcal{O}_X}^\mathbf{L} \mathcal{G})
$$ avec produit tensoriel dérivé comme défini ci-dessus.
\end{definition}

\noindent
Cette définition implique que, pour toute suite exacte courte de $\mathcal{O}_X$-modules $0 \to \mathcal{F}_1 \to \mathcal{F}_2 \to \mathcal{F}_3 \to 0$, nous avons une suite exacte longue de cohomologie $$
\xymatrix{
\mathcal{F}_1 \otimes_{\mathcal{O}_X} \mathcal{G} \ar[r] &
\mathcal{F}_2 \otimes_{\mathcal{O}_X} \mathcal{G} \ar[r] &
\mathcal{F}_3 \otimes_{\mathcal{O}_X} \mathcal{G} \ar[r] & 0 \\
\text{Tor}_1^{\mathcal{O}_X}(\mathcal{F}_1, \mathcal{G}) \ar[r] &
\text{Tor}_1^{\mathcal{O}_X}(\mathcal{F}_2, \mathcal{G}) \ar[r] &
\text{Tor}_1^{\mathcal{O}_X}(\mathcal{F}_3, \mathcal{G}) \ar[ull]
}
$$ pour chaque $\mathcal{O}_X$-module $\mathcal{G}$. Ceci sera appelé la suite exacte longue de $\text{Tor}$ associée à la situation.

\begin{lemma}

\label{lemma-flat-tor-zero}

\begin{slogan}

Tor mesure l'écart de planéité.

\end{slogan}

Soit $(X, \mathcal{O}_X)$ un espace annelé. Soit $\mathcal{F}$ un $\mathcal{O}_X$-module. Les conditions suivantes sont équivalentes :

\begin{enumerate}

\item  $\mathcal{F}$ est un $\mathcal{O}_X$-module plat, et
\item  $\text{Tor}_1^{\mathcal{O}_X}(\mathcal{F}, \mathcal{G}) = 0$
pour tout $\mathcal{O}_X$-module $\mathcal{G}$.

\end{enumerate}

\end{lemma}

\begin{proof}
Si $\mathcal{F}$ est plat, alors $\mathcal{F} \otimes_{\mathcal{O}_X} -$ est un foncteur exact et ses foncteurs dérivés s'annulent. Réciproquement, supposons (2). Si $\mathcal{G} \to \mathcal{H}$ est injectif, de conoyau $\mathcal{Q}$, la suite exacte longue de $\text{Tor}$ montre que le noyau de $\mathcal{F} \otimes_{\mathcal{O}_X} \mathcal{G} \to
\mathcal{F} \otimes_{\mathcal{O}_X} \mathcal{H}$ est un quotient de $\text{Tor}_1^{\mathcal{O}_X}(\mathcal{F}, \mathcal{Q})$, qui est nul par hypothèse. Par conséquent, $\mathcal{F}$ est plat.
\end{proof}

\begin{lemma}

\label{lemma-factor-through-K-flat}

Soit $(X, \mathcal{O}_X)$ un espace annelé. Soit $a : \mathcal{K}^\bullet \to \mathcal{L}^\bullet$ un morphisme de complexes de $\mathcal{O}_X$-modules. Si $\mathcal{K}^\bullet$ est K-plat, alors il existe un complexe $\mathcal{N}^\bullet$ et des morphismes de complexes
$b : \mathcal{K}^\bullet \to \mathcal{N}^\bullet$
et $c : \mathcal{N}^\bullet \to \mathcal{L}^\bullet$ de telle manière que

\begin{enumerate}

\item  $\mathcal{N}^\bullet$ est K-plat,
\item  $c$ est un quasi-isomorphisme,
\item  $a$ est homotopique à $c \circ b$.

\end{enumerate}

Si les termes de $\mathcal{K}^\bullet$ sont plats, alors nous pouvons choisir
$\mathcal{N}^\bullet$, $b$, $c$
tel que la même chose est vraie pour $\mathcal{N}^\bullet$.

\end{lemma}

\begin{proof}

Nous utiliserons que la catégorie homotopie $K(\textit{Mod}(\mathcal{O}_X))$
est une catégorie triangulée, voir Catégories dérivées, proposition
\ref{derived-proposition-homotopy-category-triangulated}. Choisissez un triangle distingué
$\mathcal{K}^\bullet \to \mathcal{L}^\bullet \to
\mathcal{C}^\bullet \to \mathcal{K}^\bullet[1]$. Choisissez un quasi-isomorphisme $\mathcal{M}^\bullet \to \mathcal{C}^\bullet$ avec
$\mathcal{M}^\bullet$ K-plat avec des termes plats, voir lemme \ref{lemma-K-flat-resolution}. Par les axiomes des catégories triangulées, nous pouvons correspondre à la composition
$\mathcal{M}^\bullet \to \mathcal{C}^\bullet \to \mathcal{K}^\bullet[1]$
dans un triangle distingué
$\mathcal{K}^\bullet \to \mathcal{N}^\bullet \to
\mathcal{M}^\bullet \to \mathcal{K}^\bullet[1]$. D'après le lemme \ref{lemma-K-flat-two-out-of-three} Nous voyons que
$\mathcal{N}^\bullet$ est K-plat. En utilisant de nouveau les axiomes des catégories triangulées, nous pouvons choisir une application $\mathcal{N}^\bullet \to \mathcal{L}^\bullet$ qui définit un morphisme entre les triangles distingués suivants
$$
\xymatrix{
\mathcal{K}^\bullet \ar[r] \ar[d] &
\mathcal{N}^\bullet \ar[r] \ar[d] &
\mathcal{M}^\bullet \ar[r] \ar[d] &
\mathcal{K}^\bullet[1] \ar[d] \\
\mathcal{K}^\bullet \ar[r] &
\mathcal{L}^\bullet \ar[r] &
\mathcal{C}^\bullet \ar[r] &
\mathcal{K}^\bullet[1]
}
$$
Puisque deux des trois flèches sont quasi-isomorphismes, ainsi est la troisième flèche $\mathcal{N}^\bullet \to \mathcal{L}^\bullet$
par les suites exécute des longues de cohomologie associées à ces triangles distingués (ou vous pouvez regarder l'image de ce diagramme en $D(\mathcal{O}_X)$ et utiliser Derived Catégories, lemme \ref{derived-lemma-third-isomorphism-triangle}
si vous le souhaitez). Ceci termine la preuve de (1), (2), et (3). Pour prouver l'affirmation finale, nous pouvons choisir $\mathcal{N}^\bullet$
de telle manière que $\mathcal{N}^n \cong \mathcal{M}^n \oplus \mathcal{K}^n$, voir Catégories dérivées, lemme
\ref{derived-lemma-improve-distinguished-triangle-homotopy}. Nous obtenons donc la planéité désirée si les termes de $\mathcal{K}^\bullet$ Ils sont plats.

\end{proof}











\section{Image inverse dérivée}

\label{section-derived-pullback}

\noindent
Soit $f : (X, \mathcal{O}_X) \to (Y, \mathcal{O}_Y)$ un morphisme d'espaces annelés. Les résolutions K-plates permettent de définir le foncteur image inverse dérivée $$
Lf^* : D(\mathcal{O}_Y) \to D(\mathcal{O}_X)
$$ Plus précisément, pour tout complexe de $\mathcal{O}_Y$-modules $\mathcal{G}^\bullet$, choisissons une résolution K-plate $\mathcal{K}^\bullet \to \mathcal{G}^\bullet$ et posons $Lf^*\mathcal{G}^\bullet = f^*\mathcal{K}^\bullet$. Les lemmes \ref{lemma-pullback-K-flat}, \ref{lemma-K-flat-resolution} et \ref{lemma-derived-tor-quasi-isomorphism-other-side} montrent que cette définition est indépendante des choix. Pour régler tous les détails, il est toutefois plus commode d'employer une théorie générale.

\begin{lemma}

\label{lemma-derived-base-change}
La construction ci-dessus est indépendante des choix et définit un foncteur exact des catégories triangulées $Lf^* : D(\mathcal{O}_Y) \to D(\mathcal{O}_X)$.
\end{lemma}

\begin{proof}

Pour cela, nous utilisons la théorie générale développée dans Catégories dérivées, section \ref{derived-section-derived-functors}. Posons $\mathcal{D} = K(\mathcal{O}_Y)$ et $\mathcal{D}' = D(\mathcal{O}_X)$. Notons $F : \mathcal{D} \to \mathcal{D}'$ le foncteur exact entre catégories triangulées défini par la règle
$F(\mathcal{G}^\bullet) = f^*\mathcal{G}^\bullet$. Soit $S$ l'ensemble des quasi-isomorphismes de
$\mathcal{D} = K(\mathcal{O}_Y)$. On obtient une situation comme dans Catégories dérivées, situation \ref{derived-situation-derived-functor},
de sorte que Catégories dérivées, définition
\ref{derived-definition-right-derived-functor-defined}
est applicable. Nous affirmons que $LF$ est défini partout. Cela résulte de Catégories dérivées, lemme \ref{derived-lemma-find-existence-computes},
en prenant pour $\mathcal{P} \subset \Ob(\mathcal{D})$ la collection des complexes $K$-plats : la condition (1) résulte du lemme \ref{lemma-K-flat-resolution};
pour vérifier (2), il faut montrer que, pour tout quasi-isomorphisme
$\mathcal{K}_1^\bullet  \to \mathcal{K}_2^\bullet$ entre complexes K-plats de $\mathcal{O}_Y$-modules, l'application
$f^*\mathcal{K}_1^\bullet  \to f^*\mathcal{K}_2^\bullet$ est un quasi-isomorphisme. Écrivons cette application sous la forme
$$
f^{-1}\mathcal{K}_1^\bullet \otimes_{f^{-1}\mathcal{O}_Y} \mathcal{O}_X
\longrightarrow
f^{-1}\mathcal{K}_2^\bullet \otimes_{f^{-1}\mathcal{O}_Y} \mathcal{O}_X
$$
Le foncteur $f^{-1}$ est exact; l'application
$f^{-1}\mathcal{K}_1^\bullet  \to f^{-1}\mathcal{K}_2^\bullet$ est donc un quasi-isomorphisme. Le lemme \ref{lemma-pullback-K-flat},
appliqué au morphisme $(X, f^{-1}\mathcal{O}_Y) \to (Y, \mathcal{O}_Y)$
montre que les complexes $f^{-1}\mathcal{K}_1^\bullet$ et $f^{-1}\mathcal{K}_2^\bullet$
sont des complexes K-plats de $f^{-1}\mathcal{O}_Y$-modules. Le lemme \ref{lemma-derived-tor-quasi-isomorphism-other-side}
garantit que l'application affichée est un quasi-isomorphisme.
$$
LF :
D(\mathcal{O}_Y) = S^{-1}\mathcal{D}
\longrightarrow
\mathcal{D}' = D(\mathcal{O}_X)
$$
Voir Catégories dérivées, Équation (\ref{derived-equation-everywhere}). Enfin, les catégories dérivées, lemme \ref{derived-lemma-find-existence-computes}
garantit également que
$LF(\mathcal{K}^\bullet) = F(\mathcal{K}^\bullet) = f^*\mathcal{K}^\bullet$
lorsque $\mathcal{K}^\bullet$ est K-plat, c'est-à-dire, $Lf^* = LF$ est effectivement calculé de la manière décrite ci-dessus.

\end{proof}

\begin{lemma}

\label{lemma-derived-pullback-composition}
Soient $f : X \to Y$ et $g : Y \to Z$ des morphismes d'espaces annelés. Alors $Lf^* \circ Lg^* = L(g \circ f)^*$ comme foncteurs $D(\mathcal{O}_Z) \to D(\mathcal{O}_X)$.
\end{lemma}

\begin{proof}
Soit $E$ un objet de $D(\mathcal{O}_Z)$. Par construction, $Lg^*E$ se calcule en choisissant un complexe K-plat $\mathcal{K}^\bullet$ représentant $E$ sur $Z$ et en posant $Lg^*E = g^*\mathcal{K}^\bullet$. D'après le lemme \ref{lemma-pullback-K-flat}, $g^*\mathcal{K}^\bullet$ est K-plat sur $Y$. Ainsi, $Lf^*Lg^*E$ est représenté par $f^*g^*\mathcal{K}^\bullet = (g \circ f)^*\mathcal{K}^\bullet$, qui représente également $L(g \circ f)^*E$.
\end{proof}

\begin{lemma}

\label{lemma-pullback-tensor-product}
Soit $f : (X, \mathcal{O}_X) \to (Y, \mathcal{O}_Y)$ un morphisme d'espaces annelés. Il existe un isomorphisme canonique bifonctoriel $$
Lf^*(
\mathcal{F}^\bullet \otimes_{\mathcal{O}_Y}^{\mathbf{L}} \mathcal{G}^\bullet
) =
Lf^*\mathcal{F}^\bullet 
\otimes_{\mathcal{O}_X}^{\mathbf{L}}
Lf^*\mathcal{G}^\bullet 
$$ pour $\mathcal{F}^\bullet, \mathcal{G}^\bullet \in \Ob(D(\mathcal{O}_Y))$.
\end{lemma}

\begin{proof}
On peut supposer que $\mathcal{F}^\bullet$ et $\mathcal{G}^\bullet$ sont des complexes K-plats. Dans ce cas, $\mathcal{F}^\bullet \otimes_{\mathcal{O}_Y}^{\mathbf{L}} \mathcal{G}^\bullet$ est le complexe total associé au bicomplexe $\mathcal{F}^\bullet \otimes_{\mathcal{O}_Y} \mathcal{G}^\bullet$. D'après le lemme \ref{lemma-tensor-product-K-flat}, $\text{Tot}(\mathcal{F}^\bullet \otimes_{\mathcal{O}_Y} \mathcal{G}^\bullet)$ est lui aussi K-plat. L'isomorphisme du lemme provient donc de l'isomorphisme $$
\text{Tot}(f^*\mathcal{F}^\bullet \otimes_{\mathcal{O}_X}
f^*\mathcal{G}^\bullet)
\longrightarrow
f^*\text{Tot}(\mathcal{F}^\bullet \otimes_{\mathcal{O}_Y} \mathcal{G}^\bullet)
$$ dont les constituants sont les isomorphismes $f^*\mathcal{F}^p \otimes_{\mathcal{O}_X} f^*\mathcal{G}^q \to
f^*(\mathcal{F}^p \otimes_{\mathcal{O}_Y} \mathcal{G}^q)$ des modules, lemme \ref{modules-lemma-tensor-product-pullback}.
\end{proof}

\begin{lemma}

\label{lemma-variant-derived-pullback}
Soit $f : (X, \mathcal{O}_X) \to (Y, \mathcal{O}_Y)$ un morphisme d'espaces annelés. Il y a un isomorphisme bifonctionnel canonique $$
\mathcal{F}^\bullet
\otimes_{\mathcal{O}_X}^{\mathbf{L}}
Lf^*\mathcal{G}^\bullet
=
\mathcal{F}^\bullet 
\otimes_{f^{-1}\mathcal{O}_Y}^{\mathbf{L}}
f^{-1}\mathcal{G}^\bullet 
$$ pour $\mathcal{F}^\bullet$ dans $D(\mathcal{O}_X)$ et $\mathcal{G}^\bullet$ dans $D(\mathcal{O}_Y)$.
\end{lemma}

\begin{proof}
Soit $\mathcal{F}$ un $\mathcal{O}_X$-module et soit $\mathcal{G}$ un $\mathcal{O}_Y$-module. Alors $\mathcal{F} \otimes_{\mathcal{O}_X} f^*\mathcal{G} =
\mathcal{F} \otimes_{f^{-1}\mathcal{O}_Y} f^{-1}\mathcal{G}$ parce que $f^*\mathcal{G} =
\mathcal{O}_X \otimes_{f^{-1}\mathcal{O}_Y} f^{-1}\mathcal{G}$. Le lemme suit de ceci et les définitions.
\end{proof}

\begin{lemma}

\label{lemma-tensor-pull-compatibility}
Soit $f : (X, \mathcal{O}_X) \to (Y, \mathcal{O}_Y)$ un morphisme des espaces annelés. Soit $\mathcal{K}^\bullet$ et $\mathcal{M}^\bullet$ des complexes des modules $\mathcal{O}_Y$. Le schéma $$
\xymatrix{
Lf^*(\mathcal{K}^\bullet
\otimes_{\mathcal{O}_Y}^\mathbf{L}
\mathcal{M}^\bullet) \ar[r] \ar[d] &
Lf^*\text{Tot}(\mathcal{K}^\bullet
\otimes_{\mathcal{O}_Y}
\mathcal{M}^\bullet) \ar[d] \\
Lf^*\mathcal{K}^\bullet \otimes_{\mathcal{O}_X}^\mathbf{L}
Lf^*\mathcal{M}^\bullet \ar[d] &
f^*\text{Tot}(\mathcal{K}^\bullet
\otimes_{\mathcal{O}_Y}
\mathcal{M}^\bullet) \ar[d] \\
f^*\mathcal{K}^\bullet \otimes_{\mathcal{O}_X}^\mathbf{L}
f^*\mathcal{M}^\bullet \ar[r] &
\text{Tot}(f^*\mathcal{K}^\bullet \otimes_{\mathcal{O}_X}
f^*\mathcal{M}^\bullet)
}
$$ se déplace.
\end{lemma}

\begin{proof}
Nous utiliserons l'existence de résolutions K-plates comme dans le lemme \ref{lemma-pullback-K-flat}. Si nous choisissons ces résolutions $\mathcal{P}^\bullet \to \mathcal{K}^\bullet$ et $\mathcal{Q}^\bullet \to \mathcal{M}^\bullet$, alors nous voyons que $$
\xymatrix{
Lf^*\text{Tot}(\mathcal{P}^\bullet
\otimes_{\mathcal{O}_Y}
\mathcal{Q}^\bullet) \ar[r] \ar[d] &
Lf^*\text{Tot}(\mathcal{K}^\bullet
\otimes_{\mathcal{O}_Y}
\mathcal{M}^\bullet) \ar[d] \\
f^*\text{Tot}(\mathcal{P}^\bullet
\otimes_{\mathcal{O}_Y}
\mathcal{Q}^\bullet) \ar[d] \ar[r] &
f^*\text{Tot}(\mathcal{K}^\bullet
\otimes_{\mathcal{O}_Y}
\mathcal{M}^\bullet) \ar[d] \\
\text{Tot}(f^*\mathcal{P}^\bullet \otimes_{\mathcal{O}_X}
f^*\mathcal{Q}^\bullet) \ar[r] &
\text{Tot}(f^*\mathcal{K}^\bullet \otimes_{\mathcal{O}_X}
f^*\mathcal{M}^\bullet)
}
$$ commute. Or le côté gauche de ce diagramme coïncide avec celui du diagramme de l'énoncé, par notre choix de $\mathcal{P}^\bullet$ et $\mathcal{Q}^\bullet$ et par le lemme \ref{lemma-tensor-product-K-flat}.
\end{proof}









\section{Cohomologie des complexes non bornés}

\label{section-unbounded}

\noindent

Soit $(X, \mathcal{O}_X)$ un espace annelé. La catégorie $\textit{Mod}(\mathcal{O}_X)$ est une catégorie abélienne de Grothendieck : elle admet toutes les colimites, les colimites filtrantes y sont exactes et elle possède un générateur, à savoir
$$
\bigoplus\nolimits_{U \subset X\text{ open}} j_{U!}\mathcal{O}_U,
$$
voir Modules, section \ref{modules-section-kernels} et les lemmes \ref{modules-lemma-j-shriek-flat} et
\ref{modules-lemma-module-quotient-flat}. Par Injectifs, théorème
\ref{injectives-theorem-K-injective-embedding-grothendieck}
pour tout complexe $\mathcal{F}^\bullet$ de $\mathcal{O}_X$-modules, il existe un quasi-isomorphisme
$\mathcal{F}^\bullet \to \mathcal{I}^\bullet$ vers un complexe K-injectif de $\mathcal{O}_X$-modules dont tous les termes sont injectifs;
$\mathcal{O}_X$-modules; de plus, cette inclusion peut être choisie fonctoriellement en le complexe $\mathcal{F}^\bullet$. Il résulte alors de Catégories dérivées, lemme \ref{derived-lemma-enough-K-injectives-implies},
que

\begin{enumerate}

\item tout foncteur exact $F : K(\textit{Mod}(\mathcal{O}_X)) \to \mathcal{D}$
à valeurs dans une catégorie triangulée $\mathcal{D}$ admet un foncteur dérivé à droite
$RF : D(\mathcal{O}_X) \to \mathcal{D}$,
\item pour tout foncteur additif
$F : \textit{Mod}(\mathcal{O}_X) \to \mathcal{A}$
à valeurs dans une catégorie abélienne $\mathcal{A}$, on considère le foncteur exact
$F : K(\textit{Mod}(\mathcal{O}_X)) \to K(\mathcal{A})$ induit par $F$
et nous obtenons un foncteur dérivé à droite
$RF : D(\mathcal{O}_X) \to D(\mathcal{A})$.

\end{enumerate}

Par construction, on a $RF(\mathcal{F}^\bullet) = F(\mathcal{I}^\bullet)$,
où $\mathcal{F}^\bullet \to \mathcal{I}^\bullet$ est comme ci-dessus.

\medskip\noindent
Voici quelques exemples de ce qui précède :
\begin{enumerate}
\item Le foncteur $\Gamma(X, -) : \textit{Mod}(\mathcal{O}_X) \to
\text{Mod}_{\Gamma(X, \mathcal{O}_X)}$ donne lieu à $$
R\Gamma(X, -) : D(\mathcal{O}_X) \to D(\Gamma(X, \mathcal{O}_X))
$$ Nous utiliserons la notation $H^i(X, K) = H^i(R\Gamma(X, K))$ pour la cohomologie.
\item Pour un $U \subset X$ ouvert nous considérons le foncteur $\Gamma(U, -) : \textit{Mod}(\mathcal{O}_X) \to
\text{Mod}_{\Gamma(U, \mathcal{O}_X)}$. Ceci donne lieu à $$
R\Gamma(U, -) : D(\mathcal{O}_X) \to D(\Gamma(U, \mathcal{O}_X))
$$ Nous utiliserons la notation $H^i(U, K) = H^i(R\Gamma(U, K))$ pour la cohomologie.
\item Pour un morphisme d'espaces annelés $f : (X, \mathcal{O}_X) \to (Y, \mathcal{O}_Y)$ nous considérons le foncteur $f_* : \textit{Mod}(\mathcal{O}_X) \to \textit{Mod}(\mathcal{O}_Y)$ qui donne lieu à l'image directe totale $$
Rf_* : D(\mathcal{O}_X) \longrightarrow D(\mathcal{O}_Y)
$$ sur des catégories dérivées non délimitées.
\end{enumerate}

\begin{lemma}

\label{lemma-adjoint}
Soit $f : (X, \mathcal{O}_X) \to (Y, \mathcal{O}_Y)$ un morphisme des espaces annelés. Le foncteur $Rf_*$ défini ci-dessus et le foncteur $Lf^*$ défini dans le lemme \ref{lemma-derived-base-change} sont adjoints: $$
\Hom_{D(\mathcal{O}_X)}(Lf^*\mathcal{G}^\bullet, \mathcal{F}^\bullet)
=
\Hom_{D(\mathcal{O}_Y)}(\mathcal{G}^\bullet, Rf_*\mathcal{F}^\bullet)
$$ bifonctoriellement en $\mathcal{F}^\bullet \in \Ob(D(\mathcal{O}_X))$ et $\mathcal{G}^\bullet \in \Ob(D(\mathcal{O}_Y))$.
\end{lemma}

\begin{proof}
Ceci résulte formellement du fait que $Rf_*$ et $Lf^*$ existent, voir Catégories dérivées, lemme \ref{derived-lemma-derived-adjoint-functors}.
\end{proof}

\begin{lemma}

\label{lemma-derived-pushforward-composition}
Soient $f : X \to Y$ et $g : Y \to Z$ des morphismes d'espaces annelés. Alors $Rg_* \circ Rf_* = R(g \circ f)_*$ comme foncteurs $D(\mathcal{O}_X) \to D(\mathcal{O}_Z)$.
\end{lemma}

\begin{proof}
Par lemme \ref{lemma-adjoint} nous voyons que $Rg_* \circ Rf_*$ est adjoint à $Lf^* \circ Lg^*$. Nous avons $Lf^* \circ Lg^* = L(g \circ f)^*$ par lemme \ref{lemma-derived-pullback-composition} et donc par unicité des foncteurs adjoints nous avons $Rg_* \circ Rf_* = R(g \circ f)_*$.
\end{proof}

\begin{remark}

\label{remark-base-change}

La construction des foncteurs dérivés $Lf^*$ et $Rf_*$
permet de construire l'application de changement de base en pleine généralité.
$$
\xymatrix{
X' \ar[r]_{g'} \ar[d]_{f'} &
X \ar[d]^f \\
S' \ar[r]^g &
S
}
$$
soit un diagramme commutatif d'espaces annelés. Soit $K$ un objet de
$D(\mathcal{O}_X)$. Il existe alors une application canonique de changement de base
$$
Lg^*Rf_*K \longrightarrow R(f')_*L(g')^*K
$$
dans $D(\mathcal{O}_{S'})$. Cette application est, par adjonction, associée à une application
$L(f')^*Lg^*Rf_*K \to L(g')^*K$
Puisque $L(f')^*Lg^* = L(g')^*Lf^*$, il s'agit de la même chose que de donner une application
$L(g')^*Lf^*Rf_*K \to L(g')^*K$
que nous pouvons prendre pour être $L(g')^*$ de l'application d'adjonction
$Lf^*Rf_*K \to K$.

\end{remark}

\begin{remark}

\label{remark-compose-base-change}

Considérez un diagramme commutatif
$$
\xymatrix{
X' \ar[r]_k \ar[d]_{f'} & X \ar[d]^f \\
Y' \ar[r]^l \ar[d]_{g'} & Y \ar[d]^g \\
Z' \ar[r]^m & Z
}
$$
d'espaces annelés. Les applications de changement de base de la remarque \ref{remark-base-change}
associées aux deux carrés se composent en l'application de changement de base du rectangle extérieur. Plus précisément, la composée

\begin{align*}
Lm^* \circ R(g \circ f)_*
& =
Lm^* \circ Rg_* \circ Rf_* \\
& \to Rg'_* \circ Ll^* \circ Rf_* \\
& \to Rg'_* \circ Rf'_* \circ Lk^* \\
& = R(g' \circ f')_* \circ Lk^*
\end{align*}

est l'application de changement de base pour le rectangle. Nous omettons la vérification.

\end{remark}

\begin{remark}

\label{remark-compose-base-change-horizontal}

Considérez un diagramme commutatif
$$
\xymatrix{
X'' \ar[r]_{g'} \ar[d]_{f''} & X' \ar[r]_g \ar[d]_{f'} & X \ar[d]^f \\
Y'' \ar[r]^{h'} & Y' \ar[r]^h & Y
}
$$
d'espaces annelés. Les applications de changement de base de la remarque \ref{remark-base-change}
associées aux deux carrés se composent en l'application de changement de base du rectangle extérieur. Plus précisément, la composée

\begin{align*}
L(h \circ h')^* \circ Rf_*
& =
L(h')^* \circ Lh_* \circ Rf_* \\
& \to L(h')^* \circ Rf'_* \circ Lg^* \\
& \to Rf''_* \circ L(g')^* \circ Lg^* \\
& = Rf''_* \circ L(g \circ g')^*
\end{align*}

est l'application de changement de base pour le rectangle. Nous omettons la vérification.

\end{remark}

\begin{lemma}

\label{lemma-adjoints-push-pull-compatibility}
Soit $f : (X, \mathcal{O}_X) \to (Y, \mathcal{O}_Y)$ un morphisme d'espaces annelés. Soit $\mathcal{K}^\bullet$ un complexe de $\mathcal{O}_X$-modules. Le diagramme $$
\xymatrix{
Lf^*f_*\mathcal{K}^\bullet \ar[r] \ar[d] &
f^*f_*\mathcal{K}^\bullet \ar[d] \\
Lf^*Rf_*\mathcal{K}^\bullet \ar[r] &
\mathcal{K}^\bullet
}
$$ provenant de $Lf^* \to f^*$ sur les complexes, $f_* \to Rf_*$ sur les complexes, et l'adjonction $Lf^* \circ Rf_* \to \text{id}$ se déplace dans $D(\mathcal{O}_X)$.
\end{lemma}

\begin{proof}

Nous utiliserons l'existence de résolutions K-plates et de résolutions K-injectives; voir le lemme \ref{lemma-pullback-K-flat}
et la discussion ci-dessus. Choisissez un quasi-isomorphisme
$\mathcal{K}^\bullet \to \mathcal{I}^\bullet$ où $\mathcal{I}^\bullet$
est K-injective comme un complexe de $\mathcal{O}_X$-modules. Choisissez un quasi-isomorphisme $\mathcal{Q}^\bullet \to f_*\mathcal{I}^\bullet$
où $\mathcal{Q}^\bullet$ est K-plat comme un complexe de
$\mathcal{O}_Y$-modules. Nous pouvons choisir un complexe K-plat de
$\mathcal{O}_Y$-modules $\mathcal{P}^\bullet$
et un diagramme de morphismes de complexes
$$
\xymatrix{
\mathcal{P}^\bullet \ar[r] \ar[d] &
f_*\mathcal{K}^\bullet \ar[d] \\
\mathcal{Q}^\bullet \ar[r] & f_*\mathcal{I}^\bullet
}
$$
commutatif à homotopie près, dont la flèche horizontale supérieure est un quasi-isomorphisme. Un tel $\mathcal{P}^\bullet$ existe parce que les quasi-isomorphismes forment un système multiplicatif dans la catégorie homotopique des complexes; on peut ensuite remplacer
$\mathcal{P}^\bullet$ par un complexe K-plat. En appliquant l'image inverse, on obtient un diagramme de morphismes de complexes
$$
\xymatrix{
f^*\mathcal{P}^\bullet \ar[r] \ar[d] &
f^*f_*\mathcal{K}^\bullet \ar[d] \ar[r] &
\mathcal{K}^\bullet \ar[d] \\
f^*\mathcal{Q}^\bullet \ar[r] &
f^*f_*\mathcal{I}^\bullet \ar[r] &
\mathcal{I}^\bullet
}
$$
Le rectangle extérieur prouve l'assertion du lemme.

\end{proof}

\begin{remark}

\label{remark-cup-product}
Soit $f : (X, \mathcal{O}_X) \to (Y, \mathcal{O}_Y)$ un morphisme d'espaces annelés. L'adjonction entre $Lf^*$ et $Rf_*$ permet de construire un produit cup relatif $$
Rf_*K \otimes_{\mathcal{O}_Y}^\mathbf{L} Rf_*L
\longrightarrow
Rf_*(K \otimes_{\mathcal{O}_X}^\mathbf{L} L)
$$ dans $D(\mathcal{O}_Y)$ pour tous $K, L$ dans $D(\mathcal{O}_X)$. Autrement dit, cette application est adjointe à une application $Lf^*(Rf_*K \otimes_{\mathcal{O}_Y}^\mathbf{L} Rf_*L) \to
K \otimes_{\mathcal{O}_X}^\mathbf{L} L$ pour laquelle nous pouvons prendre la composition de l'isomorphisme $Lf^*(Rf_*K \otimes_{\mathcal{O}_Y}^\mathbf{L} Rf_*L) =
Lf^*Rf_*K \otimes_{\mathcal{O}_X}^\mathbf{L} Lf^*Rf_*L$ (lemme \ref{lemma-pullback-tensor-product}) avec l'application $Lf^*Rf_*K \otimes_{\mathcal{O}_X}^\mathbf{L} Lf^*Rf_*L
\to K \otimes_{\mathcal{O}_X}^\mathbf{L} L$ provenant de la co-unité $Lf^* \circ Rf_* \to \text{id}$.
\end{remark}





\section{Cohomologie des complexes filtrés}

\label{section-cohomology-filtered-object}

\noindent
Les complexes filtrés de faisceaux apparaissent souvent de manière naturelle dans l'étude de la cohomologie des variétés algébriques; par exemple, le complexe de de Rham est muni de sa filtration de Hodge. Dans cette section, nous utilisons le très général lemme \ref{injectives-lemma-K-injective-embedding-filtration} d'Injectifs pour construire des suites spectrales en cohomologie et les relier à celles construites précédemment.

\begin{lemma}

\label{lemma-spectral-sequence-filtered-object}

Soit $(X, \mathcal{O}_X)$ un espace annelé. Soit $\mathcal{F}^\bullet$ un complexe filtré de $\mathcal{O}_X$-modules. Il existe une suite spectrale canonique $(E_r, \text{d}_r)_{r \geq 1}$ de
$\Gamma(X, \mathcal{O}_X)$-modules bigradués, où $d_r$ est de bidegré $(r, -r + 1)$, et
$$
E_1^{p, q} = H^{p + q}(X, \text{gr}^p\mathcal{F}^\bullet)
$$
Si, pour tout $n$, on a
$$
H^n(X, F^p\mathcal{F}^\bullet) = 0\text{ pour }p \gg 0
\quad\text{et}\quad
H^n(X, F^p\mathcal{F}^\bullet) = H^n(X, \mathcal{F}^\bullet)\text{ pour }p \ll 0
$$
alors la suite spectrale est bornée et converge vers
$H^*(X, \mathcal{F}^\bullet)$.

\end{lemma}

\begin{proof}
(Pour une preuve dans le cas d'un complexe borné inférieurement de modules munis de filtrations finies, voir la remarque ci-dessous.) Choisissons un morphisme de complexes filtrés $j : \mathcal{F}^\bullet \to \mathcal{J}^\bullet$ comme dans Injectifs, lemme \ref{injectives-lemma-K-injective-embedding-filtration}. La suite spectrale est celle d'Homologie, section \ref{homology-section-filtered-complex}, associée au complexe filtré $$
\Gamma(X, \mathcal{J}^\bullet)
\quad\text{avec}\quad
F^p\Gamma(X, \mathcal{J}^\bullet) = \Gamma(X, F^p\mathcal{J}^\bullet)
$$ Comme la cohomologie se calcule sur des représentants K-injectifs, la page $E_1$ est celle indiquée dans le lemme. Sous les hypothèses énoncées, la convergence et le caractère borné résultent d'Homologie, lemme \ref{homology-lemma-ss-converges-trivial}.
\end{proof}

\begin{remark}

\label{remark-spectral-sequence-filtered-object}

Soit $(X, \mathcal{O}_X)$ un espace annelé. Soit $\mathcal{F}^\bullet$ un complexe filtré de $\mathcal{O}_X$-modules. Si $\mathcal{F}^\bullet$
est borné inférieurement et si, pour tout $n$, la filtration de $\mathcal{F}^n$
est finie, on peut construire la suite spectrale du lemme \ref{lemma-spectral-sequence-filtered-object}
sans recourir à Injectifs, lemme
\ref{injectives-lemma-K-injective-embedding-filtration}. En effet, d'après Catégories dérivées, lemme
\ref{derived-lemma-right-resolution-by-filtered-injectives}
il y a un quasi-isomorphisme filtré
$i : \mathcal{F}^\bullet \to \mathcal{I}^\bullet$
de complexes filtrés, où $\mathcal{I}^\bullet$ est borné inférieurement, où la filtration de $\mathcal{I}^n$ est finie pour tout $n$ et où chaque $\text{gr}^p\mathcal{I}^n$ est un $\mathcal{O}_X$-module injectif. Nous prenons alors la suite spectrale associée à
$$
\Gamma(X, \mathcal{I}^\bullet)
\quad\text{avec}\quad
F^p\Gamma(X, \mathcal{I}^\bullet) = \Gamma(X, F^p\mathcal{I}^\bullet)
$$
Comme la cohomologie se calcule sur des complexes bornés inférieurement d'injectifs, la page $E_1$ est celle indiquée dans le lemme. Sous les hypothèses énoncées, la convergence et le caractère borné résultent d'Homologie, lemme \ref{homology-lemma-biregular-ss-converges}. Il s'agit en fait d'un cas particulier de la suite spectrale de Catégories dérivées, lemme \ref{derived-lemma-ss-filtered-derived}.

\end{remark}

\begin{example}

\label{example-spectral-sequence}

Soit $(X, \mathcal{O}_X)$ un espace annelé. Soit $\mathcal{F}^\bullet$ un complexe de $\mathcal{O}_X$-modules. Nous pouvons appliquer le lemme \ref{lemma-spectral-sequence-filtered-object}
avec $F^p\mathcal{F}^\bullet = \tau_{\leq -p}\mathcal{F}^\bullet$. (Si $\mathcal{F}^\bullet$ est borné inférieurement, nous pouvons utiliser la remarque \ref{remark-spectral-sequence-filtered-object}.) Nous obtenons alors une suite spectrale
$$
E_1^{p, q} = H^{p + q}(X, H^{-p}(\mathcal{F}^\bullet)[p]) =
H^{2p + q}(X, H^{-p}(\mathcal{F}^\bullet))
$$
Après renumérotation $p = -j$ et $q = i + 2j$, on voit que, pour tout
$K \in D(\mathcal{O}_X)$, il existe une suite spectrale
$(E'_r, d'_r)_{r \geq 2}$ de modules bigradués, où
$d'_r$ est de bidegré $(r, -r + 1)$, et telle que
$$
(E'_2)^{i, j} = H^i(X, H^j(K))
$$
Si $K$ est borné inférieurement, par exemple, cette suite spectrale est bornée et converge vers $H^{i + j}(X, K)$. Dans le cas borné inférieurement, elle est un exemple de la seconde suite spectrale de Catégories dérivées, lemme \ref{derived-lemma-two-ss-complex-functor}
(construit en utilisant les résolutions de Cartan-Eilenberg).

\end{example}

\begin{example}

\label{example-spectral-sequence-bis}

Soit $(X, \mathcal{O}_X)$ un espace annelé. Soit $\mathcal{F}^\bullet$ un complexe de $\mathcal{O}_X$-modules. Nous pouvons appliquer le lemme \ref{lemma-spectral-sequence-filtered-object}
avec $F^p\mathcal{F}^\bullet = \sigma_{\geq p}\mathcal{F}^\bullet$. Nous obtenons alors une suite spectrale
$$
E_1^{p, q} = H^{p + q}(X, \mathcal{F}^p[-p]) = H^q(X, \mathcal{F}^p)
$$
Si $\mathcal{F}^\bullet$ est borné inférieurement, alors

\begin{enumerate}

\item Nous pouvons utiliser la remarque \ref{remark-spectral-sequence-filtered-object}
pour construire cette suite spectrale,
\item la suite spectrale est bornée et converge vers
$H^{i + j}(X, \mathcal{F}^\bullet)$,
\item la suite spectrale est égale à la première suite spectrale de Catégories dérivées, lemme \ref{derived-lemma-two-ss-complex-functor}
(construit en utilisant les résolutions de Cartan-Eilenberg).

\end{enumerate}

\end{example}

\begin{lemma}

\label{lemma-relative-spectral-sequence-filtered-object}
Soit $f : (X, \mathcal{O}_X) \to (Y, \mathcal{O}_Y)$ un morphisme d'espaces annelés, et soit $\mathcal{F}^\bullet$ un complexe filtré de $\mathcal{O}_X$-modules. Il existe une suite spectrale canonique $(E_r, \text{d}_r)_{r \geq 1}$ de $\mathcal{O}_Y$-modules bigradués, où $d_r$ est de bidegré $(r, -r + 1)$, et $$
E_1^{p, q} = R^{p + q}f_*\text{gr}^p\mathcal{F}^\bullet
$$ Si pour chaque $n$ nous avons $$
R^nf_*F^p\mathcal{F}^\bullet = 0 \text{ pour }p \gg 0
\quad\text{et}\quad
R^nf_*F^p\mathcal{F}^\bullet = R^nf_*\mathcal{F}^\bullet \text{ pour }p \ll 0
$$, alors la suite spectrale est bornée et converge vers $Rf_*\mathcal{F}^\bullet$.
\end{lemma}

\begin{proof}

La preuve est exactement la même que celle du lemme \ref{lemma-spectral-sequence-filtered-object}.

\end{proof}





\section{Résolution de Godement}

\label{section-godement}

\noindent
Une référence est \cite{Godement}.

\medskip\noindent

Soit $(X, \mathcal{O}_X)$ un espace annelé. Notons $X_{disc}$ l'espace topologique discret ayant les mêmes points que $X$, et $f : X_{disc} \to X$
l'application continue évidente. Posons $\mathcal{O}_{X_{disc}} = f^{-1}\mathcal{O}_X$. Alors $f : (X_{disc}, \mathcal{O}_{X_{disc}}) \to (X, \mathcal{O}_X)$
est un morphisme plat d'espaces annelés. Appliquons la construction
{\it duale} de Simplicial, section \ref{simplicial-section-standard}, à la paire de foncteurs adjoints $f^*, f_*$ sur les faisceaux de modules. On obtient ainsi un objet cosimplicial augmenté
$$
\xymatrix{
\text{id} \ar[r] &
f_*f^* \ar@<1ex>[r] \ar@<-1ex>[r] &
f_*f^*f_*f^* \ar@<0ex>[l] \ar@<-2ex>[r] \ar@<0ex>[r] \ar@<2ex>[r] &
f_*f^*f_*f^*f_*f^* \ar@<1ex>[l] \ar@<-1ex>[l]
}
$$
dans la catégorie des endofoncteurs de $\textit{Mod}(\mathcal{O}_X)$;
voir Simplicial, lemme \ref{simplicial-lemma-standard-simplicial}. De plus, l'augmentation
$$
\xymatrix{
f^* \ar[r] &
f^*f_*f^* \ar@<1ex>[r] \ar@<-1ex>[r] &
f^*f_*f^*f_*f^* \ar@<0ex>[l] \ar@<-2ex>[r] \ar@<0ex>[r] \ar@<2ex>[r] &
f^*f_*f^*f_*f^*f_*f^* \ar@<1ex>[l] \ar@<-1ex>[l]
}
$$
est une équivalence d'homotopie, voir Simplicial, lemme \ref{simplicial-lemma-standard-simplicial-homotopy}.

\begin{lemma}

\label{lemma-godement-resolution}
Soit $(X, \mathcal{O}_X)$ un espace annelé. Pour tout faisceau de $\mathcal{O}_X$-modules $\mathcal{F}$, il existe une résolution $$
0 \to
\mathcal{F} \to
f_*f^*\mathcal{F} \to
f_*f^*f_*f^*\mathcal{F} \to
f_*f^*f_*f^*f_*f^*\mathcal{F} \to \ldots
$$ fonctorielle en $\mathcal{F}$, telle que chaque terme $f_*f^* \ldots f_*f^*\mathcal{F}$ est un $\mathcal{O}_X$-module flasque et que, pour tout $x \in X$, l'application $$
\mathcal{F}_x[0] \to \Big(
(f_*f^*\mathcal{F})_x \to
(f_*f^*f_*f^*\mathcal{F})_x \to
(f_*f^*f_*f^*f_*f^*\mathcal{F})_x \to \ldots
\Big)
$$ est une équivalence d'homotopie dans la catégorie des complexes de $\mathcal{O}_{X, x}$-modules.
\end{lemma}

\begin{proof}

Le complexe $f_*f^*\mathcal{F} \to  f_*f^*f_*f^*\mathcal{F} \to
f_*f^*f_*f^*f_*f^*\mathcal{F} \to \ldots$ est le complexe associé à l'objet cosimplicial avec les termes
$f_*f^*\mathcal{F}, f_*f^*f_*f^*\mathcal{F},
f_*f^*f_*f^*f_*f^*\mathcal{F}, \ldots$ décrit ci-dessus; voir Simplicial, section \ref{simplicial-section-dold-kan-cosimplicial}. L'augmentation fournit l'application $\mathcal{F} \to f_*f^*\mathcal{F}$
indiquée. Pour tout faisceau abélien $\mathcal{H}$ sur $X_{disc}$, l'image directe $f_*\mathcal{H}$ est flasque, car $X_{disc}$
est discret et l'image directe d'un faisceau flasque est flasque. La même observation vaut pour les $\mathcal{O}_X$-modules.

\medskip\noindent

Si $x \in X_{disc} = X$ est un point, alors $(f^*\mathcal{G})_x = \mathcal{G}_x$
pour tout $\mathcal{O}_X$-module $\mathcal{G}$. Ainsi, $f^*$ est un foncteur exact, et un complexe de $\mathcal{O}_X$-modules
$\mathcal{G}_1 \to \mathcal{G}_2 \to \mathcal{G}_3$
est exact si et seulement si
$f^*\mathcal{G}_1 \to f^*\mathcal{G}_2 \to f^*\mathcal{G}_3$
est exact (voir Modules, lemme \ref{modules-lemma-abelian}). Le résultat mentionné au début de cette section montre que l'image inverse par $f^*$ fournit une équivalence d'homotopie de l'objet cosimplicial constant $f^*\mathcal{F}$ vers l'objet cosimplicial dont les termes sont
$f_*f^*\mathcal{F}, f_*f^*f_*f^*\mathcal{F},
f_*f^*f_*f^*f_*f^*\mathcal{F}, \ldots$. Par Simplicial, lemme \ref{simplicial-lemma-homotopy-equivalence-s-Q}
nous obtenons que
$$
f^*\mathcal{F}[0] \to \Big(
f^*f_*f^*\mathcal{F} \to
f^*f_*f^*f_*f^*\mathcal{F} \to
f^*f_*f^*f_*f^*f_*f^*\mathcal{F} \to \ldots
\Big)
$$
est une équivalence d'homotopie. Les deux dernières assertions du lemme en résultent immédiatement.

\end{proof}

\begin{lemma}

\label{lemma-godement-resolution-bounded-below}
Soit $(X, \mathcal{O}_X)$ un espace annelé. Soit $\mathcal{F}^\bullet$ un complexe borné inférieurement de $\mathcal{O}_X$-modules. Il existe un quasi-isomorphisme $\mathcal{F}^\bullet \to \mathcal{G}^\bullet$, où $\mathcal{G}^\bullet$ est un complexe borné inférieurement de $\mathcal{O}_X$-modules flasques et où, pour tout $x \in X$, l'application $\mathcal{F}^\bullet_x \to \mathcal{G}^\bullet_x$ est une équivalence d'homotopie dans la catégorie des complexes de $\mathcal{O}_{X, x}$-modules.
\end{lemma}

\begin{proof}

Soit $\mathcal{A}$ la catégorie de complexes de $\mathcal{O}_X$-modules et soit $\mathcal{B}$ la catégorie de complexes de $\mathcal{O}_X$-modules. Nous pouvons alors appliquer la discussion ci-dessus aux foncteurs
$f^*$ et $f_*$ entre $\mathcal{A}$ et $\mathcal{B}$. En raisonnant exactement comme dans la preuve du lemme \ref{lemma-godement-resolution},
nous obtenons une résolution
$$
0 \to
\mathcal{F}^\bullet \to
f_*f^*\mathcal{F}^\bullet \to
f_*f^*f_*f^*\mathcal{F}^\bullet \to
f_*f^*f_*f^*f_*f^*\mathcal{F}^\bullet \to \ldots
$$
dans la catégorie abélienne $\mathcal{A}$, de telle sorte que chaque terme de chaque
$f_*f^*\ldots f_*f^*\mathcal{F}^\bullet$ soit un
$\mathcal{O}_X$-module flasque et que, pour tout $x \in X$, l'application
$$
\mathcal{F}^\bullet_x[0] \to \Big(
(f_*f^*\mathcal{F}^\bullet)_x \to
(f_*f^*f_*f^*\mathcal{F}^\bullet)_x \to
(f_*f^*f_*f^*f_*f^*\mathcal{F}^\bullet)_x \to \ldots
\Big)
$$
est une équivalence d'homotopie dans la catégorie des complexes de complexes de $\mathcal{O}_{X, x}$-modules. Comme un complexe de complexes est un bicomplexe, nous pouvons considérer l'application induite
$$
\mathcal{F}^\bullet \to
\mathcal{G}^\bullet =
\text{Tot}(
f_*f^*\mathcal{F}^\bullet \to
f_*f^*f_*f^*\mathcal{F}^\bullet \to
f_*f^*f_*f^*f_*f^*\mathcal{F}^\bullet \to \ldots
)
$$
Comme le complexe $\mathcal{F}^\bullet$ est borné inférieurement, il en est de même de $\mathcal{G}^\bullet$; en fait, chaque terme de $\mathcal{G}^\bullet$ est une somme directe finie de termes des complexes $f_*f^*\ldots f_*f^*\mathcal{F}^\bullet$
et est donc flasque. La dernière assertion du lemme résulte alors d'Homologie, lemme \ref{homology-lemma-homotopy-complex-complexes}. Comme cela montre en particulier que
$\mathcal{F}^\bullet \to \mathcal{G}^\bullet$
est un quasi-isomorphisme, la preuve est complète.

\end{proof}
\section{Cup-produit}

\label{section-cup-product}

\noindent

Soit $(X, \mathcal{O}_X)$ un espace annelé. Soient $K, M$ des objets de $D(\mathcal{O}_X)$. Posons $A = \Gamma(X, \mathcal{O}_X)$. Le cup-produit (global) dans ce cadre est un morphisme
$$
\mu :
R\Gamma(X, K) \otimes_A^\mathbf{L} R\Gamma(X, M)
\longrightarrow
R\Gamma(X, K \otimes_{\mathcal{O}_X}^\mathbf{L} M)
$$
dans $D(A)$. Nous le définissons comme le cup-produit relatif associé au morphisme d'espaces annelés $f : (X, \mathcal{O}_X) \to (pt, A)$
de la remarque \ref{remark-cup-product}, via $D(pt, A) = D(A)$. Ce morphisme définit en particulier les accouplements
$$
\cup :
H^i(X, K) \times H^j(X, M)
\longrightarrow
H^{i + j}(X, K \otimes_{\mathcal{O}_X}^\mathbf{L} M)
$$
Plus précisément, étant donnés $\xi \in H^i(X, K) = H^i(R\Gamma(X, K))$ et
$\eta \in H^j(X, M) = H^j(R\Gamma(X, M))$, nous pouvons d'abord les ``tensoriser'' pour obtenir un élément $\xi \otimes \eta$ de
$H^{i + j}(R\Gamma(X, K) \otimes_A^\mathbf{L} R\Gamma(X, M))$; voir Compléments sur l'algèbre, section \ref{more-algebra-section-products-tor}. Nous pouvons ensuite appliquer $\mu$ pour obtenir l'élément souhaité
$\xi \cup \eta = \mu(\xi \otimes \eta)$
de $H^{i + j}(X, K \otimes_{\mathcal{O}_X}^\mathbf{L} M)$.

\medskip\noindent
Voici une autre manière d'interpréter le cup-produit de $\xi$ et $\eta$. En effet, nous pouvons écrire $$
H^i(R\Gamma(X, K)) = \Hom_{D(\mathcal{O}_X)}(\mathcal{O}_X[-i], K)
\quad\text{et}\quad
H^j(R\Gamma(X, M)) = \Hom_{D(\mathcal{O}_X)}(\mathcal{O}_X[-j], M)
$$ puisque $\Hom(\mathcal{O}_X, -) = \Gamma(X, -)$. Ainsi, $\xi$ et $\eta$ sont la ``même'' chose que les morphismes $$
\tilde \xi : \mathcal{O}_X[-i] \to K
\quad\text{et}\quad
\tilde \eta : \mathcal{O}_X[-j] \to M
$$ En combinant ceci avec la fonctorialité du produit tensoriel dérivé, nous obtenons $$
\mathcal{O}_X[-i - j] =
\mathcal{O}_X[-i] \otimes_{\mathcal{O}_X}^\mathbf{L} \mathcal{O}_X[-j]
\xrightarrow{\tilde \xi \otimes \tilde \eta}
K \otimes_{\mathcal{O}_X}^\mathbf{L} M
$$ qui, par le même raisonnement que ci-dessus, est un élément de $H^{i + j}(X, K \otimes_{\mathcal{O}_X}^\mathbf{L} M)$.

\begin{lemma}

\label{lemma-second-cup-equals-first}

Cette construction donne le cup-produit.

\end{lemma}

\begin{proof}

Avec $f : (X, \mathcal{O}_X) \to (pt, A)$ comme ci-dessus nous avons
$Rf_*(-) = R\Gamma(X, -)$ et notre morphisme $\mu$ est adjoint du morphisme
$$
Lf^*(Rf_*K \otimes_A^\mathbf{L} Rf_*M) =
Lf^*Rf_*K \otimes_{\mathcal{O}_X}^\mathbf{L} Lf^*Rf_*M
\xrightarrow{\epsilon_K \otimes \epsilon_M}
K \otimes_{\mathcal{O}_X}^\mathbf{L} M
$$
où $\epsilon$ est la counité de l'adjonction entre
$Lf^*$ et $Rf_*$. Si nous considérons $\xi$ et $\eta$ comme des morphismes $\xi : A[-i] \to R\Gamma(X, K)$
et $\eta : A[-j] \to R\Gamma(X, M)$, alors le tenseur $\xi \otimes \eta$ correspond au morphisme\footnote{Un signe est caché ici : l'égalité est définie par la composition
$$
A[-i - j] \to (A \otimes_A^\mathbf{L} A)[-i - j] \to
A[-i] \otimes_A^\mathbf{L} A[-j]
$$
où, dans la deuxième étape, nous utilisons l'identification de Compléments sur l'algèbre, point (\ref{more-algebra-item-shift-tensor}), qui fait en principe intervenir un signe. Dans ce cas, ce signe vaut toutefois $+1$ par convention; même s'il ne valait pas
$+1$, cela n'aurait aucune importance puisque nous avons utilisé le même signe dans l'identification
$\mathcal{O}_X[-i - j] =
\mathcal{O}_X[-i] \otimes_{\mathcal{O}_X}^\mathbf{L} \mathcal{O}_X[-j]$.}
$$
A[-i - j] = A[-i] \otimes_A^\mathbf{L} A[-j]
\xrightarrow{\xi \otimes \eta}
R\Gamma(X, K) \otimes_A^\mathbf{L} R\Gamma(X, M)
$$
Par définition, le cup-produit $\xi \cup \eta$ est le morphisme
$A[-i - j] \to R\Gamma(X, K \otimes_{\mathcal{O}_X}^\mathbf{L} M)$
dont l'adjoint est
$$
(\epsilon_K \otimes \epsilon_M) \circ Lf^*(\xi \otimes \eta) =
(\epsilon_K \circ Lf^*\xi) \otimes (\epsilon_M \circ Lf^*\eta)
$$
Cependant, il est facile de voir que
$\epsilon_K \circ Lf^*\xi = \tilde \xi$ et
$\epsilon_M \circ Lf^*\eta = \tilde \eta$. Nous concluons que $\widetilde{\xi \cup \eta} = \tilde \xi \otimes \tilde \eta$,
ce qui donne l'identification souhaitée.

\end{proof}

\begin{remark}

\label{remark-cup-with-element-map-total-cohomology}
Soit $(X, \mathcal{O}_X)$ un espace annelé. Soient $K, M$ des objets de $D(\mathcal{O}_X)$. Posons $A = \Gamma(X, \mathcal{O}_X)$. Pour $\xi \in H^i(X, K)$, nous obtenons le morphisme associé $$
\xi = ``\xi \cup -'' :
R\Gamma(X, M)[-i]
\to
R\Gamma(X, K \otimes_{\mathcal{O}_X}^\mathbf{L} M)
$$ en représentant $\xi$ par un morphisme $\xi : A[-i] \to R\Gamma(X, K)$ comme dans la preuve du lemme \ref{lemma-second-cup-equals-first}, puis en utilisant la composition $$
R\Gamma(X, M)[-i] = A[-i] \otimes_A^\mathbf{L} R\Gamma(X, M)
\xrightarrow{\xi \otimes 1}
R\Gamma(X, K) \otimes_A^\mathbf{L} R\Gamma(X, M)
\to
R\Gamma(X, K \otimes_{\mathcal{O}_X}^\mathbf{L} M)
$$ où la deuxième flèche est le cup-produit global $\mu$ ci-dessus. En cohomologie, on retrouve ainsi le cup-produit par $\xi$, comme il ressort du lemme \ref{lemma-second-cup-equals-first} et de sa preuve.
\end{remark}

\noindent
Formulons et démontrons une compatibilité naturelle du cup-produit relatif. Supposons donné un morphisme $f : (X, \mathcal{O}_X) \to (Y, \mathcal{O}_Y)$ d'espaces annelés. Soient $\mathcal{K}^\bullet$ et $\mathcal{M}^\bullet$ des complexes de $\mathcal{O}_X$-modules. Il existe un cup-produit naïf $$
\text{Tot}(
f_*\mathcal{K}^\bullet
\otimes_{\mathcal{O}_Y}
f_*\mathcal{M}^\bullet)
\longrightarrow
f_*\text{Tot}(\mathcal{K}^\bullet
\otimes_{\mathcal{O}_X}
\mathcal{M}^\bullet)
$$ Nous affirmons que celui-ci est relié au cup-produit relatif.

\begin{lemma}

\label{lemma-cup-compatible-with-naive}
Dans la situation ci-dessus, le diagramme ci-dessous commute $$
\xymatrix{
f_*\mathcal{K}^\bullet
\otimes_{\mathcal{O}_Y}^\mathbf{L}
f_*\mathcal{M}^\bullet \ar[r] \ar[d]
&
Rf_*\mathcal{K}^\bullet
\otimes_{\mathcal{O}_Y}^\mathbf{L}
Rf_*\mathcal{M}^\bullet \ar[d]^{\text{Remarque \ref{remark-cup-product}}} \\
\text{Tot}(
f_*\mathcal{K}^\bullet
\otimes_{\mathcal{O}_Y}
f_*\mathcal{M}^\bullet) \ar[d]_{\text{cup-produit naïf}} &
Rf_*(\mathcal{K}^\bullet
\otimes_{\mathcal{O}_X}^\mathbf{L}
\mathcal{M}^\bullet) \ar[d] \\
f_*\text{Tot}(\mathcal{K}^\bullet
\otimes_{\mathcal{O}_X}
\mathcal{M}^\bullet) \ar[r] &
Rf_*\text{Tot}(\mathcal{K}^\bullet
\otimes_{\mathcal{O}_X}
\mathcal{M}^\bullet)
}
$$
\end{lemma}

\begin{proof}

D'après la construction de la remarque \ref{remark-cup-product}, le parcours du diagramme dans le sens horaire donne un morphisme
$$
f_*\mathcal{K}^\bullet
\otimes_{\mathcal{O}_Y}^\mathbf{L}
f_*\mathcal{M}^\bullet 
\longrightarrow
Rf_*\text{Tot}(\mathcal{K}^\bullet
\otimes_{\mathcal{O}_X}
\mathcal{M}^\bullet)
$$
dont l'adjoint est le morphisme

\begin{align*}
Lf^*(f_*\mathcal{K}^\bullet
\otimes_{\mathcal{O}_Y}^\mathbf{L}
f_*\mathcal{M}^\bullet)
& =
Lf^*f_*\mathcal{K}^\bullet
\otimes_{\mathcal{O}_Y}^\mathbf{L}
Lf^*f_*\mathcal{M}^\bullet \\
& \to
Lf^*Rf_*\mathcal{K}^\bullet
\otimes_{\mathcal{O}_Y}^\mathbf{L}
Lf^*Rf_*\mathcal{M}^\bullet \\
& \to
\mathcal{K}^\bullet
\otimes_{\mathcal{O}_Y}^\mathbf{L}
\mathcal{M}^\bullet \\
& \to
\text{Tot}(\mathcal{K}^\bullet
\otimes_{\mathcal{O}_X}
\mathcal{M}^\bullet)
\end{align*}

D'après le lemme \ref{lemma-adjoints-push-pull-compatibility} c'est aussi égal à

\begin{align*}
Lf^*(f_*\mathcal{K}^\bullet
\otimes_{\mathcal{O}_Y}^\mathbf{L}
f_*\mathcal{M}^\bullet)
& =
Lf^*f_*\mathcal{K}^\bullet
\otimes_{\mathcal{O}_Y}^\mathbf{L}
Lf^*f_*\mathcal{M}^\bullet \\
& \to
f^*f_*\mathcal{K}^\bullet
\otimes_{\mathcal{O}_Y}^\mathbf{L}
f^*f_*\mathcal{M}^\bullet \\
& \to
\mathcal{K}^\bullet
\otimes_{\mathcal{O}_Y}^\mathbf{L}
\mathcal{M}^\bullet \\
& \to
\text{Tot}(\mathcal{K}^\bullet
\otimes_{\mathcal{O}_X}
\mathcal{M}^\bullet)
\end{align*}

En parcourant le diagramme dans le sens antihoraire, nous obtenons le morphisme adjoint du morphisme

\begin{align*}
Lf^*(f_*\mathcal{K}^\bullet
\otimes_{\mathcal{O}_Y}^\mathbf{L}
f_*\mathcal{M}^\bullet)
& \to
Lf^*\text{Tot}(
f_*\mathcal{K}^\bullet
\otimes_{\mathcal{O}_Y}
f_*\mathcal{M}^\bullet) \\
& \to
Lf^*f_*\text{Tot}(\mathcal{K}^\bullet
\otimes_{\mathcal{O}_X}
\mathcal{M}^\bullet) \\
& \to
Lf^*Rf_*\text{Tot}(\mathcal{K}^\bullet
\otimes_{\mathcal{O}_X}
\mathcal{M}^\bullet) \\
& \to
\text{Tot}(\mathcal{K}^\bullet
\otimes_{\mathcal{O}_X}
\mathcal{M}^\bullet)
\end{align*}

D'après le lemme \ref{lemma-adjoints-push-pull-compatibility} c'est aussi égal à

\begin{align*}
Lf^*(f_*\mathcal{K}^\bullet
\otimes_{\mathcal{O}_Y}^\mathbf{L}
f_*\mathcal{M}^\bullet)
& \to
Lf^*\text{Tot}(
f_*\mathcal{K}^\bullet
\otimes_{\mathcal{O}_Y}
f_*\mathcal{M}^\bullet) \\
& \to
Lf^*f_*\text{Tot}(\mathcal{K}^\bullet
\otimes_{\mathcal{O}_X}
\mathcal{M}^\bullet) \\
& \to
f^*f_*\text{Tot}(\mathcal{K}^\bullet
\otimes_{\mathcal{O}_X}
\mathcal{M}^\bullet) \\
& \to
\text{Tot}(\mathcal{K}^\bullet
\otimes_{\mathcal{O}_X}
\mathcal{M}^\bullet)
\end{align*}

La démonstration s'achève alors par l'examen du diagramme
$$
\xymatrix{
Lf^*(f_*\mathcal{K}^\bullet
\otimes_{\mathcal{O}_Y}^\mathbf{L}
f_*\mathcal{M}^\bullet) \ar[d] \ar[rr] & &
Lf^*f_*\mathcal{K}^\bullet \otimes_{\mathcal{O}_X}^\mathbf{L}
Lf^*f_*\mathcal{M}^\bullet \ar[d] \\
Lf^*\text{Tot}(
f_*\mathcal{K}^\bullet
\otimes_{\mathcal{O}_Y}
f_*\mathcal{M}^\bullet) \ar[d]_{naive} \ar[r] &
f^*\text{Tot}(
f_*\mathcal{K}^\bullet
\otimes_{\mathcal{O}_Y}
f_*\mathcal{M}^\bullet) \ar[ldd]^{naive} \ar[dd] &
f^*f_*\mathcal{K}^\bullet \otimes_{\mathcal{O}_X}^\mathbf{L}
f^*f_*\mathcal{M}^\bullet \ar[dd] \ar[ldd] \\
Lf^*f_*\text{Tot}(\mathcal{K}^\bullet
\otimes_{\mathcal{O}_X}
\mathcal{M}^\bullet) \ar[d] \\
f^*f_*\text{Tot}(\mathcal{K}^\bullet \otimes_{\mathcal{O}_X}
\mathcal{M}^\bullet) \ar[rd] &
\text{Tot}(f^*f_*\mathcal{K}^\bullet \otimes_{\mathcal{O}_X}
f^*f_*\mathcal{M}^\bullet) \ar[d] &
\mathcal{K}^\bullet \otimes_{\mathcal{O}_X}^\mathbf{L}
\mathcal{M}^\bullet \ar[ld] \\
& \text{Tot}(\mathcal{K}^\bullet
\otimes_{\mathcal{O}_X}
\mathcal{M}^\bullet)
}
$$
Tous les polygones de ce diagramme commutent. Cela résulte du lemme \ref{lemma-tensor-pull-compatibility}. Le carré comportant les deux cup-produits naïfs commute parce que
$Lf^* \to f^*$ est fonctoriel dans le complexe de modules. De même avec le carré impliquant les deux cartes
$\mathcal{A}^\bullet \otimes^\mathbf{L} \mathcal{B}^\bullet \to
\text{Tot}(\mathcal{A}^\bullet \otimes \mathcal{B}^\bullet)$. Enfin, la commutativité du reste du carré vaut au niveau des complexes et peut être considérée comme la définition du cup-produit naïf (par l'adjonction de $f^*$ et $f_*$). La preuve est achevée, car les deux morphismes décrits ci-dessus font le tour du bord extérieur du diagramme.

\end{proof}

\noindent
Soit $(X, \mathcal{O}_X)$ un espace annelé. Soient $\mathcal{K}^\bullet$ et $\mathcal{M}^\bullet$ des complexes de $\mathcal{O}_X$-modules. Nous avons alors un cup-produit naïf $$
\mu' :
\text{Tot}(
\Gamma(X, \mathcal{K}^\bullet) \otimes_A \Gamma(X, \mathcal{M}^\bullet))
\longrightarrow
\Gamma(X, \text{Tot}(
\mathcal{K}^\bullet \otimes_{\mathcal{O}_X} \mathcal{M}^\bullet))
$$ D'après le lemme \ref{lemma-cup-compatible-with-naive}, appliqué au morphisme $(X, \mathcal{O}_X) \to (pt, A)$, ce cup-produit naïf est relié au cup-produit $\mu$ défini au premier paragraphe de cette section par le diagramme commutatif suivant $$
\xymatrix{
\Gamma(X, \mathcal{K}^\bullet)
\otimes_A^\mathbf{L}
\Gamma(X, \mathcal{M}^\bullet) \ar[d] \ar[r] &
R\Gamma(X, \mathcal{K}^\bullet)
\otimes_A^\mathbf{L}
R\Gamma(X, \mathcal{M}^\bullet) \ar[d]^-\mu \\
\text{Tot}(\Gamma(X, \mathcal{K}^\bullet)
\otimes_A
\Gamma(X, \mathcal{M}^\bullet)) \ar[d]_-{\mu'} &
R\Gamma(X, \mathcal{K}^\bullet
\otimes_{\mathcal{O}_X}^\mathbf{L}
\mathcal{M}^\bullet) \ar[d] \\
\Gamma(X, \text{Tot}(\mathcal{K}^\bullet
\otimes_{\mathcal{O}_X}
\mathcal{M}^\bullet)) \ar[r] &
R\Gamma(X, \text{Tot}(\mathcal{K}^\bullet
\otimes_{\mathcal{O}_X}
\mathcal{M}^\bullet))
}
$$ dans $D(A)$. Sur la cohomologie, nous obtenons le diagramme commutatif $$
\xymatrix{
H^i(\Gamma(X, \mathcal{K}^\bullet)) \times
H^j(\Gamma(X, \mathcal{M}^\bullet)) \ar[d] \ar[r] &
H^{i + j}(X,
\text{Tot}(\mathcal{K}^\bullet \otimes_{\mathcal{O}_X} \mathcal{M}^\bullet)) \\
H^i(X, \mathcal{K}^\bullet) \times
H^j(X, \mathcal{M}^\bullet) \ar[r]^\cup &
H^{i + j}(X, \mathcal{K}^\bullet \otimes_{\mathcal{O}_X}^\mathbf{L}
\mathcal{M}^\bullet) \ar[u]
}
$$ qui relie le cup-produit naïf au véritable cup-produit.

\begin{lemma}

\label{lemma-diagrams-commute}
Soit $(X, \mathcal{O}_X)$ un espace annelé. Soient $\mathcal{K}^\bullet$ et $\mathcal{M}^\bullet$ des complexes bornés inférieurement de $\mathcal{O}_X$-modules. Soit $\mathcal{U} : X = \bigcup_{i \in I} U_i$ un recouvrement ouvert. Alors $$
\xymatrix{
\text{Tot}(\check{\mathcal{C}}^\bullet(\mathcal{U}, \mathcal{K}^\bullet))
\otimes_A^\mathbf{L}
\text{Tot}(\check{\mathcal{C}}^\bullet(\mathcal{U}, \mathcal{M}^\bullet))
\ar[d] \ar[r] &
R\Gamma(X, \mathcal{K}^\bullet)
\otimes_A^\mathbf{L}
R\Gamma(X, \mathcal{M}^\bullet) \ar[d]^\mu \\
\text{Tot}(
\text{Tot}(\check{\mathcal{C}}^\bullet(\mathcal{U}, \mathcal{K}^\bullet))
\otimes_A
\text{Tot}(\check{\mathcal{C}}^\bullet(\mathcal{U}, \mathcal{M}^\bullet)))
\ar[d]^{(\ref{equation-needs-signs})} &
R\Gamma(X,
\mathcal{K}^\bullet \otimes_{\mathcal{O}_X}^\mathbf{L} \mathcal{M}^\bullet)
\ar[d] \\
\text{Tot}(
\check{\mathcal{C}}^\bullet({\mathcal U},
\text{Tot}(\mathcal{K}^\bullet \otimes_{\mathcal{O}_X} \mathcal{M}^\bullet)
)) \ar[r] &
R\Gamma(X,
\text{Tot}(\mathcal{K}^\bullet \otimes_{\mathcal{O}_X} \mathcal{M}^\bullet))
}
$$ où les flèches horizontales sont celles du lemme \ref{lemma-cech-complex-complex}, commute dans $D(A)$.
\end{lemma}

\begin{proof}

Choisissez des quasi-isomorphismes de complexes
$a : \mathcal{K}^\bullet \to \mathcal{K}_1^\bullet$ et
$b : \mathcal{M}^\bullet \to \mathcal{M}_1^\bullet$
comme dans le lemme \ref{lemma-godement-resolution-bounded-below}. Puisque les morphismes $a$ et $b$ induisent des équivalences d'homotopie sur les germes, nous voyons que le morphisme induit
$$
\text{Tot}(\mathcal{K}^\bullet \otimes_{\mathcal{O}_X} \mathcal{M}^\bullet)
\to
\text{Tot}(\mathcal{K}_1^\bullet \otimes_{\mathcal{O}_X} \mathcal{M}_1^\bullet)
$$
est lui aussi une équivalence d'homotopie sur les germes (Compléments sur l'algèbre, lemme
\ref{more-algebra-lemma-derived-tor-homotopy}) et donc un quasi-isomorphisme. Ainsi les cibles
$$
R\Gamma(X,
\text{Tot}(\mathcal{K}^\bullet
\otimes_{\mathcal{O}_X} \mathcal{M}^\bullet)) =
R\Gamma(X,
\text{Tot}(\mathcal{K}_1^\bullet
\otimes_{\mathcal{O}_X} \mathcal{M}_1^\bullet))
$$
des deux diagrammes coïncident dans $D(A)$. Il suffit donc de démontrer l'assertion après avoir remplacé $\mathcal{K}$ et $\mathcal{M}$
par $\mathcal{K}_1$ et $\mathcal{M}_1$. Nous sommes ainsi ramenés au cas traité dans le paragraphe suivant.

\medskip\noindent

Supposons que $\mathcal{K}^\bullet$ et $\mathcal{M}^\bullet$ soient des complexes bornés inférieurement de $\mathcal{O}_X$-modules flasques et considérons le diagramme qui relie le cup-produit au cup-produit (\ref{equation-needs-signs}) sur les complexes de {\v C}ech. Nous pouvons alors considérer le diagramme commutatif
$$
\xymatrix{
\Gamma(X, \mathcal{K}^\bullet)
\otimes_A^\mathbf{L}
\Gamma(X, \mathcal{M}^\bullet) \ar[d] \ar[r] &
\text{Tot}(\check{\mathcal{C}}^\bullet(\mathcal{U}, \mathcal{K}^\bullet))
\otimes_A^\mathbf{L}
\text{Tot}(\check{\mathcal{C}}^\bullet(\mathcal{U}, \mathcal{M}^\bullet))
\ar[d] \\
\text{Tot}(\Gamma(X, \mathcal{K}^\bullet)
\otimes_A
\Gamma(X, \mathcal{M}^\bullet)) \ar[d] \ar[r] &
\text{Tot}(
\text{Tot}(\check{\mathcal{C}}^\bullet(\mathcal{U}, \mathcal{K}^\bullet))
\otimes_A
\text{Tot}(\check{\mathcal{C}}^\bullet(\mathcal{U}, \mathcal{M}^\bullet)))
\ar[d]^{(\ref{equation-needs-signs})} \\
\Gamma(X, \text{Tot}(\mathcal{K}^\bullet
\otimes_{\mathcal{O}_X}
\mathcal{M}^\bullet)) \ar[r] &
\text{Tot}(
\check{\mathcal{C}}^\bullet({\mathcal U},
\text{Tot}(\mathcal{K}^\bullet \otimes_{\mathcal{O}_X} \mathcal{M}^\bullet)
))
}
$$
Dans ce diagramme, les flèches horizontales sont des isomorphismes dans $D(A)$, car, pour un complexe borné inférieurement de modules flasques tel que $\mathcal{K}^\bullet$,
nous avons
$$
\Gamma(X, \mathcal{K}^\bullet) =
\text{Tot}(\check{\mathcal{C}}^\bullet(\mathcal{U}, \mathcal{K}^\bullet)) =
R\Gamma(X, \mathcal{K}^\bullet)
$$
dans $D(A)$. Cela résulte du lemme \ref{lemma-flasque-acyclic}, de Catégories dérivées, lemme \ref{derived-lemma-leray-acyclicity}, et du lemme \ref{lemma-cech-complex-complex-computes}. La commutativité du diagramme du lemme faisant intervenir (\ref{equation-needs-signs}) résulte donc de celle, déjà démontrée, du lemme \ref{lemma-cup-compatible-with-naive},
où $f$ est le morphisme vers un point (voir la discussion qui suit \ref{lemma-cup-compatible-with-naive}).

\end{proof}

\begin{lemma}

\label{lemma-cup-product-associative}
Soit $f : (X, \mathcal{O}_X) \to (Y, \mathcal{O}_Y)$ un morphisme d'espaces annelés. Le cup-produit relatif de la remarque \ref{remark-cup-product} est associatif au sens où le diagramme $$
\xymatrix{
Rf_*K \otimes_{\mathcal{O}_Y}^\mathbf{L}
Rf_*L \otimes_{\mathcal{O}_Y}^\mathbf{L}
Rf_*M \ar[r] \ar[d] &
Rf_*(K \otimes_{\mathcal{O}_X}^\mathbf{L} L)
\otimes_{\mathcal{O}_Y}^\mathbf{L} Rf_*M \ar[d] \\
Rf_*K \otimes_{\mathcal{O}_Y}^\mathbf{L}
Rf_*(L \otimes_{\mathcal{O}_X}^\mathbf{L} M) \ar[r] &
Rf_*(K \otimes_{\mathcal{O}_X}^\mathbf{L} 
L \otimes_{\mathcal{O}_X}^\mathbf{L} M)
}
$$ est commutatif dans $D(\mathcal{O}_Y)$ pour tous $K, L, M$ dans $D(\mathcal{O}_X)$.
\end{lemma}

\begin{proof}

En suivant chacun des deux chemins, nous obtenons le morphisme adjoint du morphisme évident

\begin{align*}
Lf^*(Rf_*K \otimes_{\mathcal{O}_Y}^\mathbf{L}
Rf_*L \otimes_{\mathcal{O}_Y}^\mathbf{L}
Rf_*M) & =
Lf^*(Rf_*K) \otimes_{\mathcal{O}_X}^\mathbf{L}
Lf^*(Rf_*L) \otimes_{\mathcal{O}_X}^\mathbf{L}
Lf^*(Rf_*M) \\
& \to
K \otimes_{\mathcal{O}_X}^\mathbf{L} 
L \otimes_{\mathcal{O}_X}^\mathbf{L} M
\end{align*}

en $D(\mathcal{O}_X)$.

\end{proof}

\begin{lemma}

\label{lemma-cup-product-commutative}

Soit $f : (X, \mathcal{O}_X) \to (Y, \mathcal{O}_Y)$
un morphisme d'espaces annelés. Le cup-produit relatif de la remarque \ref{remark-cup-product} est commutatif au sens où le diagramme
$$
\xymatrix{
Rf_*K \otimes_{\mathcal{O}_Y}^\mathbf{L} Rf_*L \ar[r] \ar[d]_\psi &
Rf_*(K \otimes_{\mathcal{O}_X}^\mathbf{L} L) \ar[d]^{Rf_*\psi} \\
Rf_*L \otimes_{\mathcal{O}_Y}^\mathbf{L} Rf_*K \ar[r] &
Rf_*(L \otimes_{\mathcal{O}_X}^\mathbf{L} K)
}
$$
est commutatif dans $D(\mathcal{O}_Y)$ pour tous $K, L$ dans $D(\mathcal{O}_X)$. Ici, $\psi$ est la contrainte de commutativité de la catégorie dérivée (lemme \ref{lemma-symmetric-monoidal-derived}).

\end{lemma}

\begin{proof}
Démonstration omise.
\end{proof}

\begin{lemma}

\label{lemma-compose-cup-product}
Soient $f : (X, \mathcal{O}_X) \to (Y, \mathcal{O}_Y)$ et $g : (Y, \mathcal{O}_Y) \to (Z, \mathcal{O}_Z)$ des morphismes d'espaces annelés. Le cup-produit relatif de la remarque \ref{remark-cup-product} est compatible aux compositions au sens où le diagramme $$
\xymatrix{
R(g \circ f)_*K \otimes_{\mathcal{O}_Z}^\mathbf{L} R(g \circ f)_*L
\ar@{=}[rr] \ar[d] & &
Rg_*Rf_*K \otimes_{\mathcal{O}_Z}^\mathbf{L} Rg_*Rf_*L \ar[d] \\
R(g \circ f)_*(K \otimes_{\mathcal{O}_X}^\mathbf{L} L) \ar@{=}[r] &
Rg_*Rf_*(K \otimes_{\mathcal{O}_X}^\mathbf{L} L) &
Rg_*(Rf_*K \otimes_{\mathcal{O}_Y}^\mathbf{L}  Rf_*L) \ar[l]
}
$$ est commutatif en $D(\mathcal{O}_Z)$ pour tous $K, L$ en $D(\mathcal{O}_X)$.
\end{lemma}

\begin{proof}

En effet, chacun des deux parcours du diagramme donne le morphisme adjoint du morphisme

\begin{align*}
& L(g \circ f)^*\left(R(g \circ f)_*K
\otimes_{\mathcal{O}_Z}^\mathbf{L}
R(g \circ f)_*L\right) \\
& =
L(g \circ f)^*R(g \circ f)_*K
\otimes_{\mathcal{O}_X}^\mathbf{L}
L(g \circ f)^*R(g \circ f)_*L) \\
& \to
K \otimes_{\mathcal{O}_X}^\mathbf{L} L
\end{align*}

dans $D(\mathcal{O}_X)$. Pour le voir, on utilise que la composée des counités
$$
L(g \circ f)^*R(g \circ f)_* =
Lf^* Lg^* Rg_* Rf_*  \to
Lf^* Rf_* \to \text{id}
$$
est la counité pour $L(g \circ f)^*$ et $R(g \circ f)_*$. Voir Catégories, lemme \ref{categories-lemma-compose-counits}.

\end{proof}

\begin{lemma}

\label{lemma-base-change-cup-product}
Considérez un carré commutatif $$
\xymatrix{
(X', \mathcal{O}_{X'}) \ar[r]_{g'} \ar[d]_{f'} &
(X, \mathcal{O}_X) \ar[d]^f \\
(Y', \mathcal{O}_{Y'}) \ar[r]^g & (Y, \mathcal{O}_Y) 
}
$$ d'espaces annelés. Soient $K, L$ dans $D(\mathcal{O}_X)$. Le cup-produit relatif est compatible avec ce carré au sens où le diagramme $$
\xymatrix{
Lg^*(Rf_*K \otimes_{\mathcal{O}_Y}^\mathbf{L} Rf_*L)
\ar[r] \ar@{=}[d] &
Lg^*(Rf_*(K \otimes_{\mathcal{O}_X}^\mathbf{L} L)) \ar[d] \\
Lg^*Rf_*K \otimes_{\mathcal{O}_{Y'}}^\mathbf{L} Lg^*Rf_*L \ar[d] &
R(f')_*L(g')^*(K \otimes_{\mathcal{O}_X}^\mathbf{L} L) \ar@{=}[d] \\
R(f')_*(L(g')^*K \otimes_{\mathcal{O}_{Y'}} R(f')_*(L(g')^*L \ar[r] &
R(f')_*(L(g')^*K \otimes_{\mathcal{O}_{X'}}^\mathbf{L} L(g')^*L)
}
$$ est commutatif dans $D(\mathcal{O}_{Y'})$. Les flèches horizontales sont données par le cup-produit relatif (remarque \ref{remark-cup-product}) et les flèches verticales par le morphisme de changement de base (remarque \ref{remark-base-change}) et le lemme \ref{lemma-pullback-tensor-product}.
\end{lemma}

\begin{proof}
Démonstration omise.
\end{proof}













\section{Quelques propriétés des complexes K-injectifs}

\label{section-properties-K-injective}

\noindent

Soit $(X, \mathcal{O}_X)$ un espace annelé. Soit $U \subset X$ un sous-ensemble ouvert. Notons $j : (U, \mathcal{O}_U) \to (X, \mathcal{O}_X)$
l'immersion ouverte correspondante. Le foncteur d'image inverse $j^*$ est exact, puisqu'il n'est autre que le foncteur de restriction. Ainsi, l'image inverse dérivée $Lj^*$ se calcule sur tout complexe en restreignant simplement ce complexe. Nous notons souvent simplement le foncteur correspondant
$$
D(\mathcal{O}_X) \to D(\mathcal{O}_U),
\quad
E \mapsto j^*E = E|_U
$$
De même, l'extension de zéro
$j_! : \textit{Mod}(\mathcal{O}_U) \to \textit{Mod}(\mathcal{O}_X)$
(voir Faisceaux, section \ref{sheaves-section-open-immersions}) est un foncteur exact (Modules, lemme \ref{modules-lemma-j-shriek-exact}). Il induit donc un foncteur
$$
j_! : D(\mathcal{O}_U) \to D(\mathcal{O}_X),\quad
F \mapsto j_!F
$$
en appliquant simplement $j_!$ à tout complexe représentant l'objet $F$.

\begin{lemma}

\label{lemma-restrict-K-injective-to-open}

Soit $X$ un espace annelé. Soit $U \subset X$ un sous-espace ouvert. La restriction d'un complexe K-injectif de $\mathcal{O}_X$-modules à $U$ est un complexe K-injectif de $\mathcal{O}_U$-modules.

\end{lemma}

\begin{proof}
Cela résulte de Catégories dérivées, lemme \ref{derived-lemma-adjoint-preserve-K-injectives}, et du fait que le foncteur de restriction admet pour adjoint à gauche le foncteur exact $j_!$. Pour la construction de $j_!$, voir Faisceaux, section \ref{sheaves-section-open-immersions}, et, pour son exactitude, voir Modules, lemme \ref{modules-lemma-j-shriek-exact}.
\end{proof}

\begin{lemma}

\label{lemma-unbounded-cohomology-of-open}
Soit $X$ un espace annelé. Soit $U \subset X$ un sous-espace ouvert. Pour $K$ dans $D(\mathcal{O}_X)$, nous avons $H^p(U, K) = H^p(U, K|_U)$.
\end{lemma}

\begin{proof}

Soit $\mathcal{I}^\bullet$ un complexe K-injectif de $\mathcal{O}_X$-modules représentant $K$. Alors
$$
H^q(U, K) = H^q(\Gamma(U, \mathcal{I}^\bullet)) =
H^q(\Gamma(U, \mathcal{I}^\bullet|_U))
$$
par construction de la cohomologie. D'après le lemme \ref{lemma-restrict-K-injective-to-open},
le complexe $\mathcal{I}^\bullet|_U$ est un complexe K-injectif représentant $K|_U$ et le lemme suit.

\end{proof}

\begin{lemma}

\label{lemma-sheafification-cohomology}
Soit $(X, \mathcal{O}_X)$ un espace annelé. Soit $K$ un objet de $D(\mathcal{O}_X)$. La faisceautisation de $$
U \mapsto H^q(U, K) = H^q(U, K|_U)
$$ est le $q$-ième faisceau de cohomologie $H^q(K)$ de $K$.
\end{lemma}

\begin{proof}

L'égalité $H^q(U, K) = H^q(U, K|_U)$ résulte du lemme \ref{lemma-unbounded-cohomology-of-open}. Choisissons un complexe K-injectif $\mathcal{I}^\bullet$ représentant $K$. Alors
$$
H^q(U, K) =
\frac{\Ker(\mathcal{I}^q(U) \to \mathcal{I}^{q + 1}(U))}
{\Im(\mathcal{I}^{q - 1}(U) \to \mathcal{I}^q(U))}.
$$
par notre construction de la cohomologie. Comme
$H^q(K) = \Ker(\mathcal{I}^q \to \mathcal{I}^{q + 1})/
\Im(\mathcal{I}^{q - 1} \to \mathcal{I}^q)$, le résultat est clair.

\end{proof}

\begin{lemma}

\label{lemma-restrict-direct-image-open}
Soit $f : (X, \mathcal{O}_X) \to (Y, \mathcal{O}_Y)$ un morphisme d'espaces annelés. Étant donné un sous-espace ouvert $V \subset Y$, posons $U = f^{-1}(V)$ et notons $g : U \to V$ le morphisme induit. Alors $(Rf_*E)|_V = Rg_*(E|_U)$ pour $E$ dans $D(\mathcal{O}_X)$.
\end{lemma}

\begin{proof}
Représentons $E$ par un complexe K-injectif $\mathcal{I}^\bullet$ de $\mathcal{O}_X$-modules. Alors $Rf_*(E) = f_*\mathcal{I}^\bullet$ et $Rg_*(E|_U) = g_*(\mathcal{I}^\bullet|_U)$ d'après le lemme \ref{lemma-restrict-K-injective-to-open}. Comme il est clair que $(f_*\mathcal{F})|_V = g_*(\mathcal{F}|_U)$ pour tout faisceau $\mathcal{F}$ sur $X$, le résultat suit.
\end{proof}

\begin{lemma}

\label{lemma-Leray-unbounded}
Soit $f : X \to Y$ un morphisme d'espaces annelés. Alors $R\Gamma(Y, -) \circ Rf_* = R\Gamma(X, -)$ comme foncteurs $D(\mathcal{O}_X) \to D(\Gamma(Y, \mathcal{O}_Y))$. Plus généralement, pour $V \subset Y$ ouvert et $U = f^{-1}(V)$, nous avons $R\Gamma(U, -) = R\Gamma(V, -) \circ Rf_*$.
\end{lemma}

\begin{proof}
Soit $Z$ l'espace annelé dont l'espace sous-jacent est réduit à un point et tel que $\Gamma(Z, \mathcal{O}_Z) = \Gamma(Y, \mathcal{O}_Y)$. Il existe un morphisme canonique $Y \to Z$ d'espaces annelés qui induit l'identification sur les sections globales des faisceaux structuraux. Alors $D(\mathcal{O}_Z) = D(\Gamma(Y, \mathcal{O}_Y))$. L'égalité $R\Gamma(Y, -) \circ Rf_* = R\Gamma(X, -)$ découle donc du lemme \ref{lemma-derived-pushforward-composition} appliqué à $X \to Y \to Z$.

\medskip\noindent

Le second énoncé, plus général, découle du premier après application du lemme \ref{lemma-restrict-direct-image-open}.

\end{proof}

\begin{lemma}

\label{lemma-unbounded-describe-higher-direct-images}
Soit $f : (X, \mathcal{O}_X) \to (Y, \mathcal{O}_Y)$ un morphisme d'espaces annelés. Soit $K$ dans $D(\mathcal{O}_X)$. Alors $H^i(Rf_*K)$ est le faisceau associé au préfaisceau $$
V \mapsto H^i(f^{-1}(V), K) = H^i(V, Rf_*K)
$$
\end{lemma}

\begin{proof}
L'égalité $H^i(f^{-1}(V), K) = H^i(V, Rf_*K)$ s'obtient en prenant la cohomologie du second énoncé du lemme \ref{lemma-Leray-unbounded}. L'énoncé sur la faisceautisation découle alors du lemme \ref{lemma-sheafification-cohomology}.
\end{proof}

\begin{lemma}

\label{lemma-modules-abelian-unbounded}
Soit $X$ un espace annelé. Soit $K$ un objet de $D(\mathcal{O}_X)$ et notons $K_{ab}$ son image dans $D(\underline{\mathbf{Z}}_X)$.
\begin{enumerate}
\item Pour tout $U \subset X$ ouvert, il y a une application canonique $R\Gamma(U, K) \to R\Gamma(U, K_{ab})$ qui est un isomorphisme dans $D(\textit{Ab})$.
\item Soit $f : X \to Y$ un morphisme d'espaces annelés. Il y a une application canonique $Rf_*K \to Rf_*(K_{ab})$ qui est un isomorphisme dans $D(\underline{\mathbf{Z}}_Y)$.
\end{enumerate}

\end{lemma}

\begin{proof}

Le morphisme est construit comme suit. Choisissons un complexe K-injectif
$\mathcal{I}^\bullet$ représentant $K$. Choisissez un quasi-isomorphisme
$\mathcal{I}^\bullet \to \mathcal{J}^\bullet$ où $\mathcal{J}^\bullet$
est un complexe K-injectif de groupes abéliens. Alors le morphisme de (1) est donné par $\Gamma(U, \mathcal{I}^\bullet) \to \Gamma(U, \mathcal{J}^\bullet)$
et celui de (2) par
$f_*\mathcal{I}^\bullet \to f_*\mathcal{J}^\bullet$. Pour montrer que ces morphismes sont des isomorphismes, il suffit de prouver qu'ils induisent des isomorphismes sur les groupes et les faisceaux de cohomologie. D'après les lemmes \ref{lemma-unbounded-cohomology-of-open} et
\ref{lemma-unbounded-describe-higher-direct-images}
il suffit de montrer que l'application
$$
H^0(X, K) \longrightarrow H^0(X, K_{ab})
$$
est un isomorphisme. Observez que
$$
H^0(X, K) = \Hom_{D(\mathcal{O}_X)}(\mathcal{O}_X, K)
$$
et de même pour l'autre groupe. Choisissons un complexe quelconque $\mathcal{K}^\bullet$
de $\mathcal{O}_X$-modules représentant $K$. Par la construction de la catégorie dérivée comme une localisation nous avons
$$
\Hom_{D(\mathcal{O}_X)}(\mathcal{O}_X, K) =
\colim_{s : \mathcal{F}^\bullet \to \mathcal{O}_X}
\Hom_{K(\mathcal{O}_X)}(\mathcal{F}^\bullet, \mathcal{K}^\bullet)
$$
où la colimite est prise sur les quasi-isomorphismes $s$ de complexes de
$\mathcal{O}_X$-modules. De même, nous avons
$$
\Hom_{D(\underline{\mathbf{Z}}_X)}(\underline{\mathbf{Z}}_X, K) =
\colim_{s : \mathcal{G}^\bullet \to \underline{\mathbf{Z}}_X}
\Hom_{K(\underline{\mathbf{Z}}_X)}(\mathcal{G}^\bullet, \mathcal{K}^\bullet)
$$
Ensuite, nous observons que les quasi-isomorphismes
$s : \mathcal{G}^\bullet \to \underline{\mathbf{Z}}_X$
avec $\mathcal{G}^\bullet$ complexe borné supérieurement de
$\underline{\mathbf{Z}}_X$-modules plats forment une famille cofinale dans le système. (Ceci découle de Modules, lemme \ref{modules-lemma-module-quotient-flat}, et de Catégories dérivées, lemme \ref{derived-lemma-subcategory-left-resolution}; voir la discussion de la section \ref{section-flat}.) Nous pouvons donc construire un inverse au morphisme
$H^0(X, K) \longrightarrow H^0(X, K_{ab})$
en représentant un élément $\xi \in H^0(X, K_{ab})$ par une paire
$$
(s : \mathcal{G}^\bullet \to \underline{\mathbf{Z}}_X,
a : \mathcal{G}^\bullet \to \mathcal{K}^\bullet)
$$
où $\mathcal{G}^\bullet$ est un complexe borné supérieurement de
$\underline{\mathbf{Z}}_X$-modules plats, et envoyer cet élément sur
$$
(\mathcal{G}^\bullet \otimes_{\underline{\mathbf{Z}}_X} \mathcal{O}_X
\to \mathcal{O}_X,
\mathcal{G}^\bullet  \otimes_{\underline{\mathbf{Z}}_X} \mathcal{O}_X
\to \mathcal{K}^\bullet)
$$
Le seul point à noter est que la première flèche est un quasi-isomorphisme d'après les lemmes \ref{lemma-derived-tor-quasi-isomorphism-other-side} et
\ref{lemma-bounded-flat-K-flat}. Nous omettons la vérification détaillée que cette construction est effectivement un inverse.

\end{proof}

\begin{lemma}

\label{lemma-adjoint-lower-shriek-restrict}
Soit $(X, \mathcal{O}_X)$ un espace annelé. Soit $U \subset X$ un sous-ensemble ouvert. Notons $j : (U, \mathcal{O}_U) \to (X, \mathcal{O}_X)$ l'immersion ouverte correspondante. Le foncteur de restriction $D(\mathcal{O}_X) \to D(\mathcal{O}_U)$ est adjoint à droite du foncteur d'extension par zéro $j_! : D(\mathcal{O}_U) \to D(\mathcal{O}_X)$.
\end{lemma}

\begin{proof}
Ceci résulte formellement du fait que $j_!$ et $j^*$ sont adjoints et exacts (et donc $Lj_! = j_!$ et $Rj^* = j^*$ existent), voir Catégories dérivées, lemme \ref{derived-lemma-derived-adjoint-functors}.
\end{proof}

\begin{lemma}

\label{lemma-K-injective-flat}

Soit $f : X \to Y$ un morphisme plat d'espaces annelés. Si $\mathcal{I}^\bullet$ est un complexe K-injectif de $\mathcal{O}_X$-modules, alors $f_*\mathcal{I}^\bullet$ est K-injectif comme complexe de
$\mathcal{O}_Y$-modules.

\end{lemma}

\begin{proof}
Ceci est vrai parce que $$
\Hom_{K(\mathcal{O}_Y)}(\mathcal{F}^\bullet, f_*\mathcal{I}^\bullet)
=
\Hom_{K(\mathcal{O}_X)}(f^*\mathcal{F}^\bullet, \mathcal{I}^\bullet)
$$ par Faisceaux, lemme \ref{sheaves-lemma-adjoint-pullback-pushforward-modules} et le fait que $f^*$ est exact comme $f$ est supposé être plat.
\end{proof}





\section{Mayer--Vietoris non borné}

\label{section-unbounded-mayer-vietoris}

\noindent
Il existe également une suite de Mayer--Vietoris pour la cohomologie non bornée.

\begin{lemma}

\label{lemma-exact-sequence-lower-shriek}
Soit $(X, \mathcal{O}_X)$ un espace annelé. Soit $X = U \cup V$ l'union de deux sous-espaces ouverts. Pour tout objet $E$ de $D(\mathcal{O}_X)$ nous avons un triangle distingué $$
j_{U \cap V!}E|_{U \cap V} \to
j_{U!}E|_U \oplus j_{V!}E|_V \to E \to 
j_{U \cap V!}E|_{U \cap V}[1]
$$ dans $D(\mathcal{O}_X)$.
\end{lemma}

\begin{proof}

Nous avons vu dans la section \ref{section-properties-K-injective}
que les foncteurs de restriction et d'extension par zéro se calculent en appliquant simplement les foncteurs correspondants à chaque complexe. Soit $\mathcal{E}^\bullet$ un complexe de $\mathcal{O}_X$-modules représentant $E$. Le triangle distingué du lemme est le triangle distingué associé (par Catégories dérivées, section
\ref{derived-section-canonical-delta-functor} et en particulier \ref{derived-lemma-derived-canonical-delta-functor}) à la suite exacte courte de complexes de $\mathcal{O}_X$-modules
$$
0 \to j_{U \cap V!}\mathcal{E}^\bullet|_{U \cap V} \to
j_{U!}\mathcal{E}^\bullet|_U \oplus j_{V!}\mathcal{E}^\bullet|_V
\to \mathcal{E}^\bullet \to 0
$$
L'exactitude de cette suite se vérifie sur les germes à l'aide de Faisceaux, lemme \ref{sheaves-lemma-j-shriek-modules}
(calcul omis).

\end{proof}

\begin{lemma}

\label{lemma-exact-sequence-j-star}
Soit $(X, \mathcal{O}_X)$ un espace annelé. Soit $X = U \cup V$ l'union de deux sous-espaces ouverts. Pour tout objet $E$ de $D(\mathcal{O}_X)$ nous avons un triangle distingué $$
E \to 
Rj_{U, *}E|_U \oplus Rj_{V, *}E|_V \to
Rj_{U \cap V, *}E|_{U \cap V} \to
E[1]
$$ dans $D(\mathcal{O}_X)$.
\end{lemma}

\begin{proof}

Choisissons un complexe K-injectif $\mathcal{I}^\bullet$ représentant $E$
dont les termes $\mathcal{I}^n$ sont des objets injectifs de
$\textit{Mod}(\mathcal{O}_X)$; voir Injectifs, théorème
\ref{injectives-theorem-K-injective-embedding-grothendieck}. Nous avons vu que $\mathcal{I}^\bullet|U$ est également un complexe K-injectif (lemme \ref{lemma-restrict-K-injective-to-open}). D'où :
$Rj_{U, *}E|_U$ est représenté par $j_{U, *}\mathcal{I}^\bullet|_U$. De même pour $V$ et $U \cap V$. Par conséquent, le triangle distingué du lemme est le triangle distingué associé (par Catégories dérivées, section
\ref{derived-section-canonical-delta-functor} et en particulier \ref{derived-lemma-derived-canonical-delta-functor}) à la suite exacte courte des complexes
$$
0 \to
\mathcal{I}^\bullet \to
j_{U, *}\mathcal{I}^\bullet|_U \oplus j_{V, *}\mathcal{I}^\bullet|_V \to
j_{U \cap V, *}\mathcal{I}^\bullet|_{U \cap V} \to
0.
$$
Cette suite est exacte parce que, pour tout ouvert $W \subset X$ et tout $n$, la suite
$$
0 \to
\mathcal{I}^n(W) \to
\mathcal{I}^n(W \cap U) \oplus \mathcal{I}^n(W \cap V) \to
\mathcal{I}^n(W \cap U \cap V) \to
0
$$
est exacte (voir la preuve du lemme \ref{lemma-mayer-vietoris}).

\end{proof}

\begin{lemma}

\label{lemma-mayer-vietoris-hom}
Soit $(X, \mathcal{O}_X)$ un espace annelé. Soit $X = U \cup V$ l'union de deux sous-espaces ouverts de $X$. Pour les objets $E$, $F$ de $D(\mathcal{O}_X)$ nous avons une suite de Mayer--Vietoris $$
\xymatrix{
& \ldots \ar[r] & \Ext^{-1}(E_{U \cap V}, F_{U \cap V}) \ar[lld] \\
\Hom(E, F) \ar[r] &
\Hom(E_U, F_U) \oplus
\Hom(E_V, F_V) \ar[r] &
\Hom(E_{U \cap V}, F_{U \cap V})
}
$$ où les indices désignent les restrictions aux ouverts concernés, et où les $\Hom$ et $\Ext$ sont pris dans les catégories dérivées correspondantes.
\end{lemma}

\begin{proof}
Appliquons le triangle distingué du lemme \ref{lemma-exact-sequence-lower-shriek} pour obtenir une suite exacte longue de $\Hom$ (Catégories dérivées, lemme \ref{derived-lemma-representable-homological}) et utilisons l'égalité $$
\Hom_{D(\mathcal{O}_X)}(j_{U!}E|_U, F) =
\Hom_{D(\mathcal{O}_U)}(E|_U, F|_U)
$$ d'après le lemme \ref{lemma-adjoint-lower-shriek-restrict}.
\end{proof}

\begin{lemma}

\label{lemma-unbounded-mayer-vietoris}
Soit $(X, \mathcal{O}_X)$ un espace annelé. Supposons que $X = U \cup V$ soit la réunion de deux sous-ensembles ouverts. Pour un objet $E$ de $D(\mathcal{O}_X)$, nous avons un triangle distingué $$
R\Gamma(X, E) \to R\Gamma(U, E) \oplus R\Gamma(V, E) \to
R\Gamma(U \cap V, E) \to R\Gamma(X, E)[1]
$$ et, en particulier, une suite exacte longue de cohomologie $$
\ldots \to
H^n(X, E) \to
H^n(U, E) \oplus H^0(V, E) \to
H^n(U \cap V, E) \to
H^{n + 1}(X, E) \to \ldots
$$ La construction du triangle distingué et de la suite exacte longue est fonctorielle en $E$.
\end{lemma}

\begin{proof}

Choisissez un complexe K-injectif $\mathcal{I}^\bullet$
représentant $E$. Nous pouvons supposer que $\mathcal{I}^n$ est un objet injectif de $\textit{Mod}(\mathcal{O}_X)$ pour tout $n$; voir Injectifs, théorème
\ref{injectives-theorem-K-injective-embedding-grothendieck}. Alors $R\Gamma(X, E)$ est calculé par $\Gamma(X, \mathcal{I}^\bullet)$. De même pour $U$, $V$, $U \cap V$ d'après le lemme \ref{lemma-restrict-K-injective-to-open}. Par conséquent, le triangle distingué du lemme est le triangle distingué associé (par Catégories dérivées, section
\ref{derived-section-canonical-delta-functor} et en particulier \ref{derived-lemma-derived-canonical-delta-functor}) à la suite exacte courte des complexes
$$
0 \to
\mathcal{I}^\bullet(X) \to
\mathcal{I}^\bullet(U) \oplus \mathcal{I}^\bullet(V) \to
\mathcal{I}^\bullet(U \cap V) \to
0.
$$
Nous avons vu qu'il s'agit d'une suite exacte courte dans la preuve du lemme \ref{lemma-mayer-vietoris}. L'assertion finale découle de la fonctorielité de la construction dans Injectifs, théorème
\ref{injectives-theorem-K-injective-embedding-grothendieck}.

\end{proof}

\begin{lemma}

\label{lemma-unbounded-relative-mayer-vietoris}
Soit $f : X \to Y$ un morphisme d'espaces annelés. Supposons que $X = U \cup V$ soit la réunion de deux sous-ensembles ouverts. Notons $a = f|_U : U \to Y$, $b = f|_V : V \to Y$ et $c = f|_{U \cap V} : U \cap V \to Y$. Pour chaque objet $E$ de $D(\mathcal{O}_X)$, il existe un triangle distingué $$
Rf_*E \to
Ra_*(E|_U) \oplus Rb_*(E|_V) \to
Rc_*(E|_{U \cap V}) \to
Rf_*E[1]
$$ Ce triangle est fonctoriel dans $E$.
\end{lemma}

\begin{proof}

Choisissez un complexe K-injectif $\mathcal{I}^\bullet$
représentant $E$. Nous pouvons supposer que $\mathcal{I}^n$ est un objet injectif de $\textit{Mod}(\mathcal{O}_X)$ pour tout $n$; voir Injectifs, théorème
\ref{injectives-theorem-K-injective-embedding-grothendieck}. Alors $Rf_*E$ est calculé par $f_*\mathcal{I}^\bullet$. De même pour $U$, $V$, $U \cap V$ d'après le lemme \ref{lemma-restrict-K-injective-to-open}. Par conséquent, le triangle distingué du lemme est le triangle distingué associé (par Catégories dérivées, section
\ref{derived-section-canonical-delta-functor} et en particulier \ref{derived-lemma-derived-canonical-delta-functor}) à la suite exacte courte des complexes
$$
0 \to
f_*\mathcal{I}^\bullet \to
a_*\mathcal{I}^\bullet|_U \oplus b_*\mathcal{I}^\bullet|_V \to
c_*\mathcal{I}^\bullet|_{U \cap V} \to
0.
$$
C'est une suite exacte courte de complexes d'après le lemme \ref{lemma-relative-mayer-vietoris}
et le fait que $R^1f_*\mathcal{I} = 0$
pour un objet injectif $\mathcal{I}$ de $\textit{Mod}(\mathcal{O}_X)$. L'assertion finale découle de la fonctorielité de la construction dans Injectifs, théorème
\ref{injectives-theorem-K-injective-embedding-grothendieck}.

\end{proof}

\begin{lemma}

\label{lemma-pushforward-restriction}
Soit $(X, \mathcal{O}_X)$ un espace annelé. Soit $j : U \to X$ l'inclusion d'un sous-espace ouvert. Soit $T \subset X$ une partie fermée contenue dans $U$.
\begin{enumerate}
\item Si $E$ est un objet de $D(\mathcal{O}_X)$ dont les faisceaux de cohomologie sont à support dans $T$, alors $E \to Rj_*(E|_U)$ est un isomorphisme.
\item Si $F$ est un objet de $D(\mathcal{O}_U)$ dont les faisceaux de cohomologie sont à support dans $T$, alors $j_!F \to Rj_*F$ est un isomorphisme.
\end{enumerate}

\end{lemma}

\begin{proof}

Posons $V = X \setminus T$ et $W = U \cap V$. Remarquons que $X = U \cup V$ est un recouvrement ouvert de $X$. Notons $j_W : W \to V$ l'immersion ouverte. Soit $E$ un objet de $D(\mathcal{O}_X)$ dont les faisceaux de cohomologie sont à support dans $T$. D'après le lemme \ref{lemma-restrict-direct-image-open}, nous avons
$(Rj_*E|_U)|_V = Rj_{W, *}(E|_W) = 0$, car $E|_W = 0$ par hypothèse. D'autre part, $Rj_*(E|_U)|_U = E|_U$. Ainsi, (1) est clair. Soit $F$ un objet de $D(\mathcal{O}_U)$ dont les faisceaux de cohomologie sont à support dans $T$. D'après le lemme \ref{lemma-restrict-direct-image-open}, nous avons
$(Rj_*F)|_V = Rj_{W, *}(F|_W) = 0$, car $F|_W = 0$ par hypothèse. Nous avons aussi $(j_!F)|_V = j_{W!}(F|_W) = 0$ (la première égalité résulte immédiatement de la définition de l'extension par zéro).
Comme $(Rj_*F)|_U = F$ et $(j_!F)|_U = F$, nous voyons que (2) est vérifié.

\end{proof}

\begin{lemma}

\label{lemma-mayer-vietoris-cup}

Soit $(X, \mathcal{O}_X)$ un espace annelé. Posons $A = \Gamma(X, \mathcal{O}_X)$. Supposons que $X = U \cup V$ soit la réunion de deux sous-ensembles ouverts.
Pour des objets $K$ et $M$ de $D(\mathcal{O}_X)$, nous avons un morphisme de triangles distingués
$$
\xymatrix{
R\Gamma(X, K) \otimes_A^\mathbf{L} R\Gamma(X, M) \ar[r] \ar[d] &
R\Gamma(X, K \otimes_{\mathcal{O}_X}^\mathbf{L} M) \ar[d] \\
R\Gamma(X, K) \otimes_A^\mathbf{L}
(R\Gamma(U, M) \oplus R\Gamma(V, M)) \ar[r] \ar[d] &
R\Gamma(U, K \otimes_{\mathcal{O}_X}^\mathbf{L} M)
\oplus R\Gamma(V, K \otimes_{\mathcal{O}_X}^\mathbf{L} M)) \ar[d] \\
R\Gamma(X, K) \otimes_A^\mathbf{L} R\Gamma(U \cap V, M) \ar[r] \ar[d] &
R\Gamma(U \cap V, K \otimes_{\mathcal{O}_X}^\mathbf{L} M) \ar[d] \\
R\Gamma(X, K) \otimes_A^\mathbf{L} R\Gamma(X, M)[1] \ar[r] &
R\Gamma(X, K \otimes_{\mathcal{O}_X}^\mathbf{L} M)[1]
}
$$
où

\begin{enumerate}

\item les flèches horizontales sont données par le cup-produit,
\item sur le côté droit nous avons le triangle distingué du lemme \ref{lemma-unbounded-mayer-vietoris} pour
$K \otimes_{\mathcal{O}_X}^\mathbf{L} M$,
\item sur le côté gauche nous avons le foncteur exact
$R\Gamma(X, K) \otimes_A^\mathbf{L} - $ appliqué au triangle distingué du lemme \ref{lemma-unbounded-mayer-vietoris} pour $M$.

\end{enumerate}

\end{lemma}

\begin{proof}

Choisissons un complexe K-plat $T^\bullet$ de $A$-modules plats représentant
$R\Gamma(X, K)$, voir Compléments sur l'algèbre, lemme \ref{more-algebra-lemma-K-flat-resolution}. Notons $T^\bullet \otimes_A \mathcal{O}_X$ l'image inverse de $T^\bullet$
par le morphisme d'espaces annelés $(X, \mathcal{O}_X) \to (pt, A)$. Il existe un morphisme d'adjonction naturel
$\epsilon : T^\bullet \otimes_A \mathcal{O}_X \to K$ dans $D(\mathcal{O}_X)$. Remarquons que $T^\bullet \otimes_A \mathcal{O}_X$ est un complexe K-plat de $\mathcal{O}_X$-modules à termes plats; voir le lemme \ref{lemma-pullback-K-flat} et Modules, lemme \ref{modules-lemma-pullback-flat}. D'après le lemme \ref{lemma-factor-through-K-flat}, nous pouvons trouver un morphisme de complexes
$$
T^\bullet \otimes_A \mathcal{O}_X \longrightarrow \mathcal{K}^\bullet
$$
de $\mathcal{O}_X$-modules représentant $\epsilon$
tel que $\mathcal{K}^\bullet$ soit un complexe K-plat à termes plats. En effet, par la construction de
$D(\mathcal{O}_X)$, nous pouvons d'abord représenter $\epsilon$ par un morphisme de complexes
$e : T^\bullet \otimes_A \mathcal{O}_X \to \mathcal{L}^\bullet$
de $\mathcal{O}_X$-modules représentant $\epsilon$
puis appliquer le lemme à $e$. Choisissons un complexe K-injectif $\mathcal{I}^\bullet$ dont les termes sont des $\mathcal{O}_X$-modules injectifs et qui représente $M$. Enfin, choisissons un quasi-isomorphisme
$$
\text{Tot}(\mathcal{K}^\bullet \otimes_\mathcal{O} \mathcal{I}^\bullet)
\longrightarrow
\mathcal{J}^\bullet
$$
vers un complexe K-injectif dont les termes sont des $\mathcal{O}_X$-modules injectifs. Remarquons que la source et la cible de cette flèche représentent
$K \otimes_{\mathcal{O}_X}^\mathbf{L} M$ dans $D(\mathcal{O}_X)$. À ce stade, pour tout ouvert $W \subset X$, nous obtenons un morphisme de complexes
$$
\text{Tot}(T^\bullet \otimes_A \mathcal{I}^\bullet(W))
\to
\text{Tot}(\mathcal{K}^\bullet(W) \otimes_A \mathcal{I}^\bullet(W))
\to
\mathcal{J}^\bullet(W)
$$
de $A$-modules dont la composition représente l'application
$$
R\Gamma(X, K) \otimes_A^\mathbf{L} R\Gamma(W, M)
\longrightarrow
R\Gamma(W, K \otimes_{\mathcal{O}_X}^\mathbf{L} M)
$$
dans $D(A)$. Ces morphismes sont manifestement compatibles aux applications de restriction. Nous pouvons donc considérer le diagramme commutatif suivant de complexes de $A$-modules
$$
\xymatrix{
0 \ar[d] & 0 \ar[d] \\
\text{Tot}(T^\bullet \otimes_A \mathcal{I}^\bullet(X)) \ar[d] \ar[r] &
\mathcal{J}^\bullet(X) \ar[d] \\
\text{Tot}(T^\bullet \otimes_A
(\mathcal{I}^\bullet(U) \oplus \mathcal{I}^\bullet(V)) \ar[d] \ar[r] &
\mathcal{J}^\bullet(U) \oplus \mathcal{J}^\bullet(V) \ar[d] \\
\text{Tot}(T^\bullet \otimes_A \mathcal{I}^\bullet(U \cap V)) \ar[r] \ar[d] &
\mathcal{J}^\bullet(U \cap V) \ar[d] \\
0 & 0
}
$$
D'après la preuve du lemme \ref{lemma-mayer-vietoris}, les colonnes sont des suites exactes de complexes de $A$-modules (ce qui utilise aussi le fait que
$\text{Tot}(T^\bullet \otimes_A -)$ transforme les suites exactes courtes de complexes de $A$-modules en suites exactes courtes, puisque les termes de $T^\bullet$ sont des $A$-modules plats). Comme les triangles distingués du lemme \ref{lemma-unbounded-mayer-vietoris}
sont les triangles distingués associés à ces suites exactes courtes de complexes, le résultat souhaité découle de la fonctorialité de la construction du triangle distingué associé, étudiée dans Catégories dérivées, section \ref{derived-section-canonical-delta-functor}.

\end{proof}











\section{Cohomologie à support dans une partie fermée, II}

\label{section-cohomology-support-bis}

\noindent
Nous poursuivons la discussion commencée dans la section \ref{section-cohomology-support}.

\medskip\noindent
Soit $(X, \mathcal{O}_X)$ un espace annelé. Soit $Z \subset X$ une partie fermée. Dans cette situation, nous pouvons considérer le foncteur $\textit{Mod}(\mathcal{O}_X) \to \textit{Mod}(\mathcal{O}_X(X))$ donné par $\mathcal{F} \mapsto \Gamma_Z(X, \mathcal{F})$. Voir Modules, définition \ref{modules-definition-support} et Modules, lemme \ref{modules-lemma-support-section-closed}. En utilisant des résolutions K-injectives (voir la section \ref{section-unbounded}), nous obtenons le foncteur dérivé à droite $$
R\Gamma_Z(X, - ) : D(\mathcal{O}_X) \to D(\mathcal{O}_X(X))
$$ Pour un objet $K$ de $D(\mathcal{O}_X)$, nous notons $H^q_Z(X, K) = H^q(R\Gamma_Z(X, K))$ le module de cohomologie à support dans $Z$. Nous verrons plus loin (lemme \ref{lemma-sections-support-abelian-unbounded}) que cette notation coïncide avec la construction de la section \ref{section-cohomology-support}.

\medskip\noindent
Pour un $\mathcal{O}_X$-module $\mathcal{F}$, nous pouvons considérer le {\it sous-faisceau des sections à support dans $Z$}, noté $\mathcal{H}_Z(\mathcal{F})$, défini par la formule $$
\mathcal{H}_Z(\mathcal{F})(U) =
\{s \in \mathcal{F}(U) \mid \text{Supp}(s) \subset U \cap Z\} =
\Gamma_{Z \cap U}(U, \mathcal{F}|_U)
$$ Comme expliqué dans Modules, remarque \ref{modules-remark-sections-support-in-closed-modules}, nous pouvons voir $\mathcal{H}_Z(\mathcal{F})$ comme un $\mathcal{O}_X|_Z$-module sur $Z$, et nous obtenons un foncteur $$
\textit{Mod}(\mathcal{O}_X) \longrightarrow \textit{Mod}(\mathcal{O}_X|_Z),
\quad
\mathcal{F} \longmapsto \mathcal{H}_Z(\mathcal{F})
\text{ considéré comme }\mathcal{O}_X|_Z\text{-module sur }Z
$$ Ce foncteur est exact à gauche, mais n'est pas exact en général. Exactement comme ci-dessus, nous obtenons un foncteur dérivé à droite $$
R\mathcal{H}_Z : D(\mathcal{O}_X) \longrightarrow D(\mathcal{O}_X|_Z)
$$ Nous posons $\mathcal{H}^q_Z(K) = H^q(R\mathcal{H}_Z(K))$, de sorte que $\mathcal{H}^0_Z(\mathcal{F}) = \mathcal{H}_Z(\mathcal{F})$ pour tout faisceau de $\mathcal{O}_X$-modules $\mathcal{F}$.

\begin{lemma}

\label{lemma-cohomology-with-support-sheaf-on-support}
Soit $(X, \mathcal{O}_X)$ un espace annelé. Soit $i : Z \to X$ l'inclusion d'une partie fermée.
\begin{enumerate}
\item $R\mathcal{H}_Z : D(\mathcal{O}_X) \to D(\mathcal{O}_X|_Z)$ est adjoint à droite de $i_* : D(\mathcal{O}_X|_Z) \to D(\mathcal{O}_X)$.
\item Pour $K$ dans $D(\mathcal{O}_X|_Z)$ nous avons $R\mathcal{H}_Z(i_*K) = K$.
\item Soit $\mathcal{G}$ un faisceau de $\mathcal{O}_X|_Z$-modules sur $Z$. Alors $\mathcal{H}^p_Z(i_*\mathcal{G}) = 0$ pour $p > 0$.
\end{enumerate}

\end{lemma}

\begin{proof}

Le foncteur $i_*$ est exact, donc $i_* = Ri_* = Li_*$. Par conséquent, la partie (1) du lemme découle de Modules, lemme \ref{modules-lemma-adjoint-section-with-support},
et de Catégories dérivées, lemme \ref{derived-lemma-derived-adjoint-functors}. Soit $K$ comme dans (2). Nous pouvons représenter $K$ par un complexe K-injectif
$\mathcal{I}^\bullet$ de $\mathcal{O}_X|_Z$-modules. D'après le lemme \ref{lemma-K-injective-flat},
le complexe $i_*\mathcal{I}^\bullet$, qui représente $i_*K$, est un complexe K-injectif de $\mathcal{O}_X$-modules. Ainsi,
$R\mathcal{H}_Z(i_*K)$ est calculé par
$\mathcal{H}_Z(i_*\mathcal{I}^\bullet) = \mathcal{I}^\bullet$
ce qui prouve (2). La partie (3) est un cas particulier de (2).

\end{proof}

\noindent
Soit $(X, \mathcal{O}_X)$ un espace annelé et soit $Z \subset X$ une partie fermée. La catégorie des $\mathcal{O}_X$-modules dont le support est contenu dans $Z$ est une sous-catégorie de Serre de la catégorie de tous les $\mathcal{O}_X$-modules, voir Homologie, définition \ref{homology-definition-serre-subcategory} et Modules, lemme \ref{modules-lemma-support-section-closed}. Nous notons $D_Z(\mathcal{O}_X)$ la sous-catégorie triangulée strictement saturée de $D(\mathcal{O}_X)$ formée des complexes dont les faisceaux de cohomologie sont à support dans $Z$, voir Catégories dérivées, section \ref{derived-section-triangulated-sub}.

\begin{lemma}

\label{lemma-complexes-with-support-on-closed}
Soit $(X, \mathcal{O}_X)$ un espace annelé. Soit $i : Z \to X$ l'inclusion d'une partie fermée.
\begin{enumerate}
\item Pour $K$ dans $D(\mathcal{O}_X|_Z)$ nous avons $i_*K$ dans $D_Z(\mathcal{O}_X)$.
\item Le foncteur $i_* : D(\mathcal{O}_X|_Z) \to D_Z(\mathcal{O}_X)$ est une équivalence, de quasi-inverse $i^{-1}|_{D_Z(\mathcal{O}_X)} = R\mathcal{H}_Z|_{D_Z(\mathcal{O}_X)}$.
\item Le foncteur $i_* \circ R\mathcal{H}_Z : D(\mathcal{O}_X) \to D_Z(\mathcal{O}_X)$ est adjoint à droite du foncteur d'inclusion $D_Z(\mathcal{O}_X) \to D(\mathcal{O}_X)$.
\end{enumerate}

\end{lemma}

\begin{proof}
La partie (1) résulte immédiatement des définitions. La partie (3) est une conséquence formelle de la partie (2) et du lemme \ref{lemma-cohomology-with-support-sheaf-on-support}. Dans la suite de la preuve, nous démontrons la partie (2).

\medskip\noindent

Considérons $i$ comme le morphisme d'espaces annelés
$i : (Z, \mathcal{O}_X|_Z) \to (X, \mathcal{O}_X)$. Rappelons que $i^*$ et $i_*$ forment une paire de foncteurs adjoints. Puisque $i$ est une immersion fermée, $i_*$ est exact. Comme $i^{-1}\mathcal{O}_X = \mathcal{O}_X|_Z$ est le faisceau structural de $(Z, \mathcal{O}_X|_Z)$, nous voyons que $i^* = i^{-1}$
est exact et que $i^*i_* = i^{-1}i_*$
est isomorphe au foncteur identité. Voir Modules, lemmes \ref{modules-lemma-exactness-pushforward-pullback} et
\ref{modules-lemma-i-star-exact}. Ainsi
$i_* : D(\mathcal{O}_X|_Z) \to D_Z(\mathcal{O}_X)$
est pleinement fidèle et $i^{-1}$ fournit un inverse à gauche. D'autre part, supposons que $K$ soit un objet de
$D_Z(\mathcal{O}_X)$ et considérons le morphisme d'adjonction
$K \to i_*i^{-1}K$. L'exactitude de $i_*$ et de $i^{-1}$ montre que celui-ci induit les morphismes d'adjonction
$H^n(K) \to i_*i^{-1}H^n(K)$ sur les faisceaux de cohomologie. Comme ces faisceaux sont à support dans $Z$, ces morphismes d'adjonction sont des isomorphismes; nous concluons donc que $i_* : D(\mathcal{O}_X|_Z) \to D_Z(\mathcal{O}_X)$
est une équivalence.

\medskip\noindent

Pour terminer la preuve, il suffit de montrer que $R\mathcal{H}_Z(K) = i^{-1}K$ lorsque
$K$ est un objet de $D_Z(\mathcal{O}_X)$. Pour cela, nous pouvons utiliser l'égalité
$K = i_*i^{-1}K$, que nous venons de démontrer. Le lemme \ref{lemma-cohomology-with-support-sheaf-on-support}
donne alors le résultat souhaité.

\end{proof}

\begin{lemma}

\label{lemma-sections-with-support-K-injective}
Soit $(X, \mathcal{O}_X)$ un espace annelé. Soit $i : Z \to X$ l'inclusion d'une partie fermée. Si $\mathcal{I}^\bullet$ est un complexe K-injectif de $\mathcal{O}_X$-modules, alors $\mathcal{H}_Z(\mathcal{I}^\bullet)$ est un complexe K-injectif de $\mathcal{O}_X|_Z$-modules.
\end{lemma}

\begin{proof}
Puisque $i_* : \textit{Mod}(\mathcal{O}_X|_Z) \to \textit{Mod}(\mathcal{O}_X)$ est exact et adjoint à gauche de $\mathcal{H}_Z$ (Modules, lemme \ref{modules-lemma-adjoint-section-with-support}), ceci découle de Catégories dérivées, lemme \ref{derived-lemma-adjoint-preserve-K-injectives}.
\end{proof}

\begin{lemma}

\label{lemma-local-to-global-sections-with-support}
Soit $(X, \mathcal{O}_X)$ un espace annelé. Soit $i : Z \to X$ l'inclusion d'une partie fermée. Alors $R\Gamma(Z, - ) \circ R\mathcal{H}_Z = R\Gamma_Z(X, - )$ comme foncteurs $D(\mathcal{O}_X) \to D(\mathcal{O}_X(X))$.
\end{lemma}

\begin{proof}
Cela résulte de la construction des foncteurs dérivés à droite au moyen de résolutions K-injectives, du lemme \ref{lemma-sections-with-support-K-injective} et de l'égalité $\Gamma_Z(X, -) = \Gamma(Z, -) \circ \mathcal{H}_Z$.
\end{proof}

\begin{lemma}

\label{lemma-triangle-sections-with-support}
Soit $(X, \mathcal{O}_X)$ un espace annelé. Soit $i : Z \to X$ l'inclusion d'une partie fermée. Posons $U = X \setminus Z$. Il existe un triangle distingué $$
R\Gamma_Z(X, K) \to R\Gamma(X, K) \to R\Gamma(U, K) \to
R\Gamma_Z(X, K)[1]
$$ dans $D(\mathcal{O}_X(X))$, fonctoriellement en $K$ dans $D(\mathcal{O}_X)$.
\end{lemma}

\begin{proof}

Choisissons un complexe K-injectif $\mathcal{I}^\bullet$ dont tous les termes sont des $\mathcal{O}_X$-modules injectifs et qui représente $K$. Voir la section \ref{section-unbounded}. Rappelons que $\mathcal{I}^\bullet|_U$
est un complexe K-injectif de $\mathcal{O}_U$-modules; voir le lemme \ref{lemma-restrict-K-injective-to-open}. Par conséquent, chacun des foncteurs dérivés du triangle distingué s'obtient en appliquant le foncteur sous-jacent à $\mathcal{I}^\bullet$. Il suffit donc de prouver que, pour tout $\mathcal{O}_X$-module injectif $\mathcal{I}$, nous avons une suite exacte courte
$$
0 \to \Gamma_Z(X, \mathcal{I}) \to \Gamma(X, \mathcal{I})
\to \Gamma(U, \mathcal{I}) \to 0
$$
Cela découle du lemme \ref{lemma-injective-restriction-surjective}
et des définitions.

\end{proof}

\begin{lemma}

\label{lemma-triangle-sections-with-support-sheaves}
Soit $(X, \mathcal{O}_X)$ un espace annelé. Soit $i : Z \to X$ l'inclusion d'une partie fermée. Notons $j : U = X \setminus Z \to X$ l'inclusion du complémentaire. Il existe un triangle distingué $$
i_*R\mathcal{H}_Z(K) \to K \to Rj_*(K|_U) \to
i_*R\mathcal{H}_Z(K)[1]
$$ dans $D(\mathcal{O}_X)$, fonctoriellement en $K$ dans $D(\mathcal{O}_X)$.
\end{lemma}

\begin{proof}

Choisissons un complexe K-injectif $\mathcal{I}^\bullet$ dont tous les termes sont des $\mathcal{O}_X$-modules injectifs et qui représente $K$. Voir la section \ref{section-unbounded}. Rappelons que $\mathcal{I}^\bullet|_U$
est un complexe K-injectif de $\mathcal{O}_U$-modules; voir le lemme \ref{lemma-restrict-K-injective-to-open}. Par conséquent, chacun des foncteurs dérivés du triangle distingué s'obtient en appliquant le foncteur sous-jacent à $\mathcal{I}^\bullet$. Il suffit donc de prouver que, pour tout $\mathcal{O}_X$-module injectif $\mathcal{I}$, nous avons une suite exacte courte
$$
0 \to i_*\mathcal{H}_Z(\mathcal{I}) \to \mathcal{I}
\to j_*(\mathcal{I}|_U) \to 0
$$
Cela découle du lemme \ref{lemma-injective-restriction-surjective}
et des définitions.

\end{proof}

\begin{lemma}

\label{lemma-sections-support-in-closed-disjoint-open}
Soit $(X, \mathcal{O}_X)$ un espace annelé. Soit $Z \subset X$ une partie fermée. Soit $j : U \to X$ l'inclusion d'un sous-ensemble ouvert tel que $U \cap Z = \emptyset$. Alors $R\mathcal{H}_Z(Rj_*K) = 0$ pour tout $K$ dans $D(\mathcal{O}_U)$.
\end{lemma}

\begin{proof}

Choisissons un complexe K-injectif $\mathcal{I}^\bullet$ de $\mathcal{O}_U$-modules représentant $K$. Alors $j_*\mathcal{I}^\bullet$ représente $Rj_*K$. D'après le lemme \ref{lemma-K-injective-flat}, le complexe $j_*\mathcal{I}^\bullet$ est un complexe K-injectif de $\mathcal{O}_X$-modules. Ainsi,
$\mathcal{H}_Z(j_*\mathcal{I}^\bullet)$ représente $R\mathcal{H}_Z(Rj_*K)$. Il suffit donc de montrer que $\mathcal{H}_Z(j_*\mathcal{G}) = 0$
pour tout faisceau abélien $\mathcal{G}$ sur $U$. Il faut donc montrer qu'une section $s$ de $j_*\mathcal{G}$ sur un ouvert $W$, à support dans $W \cap Z$, est nulle. La condition sur le support signifie que
$s|_{W \setminus W \cap Z} = 0$. Comme $j_*\mathcal{G}(W) =
\mathcal{G}(U \cap W) = j_*\mathcal{G}(W \setminus W \cap Z)$
cela implique que $s$ est nulle, comme souhaité.

\end{proof}

\begin{lemma}

\label{lemma-sections-support-abelian-unbounded}
Soit $(X, \mathcal{O}_X)$ un espace annelé. Soit $Z \subset X$ une partie fermée. Soit $K$ un objet de $D(\mathcal{O}_X)$ et notons $K_{ab}$ son image dans $D(\underline{\mathbf{Z}}_X)$.
\begin{enumerate}
\item Il y a une application canonique $R\Gamma_Z(X, K) \to R\Gamma_Z(X, K_{ab})$ qui est un isomorphisme dans $D(\textit{Ab})$.
\item Il y a une application canonique $R\mathcal{H}_Z(K) \to R\mathcal{H}_Z(K_{ab})$ qui est un isomorphisme dans $D(\underline{\mathbf{Z}}_Z)$.
\end{enumerate}

\end{lemma}

\begin{proof}

Démonstration de (1). Le morphisme est construit comme suit. Choisissons un complexe K-injectif de $\mathcal{O}_X$-modules $\mathcal{I}^\bullet$ représentant $K$. Choisissons un quasi-isomorphisme
$\mathcal{I}^\bullet \to \mathcal{J}^\bullet$ où $\mathcal{J}^\bullet$
est un complexe K-injectif de groupes abéliens. Alors le morphisme de (1) est donné par
$$
\Gamma_Z(X, \mathcal{I}^\bullet) \to \Gamma_Z(X, \mathcal{J}^\bullet)
$$
déterminé par le fait que $\Gamma_Z$ est un foncteur sur les faisceaux abéliens. Une vérification immédiate montre que le morphisme obtenu et les morphismes canoniques du lemme \ref{lemma-modules-abelian-unbounded}
s'insèrent dans un morphisme de triangles distingués
$$
\xymatrix{
R\Gamma_Z(X, K) \ar[r] \ar[d] &
R\Gamma(X, K) \ar[r] \ar[d] &
R\Gamma(U, K) \ar[d] \\
R\Gamma_Z(X, K_{ab}) \ar[r] &
R\Gamma(X, K_{ab}) \ar[r] &
R\Gamma(U, K_{ab})
}
$$
du lemme \ref{lemma-triangle-sections-with-support}. Comme deux des trois flèches sont des isomorphismes d'après le lemme cité, nous concluons par Catégories dérivées, lemme
\ref{derived-lemma-third-isomorphism-triangle}.

\medskip\noindent

La preuve de (2) est omise. Indication : utiliser le même argument avec le lemme \ref{lemma-triangle-sections-with-support-sheaves}
pour le triangle distingué.

\end{proof}

\begin{remark}

\label{remark-support-cup-product}

Soit $(X, \mathcal{O}_X)$ un espace annelé. Soit $i : Z \to X$
l'inclusion d'une partie fermée. Étant donnés $K$ et $M$ dans
$D(\mathcal{O}_X)$, il existe un morphisme canonique
$$
K|_Z \otimes_{\mathcal{O}_X|_Z}^\mathbf{L} R\mathcal{H}_Z(M)
\longrightarrow
R\mathcal{H}_Z(K \otimes_{\mathcal{O}_X}^\mathbf{L} M)
$$
dans $D(\mathcal{O}_X|_Z)$. Ici, $K|_Z = i^{-1}K$ est la restriction de
$K$ à $Z$, considérée comme un objet de $D(\mathcal{O}_X|_Z)$. Par l'adjonction entre $i_*$ et $R\mathcal{H}_Z$ du lemme \ref{lemma-cohomology-with-support-sheaf-on-support}, il suffit, pour construire ce morphisme, de produire un morphisme canonique
$$
i_*\left(K|_Z \otimes_{\mathcal{O}_X|_Z}^\mathbf{L} R\mathcal{H}_Z(M)\right)
\longrightarrow
K \otimes_{\mathcal{O}_X}^\mathbf{L} M
$$
Pour construire ce morphisme, choisissons un complexe K-injectif $\mathcal{I}^\bullet$
de $\mathcal{O}_X$-modules représentant $M$ et un complexe K-plat
$\mathcal{K}^\bullet$ de $\mathcal{O}_X$-modules représentant $K$. Remarquons que $\mathcal{K}^\bullet|_Z$ est un complexe K-plat de
$\mathcal{O}_X|_Z$-modules représentant $K|_Z$; voir le lemme \ref{lemma-pullback-K-flat}. Nous devons donc produire un morphisme de complexes
$$
i_*\text{Tot}\left(
\mathcal{K}^\bullet|_Z \otimes_{\mathcal{O}_X|_Z}
\mathcal{H}_Z(\mathcal{I}^\bullet)\right)
\longrightarrow
\text{Tot}(\mathcal{K}^\bullet \otimes_{\mathcal{O}_X} \mathcal{I}^\bullet)
$$
de $\mathcal{O}_X$-modules. Pour cela, il suffit de produire des morphismes
$$
i_*(\mathcal{K}^a|_Z \otimes_{\mathcal{O}_X|_Z}
\mathcal{H}_Z(\mathcal{I}^b))
\longrightarrow
\mathcal{K}^a \otimes_{\mathcal{O}_X} \mathcal{I}^b
$$
En regardant les germes (par exemple), nous voyons que le côté gauche de cette formule est égal à
$\mathcal{K}^a \otimes_{\mathcal{O}_X} i_*\mathcal{H}_Z(\mathcal{I}^b)$
et nous pouvons utiliser l'inclusion
$\mathcal{H}_Z(\mathcal{I}^b) \to \mathcal{I}^b$ pour obtenir notre morphisme.

\end{remark}

\begin{remark}

\label{remark-support-cup-product-global}

Avec les notations de la remarque \ref{remark-support-cup-product},
nous obtenons un cup-produit canonique

\begin{align*}
H^a(X, K) \times H^b_Z(X, M)
& =
H^a(X, K) \times H^b(Z, R\mathcal{H}_Z(M)) \\
& \to
H^a(Z, K|_Z) \times H^b(Z, R\mathcal{H}_Z(M)) \\
& \to
H^{a + b}(Z, K|_Z \otimes_{\mathcal{O}_X|_Z}^\mathbf{L} R\mathcal{H}_Z(M)) \\
& \to
H^{a + b}(Z, R\mathcal{H}_Z(K \otimes_{\mathcal{O}_X}^\mathbf{L} M)) \\
& =
H^{a + b}_Z(X, K \otimes_{\mathcal{O}_X}^\mathbf{L} M)
\end{align*}

Ici, les signes d'égalité sont donnés par le lemme \ref{lemma-local-to-global-sections-with-support}; la première flèche est la restriction à $Z$, la deuxième est le cup-produit (section \ref{section-cup-product}) et la troisième est le morphisme de la remarque \ref{remark-support-cup-product}.

\end{remark}

\begin{lemma}

\label{lemma-support-cup-product}
Avec les notations de la remarque \ref{remark-support-cup-product}, le diagramme $$
\xymatrix{
H^i(X, K) \times H^j_Z(X, M) \ar[r] \ar[d] &
H^{i + j}_Z(X, K \otimes_{\mathcal{O}_X}^\mathbf{L} M) \ar[d] \\
H^i(X, K) \times H^j(X, M) \ar[r] &
H^{i + j}(X, K \otimes_{\mathcal{O}_X}^\mathbf{L} M)
}
$$ commute, la flèche horizontale supérieure étant le cup-produit de la remarque \ref{remark-support-cup-product-global}.
\end{lemma}

\begin{proof}
Démonstration omise.
\end{proof}

\begin{remark}

\label{remark-support-functorial}

Soit $f : (X', \mathcal{O}_{X'}) \to (X, \mathcal{O}_X)$ un morphisme d'espaces annelés. Soit $Z \subset X$ une partie fermée et posons $Z' = f^{-1}(Z)$. Notons $f|_{Z'} : (Z', \mathcal{O}_{X'}|_{Z'}) \to (Z, \mathcal{O}_X|Z)$ le morphisme induit d'espaces annelés. Pour tout $K$ dans $D(\mathcal{O}_X)$, il existe un morphisme canonique
$$
L(f|_{Z'})^*R\mathcal{H}_Z(K) \longrightarrow R\mathcal{H}_{Z'}(Lf^*K)
$$
dans $D(\mathcal{O}_{X'}|_{Z'})$. Notons $i : Z \to X$ et $i' : Z' \to X'$
les morphismes d'inclusion. D'après le lemme \ref{lemma-complexes-with-support-on-closed}, partie (2), appliqué à $i'$, cela revient à donner un morphisme
$$
i'_* L(f|_{Z'})^* R\mathcal{H}_Z(K)
\longrightarrow
i'_*R\mathcal{H}_{Z'}(Lf^*K)
$$
dans $D_{Z'}(\mathcal{O}_{X'})$. Le morphisme de foncteurs
$Lf^* \circ i_* \to i'_* \circ L(f|_{Z'})^*$ de la remarque \ref{remark-base-change} est un isomorphisme dans ce cas (cela se vérifie sur les germes en utilisant que $i_*$ et $i'_*$
sont exactes et $\mathcal{O}_{Z, z} = \mathcal{O}_{X, z}$
et de même pour $Z'$). Il suffit donc de construire la flèche horizontale supérieure dans le diagramme suivant
$$
\xymatrix{
Lf^* i_* R\mathcal{H}_Z(K) \ar[rr] \ar[rd] & &
i'_* R\mathcal{H}_{Z'}(Lf^*K) \ar[ld] \\
& Lf^*K
}
$$
Le complexe $Lf^* i_* R\mathcal{H}_Z(K)$ est à support dans $Z'$. La flèche sud-est provient du morphisme d'adjonction $i_*R\mathcal{H}_Z(K) \to K$
(lemme \ref{lemma-cohomology-with-support-sheaf-on-support}). Puisque le morphisme d'adjonction $i'_* R\mathcal{H}_{Z'}(Lf^*K) \to Lf^*K$ est universel d'après le lemme \ref{lemma-complexes-with-support-on-closed}, partie (3), la flèche sud-est se factorise de manière unique par la flèche sud-ouest, ce qui fournit la flèche souhaitée.

\end{remark}

\begin{lemma}

\label{lemma-support-functorial}
Avec les notations et les hypothèses de la remarque \ref{remark-support-functorial}, le diagramme $$
\xymatrix{
H^p_Z(X, K) \ar[r] \ar[d] & H^p_{Z'}(X', Lf^*K) \ar[d] \\
H^p(X, K) \ar[r] & H^p(X', Lf^*K)
}
$$ commute. Ici, la flèche horizontale supérieure provient des identifications $H^p_Z(X, K) = H^p(Z, R\mathcal{H}_Z(K))$ et $H^p_{Z'}(X', Lf^*K) = H^p(Z', R\mathcal{H}_{Z'}(K'))$, du morphisme d'image inverse $H^p(Z, R\mathcal{H}_Z(K)) \to H^p(Z', L(f|_{Z'})^*R\mathcal{H}_Z(K))$ et du morphisme construit dans la remarque \ref{remark-support-functorial}.
\end{lemma}

\begin{proof}
Démonstration omise. Indication : en utilisant $H^p(Z, R\mathcal{H}_Z(K)) = H^p(X, i_*R\mathcal{H}_Z(K))$, et de même pour $R\mathcal{H}_{Z'}(Lf^*K)$, le résultat découle de la fonctorialité des morphismes d'image inverse et du diagramme commutatif utilisé pour définir le morphisme de la remarque \ref{remark-support-functorial}.
\end{proof}





\section{Systèmes projectifs et cohomologie, I}

\label{section-inverse-systems}

\noindent
Soit $A$ un anneau et soit $I \subset A$ un idéal. Nous démontrons quelques résultats sur des systèmes projectifs de faisceaux de $A/I^n$-modules.

\begin{lemma}

\label{lemma-ML-general}

Soit $I$ un idéal d'un anneau $A$. Soit $X$ un espace topologique. Soit
$$
\ldots \to \mathcal{F}_3 \to \mathcal{F}_2 \to \mathcal{F}_1
$$
soit un système projectif de faisceaux de $A$-modules sur $X$
tel que $\mathcal{F}_n = \mathcal{F}_{n + 1}/I^n\mathcal{F}_{n + 1}$. Soit $p \geq 0$. Supposons que
$$
\bigoplus\nolimits_{n \geq 0} H^{p + 1}(X, I^n\mathcal{F}_{n + 1})
$$
satisfait à la condition de chaîne ascendante comme module gradué sur
$\bigoplus_{n \geq 0} I^n/I^{n + 1}$. Alors le système projectif $M_n = H^p(X, \mathcal{F}_n)$ satisfait à la condition de Mittag--Leffler\footnote{En fait, il existe un $c \geq 0$ tel que $\Im(M_n \to M_{n - c})$ soit l'image stable pour tout $n \geq c$.}.

\end{lemma}

\begin{proof}

Posons $N_n = H^{p + 1}(X, I^n\mathcal{F}_{n + 1})$ et soit
$\delta_n : M_n \to N_n$ l'application de bord en cohomologie provenant de la suite exacte courte
$0 \to I^n\mathcal{F}_{n + 1} \to \mathcal{F}_{n + 1} \to \mathcal{F}_n \to 0$. Alors $\bigoplus \Im(\delta_n) \subset \bigoplus N_n$ est un sous-module gradué. Soient $s \in M_n$ et $f \in I^m$; nous avons alors un diagramme commutatif
$$
\xymatrix{
0 \ar[r] &
I^n\mathcal{F}_{n + 1} \ar[d]_f \ar[r] &
\mathcal{F}_{n + 1} \ar[d]_f \ar[r] &
\mathcal{F}_n \ar[d]_f \ar[r] & 0 \\
0 \ar[r] &
I^{n + m}\mathcal{F}_{n + m + 1} \ar[r] &
\mathcal{F}_{n + m + 1} \ar[r] &
\mathcal{F}_{n + m} \ar[r] & 0
}
$$
Le morphisme vertical du milieu s'obtient en relevant une section locale de
$\mathcal{F}_{n + 1}$ en une section de $\mathcal{F}_{n + m + 1}$,
puis en la multipliant par $f$; il en va de même pour les autres flèches verticales. Nous en déduisons que $\delta_{n + m}(fs) = f \delta_n(s)$. Par hypothèse, nous pouvons trouver $s_j \in M_{n_j}$, $j = 1, \ldots, N$,
tels que les $\delta_{n_j}(s_j)$
engendrent $\bigoplus \Im(\delta_n)$ comme module gradué. Soit $n > c = \max(n_j)$ et soit $s \in M_n$. Nous pouvons alors trouver $f_j \in I^{n - n_j}$ tels que
$\delta_n(s) = \sum f_j \delta_{n_j}(s_j)$. Nous concluons que
$\delta(s - \sum f_j s_j) = 0$, c'est-à-dire, nous pouvons trouver $s' \in M_{n + 1}$
s'envoyant sur $s - \sum f_js_j$ dans $M_n$. Il s'ensuit que
$$
\Im(M_{n + 1} \to M_{n - c}) = \Im(M_n \to M_{n - c})
$$
En effet, les éléments $f_js_j$ s'envoient sur zéro dans $M_{n - c}$. Cela prouve le lemme.

\end{proof}

\begin{lemma}

\label{lemma-ML-general-better}

Soit $I$ un idéal d'un anneau $A$. Soit $X$ un espace topologique. Soit
$$
\ldots \to \mathcal{F}_3 \to \mathcal{F}_2 \to \mathcal{F}_1
$$
un système projectif de faisceaux de $A$-modules sur $X$
tel que $\mathcal{F}_n = \mathcal{F}_{n + 1}/I^n\mathcal{F}_{n + 1}$. Soit $p \geq 0$. Pour $n$, posons
$$
N_n =
\bigcap\nolimits_{m \geq n}
\Im\left(
H^{p + 1}(X, I^n\mathcal{F}_{m + 1}) \to H^{p + 1}(X, I^n\mathcal{F}_{n + 1})
\right)
$$
Si $\bigoplus N_n$ satisfait à la condition de chaîne ascendante comme module gradué sur
$\bigoplus_{n \geq 0} I^n/I^{n + 1}$, alors le système projectif
$M_n = H^p(X, \mathcal{F}_n)$ satisfait à la condition de Mittag--Leffler\footnote{En fait, il existe un $c \geq 0$ tel que $\Im(M_n \to M_{n - c})$ soit l'image stable pour tout $n \geq c$.}.

\end{lemma}

\begin{proof}
La preuve est exactement la même que celle du lemme \ref{lemma-ML-general}. En effet, les arguments donnés là prouvent le résultat dès que l'on montre que $\bigoplus N_n$ est un sous-module gradué sur $\bigoplus_{n \geq 0} I^n/I^{n + 1}$ de $\bigoplus H^{p + 1}(X, I^n\mathcal{F}_{n + 1})$ et que les applications de bord $\delta_n : M_n \to H^{p + 1}(X, I^n\mathcal{F}_{n + 1})$ ont leur image contenue dans $N_n$.

\medskip\noindent
Supposons que $\xi \in N_n$ et $f \in I^k$. Choisissons $m \gg n + k$, puis $\xi' \in H^{p + 1}(X, I^n\mathcal{F}_{m + 1})$ relevant $\xi$. Considérons le diagramme $$
\xymatrix{
0 \ar[r] &
I^n\mathcal{F}_{m + 1} \ar[d]_f \ar[r] &
\mathcal{F}_{m + 1} \ar[d]_f \ar[r] &
\mathcal{F}_n \ar[d]_f \ar[r] & 0 \\
0 \ar[r] &
I^{n + k}\mathcal{F}_{m + 1} \ar[r] &
\mathcal{F}_{m + 1} \ar[r] &
\mathcal{F}_{n + k} \ar[r] & 0
}
$$ construit comme dans la preuve du lemme \ref{lemma-ML-general}. Nous obtenons un morphisme induit en cohomologie et nous voyons que $f \xi' \in H^{p + 1}(X, I^{n + k}\mathcal{F}_{m + 1})$ s'envoie sur $f \xi$. Puisque ceci vaut pour tout $m \gg n + k$, nous voyons que $f\xi$ appartient à $N_{n + k}$, comme souhaité.

\medskip\noindent

Pour voir que les applications de bord $\delta_n$ ont leur image contenue dans $N_n$,
nous considérons les diagrammes
$$
\xymatrix{
0 \ar[r] &
I^n\mathcal{F}_{m + 1} \ar[d] \ar[r] &
\mathcal{F}_{m + 1} \ar[d] \ar[r] &
\mathcal{F}_n \ar[d] \ar[r] & 0 \\
0 \ar[r] &
I^n\mathcal{F}_{n + 1} \ar[r] &
\mathcal{F}_{n + 1} \ar[r] &
\mathcal{F}_n \ar[r] & 0
}
$$
pour $m \geq n$. Le résultat s'obtient en examinant les morphismes induits en cohomologie.

\end{proof}

\begin{lemma}

\label{lemma-topology-I-adic-general}
Soit $I$ un idéal d'un anneau $A$. Soit $X$ un espace topologique. Soit $$
\ldots \to \mathcal{F}_3 \to \mathcal{F}_2 \to \mathcal{F}_1
$$ un système projectif de faisceaux de $A$-modules sur $X$ tel que $\mathcal{F}_n = \mathcal{F}_{n + 1}/I^n\mathcal{F}_{n + 1}$. Soit $p \geq 0$. Supposons que $$
\bigoplus\nolimits_{n \geq 0} H^p(X, I^n\mathcal{F}_{n + 1})
$$ satisfasse à la condition de chaîne ascendante comme module gradué sur $\bigoplus_{n \geq 0} I^n/I^{n + 1}$. Alors la topologie limite sur $M = \lim H^p(X, \mathcal{F}_n)$ est la topologie $I$-adique.
\end{lemma}

\begin{proof}

Posons $F^n = \Ker(M \to H^p(X, \mathcal{F}_n))$ pour $n \geq 1$ et $F^0 = M$. Observons que $I F^n \subset F^{n + 1}$. En particulier, $I^n M \subset F^n$. Ainsi, la topologie $I$-adique est plus fine que la topologie limite. Réciproquement, nous montrerons que, pour tout $n$,
il existe un $m \geq n$ tel que $F^m \subset I^nM$.\footnote{En fait, il existe un $c \geq 0$ tel que $F^{n + c} \subset I^nM$ pour tout $n$.} Nous avons des applications injectives
$$
F^n/F^{n + 1} \longrightarrow H^p(X, \mathcal{F}_{n + 1})
$$
dont l'image est contenue dans
$H^p(X, I^n\mathcal{F}_{n + 1}) \to H^p(X, \mathcal{F}_{n + 1})$. Notons
$$
E_n \subset H^p(X, I^n\mathcal{F}_{n + 1})
$$
l'image inverse de $F^n/F^{n + 1}$. Alors $\bigoplus E_n$ est un sous-module gradué sur $\bigoplus I^n/I^{n + 1}$ de
$\bigoplus H^p(X, I^n\mathcal{F}_{n + 1})$ et
$\bigoplus E_n \to \bigoplus F^n/F^{n + 1}$ est un homomorphisme de modules gradués; nous omettons les détails. Par hypothèse, $\bigoplus E_n$ est engendré par un nombre fini d'éléments homogènes sur $\bigoplus I^n/I^{n + 1}$. Comme $E_n \to F^n/F^{n + 1}$ est surjectif, il en va de même pour $\bigoplus F^n/F^{n + 1}$. Nous pouvons donc trouver $r$, des entiers $c_1, \ldots, c_r \geq 0$ et
des éléments $a_i \in F^{c_i}$ dont les images engendrent $\bigoplus F^n/F^{n + 1}$. Posons $c = \max(c_i)$.

\medskip\noindent

Pour $n \geq c$, nous affirmons que $I F^n = F^{n + 1}$. Cette assertion montre que
$F^{n + c} = I^nF^c \subset I^nM$, comme souhaité. Pour la démontrer, supposons que $a \in F^{n + 1}$. L'image de
$a$ en $F^{n + 1}/F^{n + 2}$ est une combinaison linéaire de notre $a_i$. Par conséquent $a  - \sum f_i a_i \in F^{n + 2}$
pour certains $f_i \in I^{n + 1 - c_i}$. Comme
$I^{n + 1 - c_i} = I \cdot I^{n - c_i}$ et $n \geq c_i$, nous pouvons écrire
$f_i = \sum g_{i, j} h_{i, j}$ avec $g_{i, j} \in I$
et $h_{i, j}a_i \in F^n$. Ainsi, nous voyons que
$F^{n + 1} = F^{n + 2} + IF^n$. Un simple argument d'induction donne $F^{n + 1} = F^{n + e} + IF^n$
pour tout $e > 0$. Il s'ensuit que $IF^n$ est dense dans $F^{n + 1}$. Choisissons des générateurs $k_1, \ldots, k_r$ de $I$ et considérons l'application continue
$$
u : (F^n)^{\oplus r} \longrightarrow F^{n + 1},\quad
(x_1, \ldots, x_r) \mapsto \sum k_i x_i
$$
(pour la topologie limite). D'après ce qui précède, l'image de $(F^m)^{\oplus r}$ par $u$ est dense dans
$F^{m + 1}$ pour tout $m \geq n$. Le lemme de l'application ouverte (Compléments sur l'algèbre, lemme \ref{more-algebra-lemma-open-mapping}) montre que $u$ est ouverte. L'application $u$ est donc surjective. Par conséquent, $IF^n = F^{n + 1}$
pour $n \geq c$. Ceci conclut la preuve.

\end{proof}

\begin{lemma}

\label{lemma-topology-I-adic-general-better}
Soit $I$ un idéal d'un anneau $A$. Soit $X$ un espace topologique. Soit $$
\ldots \to \mathcal{F}_3 \to \mathcal{F}_2 \to \mathcal{F}_1
$$ un système projectif de faisceaux de $A$-modules sur $X$ tel que $\mathcal{F}_n = \mathcal{F}_{n + 1}/I^n\mathcal{F}_{n + 1}$. Soit $p \geq 0$. Pour $n$, posons $$
N_n =
\bigcap\nolimits_{m \geq n}
\Im\left(
H^p(X, I^n\mathcal{F}_{m + 1}) \to H^p(X, I^n\mathcal{F}_{n + 1})
\right)
$$ Si $\bigoplus N_n$ satisfait à la condition de chaîne ascendante comme module gradué sur $\bigoplus_{n \geq 0} I^n/I^{n + 1}$, alors la topologie limite sur $M = \lim H^p(X, \mathcal{F}_n)$ est la topologie $I$-adique.
\end{lemma}

\begin{proof}
La preuve est exactement la même que celle du lemme \ref{lemma-topology-I-adic-general}. En effet, les arguments donnés là prouvent le résultat dès que l'on montre que $\bigoplus N_n$ est un sous-module gradué sur $\bigoplus_{n \geq 0} I^n/I^{n + 1}$ de $\bigoplus H^{p + 1}(X, I^n\mathcal{F}_{n + 1})$ et que $F^n/F^{n + 1} \subset H^p(X, \mathcal{F}_{n + 1})$ est contenu dans l'image de $N_n \to H^p(X, \mathcal{F}_{n + 1})$. Dans la preuve du lemme \ref{lemma-ML-general-better}, nous avons établi l'énoncé sur la structure de module.

\medskip\noindent

Soit $t \in F^n$. Choisissons un élément $s \in H^p(X, I^n\mathcal{F}_{n + 1})$
qui s'envoie sur l'image de $t$ dans $H^p(X, \mathcal{F}_{n + 1})$. Nous devons montrer que $s$ appartient à $N_n$. Or $F^n$ est le noyau de l'application
$M \to H^p(X, \mathcal{F}_n)$; par conséquent, pour tout $m \geq n$, nous pouvons envoyer $t$
vers un élément $t_m \in H^p(X, \mathcal{F}_{m + 1})$ qui s'envoie sur zéro dans $H^p(X, \mathcal{F}_n)$. Considérez la séquence de cohomologie
$$
H^{p - 1}(X, \mathcal{F}_n) \to
H^p(X, I^n\mathcal{F}_{m + 1}) \to
H^p(X, \mathcal{F}_{m + 1}) \to
H^p(X, \mathcal{F}_n)
$$
provenant de la suite exacte courte
$0 \to I^n\mathcal{F}_{m + 1} \to \mathcal{F}_{m + 1} \to \mathcal{F}_n \to 0$. Nous pouvons choisir $s_m \in H^p(X, I^n\mathcal{F}_{m + 1})$ s'envoyant sur $t_m$. En comparant la suite ci-dessus à celle correspondant à $m = n$, nous voyons que
$s_m$ s'envoie sur $s$, à un élément de l'image de
$H^{p - 1}(X, \mathcal{F}_n) \to H^p(X, I^n\mathcal{F}_{n + 1})$. Cependant, cette application se factorise par l'application
$H^p(X, I^n\mathcal{F}_{m + 1}) \to H^p(X, I^n\mathcal{F}_{n + 1})$
et nous voyons que $s$ appartient à l'image, comme souhaité.

\end{proof}





\section{Systèmes projectifs et cohomologie, II}

\label{section-inverse-systems-bis}

\noindent
Cette section poursuit la discussion dans la section \ref{section-inverse-systems} dans le cadre où l'idéal est principal.

\begin{lemma}

\label{lemma-equivalent-f-good}

Soit $(X, \mathcal{O}_X)$ un espace annelé et soit $f \in \Gamma(X, \mathcal{O}_X)$. Soit
$$
\ldots \to \mathcal{F}_3 \to \mathcal{F}_2 \to \mathcal{F}_1
$$
un système projectif de $\mathcal{O}_X$-modules. Considérons les conditions

\begin{enumerate}

\item pour tout $n \geq 1$, le morphisme
$f : \mathcal{F}_{n + 1} \to \mathcal{F}_{n + 1}$ se factorise par $\mathcal{F}_{n + 1} \to \mathcal{F}_n$ de manière à donner une suite exacte courte
$0 \to \mathcal{F}_n \to \mathcal{F}_{n + 1} \to \mathcal{F}_1 \to 0$,
\item pour tout $n \geq 1$, le morphisme
$f^n : \mathcal{F}_{n + 1} \to \mathcal{F}_{n + 1}$
se factorise par $\mathcal{F}_{n + 1} \to \mathcal{F}_1$
de manière à donner une suite exacte courte
$0 \to \mathcal{F}_1 \to \mathcal{F}_{n + 1} \to \mathcal{F}_n \to 0$

\item il existe un $\mathcal{O}_X$-module $\mathcal{G}$
qui est $f$-divisible et tel que $\mathcal{F}_n = \mathcal{G}[f^n]$,
\item il existe un $\mathcal{O}_X$-module $\mathcal{F}$
sans $f$-torsion et tel que
$\mathcal{F}_n = \mathcal{F}/f^n\mathcal{F}$.

\end{enumerate}

Alors (4) $\Rightarrow$ (3) $\Leftrightarrow$ (2) $\Leftrightarrow$ (1).

\end{lemma}

\begin{proof}
Nous omettons la preuve de l'équivalence de (1) et (2), ainsi que celle du fait que (3) implique (1). Étant donné $\mathcal{F}_n$ comme dans (1), posons, pour démontrer (3),
$\mathcal{G} = \colim \mathcal{F}_n$, où les morphismes
$\mathcal{F}_1 \to \mathcal{F}_2 \to \mathcal{F}_3 \to \ldots$
sont ceux de (1). Le morphisme $f : \mathcal{G} \to \mathcal{G}$
est surjectif, car l'image de $\mathcal{F}_{n + 1} \subset \mathcal{G}$
est $\mathcal{F}_n \subset \mathcal{G}$ d'après la suite exacte courte de (1).
Ainsi, $\mathcal{G}$ est $f$-divisible comme $\mathcal{O}_X$-module et
$\mathcal{F}_n = \mathcal{G}[f^n]$.

\medskip\noindent
Supposons donné $\mathcal{F}$ comme dans (4). Le morphisme
$\mathcal{F}/f^{n + 1}\mathcal{F} \to \mathcal{F}/f^n\mathcal{F}$
est toujours surjectif, et son noyau est l'image du morphisme
$\mathcal{F}/f\mathcal{F} \to \mathcal{F}/f^{n + 1}\mathcal{F}$
induit par la multiplication par $f^n$. Pour vérifier (2), il suffit de voir que le noyau de
$f^n : \mathcal{F} \to \mathcal{F}/f^{n + 1}\mathcal{F}$
est $f\mathcal{F}$. Il suffit donc de montrer que, pour des sections
$s, t$ de $\mathcal{F}$ sur un ouvert $U \subset X$
telles que $f^ns = f^{n + 1}t$, nous avons $s = ft$. Cela est clair, car
$f : \mathcal{F} \to \mathcal{F}$ est injectif puisque $\mathcal{F}$
est sans $f$-torsion.
\end{proof}

\begin{lemma}

\label{lemma-topology-I-adic-f}
Supposons que $X$, $f$ et $(\mathcal{F}_n)$ satisfassent à la condition (1) du lemme \ref{lemma-equivalent-f-good}. Soit $p \geq 0$ et posons
$H^p = \lim H^p(X, \mathcal{F}_n)$. Alors $f^cH^p$ est le noyau de $H^p \to H^p(X, \mathcal{F}_c)$ pour tout
$c \geq 1$. Ainsi, la topologie limite sur $H^p$ est la topologie $f$-adique.
\end{lemma}

\begin{proof}

Soit $c \geq 1$. Il est clair que $f^c H^p$ s'envoie sur zéro dans
$H^p(X, \mathcal{F}_c)$. Si $\xi = (\xi_n) \in H^p$ est suffisamment petit pour la topologie limite, alors $\xi_c = 0$, et donc $\xi_n$
s'envoie sur zéro dans $H^p(X, \mathcal{F}_c)$ pour $n \geq c$. Considérons le système projectif de suites exactes courtes
$$
0 \to \mathcal{F}_{n - c} \xrightarrow{f^c} \mathcal{F}_n \to
\mathcal{F}_c \to 0
$$
et le système projectif correspondant de suites exactes longues de cohomologie
$$
H^{p - 1}(X, \mathcal{F}_c) \to
H^p(X, \mathcal{F}_{n - c}) \to
H^p(X, \mathcal{F}_n) \to
H^p(X, \mathcal{F}_c)
$$
Puisque le terme $H^{p - 1}(X, \mathcal{F}_c)$ est indépendant de
$n$, nous pouvons choisir une suite compatible d'éléments
$\xi'_n \in H^p(X, \mathcal{F}_{n - c})$
relevant $\xi_n$. En posant $\xi' = (\xi'_n)$, nous obtenons $\xi = f^c \xi'$,
comme souhaité.

\end{proof}

\begin{lemma}

\label{lemma-limit-finite}

Soit $A$ un anneau noethérien complet pour un idéal principal $(f)$. Soit $X$ un espace topologique. Soit
$$
\ldots \to \mathcal{F}_3 \to \mathcal{F}_2 \to \mathcal{F}_1
$$
un système projectif de faisceaux de $A$-modules. Supposons que

\begin{enumerate}

\item $\Gamma(X, \mathcal{F}_1)$ est un $A$-module de type fini,
\item $X$, $f$ et $(\mathcal{F}_n)$ satisfont à la condition (1) du lemme \ref{lemma-equivalent-f-good}.

\end{enumerate}

Alors
$$
M = \lim \Gamma(X, \mathcal{F}_n)
$$
est un $A$-module de type fini, $f$ est un non-diviseur de zéro sur $M$, et
$M/fM$ est l'image de $M$ dans $\Gamma(X, \mathcal{F}_1)$.

\end{lemma}

\begin{proof}

D'après le lemme \ref{lemma-topology-I-adic-f}, nous avons
$M/fM \subset H^0(X, \mathcal{F}_1)$. D'après (1) et la noethérianité de $A$, nous obtenons que $M/fM$ est un $A$-module de type fini. De plus, $\bigcap f^nM = 0$, car $f^nM$ s'envoie sur zéro dans
$H^0(X, \mathcal{F}_n)$. Par Algèbre, lemme \ref{algebra-lemma-finite-over-complete-ring},
nous concluons que $M$ est de type fini sur $A$. Enfin, supposons que $s = (s_n) \in M = \lim H^0(X, \mathcal{F}_n)$
vérifie $fs = 0$. Alors $s_{n + 1}$ appartient au noyau de
$\mathcal{F}_{n + 1} \to \mathcal{F}_n$ d'après la condition (1) du lemme \ref{lemma-equivalent-f-good}. Ainsi, $s_n = 0$. Comme $n$ est arbitraire, nous obtenons $s = 0$. Par conséquent, $f$ est un non-diviseur de zéro sur $M$.

\end{proof}

\begin{lemma}

\label{lemma-ML}

Soit $A$ un anneau. Soit $f \in A$. Soit $X$ un espace topologique. Soit
$$
\ldots \to \mathcal{F}_3 \to \mathcal{F}_2 \to \mathcal{F}_1
$$
un système projectif de faisceaux de $A$-modules. Soit $p \geq 0$. Supposons que

\begin{enumerate}

\item soit $H^{p + 1}(X, \mathcal{F}_1)$ est un $A$-module de longueur finie, soit $A$ est noethérien et $H^{p + 1}(X, \mathcal{F}_1)$ est un $A$-module de type fini,
\item $X$, $f$ et $(\mathcal{F}_n)$ satisfont à la condition (1) du lemme \ref{lemma-equivalent-f-good}.

\end{enumerate}

Alors le système projectif $M_n = H^p(X, \mathcal{F}_n)$ satisfait à la condition de Mittag--Leffler.

\end{lemma}

\begin{proof}

Posons $I = (f)$. Nous utiliserons le critère du lemme \ref{lemma-ML-general}. Remarquons que
$f^n : \mathcal{F}_1 \to I^n\mathcal{F}_{n + 1}$
est un isomorphisme pour tout $n \geq 0$. Il suffit donc de montrer que
$$
\bigoplus\nolimits_{n \geq 1} H^{p + 1}(X, \mathcal{F}_1) \cdot f^{n + 1}
$$
est un module gradué sur $S = \bigoplus_{n \geq 0} A/(f) \cdot f^n$ satisfaisant à la condition de chaîne ascendante. Si $A$ n'est pas noethérien, alors
$H^{p + 1}(X, \mathcal{F}_1)$ est de longueur finie et le résultat est acquis. Si $A$ est noethérien, alors $S$ est un anneau noethérien et le résultat découle de ce que le module est de type fini sur $S$, par l'hypothèse de finitude sur $H^{p + 1}(X, \mathcal{F}_1)$. Nous omettons certains détails.

\end{proof}

\begin{lemma}

\label{lemma-ML-better}

Soit $A$ un anneau. Soit $f \in A$. Soit $X$ un espace topologique. Soit
$$
\ldots \to \mathcal{F}_3 \to \mathcal{F}_2 \to \mathcal{F}_1
$$
un système projectif de faisceaux de $A$-modules. Soit $p \geq 0$. Supposons que

\begin{enumerate}

\item soit il existe $m \geq 1$ tel que l'image de
$H^{p + 1}(X, \mathcal{F}_m) \to H^{p + 1}(X, \mathcal{F}_1)$
soit un $A$-module de longueur finie, soit $A$ est noethérien et l'intersection des images de
$H^{p + 1}(X, \mathcal{F}_m) \to H^{p + 1}(X, \mathcal{F}_1)$
est un $A$-module de type fini,
\item $X$, $f$ et $(\mathcal{F}_n)$ satisfont à la condition (1) du lemme \ref{lemma-equivalent-f-good}.

\end{enumerate}

Alors le système projectif $M_n = H^p(X, \mathcal{F}_n)$ satisfait à la condition de Mittag--Leffler.

\end{lemma}

\begin{proof}
Posons $I = (f)$. Nous appliquerons le critère du lemme \ref{lemma-ML-general-better} aux modules $N_n$. Pour $m \geq n$, nous avons
$I^n\mathcal{F}_{m + 1} = \mathcal{F}_{m + 1 - n}$. Ainsi, nous voyons que
$$
N_n = \bigcap\nolimits_{m \geq 1} \Im\left(
H^{p + 1}(X, \mathcal{F}_m) \to H^{p + 1}(X, \mathcal{F}_1)
\right)
$$
est indépendant de $n$ et que $\bigoplus N_n = \bigoplus N_1 \cdot f^{n + 1}$. Nous concluons donc exactement comme dans la preuve du lemme \ref{lemma-ML}.
\end{proof}

\begin{remark}

\label{remark-compare-ML}

Soit $(X, \mathcal{O}_X)$ un espace annelé. Soit $f \in \Gamma(X, \mathcal{O}_X)$. Soit $\mathcal{F}$ un $\mathcal{O}_X$-module. Si $\mathcal{F}$ est sans $f$-torsion, alors, pour tout $p \geq 0$, nous avons une suite exacte courte de systèmes projectifs
$$
0 \to \{H^p(X, \mathcal{F})/f^nH^p(X, \mathcal{F})\}
\to \{H^p(X, \mathcal{F}/f^n\mathcal{F})\}
\to \{H^{p + 1}(X, \mathcal{F})[f^n]\} \to 0
$$
Comme le premier système projectif satisfait à la condition de Mittag--Leffler (ML), nous en déduisons trois faits :

\begin{enumerate}

\item Il existe une suite exacte courte
$$
0 \to \widehat{H^p(X, \mathcal{F})} \to
\lim H^p(X, \mathcal{F}/f^n\mathcal{F}) \to
T_f(H^{p + 1}(X, \mathcal{F})) \to 0
$$
où $\widehat{\ }$ désigne la complétion $f$-adique usuelle et
$T_f( - )$ le module de Tate $f$-adique de Compléments sur l'algèbre, exemple
\ref{more-algebra-example-spectral-sequence-principal}.
\item Nous avons
$R^1\lim H^p(X, \mathcal{F}/f^n\mathcal{F}) =
R^1\lim H^{p + 1}(X, \mathcal{F})[f^n]$.
\item Le système $\{H^{p + 1}(X, \mathcal{F})[f^n]\}$ est ML si et seulement si
$\{H^p(X, \mathcal{F}/f^n\mathcal{F})\}$ est ML.

\end{enumerate}

Voir Homologie, lemme \ref{homology-lemma-Mittag-Leffler}, et Compléments sur l'algèbre, lemmes \ref{more-algebra-lemma-six-term-Rlim} et
\ref{more-algebra-lemma-Mittag-Leffler}.

\end{remark}









\section{Limites dérivées}

\label{section-derived-limits}

\noindent
Soit $(X, \mathcal{O}_X)$ un espace annelé. Puisque la catégorie triangulée $D(\mathcal{O}_X)$ admet des produits (Injectifs, lemme \ref{injectives-lemma-derived-products}), il s'ensuit que $D(\mathcal{O}_X)$ admet des limites dérivées; voir Catégories dérivées, définition \ref{derived-definition-derived-limit}. Si $(K_n)$ est un système projectif dans $D(\mathcal{O}_X)$, nous notons $R\lim K_n$ sa limite dérivée.

\begin{lemma}

\label{lemma-RGamma-commutes-with-Rlim}
Soit $(X, \mathcal{O}_X)$ un espace annelé. Pour tout ouvert $U \subset X$, le foncteur $R\Gamma(U, -)$ commute avec $R\lim$. De plus, il existe des suites exactes courtes $$
0 \to
R^1\lim H^{m - 1}(U, K_n) \to H^m(U, R\lim K_n) \to
\lim H^m(U, K_n) \to 0
$$ pour tout système projectif $(K_n)$ dans $D(\mathcal{O}_X)$ et tout $m \in \mathbf{Z}$.
\end{lemma}

\begin{proof}

Le premier énoncé découle d'Injectifs, lemme \ref{injectives-lemma-RF-commutes-with-Rlim}. Nous pouvons ensuite appliquer Compléments sur l'algèbre, remarque \ref{more-algebra-remark-compare-derived-limit},
à $R\lim R\Gamma(U, K_n) = R\Gamma(U, R\lim K_n)$ pour obtenir les suites exactes courtes.

\end{proof}

\begin{lemma}

\label{lemma-Rf-commutes-with-Rlim}

Soit $f : (X, \mathcal{O}_X) \to (Y, \mathcal{O}_Y)$ un morphisme d'espaces annelés. Alors $Rf_*$ commute avec $R\lim$, c'est-à-dire que $Rf_*$ commute avec les limites dérivées.

\end{lemma}

\begin{proof}

Soit $(K_n)$ un système projectif dans $D(\mathcal{O}_X)$. Considérons le triangle distingué qui définit la limite dérivée
$$
R\lim K_n \to \prod K_n \to \prod K_n
$$
dans $D(\mathcal{O}_X)$. En appliquant le foncteur exact $Rf_*$, nous obtenons le triangle distingué
$$
Rf_*(R\lim K_n) \to Rf_*\left(\prod K_n\right) \to Rf_*\left(\prod K_n\right)
$$
dans $D(\mathcal{O}_Y)$. Il suffit donc de montrer que
$Rf_*$ commute aux produits de la catégorie dérivée (lesquels ne sont pas simplement donnés par les produits de complexes; voir Injectifs, lemme \ref{injectives-lemma-derived-products}). Or, puisque $Rf_*$ est adjoint à droite d'après le lemme \ref{lemma-adjoint},
cela résulte formellement de Catégories, lemme \ref{categories-lemma-adjoint-exact}. Attention : on ne peut pas appliquer directement Catégories, lemme \ref{categories-lemma-adjoint-exact},
car $R\lim K_n$ n'est pas une limite dans $D(\mathcal{O}_X)$.

\end{proof}

\begin{remark}

\label{remark-discuss-derived-limit}

Soit $(X, \mathcal{O}_X)$ un espace annelé. Soit $(K_n)$ un système projectif dans $D(\mathcal{O}_X)$. Posons $K = R\lim K_n$. Pour tout $n$ et tout $m$,
soit $\mathcal{H}^m_n = H^m(K_n)$ le $m$-ième faisceau de cohomologie de
$K_n$, et posons de même $\mathcal{H}^m = H^m(K)$. Notons
$\underline{\mathcal{H}}^m_n$ le préfaisceau
$$
U \longmapsto \underline{\mathcal{H}}^m_n(U) = H^m(U, K_n)
$$
De même, posons $\underline{\mathcal{H}}^m(U) = H^m(U, K)$. D'après le lemme \ref{lemma-sheafification-cohomology}, nous voyons que
$\mathcal{H}^m_n$ est la faisceautisation de
$\underline{\mathcal{H}}^m_n$ et que $\mathcal{H}^m$ est la faisceautisation de $\underline{\mathcal{H}}^m$. Voici un diagramme
$$
\xymatrix{
K \ar@{=}[d] &
\underline{\mathcal{H}}^m \ar[d] \ar[r] & 
\mathcal{H}^m \ar[d] \\
R\lim K_n &
\lim \underline{\mathcal{H}}^m_n \ar[r] & 
\lim \mathcal{H}^m_n
}
$$
En général, il n'est pas vrai que
$\lim \mathcal{H}^m_n$ soit le faisceau associé à
$\lim \underline{\mathcal{H}}^m_n$. Si $U \subset X$ est un ouvert, nous avons alors des suites exactes courtes

\begin{equation}
\label{equation-ses-Rlim-over-U}
0 \to
R^1\lim \underline{\mathcal{H}}^{m - 1}_n(U) \to
\underline{\mathcal{H}}^m(U) \to
\lim \underline{\mathcal{H}}^m_n(U) \to 0
\end{equation}

d'après le lemme \ref{lemma-RGamma-commutes-with-Rlim}.

\end{remark}

\noindent
Le lemme suivant s'applique à un système projectif de modules quasi-cohérents dont les morphismes de transition sont surjectifs sur un schéma.

\begin{lemma}

\label{lemma-inverse-limit-is-derived-limit}

Soit $(X, \mathcal{O}_X)$ un espace annelé. Soit $(\mathcal{F}_n)$ un système projectif de $\mathcal{O}_X$-modules. Soit $\mathcal{B}$ un ensemble d'ouverts de $X$. Supposons que

\begin{enumerate}

\item tout ouvert de $X$ admet un recouvrement dont les membres appartiennent à
$\mathcal{B}$,
\item  $H^p(U, \mathcal{F}_n) = 0$ pour $p > 0$ et $U \in \mathcal{B}$,
\item le système projectif $\mathcal{F}_n(U)$ a un $R^1\lim$ nul
pour $U \in \mathcal{B}$.

\end{enumerate}

Alors $R\lim \mathcal{F}_n = \lim \mathcal{F}_n$ et nous avons
$H^p(U, \lim \mathcal{F}_n) = 0$ pour $p > 0$ et $U \in \mathcal{B}$.

\end{lemma}

\begin{proof}

Posons $K_n = \mathcal{F}_n$ et $K = R\lim \mathcal{F}_n$. Avec les notations de la remarque \ref{remark-discuss-derived-limit} et l'hypothèse (2), nous voyons que, pour
$U \in \mathcal{B}$, nous avons $\underline{\mathcal{H}}_n^m(U) = 0$
lorsque $m \not = 0$ et $\underline{\mathcal{H}}_n^0(U) = \mathcal{F}_n(U)$. De l'équation (\ref{equation-ses-Rlim-over-U}) et l'hypothèse (3) nous voyons que $\underline{\mathcal{H}}^m(U) = 0$
lorsque $m \not = 0$ et égal à $\lim \mathcal{F}_n(U)$
lorsque $m = 0$. En faisceautisant et en utilisant (1), nous obtenons
$\mathcal{H}^m = 0$ lorsque $m \not = 0$ et égal à
$\lim \mathcal{F}_n$ lorsque $m = 0$. Ainsi, $K = \lim \mathcal{F}_n$. Comme $H^m(U, K) = \underline{\mathcal{H}}^m(U) = 0$ pour $m > 0$
(voir ci-dessus), la seconde assertion est vérifiée.

\end{proof}

\begin{lemma}

\label{lemma-cohomology-derived-limit-injective}

Soit $(X, \mathcal{O}_X)$ un espace annelé. Soit $(K_n)$ un système projectif dans $D(\mathcal{O}_X)$. Soient $x \in X$ et $m \in \mathbf{Z}$. Supposons qu'il existe un entier $n(x)$ et un système fondamental $\mathfrak{U}_x$
de voisinages ouverts de $x$ tels que, pour $U \in \mathfrak{U}_x$,

\begin{enumerate}

\item  $R^1\lim H^{m - 1}(U, K_n) = 0$,
\item  $H^m(U, K_n) \to H^m(U, K_{n(x)})$ est injectif pour $n \geq n(x)$.

\end{enumerate}

Alors l'application sur les germes $H^m(R\lim K_n)_x \to H^m(K_{n(x)})_x$ est injective.

\end{lemma}

\begin{proof}

Soit $\gamma$ un élément de $H^m(R\lim K_n)_x$ qui s'envoie sur zéro dans $H^m(K_{n(x)})_x$. Comme $H^m(R\lim K_n)$ est la faisceautisation de $U \mapsto H^m(U, R\lim K_n)$,
(d'après le lemme \ref{lemma-sheafification-cohomology}) nous pouvons choisir $U \in \mathfrak{U}_x$
et un élément $\tilde \gamma \in H^m(U, R\lim K_n)$ s'envoyant sur $\gamma$. Alors $\tilde\gamma$ s'envoie sur $\tilde\gamma_{n(x)} \in H^m(U, K_{n(x)})$. Puisque $H^m(K_{n(x)})$ est la faisceautisation de
$U \mapsto H^m(U, K_{n(x)})$
(à nouveau d'après le lemme \ref{lemma-sheafification-cohomology}), nous voyons qu'après avoir rétréci $U$, nous pouvons supposer que $\tilde\gamma_{n(x)} = 0$. Pour cet ouvert $U$, considérons la suite exacte courte
$$
0 \to
R^1\lim H^{m - 1}(U, K_n) \to H^m(U, R\lim K_n) \to
\lim H^m(U, K_n) \to 0
$$
du lemme \ref{lemma-RGamma-commutes-with-Rlim}. Par l'hypothèse (1), le groupe de gauche est nul et, par l'hypothèse (2), celui de droite s'injecte dans $H^m(U, K_{n(x)})$. Nous concluons que $\tilde\gamma = 0$
et donc que $\gamma = 0$, comme souhaité.

\end{proof}

\begin{lemma}

\label{lemma-is-limit-per-point}

Soit $(X, \mathcal{O}_X)$ un espace annelé et soit $E \in D(\mathcal{O}_X)$. Supposons que, pour tout $x \in X$, il existe une fonction $p(x, -) : \mathbf{Z} \to \mathbf{Z}$ et un système fondamental $\mathfrak{U}_x$ de voisinages ouverts de $x$
tels que
$$
H^p(U, H^{m - p}(E)) = 0 \text{ pour }
U \in \mathfrak{U}_x \text{ et } p > p(x, m)
$$
Alors l'application $E \to R\lim \tau_{\geq -n} E$ de Catégories dérivées, remarque
\ref{derived-remark-map-into-derived-limit-truncations}
est un isomorphisme dans $D(\mathcal{O}_X)$.

\end{lemma}

\begin{proof}

Posons $K_n = \tau_{\geq -n}E$ et $K = R\lim K_n$. Le morphisme canonique $E \to K$
provient des morphismes canoniques $E \to K_n = \tau_{\geq -n}E$. Nous devons montrer que $E \to K$ induit un isomorphisme
$H^m(E) \to H^m(K)$ de faisceaux de cohomologie. Dans la suite de la preuve, fixons $m$. Si $n \geq -m$, alors le morphisme $E \to \tau_{\geq -n}E = K_n$ induit un isomorphisme
$H^m(E) \to H^m(K_n)$. Pour finir la preuve il suffit de montrer que pour chaque $x \in X$
il existe un entier $n(x) \geq -m$ de sorte que l'application
$H^m(K)_x \to H^m(K_{n(x)})_x$ est injective. En effet, la composée
$$
H^m(E)_x \to H^m(K)_x \to H^m(K_{n(x)})_x
$$
est une bijection et la deuxième flèche est injective; la première est donc bijective. Posons
$$
n(x) = 1 + \max\{-m, p(x, m - 1) - m, -1 + p(x, m) - m, -2 + p(x, m + 1) - m\}.
$$
de sorte que, dans tous les cas, $n(x) \geq -m$. Assertion : les morphismes
$$
H^{m - 1}(U, K_{n + 1}) \to H^{m - 1}(U, K_n)
\quad\text{et}\quad
H^m(U, K_{n + 1}) \to H^m(U, K_n)
$$
sont des isomorphismes pour $n \geq n(x)$ et $U \in \mathfrak{U}_x$. L'assertion implique que les conditions (1) et (2) du lemme \ref{lemma-cohomology-derived-limit-injective}
sont satisfaites et fournit donc l'injectivité souhaitée. Rappelons (Catégories dérivées, remarque
\ref{derived-remark-truncation-distinguished-triangle}) que nous avons des triangles distingués
$$
H^{-n - 1}(E)[n + 1] \to
K_{n + 1} \to K_n \to H^{-n - 1}(E)[n + 2]
$$
En examinant la suite exacte longue de cohomologie associée, l'assertion découle de la nullité des groupes
$$
H^{m + n}(U, H^{-n - 1}(E)),\quad
H^{m + n + 1}(U, H^{-n - 1}(E)),\quad
H^{m + n + 2}(U, H^{-n - 1}(E))
$$
pour $n \geq n(x)$ et $U \in \mathfrak{U}_x$. Cette nullité découle de notre choix de $n(x)$
et de l'hypothèse du lemme.

\end{proof}

\begin{lemma}

\label{lemma-is-limit-spaltenstein}

\begin{reference}
\cite[proposition 3.13]{Spaltenstein}
\end{reference}
Soit $(X, \mathcal{O}_X)$ un espace annelé. Soit $E \in D(\mathcal{O}_X)$. Supposons que, pour tout $x \in X$, il existe un entier $d_x \geq 0$ et un système fondamental $\mathfrak{U}_x$ de voisinages ouverts de $x$ tels que $$
H^p(U, H^q(E)) = 0 \text{ pour }
U \in \mathfrak{U}_x,\ p > d_x, \text{ et }q < 0
$$ Alors l'application $E \to R\lim \tau_{\geq -n} E$ de Catégories dérivées, remarque \ref{derived-remark-map-into-derived-limit-truncations}, est un isomorphisme dans $D(\mathcal{O}_X)$.
\end{lemma}

\begin{proof}
Ceci résulte du lemme \ref{lemma-is-limit-per-point} avec $p(x, m) = d_x + \max(0, m)$.
\end{proof}

\begin{lemma}

\label{lemma-is-limit}

Soit $(X, \mathcal{O}_X)$ un espace annelé. Soit $E \in D(\mathcal{O}_X)$. Supposons qu'il existe une fonction $p(-) : \mathbf{Z} \to \mathbf{Z}$
et un ensemble $\mathcal{B}$ d'ouverts de $X$ tels que

\begin{enumerate}

\item tout ouvert de $X$ admet un recouvrement dont les membres appartiennent à $\mathcal{B}$,
\item  $H^p(U, H^{m - p}(E)) = 0$ pour $p > p(m)$ et $U \in \mathcal{B}$.

\end{enumerate}

Alors l'application $E \to R\lim \tau_{\geq -n} E$ de Catégories dérivées, remarque
\ref{derived-remark-map-into-derived-limit-truncations}
est un isomorphisme dans $D(\mathcal{O}_X)$.

\end{lemma}

\begin{proof}
Appliquer le lemme \ref{lemma-is-limit-per-point} avec $p(x, m) = p(m)$ et $\mathfrak{U}_x = \{U \in \mathcal{B} \mid x \in U\}$.
\end{proof}

\begin{lemma}

\label{lemma-is-limit-dimension}
Soit $(X, \mathcal{O}_X)$ un espace annelé. Soit $E \in D(\mathcal{O}_X)$. Supposons qu'il existe un entier $d \geq 0$ et une base $\mathcal{B}$ de la topologie de $X$ telle que $$
H^p(U, H^q(E)) = 0 \text{ pour }
U \in \mathcal{B},\ p > d, \text{ et }q < 0
$$ Alors l'application $E \to R\lim \tau_{\geq -n} E$ de Catégories dérivées, remarque \ref{derived-remark-map-into-derived-limit-truncations}, est un isomorphisme dans $D(\mathcal{O}_X)$.
\end{lemma}

\begin{proof}
Appliquer le lemme \ref{lemma-is-limit-spaltenstein} avec $d_x = d$ et $\mathfrak{U}_x = \{U \in \mathcal{B} \mid x \in U\}$.
\end{proof}

\noindent
Les lemmes ci-dessus peuvent être utilisés pour calculer la cohomologie dans certaines situations.

\begin{lemma}

\label{lemma-cohomology-over-U-trivial}

Soit $(X, \mathcal{O}_X)$ un espace annelé. Soit $K$
un objet de $D(\mathcal{O}_X)$. Soit $\mathcal{B}$ un ensemble d'ouverts de $X$. Supposons que

\begin{enumerate}

\item tout ouvert de $X$ admet un recouvrement dont les membres appartiennent à $\mathcal{B}$,
\item  $H^p(U, H^q(K)) = 0$ pour tous $p > 0$, $q \in \mathbf{Z}$,
$U \in \mathcal{B}$.

\end{enumerate}

Alors $H^q(U, K) = H^0(U, H^q(K))$ pour $q \in \mathbf{Z}$
et $U \in \mathcal{B}$.

\end{lemma}

\begin{proof}

Remarquons que $K = R\lim \tau_{\geq -n} K$ d'après le lemme \ref{lemma-is-limit-dimension} avec $d = 0$. Soit $U \in \mathcal{B}$. Par l'équation (\ref{equation-ses-Rlim-over-U}), nous obtenons une suite exacte courte
$$
0 \to R^1\lim H^{q - 1}(U, \tau_{\geq -n}K) \to
H^q(U, K) \to \lim H^q(U, \tau_{\geq -n}K) \to 0
$$
La condition (2) implique
$H^q(U, \tau_{\geq -n} K) = H^0(U, H^q(\tau_{\geq -n} K))$
pour tout $q$, à l'aide de la suite spectrale de l'exemple \ref{example-spectral-sequence}. Cette suite spectrale converge parce que $\tau_{\geq -n}K$ est borné inférieurement. Si $n > -q$, alors $H^q(\tau_{\geq -n}K) = H^q(K)$. Ainsi, les systèmes de gauche et de droite de la suite exacte courte affichée sont stationnaires, de valeurs respectives
$H^0(U, H^{q - 1}(K))$ et $H^0(U, H^q(K))$. Le lemme suit.

\end{proof}

\noindent
Voici un autre cas où nous pouvons décrire la limite dérivée.

\begin{lemma}

\label{lemma-derived-limit-suitable-system}

Soit $(X, \mathcal{O}_X)$ un espace annelé. Soit $(K_n)$
un système projectif d'objets de $D(\mathcal{O}_X)$. Soit $\mathcal{B}$ un ensemble d'ouverts de $X$. Supposons que

\begin{enumerate}

\item tout ouvert de $X$ admet un recouvrement dont les membres appartiennent à $\mathcal{B}$,
\item pour tout $U \in \mathcal{B}$ et tout $q \in \mathbf{Z}$, nous avons

\begin{enumerate}

\item  $H^p(U, H^q(K_n)) = 0$ pour $p > 0$,
\item le système projectif $H^0(U, H^q(K_n))$ a un $R^1\lim$ nul.

\end{enumerate}

\end{enumerate}

Alors $H^q(R\lim K_n) = \lim H^q(K_n)$ pour $q \in \mathbf{Z}$.

\end{lemma}

\begin{proof}

Posons $K = R\lim K_n$. Soit $U \in \mathcal{B}$. D'après le lemme \ref{lemma-cohomology-over-U-trivial} et (2) a), nous avons $H^q(U, K_n) = H^0(U, H^q(K_n))$. D'après le lemme \ref{lemma-RGamma-commutes-with-Rlim} et (2) b), nous avons 
$H^q(U, K) = \lim H^0(U, H^q(K_n))$. Ainsi, $H^q(U, K)$ est la limite projective des sections des faisceaux $H^q(K_n)$ sur $U$. Comme
$\lim H^q(K_n)$ est un faisceau, l'hypothèse (1) montre que $H^q(K)$, qui est la faisceautisation du préfaisceau $U \mapsto H^q(U, K)$, est égal à $\lim H^q(K_n)$. Cela prouve le lemme.

\end{proof}







\section{Construction de résolutions K-injectives}

\label{section-K-injective}

\noindent

Soit $(X, \mathcal{O}_X)$ un espace annelé. Soit $\mathcal{F}^\bullet$ un complexe de $\mathcal{O}_X$-modules. La catégorie $\textit{Mod}(\mathcal{O}_X)$ a assez d'injectifs; nous pouvons donc utiliser Catégories dérivées, lemme \ref{derived-lemma-special-inverse-system},
pour produire un diagramme
$$
\xymatrix{
\ldots \ar[r] &
\tau_{\geq -2}\mathcal{F}^\bullet \ar[r] \ar[d] &
\tau_{\geq -1}\mathcal{F}^\bullet \ar[d] \\
\ldots \ar[r] & \mathcal{I}_2^\bullet \ar[r] & \mathcal{I}_1^\bullet
}
$$
dans la catégorie des complexes de $\mathcal{O}_X$-modules, tel que

\begin{enumerate}

\item les flèches verticales sont des quasi-isomorphismes,
\item $\mathcal{I}_n^\bullet$ est un complexe borné inférieurement d'objets injectifs,
\item les flèches $\mathcal{I}_{n + 1}^\bullet \to \mathcal{I}_n^\bullet$
sont des surjections scindées terme à terme.

\end{enumerate}

La catégorie des $\mathcal{O}_X$-modules admet des limites (elles se calculent au niveau des préfaisceaux); nous pouvons donc former la limite
$\mathcal{I}^\bullet = \lim_n \mathcal{I}_n^\bullet$. Par Catégories dérivées, lemmes
\ref{derived-lemma-bounded-below-injectives-K-injective} et
\ref{derived-lemma-limit-K-injectives}
il s'agit d'un complexe K-injectif. En général, l'application canonique

\begin{equation}
\label{equation-into-candidate-K-injective}
\mathcal{F}^\bullet \to \mathcal{I}^\bullet
\end{equation}

peut ne pas être un quasi-isomorphisme. Dans le lemme suivant, nous donnons des conditions sous lesquelles il en est un.

\begin{lemma}

\label{lemma-K-injective}

Dans la situation décrite ci-dessus, soit $\mathcal{H}^m = H^m(\mathcal{F}^\bullet)$ le faisceau de cohomologie de degré $m$. Soit $\mathcal{B}$ un ensemble de sous-ensembles ouverts de $X$. Soit $d \in \mathbf{N}$. Supposons que

\begin{enumerate}

\item tout ouvert de $X$ admet un recouvrement dont les membres appartiennent à $\mathcal{B}$,
\item pour tout $U \in \mathcal{B}$, nous avons $H^p(U, \mathcal{H}^q) = 0$
pour $p > d$ et $q < 0$\footnote{Il suffit que
$\forall m$, $\exists p(m)$, $H^p(U. \mathcal{H}^{m - p}) = 0$ pour
$p > p(m)$, voir lemme \ref{lemma-is-limit}.}.

\end{enumerate}

Alors (\ref{equation-into-candidate-K-injective}) est un quasi-isomorphisme.

\end{lemma}

\begin{proof}
D'après Catégories dérivées, lemme \ref{derived-lemma-difficulty-K-injectives}, il suffit de montrer que l'application $\mathcal{F}^\bullet \to R\lim \tau_{\geq -n} \mathcal{F}^\bullet$ est un isomorphisme. C'est le lemme \ref{lemma-is-limit-dimension}.
\end{proof}

\noindent
Voici un lemme technique sur les faisceaux de cohomologie de la limite projective d'un système de complexes de faisceaux. En un certain sens, ce n'est pas le bon énoncé à chercher à démontrer, car il faudrait prendre des limites dérivées plutôt que des limites projectives ordinaires.

\begin{lemma}

\label{lemma-inverse-limit-complexes}

Soit $(X, \mathcal{O}_X)$ un espace annelé. Soit $(\mathcal{F}_n^\bullet)$
un système projectif de complexes de $\mathcal{O}_X$-modules. Soit $m \in \mathbf{Z}$. Supposons qu'il existe un ensemble $\mathcal{B}$ de sous-ensembles ouverts de $X$ et un entier $n_0$ tels que

\begin{enumerate}

\item tout ouvert de $X$ admet un recouvrement dont les membres appartiennent à $\mathcal{B}$,
\item pour chaque $U \in \mathcal{B}$

\begin{enumerate}

\item les systèmes des groupes abéliens
$\mathcal{F}_n^{m - 2}(U)$ et $\mathcal{F}_n^{m - 1}(U)$
ont un $R^1\lim$ nul (par exemple, ils satisfont à la condition de Mittag--Leffler),
\item le système des groupes abéliens $H^{m - 1}(\mathcal{F}_n^\bullet(U))$
a un $R^1\lim$ nul (par exemple, il satisfait à la condition de Mittag--Leffler), et
\item Nous avons
$H^m(\mathcal{F}_n^\bullet(U)) = H^m(\mathcal{F}_{n_0}^\bullet(U))$
pour tous $n \geq n_0$.

\end{enumerate}

\end{enumerate}

Alors les morphismes
$H^m(\mathcal{F}^\bullet) \to \lim H^m(\mathcal{F}_n^\bullet) \to
H^m(\mathcal{F}_{n_0}^\bullet)$
sont des isomorphismes de faisceaux, où
$\mathcal{F}^\bullet = \lim \mathcal{F}_n^\bullet$ est la limite projective par terme.

\end{lemma}

\begin{proof}

Soit $U \in \mathcal{B}$. Notons que $H^m(\mathcal{F}^\bullet(U))$ est la cohomologie de
$$
\lim_n \mathcal{F}_n^{m - 2}(U) \to
\lim_n \mathcal{F}_n^{m - 1}(U) \to
\lim_n \mathcal{F}_n^m(U) \to
\lim_n \mathcal{F}_n^{m + 1}(U)
$$
Sous les hypothèses (2)a) et (2)b), nous pouvons appliquer Compléments sur l'algèbre, lemme \ref{more-algebra-lemma-apply-Mittag-Leffler-again},
pour conclure que
$$
H^m(\mathcal{F}^\bullet(U)) = \lim H^m(\mathcal{F}_n^\bullet(U))
$$
Par hypothèse (2) c), nous concluons
$$
H^m(\mathcal{F}^\bullet(U)) = H^m(\mathcal{F}_n^\bullet(U))
$$
pour tout $n \geq n_0$. Par l'hypothèse (1), nous concluons que la faisceautisation de
$U \mapsto H^m(\mathcal{F}^\bullet(U))$ est égale à celle de $U \mapsto H^m(\mathcal{F}_n^\bullet(U))$ pour tout $n \geq n_0$. Ainsi, le système projectif des faisceaux $H^m(\mathcal{F}_n^\bullet)$ est constant pour $n \geq n_0$, de valeur $H^m(\mathcal{F}^\bullet)$. Cela prouve le lemme.

\end{proof}











\section{Systèmes projectifs et cohomologie, III}

\label{section-inverse-systems-tri}

\noindent
Cette section poursuit la discussion dans la section \ref{section-inverse-systems-bis} en utilisant des limites dérivées.

\begin{lemma}

\label{lemma-formal-functions-principal}

Soit $(X, \mathcal{O}_X)$ un espace annelé. Soit $A \to \Gamma(X, \mathcal{O}_X)$
un homomorphisme d'anneaux et soit $f \in A$. Soit $E$ un objet de $D(\mathcal{O}_X)$. Posons
$$
E_n = E \otimes_{\mathcal{O}_X} (\mathcal{O}_X \xrightarrow{f^n} \mathcal{O}_X)
$$
et définissons $E^\wedge = R\lim E_n$. Pour $p \in \mathbf{Z}$,
il existe un diagramme commutatif canonique
$$
\xymatrix{
& 0 & 0 \\
0 \ar[r] &
\widehat{H^p(X, E)} \ar[r] \ar[u] &
\lim H^p(X, E_n) \ar[r] \ar[u] &
T_f(H^{p + 1}(X, E)) \ar[r] &
0 \\
0 \ar[r] &
H^0(H^p(X, E)^\wedge) \ar[r] \ar[u] &
H^p(X, E^\wedge) \ar[r] \ar[u] &
T_f(H^{p + 1}(X, E)) \ar[r] \ar@{=}[u] &
0 \\
&
R^1\lim H^p(X, E)[f^n] \ar[u] \ar[r]^\cong &
R^1\lim H^{p - 1}(X, E_n) \ar[u] \\
& 0 \ar[u] & 0 \ar[u]
}
$$
dont les lignes et les colonnes sont exactes, où $\widehat{H^p(X, E)} = \lim H^p(X, E)/f^n H^p(X, E)$ est la complétion
$f$-adique, $H^p(X, E)^\wedge$ est la complétion $f$-adique dérivée, et $T_f(H^{p + 1}(X, E))$ est le module de Tate $f$-adique; voir Compléments sur l'algèbre, exemple
\ref{more-algebra-example-spectral-sequence-principal}. Enfin, nous avons $H^p(X, E^\wedge) = H^p(R\Gamma(X, E)^\wedge)$.

\end{lemma}

\begin{proof}
Remarquons que $R\Gamma(X, E^\wedge) = R\lim R\Gamma(X, E_n)$ d'après le lemme \ref{lemma-Rf-commutes-with-Rlim}. D'autre part, nous avons $$
R\Gamma(X, E_n) =
R\Gamma(X, E) \otimes_A^\mathbf{L} (A \xrightarrow{f^n} A)
$$ (détails omis). Nous en déduisons que $R\Gamma(X, E^\wedge)$ est la complétion $f$-adique dérivée $R\Gamma(X, E)^\wedge$. Le diagramme découle alors de Compléments sur l'algèbre, lemme \ref{more-algebra-lemma-compare-completion-derived-completion}.
\end{proof}

\begin{lemma}

\label{lemma-stabilizes}
Soit $\mathcal{A}$ une catégorie abélienne. Soit $f : M \to M$ un morphisme de $\mathcal{A}$. Si $M[f^n] = \Ker(f^n : M \to M)$ se stabilise, alors les systèmes projectifs $$
(M \xrightarrow{f^n} M)
\quad\text{et}\quad
\Coker(f^n : M \to M)
$$ sont pro-isomorphes dans $D(\mathcal{A})$.
\end{lemma}

\begin{proof}
Il existe manifestement un morphisme du premier système projectif vers le second. Supposons que $M[f^c] = M[f^{c + 1}] = M[f^{c + 2}] = \ldots$. Nous pouvons alors définir un morphisme de systèmes projectifs dans $D(\mathcal{A})$ en sens inverse au moyen des diagrammes $$
\xymatrix{
M/M[f^c] \ar[r]_-{f^{n + c}} \ar[d]_{f^c} &
M \ar[d]^1 \\
M \ar[r]^{f^n} & M
}
$$ Puisque la flèche horizontale supérieure est injective, le complexe de la ligne supérieure est quasi-isomorphe à $\Coker(f^{n + c} : M \to M)$. Nous omettons quelques détails.
\end{proof}

\begin{example}

\label{example-formal-functions-principal}

Soit $(X, \mathcal{O}_X)$ un espace annelé. Soit $A \to \Gamma(X, \mathcal{O}_X)$
un homomorphisme d'anneaux et soit $f \in A$. Soit $\mathcal{F}$ un $\mathcal{O}_X$-module. Supposons qu'il existe un entier $c$ tel que $\mathcal{F}[f^c] = \mathcal{F}[f^n]$
pour tout $n \geq c$. Appliquons le lemme \ref{lemma-formal-functions-principal} avec
$E = \mathcal{F}$. D'après le lemme \ref{lemma-stabilizes}, le système projectif $(E_n)$ est pro-isomorphe au système projectif
$(\mathcal{F}/f^n\mathcal{F})$. Nous en concluons que, pour $p \in \mathbf{Z}$,
nous obtenons le diagramme commutatif
$$
\xymatrix{
& 0 & 0 \\
0 \ar[r] &
\widehat{H^p(X, \mathcal{F})} \ar[r] \ar[u] &
\lim H^p(X, \mathcal{F}/f^n\mathcal{F}) \ar[r] \ar[u] &
T_f(H^{p + 1}(X, \mathcal{F})) \ar[r] &
0 \\
0 \ar[r] &
H^0(H^p(X, \mathcal{F})^\wedge) \ar[r] \ar[u] &
H^p(R\Gamma(X, \mathcal{F})^\wedge) \ar[r] \ar[u] &
T_f(H^{p + 1}(X, \mathcal{F})) \ar[r] \ar@{=}[u] &
0 \\
&
R^1\lim H^p(X, \mathcal{F})[f^n] \ar[u] \ar[r]^\cong &
R^1\lim H^{p - 1}(X, \mathcal{F}/f^n\mathcal{F}) \ar[u] \\
& 0 \ar[u] & 0 \ar[u]
}
$$
dont les lignes et les colonnes sont exactes, où
$\widehat{H^p(X, \mathcal{F})} =
\lim H^p(X, \mathcal{F})/f^n H^p(X, \mathcal{F})$
est la complétion $f$-adique usuelle et où $M^\wedge$ désigne la complétion $f$-adique dérivée de $M$ dans $D(A)$.

\end{example}















\section{Cohomologie de {\v C}ech des complexes non bornés}

\label{section-cech-cohomology-of-unbounded-complexes}

\noindent
La construction de la section \ref{section-cech-cohomology-of-complexes} n'est pas la bonne pour les complexes non bornés. Le problème est que, dans le projet Champs, nous utilisons des sommes directes dans la totalisation d'un complexe double et que nous devrions les remplacer par un produit. Au lieu d'effectuer ce remplacement dans cette section, nous supposons que le recouvrement est fini et nous utilisons le complexe de {\v C}ech alterné.

\medskip\noindent

Soit $(X, \mathcal{O}_X)$ un espace annelé. Soit ${\mathcal F}^\bullet$ un complexe de préfaisceaux de
$\mathcal{O}_X$-modules. Soit ${\mathcal U} : X = \bigcup_{i \in I} U_i$
un {\bf recouvrement ouvert fini} de $X$. Comme le complexe de {\v C}ech alterné
$\check{\mathcal{C}}_{alt}^\bullet(\mathcal{U}, \mathcal{F})$
(section \ref{section-alternating-cech}) est fonctoriel en $\mathcal{F}$, nous obtenons un complexe double
$\check{\mathcal{C}}^\bullet_{alt}(\mathcal{U}, \mathcal{F}^\bullet)$. Dans cette section, nous travaillons avec le complexe total associé.
La construction de $\text{Tot}(\check{\mathcal{C}}^\bullet_{alt}({\mathcal U},
{\mathcal F}^\bullet))$
est fonctorielle en ${\mathcal F}^\bullet$. Il existe également une transformation fonctorielle

\begin{equation}
\label{equation-global-sections-to-alternating-cech}
\Gamma(X, {\mathcal F}^\bullet)
\longrightarrow
\text{Tot}(\check{\mathcal{C}}^\bullet_{alt}({\mathcal U},
{\mathcal F}^\bullet))
\end{equation}

de complexes définie par la règle suivante : une section
$s\in \Gamma(X, {\mathcal F}^n)$
est envoyée vers l'élément $\alpha = \{\alpha_{i_0\ldots i_p}\}$
avec $\alpha_{i_0} = s|_{U_{i_0}}$ et $\alpha_{i_0\ldots i_p} = 0$
pour $p > 0$.

\begin{lemma}

\label{lemma-alternating-cech-complex-complex}

Soit $(X, \mathcal{O}_X)$ un espace annelé. Soit $\mathcal{U} : X = \bigcup_{i \in I} U_i$ un recouvrement fini ouvert. Pour un complexe $\mathcal{F}^\bullet$
de $\mathcal{O}_X$-modules, il existe un morphisme canonique
$$
\text{Tot}(\check{\mathcal{C}}^\bullet_{alt}(\mathcal{U}, \mathcal{F}^\bullet))
\longrightarrow
R\Gamma(X, \mathcal{F}^\bullet)
$$
fonctoriel en $\mathcal{F}^\bullet$ et compatible avec (\ref{equation-global-sections-to-alternating-cech}).

\end{lemma}

\begin{proof}

Soit ${\mathcal I}^\bullet$ un complexe K-injectif dont les termes sont des $\mathcal{O}_X$-modules injectifs. Le morphisme (\ref{equation-global-sections-to-alternating-cech}) pour
$\mathcal{I}^\bullet$ est un morphisme
$\Gamma(X, {\mathcal I}^\bullet) \to
\text{Tot}(\check{\mathcal{C}}^\bullet_{alt}({\mathcal U},
{\mathcal I}^\bullet))$. Il s'agit d'un quasi-isomorphisme de complexes de groupes abéliens comme il résulte de Homologie, lemme \ref{homology-lemma-double-complex-gives-resolution}
appliqué au complexe double
$\check{\mathcal{C}}^\bullet_{alt}({\mathcal U},
{\mathcal I}^\bullet)$, au moyen des lemmes \ref{lemma-injective-trivial-cech} et \ref{lemma-alternating-usual}. Supposons que ${\mathcal F}^\bullet \to {\mathcal I}^\bullet$ soit un quasi-isomorphisme de ${\mathcal F}^\bullet$ dans un complexe K-injectif dont les termes sont injectifs (Injectifs, théorème
\ref{injectives-theorem-K-injective-embedding-grothendieck}). Puisque $R\Gamma(X, {\mathcal F}^\bullet)$ est représenté par le complexe
$\Gamma(X, {\mathcal I}^\bullet)$, nous obtenons le morphisme du lemme au moyen de
$$
\text{Tot}(\check{\mathcal{C}}^\bullet_{alt}({\mathcal U},
{\mathcal F}^\bullet))
\longrightarrow
\text{Tot}(\check{\mathcal{C}}^\bullet_{alt}({\mathcal U},
{\mathcal I}^\bullet)).
$$
Nous omettons la vérification de la fonctorialité et des compatibilités.

\end{proof}

\begin{lemma}

\label{lemma-alternating-cech-complex-complex-ss}

Soit $(X, \mathcal{O}_X)$ un espace annelé.
Soit $\mathcal{U} : X = \bigcup_{i \in I} U_i$ un recouvrement fini ouvert.
Soit $\mathcal{F}^\bullet$ un complexe de $\mathcal{O}_X$-modules. Soit $\mathcal{B}$ un ensemble d'ouverts de $X$. Supposons que

\begin{enumerate}

\item tout ouvert de $X$ possède un recouvrement dont les membres appartiennent à $\mathcal{B}$,
\item on a $U_{i_0\ldots i_p} \in \mathcal{B}$ pour tous
$i_0, \ldots, i_p \in I$,
\item pour tout $U \in \mathcal{B}$ et tout $p > 0$, on a

\begin{enumerate}

\item  $H^p(U, \mathcal{F}^q) = 0$,
\item  $H^p(U, \Coker(\mathcal{F}^{q - 1} \to \mathcal{F}^q)) = 0$,
\item  $H^p(U, H^q(\mathcal{F})) = 0$.

\end{enumerate}

\end{enumerate}

Alors le morphisme
$$
\text{Tot}(\check{\mathcal{C}}^\bullet_{alt}(\mathcal{U}, \mathcal{F}^\bullet))
\longrightarrow
R\Gamma(X, \mathcal{F}^\bullet)
$$
du lemme \ref{lemma-alternating-cech-complex-complex}
est un isomorphisme dans $D(\textit{Ab})$.

\end{lemma}

\begin{proof}
Supposons d'abord que $\mathcal{F}^\bullet$ soit borné inférieurement. Dans ce cas, le morphisme $$
\text{Tot}(\check{\mathcal{C}}^\bullet_{alt}(\mathcal{U}, \mathcal{F}^\bullet))
\longrightarrow
\text{Tot}(\check{\mathcal{C}}^\bullet(\mathcal{U}, \mathcal{F}^\bullet))
$$ est un quasi-isomorphisme d'après le lemme \ref{lemma-alternating-usual}. En effet, le morphisme de complexes doubles $\check{\mathcal{C}}^\bullet_{alt}(\mathcal{U}, \mathcal{F}^\bullet) \to
\check{\mathcal{C}}^\bullet(\mathcal{U}, \mathcal{F}^\bullet)$ induit un isomorphisme entre les premières pages des secondes suites spectrales associées à ces complexes (par Homologie, lemme \ref{homology-lemma-ss-double-complex}), et ces suites spectrales convergent (Homologie, lemme \ref{homology-lemma-first-quadrant-ss}). La conclusion découle donc, dans ce cas, du lemme \ref{lemma-cech-complex-complex-computes} et de l'hypothèse (3)(a).

\medskip\noindent

En général, par l'hypothèse (3)c), nous pouvons choisir une résolution
$\mathcal{F}^\bullet \to \mathcal{I}^\bullet = \lim \mathcal{I}_n^\bullet$
comme dans le lemme \ref{lemma-K-injective}. Le morphisme du lemme devient alors
$$
\lim_n 
\text{Tot}(\check{\mathcal{C}}^\bullet_{alt}(\mathcal{U},
\tau_{\geq -n}\mathcal{F}^\bullet))
\longrightarrow
\Gamma(X, \mathcal{I}^\bullet) =
\lim_n \Gamma(X, \mathcal{I}_n^\bullet)
$$
Ici, la flèche appartient à la catégorie dérivée, mais l'égalité de droite vaut au niveau des complexes. Remarquons que l'hypothèse (3)(b) montre que $\tau_{\geq -n}\mathcal{F}^\bullet$
est un complexe borné inférieurement satisfaisant aux hypothèses du lemme. Ainsi, le cas des complexes bornés inférieurement montre que chacun des morphismes
$$
\text{Tot}(\check{\mathcal{C}}^\bullet_{alt}(\mathcal{U},
\tau_{\geq -n}\mathcal{F}^\bullet))
\longrightarrow
\Gamma(X, \mathcal{I}_n^\bullet)
$$
est un quasi-isomorphisme. En chaque degré donné, les cohomologies des complexes du membre de gauche sont stationnaires (car le complexe de {\v C}ech alterné est fini). Il en va donc de même pour le membre de droite. La cohomologie de la limite du membre de droite est par conséquent égale à cette valeur stationnaire, d'après Homologie, lemme \ref{homology-lemma-apply-Mittag-Leffler-again}
(ou, plus précisément, Compléments sur l'algèbre, lemme
\ref{more-algebra-lemma-apply-Mittag-Leffler-again}), ce qui démontre le résultat.

\end{proof}





\section{Complexes de Hom}

\label{section-hom-complexes}

\noindent

Soit $(X, \mathcal{O}_X)$ un espace annelé. Soient
$\mathcal{L}^\bullet$ et $\mathcal{M}^\bullet$ deux complexes de $\mathcal{O}_X$-modules. Nous construisons un complexe de $\mathcal{O}_X$-modules
$\SheafHom^\bullet(\mathcal{L}^\bullet, \mathcal{M}^\bullet)$. Plus précisément, pour tout $n$, posons
$$
\SheafHom^n(\mathcal{L}^\bullet, \mathcal{M}^\bullet) =
\prod\nolimits_{n = p + q}
\SheafHom_{\mathcal{O}_X}(\mathcal{L}^{-q}, \mathcal{M}^p)
$$
On peut voir $\SheafHom^n$ comme le faisceau de $\mathcal{O}_X$-modules formé des applications $\mathcal{O}_X$-linéaires de $\mathcal{L}^\bullet$ dans $\mathcal{M}^\bullet$
(considérés comme des $\mathcal{O}_X$-modules gradués) qui sont homogènes de degré $n$. Dans cette terminologie, nous définissons le différentiel par la formule
$$
\text{d}(f) =
\text{d}_\mathcal{M} \circ f - (-1)^n f \circ \text{d}_\mathcal{L}
$$
pour
$f \in \SheafHom^n_{\mathcal{O}_X}(\mathcal{L}^\bullet, \mathcal{M}^\bullet)$. Nous omettons la vérification que $\text{d}^2 = 0$. Cette construction est un cas particulier d'Algèbre différentielle graduée, exemple \ref{dga-example-category-complexes}. Il résulte immédiatement de la construction que

\begin{equation}
\label{equation-cohomology-hom-complex}
H^n(\Gamma(U, \SheafHom^\bullet(\mathcal{L}^\bullet, \mathcal{M}^\bullet))) =
\Hom_{K(\mathcal{O}_U)}(\mathcal{L}^\bullet, \mathcal{M}^\bullet[n])
\end{equation}

pour tous $n \in \mathbf{Z}$ et tous les ouverts $U \subset X$.

\begin{lemma}

\label{lemma-compose}
Soit $(X, \mathcal{O}_X)$ un espace annelé. Étant donnés des complexes $\mathcal{K}^\bullet, \mathcal{L}^\bullet, \mathcal{M}^\bullet$ de $\mathcal{O}_X$-modules, il existe un isomorphisme $$
\SheafHom^\bullet(\mathcal{K}^\bullet,
\SheafHom^\bullet(\mathcal{L}^\bullet, \mathcal{M}^\bullet))
=
\SheafHom^\bullet(\text{Tot}(\mathcal{K}^\bullet \otimes_{\mathcal{O}_X}
\mathcal{L}^\bullet), \mathcal{M}^\bullet)
$$ de complexes de $\mathcal{O}_X$-modules, fonctoriel en $\mathcal{K}^\bullet, \mathcal{L}^\bullet, \mathcal{M}^\bullet$.
\end{lemma}

\begin{proof}
Démonstration omise. Indication : cela se démontre exactement de la même manière que Compléments sur l'algèbre, lemme \ref{more-algebra-lemma-compose}.
\end{proof}

\begin{lemma}

\label{lemma-composition}
Soit $(X, \mathcal{O}_X)$ un espace annelé. Étant donnés des complexes $\mathcal{K}^\bullet, \mathcal{L}^\bullet, \mathcal{M}^\bullet$ de $\mathcal{O}_X$-modules, il existe un morphisme canonique $$
\text{Tot}\left(
\SheafHom^\bullet(\mathcal{L}^\bullet, \mathcal{M}^\bullet)
\otimes_{\mathcal{O}_X}
\SheafHom^\bullet(\mathcal{K}^\bullet, \mathcal{L}^\bullet)
\right)
\longrightarrow
\SheafHom^\bullet(\mathcal{K}^\bullet, \mathcal{M}^\bullet)
$$ de complexes de $\mathcal{O}_X$-modules.
\end{lemma}

\begin{proof}
Démonstration omise. Indication : cela se démontre exactement de la même manière que Compléments sur l'algèbre, lemme \ref{more-algebra-lemma-composition}.
\end{proof}

\begin{lemma}

\label{lemma-diagonal-better}

Soit $(X, \mathcal{O}_X)$ un espace annelé. Étant donnés des complexes
$\mathcal{K}^\bullet, \mathcal{L}^\bullet, \mathcal{M}^\bullet$
de $\mathcal{O}_X$-modules, il existe un morphisme canonique
$$
\text{Tot}\left(
\mathcal{K}^\bullet \otimes_{\mathcal{O}_X}
\SheafHom^\bullet(\mathcal{M}^\bullet, \mathcal{L}^\bullet)
\right)
\longrightarrow
\SheafHom^\bullet(\mathcal{M}^\bullet,
\text{Tot}(\mathcal{K}^\bullet \otimes_{\mathcal{O}_X} \mathcal{L}^\bullet))
$$
de complexes de $\mathcal{O}_X$-modules, fonctoriel en les trois complexes.

\end{lemma}

\begin{proof}
Démonstration omise. Indication : cela se démontre exactement de la même manière que Compléments sur l'algèbre, lemme \ref{more-algebra-lemma-diagonal-better}.
\end{proof}

\begin{lemma}

\label{lemma-diagonal}

Soit $(X, \mathcal{O}_X)$ un espace annelé. Étant donnés des complexes
$\mathcal{K}^\bullet, \mathcal{L}^\bullet$
de $\mathcal{O}_X$-modules, il existe un morphisme canonique
$$
\mathcal{K}^\bullet
\longrightarrow
\SheafHom^\bullet(\mathcal{L}^\bullet,
\text{Tot}(\mathcal{K}^\bullet \otimes_{\mathcal{O}_X} \mathcal{L}^\bullet))
$$
de complexes de $\mathcal{O}_X$-modules, fonctoriel en les deux complexes.

\end{lemma}

\begin{proof}
Démonstration omise. Indication : cela se démontre exactement de la même manière que Compléments sur l'algèbre, lemme \ref{more-algebra-lemma-diagonal}.
\end{proof}

\begin{lemma}

\label{lemma-evaluate-and-more}

Soit $(X, \mathcal{O}_X)$ un espace annelé. Étant donnés des complexes
$\mathcal{K}^\bullet, \mathcal{L}^\bullet, \mathcal{M}^\bullet$
de $\mathcal{O}_X$-modules, il existe un morphisme canonique
$$
\text{Tot}(\SheafHom^\bullet(\mathcal{L}^\bullet,
\mathcal{M}^\bullet) \otimes_{\mathcal{O}_X} \mathcal{K}^\bullet)
\longrightarrow
\SheafHom^\bullet(\SheafHom^\bullet(\mathcal{K}^\bullet,
\mathcal{L}^\bullet), \mathcal{M}^\bullet)
$$
de complexes de $\mathcal{O}_X$-modules, fonctoriel en les trois complexes.

\end{lemma}

\begin{proof}
Démonstration omise. Indication : cela se démontre exactement de la même manière que Compléments sur l'algèbre, lemme \ref{more-algebra-lemma-evaluate-and-more}.
\end{proof}

\begin{lemma}

\label{lemma-RHom-into-K-injective}
Soit $(X, \mathcal{O}_X)$ un espace annelé. Soient $L$ et $M$ des objets de $D(\mathcal{O}_X)$. Soit $\mathcal{I}^\bullet$ un complexe K-injectif de $\mathcal{O}_X$-modules représentant $M$. Soit $\mathcal{L}^\bullet$ un complexe de $\mathcal{O}_X$-modules représentant $L$. Alors $$
H^0(\Gamma(U, \SheafHom^\bullet(\mathcal{L}^\bullet, \mathcal{I}^\bullet))) =
\Hom_{D(\mathcal{O}_U)}(L|_U, M|_U)
$$ pour tout ouvert $U \subset X$.
\end{lemma}

\begin{proof}

Nous avons

\begin{align*}
H^0(\Gamma(U, \SheafHom^\bullet(\mathcal{L}^\bullet, \mathcal{I}^\bullet)))
& =
\Hom_{K(\mathcal{O}_U)}(\mathcal{L}^\bullet|_U, \mathcal{I}^\bullet|_U) \\
& =
\Hom_{D(\mathcal{O}_U)}(L|_U, M|_U)
\end{align*}

La première égalité est (\ref{equation-cohomology-hom-complex}). La seconde est vraie parce que $\mathcal{I}^\bullet|_U$
est K-injectif d'après le lemme \ref{lemma-restrict-K-injective-to-open}.

\end{proof}

\begin{lemma}

\label{lemma-RHom-well-defined}

Soit $(X, \mathcal{O}_X)$ un espace annelé. Soit
$(\mathcal{I}')^\bullet \to \mathcal{I}^\bullet$
un quasi-isomorphisme de complexes K-injectifs de $\mathcal{O}_X$-modules. Soit $(\mathcal{L}')^\bullet \to \mathcal{L}^\bullet$
un quasi-isomorphisme de complexes de $\mathcal{O}_X$-modules. Alors
$$
\SheafHom^\bullet(\mathcal{L}^\bullet, (\mathcal{I}')^\bullet)
\longrightarrow
\SheafHom^\bullet((\mathcal{L}')^\bullet, \mathcal{I}^\bullet)
$$
est un quasi-isomorphisme.

\end{lemma}

\begin{proof}
Soit $M$ l'objet de $D(\mathcal{O}_X)$ représenté par $\mathcal{I}^\bullet$ et $(\mathcal{I}')^\bullet$. Soit $L$ l'objet de $D(\mathcal{O}_X)$ représenté par $\mathcal{L}^\bullet$ et $(\mathcal{L}')^\bullet$. D'après le lemme \ref{lemma-RHom-into-K-injective}, les faisceaux $$
H^0(\SheafHom^\bullet(\mathcal{L}^\bullet, (\mathcal{I}')^\bullet))
\quad\text{et}\quad
H^0(\SheafHom^\bullet((\mathcal{L}')^\bullet, \mathcal{I}^\bullet))
$$ sont tous deux égaux au faisceau associé au préfaisceau $$
U \longmapsto \Hom_{D(\mathcal{O}_U)}(L|_U, M|_U)
$$ Le morphisme est donc un quasi-isomorphisme.
\end{proof}

\begin{lemma}

\label{lemma-RHom-from-K-flat-into-K-injective}
Soit $(X, \mathcal{O}_X)$ un espace annelé. Soit $\mathcal{I}^\bullet$ un complexe K-injectif de $\mathcal{O}_X$-modules. Soit $\mathcal{L}^\bullet$ un complexe K-plat de $\mathcal{O}_X$-modules. Alors $\SheafHom^\bullet(\mathcal{L}^\bullet, \mathcal{I}^\bullet)$ est un complexe K-injectif de $\mathcal{O}_X$-modules.
\end{lemma}

\begin{proof}

En effet, si $\mathcal{K}^\bullet$ est un complexe acyclique de
$\mathcal{O}_X$-modules, alors

\begin{align*}
\Hom_{K(\mathcal{O}_X)}(\mathcal{K}^\bullet,
\SheafHom^\bullet(\mathcal{L}^\bullet, \mathcal{I}^\bullet))
& =
H^0(\Gamma(X,
\SheafHom^\bullet(\mathcal{K}^\bullet,
\SheafHom^\bullet(\mathcal{L}^\bullet, \mathcal{I}^\bullet)))) \\
& =
H^0(\Gamma(X, \SheafHom^\bullet(\text{Tot}(
\mathcal{K}^\bullet \otimes_{\mathcal{O}_X} \mathcal{L}^\bullet),
\mathcal{I}^\bullet))) \\
& =
\Hom_{K(\mathcal{O}_X)}(
\text{Tot}(\mathcal{K}^\bullet \otimes_{\mathcal{O}_X} \mathcal{L}^\bullet),
\mathcal{I}^\bullet) \\
& =
0
\end{align*}

La première égalité découle de (\ref{equation-cohomology-hom-complex}), la deuxième du lemme \ref{lemma-compose} et la troisième de (\ref{equation-cohomology-hom-complex}). La dernière égalité vaut parce que
$\text{Tot}(\mathcal{K}^\bullet \otimes_{\mathcal{O}_X} \mathcal{L}^\bullet)$
est acyclique puisque $\mathcal{L}^\bullet$ est K-plat (définition \ref{definition-K-flat}) et puisque $\mathcal{I}^\bullet$
est K-injectif.

\end{proof}







\section{Hom interne dans la catégorie dérivée}

\label{section-internal-hom}

\noindent

Soit $(X, \mathcal{O}_X)$ un espace annelé. Soient $L, M$ des objets de $D(\mathcal{O}_X)$. Nous voulons construire un objet
$R\SheafHom(L, M)$ de $D(\mathcal{O}_X)$ tel que, pour tout autre objet $K$ de $D(\mathcal{O}_X)$, il existe une bijection canonique

\begin{equation}
\label{equation-internal-hom}
\Hom_{D(\mathcal{O}_X)}(K, R\SheafHom(L, M))
=
\Hom_{D(\mathcal{O}_X)}(K \otimes_{\mathcal{O}_X}^\mathbf{L} L, M)
\end{equation}

Remarquons que cette formule définit $R\SheafHom(L, M)$ à isomorphisme unique près, par le lemme de Yoneda (Catégories, lemme \ref{categories-lemma-yoneda}).

\medskip\noindent
Pour construire un tel objet, choisissons un complexe K-injectif $\mathcal{I}^\bullet$ représentant $M$ et un complexe quelconque de $\mathcal{O}_X$-modules $\mathcal{L}^\bullet$ représentant $L$. Posons alors $$
R\SheafHom(L, M) = \SheafHom^\bullet(\mathcal{L}^\bullet, \mathcal{I}^\bullet)
$$ où le membre de droite est le complexe de $\mathcal{O}_X$-modules construit dans la section \ref{section-hom-complexes}. Cette définition est bien posée d'après le lemme \ref{lemma-RHom-well-defined}. Nous obtenons un foncteur $$
D(\mathcal{O}_X)^{opp} \times D(\mathcal{O}_X) \longrightarrow D(\mathcal{O}_X),
\quad
(K, L) \longmapsto R\SheafHom(K, L)
$$ Avant de démontrer (\ref{equation-internal-hom}), calculons les groupes de cohomologie de $R\SheafHom(K, L)$.

\begin{lemma}

\label{lemma-section-RHom-over-U}
Soit $(X, \mathcal{O}_X)$ un espace annelé. Soient $L, M$ des objets de $D(\mathcal{O}_X)$. Pour tout ouvert $U$, on a $$
H^0(U, R\SheafHom(L, M)) =
\Hom_{D(\mathcal{O}_U)}(L|_U, M|_U)
$$ et en particulier $H^0(X, R\SheafHom(L, M)) = \Hom_{D(\mathcal{O}_X)}(L, M)$.
\end{lemma}

\begin{proof}
Choisissons un complexe K-injectif $\mathcal{I}^\bullet$ de $\mathcal{O}_X$-modules représentant $M$ et un complexe K-plat $\mathcal{L}^\bullet$ représentant $L$. Alors $\SheafHom^\bullet(\mathcal{L}^\bullet, \mathcal{I}^\bullet)$ est K-injectif d'après le lemme \ref{lemma-RHom-from-K-flat-into-K-injective}. Nous pouvons donc calculer la cohomologie sur $U$ en prenant simplement les sections sur $U$, et le résultat découle du lemme \ref{lemma-RHom-into-K-injective}.
\end{proof}

\begin{lemma}

\label{lemma-internal-hom}
Soit $(X, \mathcal{O}_X)$ un espace annelé. Soient $K, L, M$ des objets de $D(\mathcal{O}_X)$. Avec la construction décrite ci-dessus, il existe un isomorphisme canonique $$
R\SheafHom(K, R\SheafHom(L, M)) =
R\SheafHom(K \otimes_{\mathcal{O}_X}^\mathbf{L} L, M)
$$ dans $D(\mathcal{O}_X)$, fonctoriel en $K, L, M$, qui redonne (\ref{equation-internal-hom}) après application de $H^0(X, -)$.
\end{lemma}

\begin{proof}

Choisissons un complexe K-injectif $\mathcal{I}^\bullet$ représentant
$M$ et un complexe K-plat de $\mathcal{O}_X$-modules $\mathcal{L}^\bullet$
représentant $L$. Soit $\mathcal{K}^\bullet$ un complexe quelconque de
$\mathcal{O}_X$-modules représentant $K$. Alors
$$
\SheafHom^\bullet(\mathcal{K}^\bullet,
\SheafHom^\bullet(\mathcal{L}^\bullet, \mathcal{I}^\bullet))
=
\SheafHom^\bullet(
\text{Tot}(\mathcal{K}^\bullet \otimes_{\mathcal{O}_X} \mathcal{L}^\bullet),
\mathcal{I}^\bullet)
$$
d'après le lemme \ref{lemma-compose}. Remarquons que le membre de gauche représente
$R\SheafHom(K, R\SheafHom(L, M))$ (utiliser le lemme \ref{lemma-RHom-from-K-flat-into-K-injective}) et que le membre de droite représente
$R\SheafHom(K \otimes_{\mathcal{O}_X}^\mathbf{L} L, M)$. Cela démontre la formule du lemme. En prenant les sections globales et en utilisant le lemme \ref{lemma-section-RHom-over-U},
nous obtenons (\ref{equation-internal-hom}).

\end{proof}

\begin{lemma}

\label{lemma-restriction-RHom-to-U}
Soit $(X, \mathcal{O}_X)$ un espace annelé. Soient $K, L$ des objets de $D(\mathcal{O}_X)$. La construction de $R\SheafHom(K, L)$ commute à la restriction aux ouverts, c'est-à-dire que, pour tout ouvert $U$, on a $R\SheafHom(K|_U, L|_U) = R\SheafHom(K, L)|_U$.
\end{lemma}

\begin{proof}
Cela résulte immédiatement de la construction et du lemme \ref{lemma-restrict-K-injective-to-open}.
\end{proof}

\begin{lemma}

\label{lemma-RHom-triangulated}

Soit $(X, \mathcal{O}_X)$ un espace annelé. Le bifoncteur $R\SheafHom(- , -)$
transforme les triangles distingués en triangles distingués dans les deux variables.

\end{lemma}

\begin{proof}

Cela résulte du fait que l'application
$$
(\mathcal{L}^\bullet, \mathcal{M}^\bullet) \longmapsto
\SheafHom^\bullet(\mathcal{L}^\bullet, \mathcal{M}^\bullet)
$$
transforme une suite exacte courte de complexes scindée terme à terme, dans l'une ou l'autre variable, en une suite exacte courte scindée terme à terme. Nous omettons les détails.

\end{proof}

\begin{lemma}

\label{lemma-internal-hom-composition}
Soit $(X, \mathcal{O}_X)$ un espace annelé. Étant donnés $K, L, M$ dans $D(\mathcal{O}_X)$, il existe un morphisme canonique $$
R\SheafHom(L, M) \otimes_{\mathcal{O}_X}^\mathbf{L} R\SheafHom(K, L)
\longrightarrow R\SheafHom(K, M)
$$ dans $D(\mathcal{O}_X)$, fonctoriel en $K, L, M$.
\end{lemma}

\begin{proof}

Choisissons un complexe K-injectif $\mathcal{I}^\bullet$ représentant $M$, un complexe K-injectif $\mathcal{J}^\bullet$ représentant $L$, et un complexe quelconque de $\mathcal{O}_X$-modules $\mathcal{K}^\bullet$ représentant $K$. D'après le lemme \ref{lemma-composition}, il existe un morphisme de complexes
$$
\text{Tot}\left(
\SheafHom^\bullet(\mathcal{J}^\bullet, \mathcal{I}^\bullet)
\otimes_{\mathcal{O}_X}
\SheafHom^\bullet(\mathcal{K}^\bullet, \mathcal{J}^\bullet)
\right)
\longrightarrow
\SheafHom^\bullet(\mathcal{K}^\bullet, \mathcal{I}^\bullet)
$$
Les complexes de $\mathcal{O}_X$-modules
$\SheafHom^\bullet(\mathcal{J}^\bullet, \mathcal{I}^\bullet)$,
$\SheafHom^\bullet(\mathcal{K}^\bullet, \mathcal{J}^\bullet)$,
$\SheafHom^\bullet(\mathcal{K}^\bullet, \mathcal{I}^\bullet)$
représentent respectivement $R\SheafHom(L, M)$, $R\SheafHom(K, L)$ et $R\SheafHom(K, M)$. Si nous choisissons un complexe K-plat $\mathcal{H}^\bullet$ et un quasi-isomorphisme
$\mathcal{H}^\bullet \to
\SheafHom^\bullet(\mathcal{K}^\bullet, \mathcal{J}^\bullet)$, alors il existe un morphisme
$$
\text{Tot}\left(
\SheafHom^\bullet(\mathcal{J}^\bullet, \mathcal{I}^\bullet)
\otimes_{\mathcal{O}_X} \mathcal{H}^\bullet
\right)
\longrightarrow
\text{Tot}\left(
\SheafHom^\bullet(\mathcal{J}^\bullet, \mathcal{I}^\bullet)
\otimes_{\mathcal{O}_X}
\SheafHom^\bullet(\mathcal{K}^\bullet, \mathcal{J}^\bullet)
\right)
$$
dont la source représente
$R\SheafHom(L, M) \otimes_{\mathcal{O}_X}^\mathbf{L} R\SheafHom(K, L)$. En composant les deux flèches affichées, on obtient le morphisme souhaité. Nous omettons la preuve de la fonctorialité de la construction.

\end{proof}

\begin{lemma}

\label{lemma-internal-hom-diagonal-better}
Soit $(X, \mathcal{O}_X)$ un espace annelé. Étant donnés $K, L, M$ dans $D(\mathcal{O}_X)$, il existe un morphisme canonique $$
K \otimes_{\mathcal{O}_X}^\mathbf{L} R\SheafHom(M, L)
\longrightarrow
R\SheafHom(M, K \otimes_{\mathcal{O}_X}^\mathbf{L} L)
$$ dans $D(\mathcal{O}_X)$, fonctoriel en $K, L, M$.
\end{lemma}

\begin{proof}
Choisissons un complexe K-plat $\mathcal{K}^\bullet$ représentant $K$, un complexe K-injectif $\mathcal{I}^\bullet$ représentant $L$, et un complexe quelconque de $\mathcal{O}_X$-modules $\mathcal{M}^\bullet$ représentant $M$. Choisissons un quasi-isomorphisme $\text{Tot}(\mathcal{K}^\bullet \otimes_{\mathcal{O}_X} \mathcal{I}^\bullet)
\to \mathcal{J}^\bullet$, où $\mathcal{J}^\bullet$ est K-injectif. Nous utilisons alors le morphisme $$
\text{Tot}\left(
\mathcal{K}^\bullet \otimes_{\mathcal{O}_X}
\SheafHom^\bullet(\mathcal{M}^\bullet, \mathcal{I}^\bullet)
\right)
\to
\SheafHom^\bullet(\mathcal{M}^\bullet,
\text{Tot}(\mathcal{K}^\bullet \otimes_{\mathcal{O}_X} \mathcal{I}^\bullet))
\to
\SheafHom^\bullet(\mathcal{M}^\bullet, \mathcal{J}^\bullet)
$$ où la première flèche est le morphisme du lemme \ref{lemma-diagonal-better}.
\end{proof}

\begin{lemma}

\label{lemma-internal-hom-diagonal}
Soit $(X, \mathcal{O}_X)$ un espace annelé. Étant donnés $K, L$ dans $D(\mathcal{O}_X)$, il existe un morphisme canonique $$
K \longrightarrow R\SheafHom(L, K \otimes_{\mathcal{O}_X}^\mathbf{L} L)
$$ dans $D(\mathcal{O}_X)$, fonctoriel en $K$ et $L$.
\end{lemma}

\begin{proof}
Choisissons un complexe K-plat $\mathcal{K}^\bullet$ représentant $K$ et un complexe quelconque de $\mathcal{O}_X$-modules $\mathcal{L}^\bullet$ représentant $L$. Choisissons un complexe K-injectif $\mathcal{J}^\bullet$ et un quasi-isomorphisme $\text{Tot}(\mathcal{K}^\bullet \otimes_{\mathcal{O}_X} \mathcal{L}^\bullet)
\to \mathcal{J}^\bullet$. Nous utilisons alors $$
\mathcal{K}^\bullet \to
\SheafHom^\bullet(\mathcal{L}^\bullet,
\text{Tot}(\mathcal{K}^\bullet \otimes_{\mathcal{O}_X} \mathcal{L}^\bullet))
\to
\SheafHom^\bullet(\mathcal{L}^\bullet, \mathcal{J}^\bullet)
$$ où la première flèche provient du lemme \ref{lemma-diagonal}.
\end{proof}

\begin{lemma}

\label{lemma-dual}

Soit $(X, \mathcal{O}_X)$ un espace annelé. Soit $L$ un objet de $D(\mathcal{O}_X)$. Posons $L^\vee = R\SheafHom(L, \mathcal{O}_X)$. Pour $M$ dans $D(\mathcal{O}_X)$, il existe un morphisme canonique

\begin{equation}
\label{equation-eval}
M \otimes^\mathbf{L}_{\mathcal{O}_X} L^\vee
\longrightarrow
R\SheafHom(L, M)
\end{equation}

qui induit un morphisme canonique
$$
H^0(X, M \otimes^\mathbf{L}_{\mathcal{O}_X} L^\vee)
\longrightarrow
\Hom_{D(\mathcal{O}_X)}(L, M)
$$
fonctoriel en $M$ dans $D(\mathcal{O}_X)$.

\end{lemma}

\begin{proof}
Le morphisme (\ref{equation-eval}) est un cas particulier du lemme \ref{lemma-internal-hom-composition}, en utilisant l'identification $M = R\SheafHom(\mathcal{O}_X, M)$.
\end{proof}

\begin{lemma}

\label{lemma-internal-hom-evaluate}
Soit $(X, \mathcal{O}_X)$ un espace annelé. Soient $K, L, M$ des objets de $D(\mathcal{O}_X)$. Il existe un morphisme canonique $$
R\SheafHom(L, M) \otimes_{\mathcal{O}_X}^\mathbf{L} K
\longrightarrow
R\SheafHom(R\SheafHom(K, L), M)
$$ dans $D(\mathcal{O}_X)$, fonctoriel en $K, L, M$.
\end{lemma}

\begin{proof}
Choisissons un complexe K-injectif $\mathcal{I}^\bullet$ représentant $M$, un complexe K-injectif $\mathcal{J}^\bullet$ représentant $L$, et un complexe K-plat $\mathcal{K}^\bullet$ représentant $K$. Le morphisme est défini au moyen du morphisme $$
\text{Tot}(\SheafHom^\bullet(\mathcal{J}^\bullet,
\mathcal{I}^\bullet) \otimes_{\mathcal{O}_X} \mathcal{K}^\bullet)
\longrightarrow
\SheafHom^\bullet(\SheafHom^\bullet(\mathcal{K}^\bullet,
\mathcal{J}^\bullet), \mathcal{I}^\bullet)
$$ du lemme \ref{lemma-evaluate-and-more}. Par notre choix particulier des complexes, le membre de gauche représente $R\SheafHom(L, M) \otimes_{\mathcal{O}_X}^\mathbf{L} K$ et le membre de droite représente $R\SheafHom(R\SheafHom(K, L), M)$. Nous omettons la preuve de la fonctorialité en les trois objets de $D(\mathcal{O}_X)$.
\end{proof}

\begin{remark}

\label{remark-tensor-internal-hom}

Soit $(X, \mathcal{O}_X)$ un espace annelé. Pour $K, K', M, M'$ dans
$D(\mathcal{O}_X)$, il existe un morphisme canonique
$$
R\SheafHom(K, K') \otimes_{\mathcal{O}_X}^\mathbf{L}
R\SheafHom(M, M')
\longrightarrow
R\SheafHom(K \otimes_{\mathcal{O}_X}^\mathbf{L} M,
K' \otimes_{\mathcal{O}_X}^\mathbf{L} M')
$$
En effet, d'après (\ref{equation-internal-hom}), cela revient à construire un morphisme
$$
R\SheafHom(K, K') \otimes_{\mathcal{O}_X}^\mathbf{L}
R\SheafHom(M, M') \otimes_{\mathcal{O}_X}^\mathbf{L}
K \otimes_{\mathcal{O}_X}^\mathbf{L} M
\longrightarrow
K' \otimes_{\mathcal{O}_X}^\mathbf{L} M'
$$
Pour cela, nous pouvons d'abord permuter les deux facteurs du milieu (avec les règles de signes de Compléments sur l'algèbre, section
\ref{more-algebra-section-sign-rules}), puis utiliser les morphismes
$$
R\SheafHom(K, K') \otimes_{\mathcal{O}_X}^\mathbf{L} K \to K'
\quad\text{et}\quad
R\SheafHom(M, M') \otimes_{\mathcal{O}_X}^\mathbf{L} M \to M'
$$
du lemme \ref{lemma-internal-hom-composition}, en considérant $K = R\SheafHom(\mathcal{O}_X, K)$ et de même pour
$K'$, $M$ et $M'$.

\end{remark}

\begin{remark}

\label{remark-projection-formula-for-internal-hom}

Soit $f : X \to Y$ un morphisme d'espaces annelés. Soient $K, L$ des objets de $D(\mathcal{O}_X)$. Nous affirmons qu'il existe un morphisme canonique
$$
Rf_*R\SheafHom(L, K) \longrightarrow R\SheafHom(Rf_*L, Rf_*K)
$$
En effet, d'après (\ref{equation-internal-hom}), cela revient à construire un morphisme
$Rf_*R\SheafHom(L, K) \otimes_{\mathcal{O}_Y}^\mathbf{L} Rf_*L \to Rf_*K$. Pour cela, nous pouvons utiliser la composition
$$
Rf_*R\SheafHom(L, K) \otimes_{\mathcal{O}_Y}^\mathbf{L} Rf_*L \to
Rf_*(R\SheafHom(L, K) \otimes_{\mathcal{O}_X}^\mathbf{L} L) \to
Rf_*K
$$
où la première flèche est le cup-produit relatif (remarque \ref{remark-cup-product}) et la seconde est $Rf_*$ appliqué au morphisme canonique
$R\SheafHom(L, K) \otimes_{\mathcal{O}_X}^\mathbf{L} L \to K$
provenant du lemme \ref{lemma-internal-hom-composition}
(avec $\mathcal{O}_X$ dans l'un des arguments).

\end{remark}

\begin{remark}

\label{remark-relative-cup-and-composition}

Soit $h : X \to Y$ un morphisme d'espaces annelés. Soient $K, M$ des objets de $D(\mathcal{O}_Y)$. Le schéma
$$
\xymatrix{
Rf_*R\SheafHom_{\mathcal{O}_X}(K, M)
\otimes_{\mathcal{O}_Y}^\mathbf{L} Rf_*K
\ar[r] \ar[d] &
Rf_*\left(R\SheafHom_{\mathcal{O}_X}(K, M)
\otimes_{\mathcal{O}_X}^\mathbf{L} K\right)
\ar[d] \\
R\SheafHom_{\mathcal{O}_Y}(Rf_*K, Rf_*M) \otimes_{\mathcal{O}_Y}^\mathbf{L}
Rf_*K \ar[r] &
Rf_*M
}
$$
est commutatif. Ici, la flèche verticale de gauche provient de la remarque \ref{remark-projection-formula-for-internal-hom}. La flèche horizontale supérieure est celle de la remarque \ref{remark-cup-product}. Les deux autres flèches sont des cas particuliers du morphisme du lemme \ref{lemma-internal-hom-composition} (où l'une des entrées est remplacée par $\mathcal{O}_X$ ou $\mathcal{O}_Y$).

\end{remark}

\begin{remark}

\label{remark-prepare-fancy-base-change}

Soit $h : X \to Y$ un morphisme d'espaces annelés. Soient $K, L$ des objets de $D(\mathcal{O}_Y)$. Nous affirmons qu'il existe un morphisme canonique
$$
Lh^*R\SheafHom(K, L) \longrightarrow R\SheafHom(Lh^*K, Lh^*L)
$$
dans $D(\mathcal{O}_X)$. En effet, d'après (\ref{equation-internal-hom}), démontré dans le lemme \ref{lemma-internal-hom},
un tel morphisme revient à construire un morphisme
$$
Lh^*R\SheafHom(K, L) \otimes^\mathbf{L} Lh^*K \longrightarrow Lh^*L
$$
La source de cette flèche est $Lh^*(\SheafHom(K, L) \otimes^\mathbf{L} K)$
d'après le lemme \ref{lemma-pullback-tensor-product}
donc il suffit de construire un morphisme canonique
$$
R\SheafHom(K, L) \otimes^\mathbf{L} K \longrightarrow L.
$$
Pour cela, nous prenons la flèche correspondant à
$$
\text{id} :
R\SheafHom(K, L)
\longrightarrow
R\SheafHom(K, L)
$$
via (\ref{equation-internal-hom}).

\end{remark}

\begin{remark}

\label{remark-fancy-base-change}

Supposons que
$$
\xymatrix{
X' \ar[r]_h \ar[d]_{f'} &
X \ar[d]^f \\
S' \ar[r]^g &
S
}
$$
est un diagramme commutatif d'espaces annelés. Soient $K, L$ des objets de $D(\mathcal{O}_X)$. Nous affirmons qu'il existe un morphisme canonique de changement de base
$$
Lg^*Rf_*R\SheafHom(K, L)
\longrightarrow
R(f')_*R\SheafHom(Lh^*K, Lh^*L)
$$
dans $D(\mathcal{O}_{S'})$. Plus précisément, nous prenons le morphisme adjoint de la composition

\begin{align*}
L(f')^*Lg^*Rf_*R\SheafHom(K, L)
& =
Lh^*Lf^*Rf_*R\SheafHom(K, L) \\
& \to
Lh^*R\SheafHom(K, L) \\
& \to
R\SheafHom(Lh^*K, Lh^*L)
\end{align*}

où la première flèche utilise le morphisme d'adjonction
$Lf^*Rf_* \to \text{id}$ et la seconde est le morphisme canonique construit dans la remarque \ref{remark-prepare-fancy-base-change}.

\end{remark}
\section{Faisceaux d'Ext}

\label{section-ext}

\noindent

Soit $(X, \mathcal{O}_X)$ un espace annelé et soient $K, L \in D(\mathcal{O}_X)$. La construction du Hom interne dans la catégorie dérivée fournit un faisceau bien défini de $\mathcal{O}_X$-modules
$$
\SheafExt^n(K, L) = H^n(R\SheafHom(K, L))
$$
obtenu en prenant le $n$-ième faisceau de cohomologie de l'objet $R\SheafHom(K, L)$
de $D(\mathcal{O}_X)$. Nous noterons parfois
$\SheafExt^n_{\mathcal{O}_X}(K, L)$ cet objet. Par le lemme \ref{lemma-section-RHom-over-U},
ce faisceau $\SheafExt^n$ est le faisceau associé au préfaisceau
$$
U \longmapsto \Ext^n_{D(\mathcal{O}_U)}(K|_U, L|_U)
$$
Par l'exemple \ref{example-spectral-sequence}, il existe toujours une suite spectrale
$$
E_2^{p, q} = H^p(X, \SheafExt^q(K, L))
$$
convergeant vers $\Ext^{p + q}_{D(\mathcal{O}_X)}(K, L)$
dans des situations favorables (par exemple si $L$ est borné inférieurement et
$K$ borné supérieurement).





\section{Hom dérivé global}

\label{section-global-RHom}

\noindent

Soit $(X, \mathcal{O}_X)$ un espace annelé et soient $K, L \in D(\mathcal{O}_X)$. La construction du Hom interne dans la catégorie dérivée fournit un objet bien défini
$$
R\Hom_X(K, L) = R\Gamma(X, R\SheafHom(K, L))
$$
dans $D(\Gamma(X, \mathcal{O}_X))$. Nous noterons parfois
$R\Hom_{\mathcal{O}_X}(K, L)$ cet objet. Par le lemme \ref{lemma-section-RHom-over-U},
nous avons
$$
H^0(R\Hom_X(K, L)) = \Hom_{D(\mathcal{O}_X)}(K, L),
\quad
H^p(R\Hom_X(K, L)) = \Ext_{D(\mathcal{O}_X)}^p(K, L)
$$
Si $f : Y \to X$ est un morphisme d'espaces annelés, il existe une application canonique
$$
R\Hom_X(K, L) \longrightarrow R\Hom_Y(Lf^*K, Lf^*L)
$$
dans $D(\Gamma(X, \mathcal{O}_X))$, obtenue en prenant les sections globales de l'application définie dans la remarque \ref{remark-prepare-fancy-base-change}.







\section{Collage de complexes}

\label{section-glueing-complexes}

\noindent
Nous pouvons recoller des complexes ! Plus précisément, dans certaines circonstances, des objets de la catégorie dérivée donnés localement se recollent en un objet global. Nous démontrons d'abord quelques cas simples, puis le très général \cite[théorème 3.2.4]{BBD} dans le cadre des espaces topologiques et des recouvrements ouverts.

\begin{lemma}

\label{lemma-glue}

Soit $(X, \mathcal{O}_X)$ un espace annelé. Supposons que $X = U \cup V$ soit la réunion de deux sous-espaces ouverts de $X$ et que soient donnés

\begin{enumerate}

\item un objet $A$ de $D(\mathcal{O}_U)$,
\item un objet $B$ de $D(\mathcal{O}_V)$, et
\item un isomorphisme $c : A|_{U \cap V} \to B|_{U \cap V}$.

\end{enumerate}

Alors il existe un objet $F$ de $D(\mathcal{O}_X)$
et des isomorphismes $f : F|_U \to A$, $g : F|_V \to B$ tels que $c = g|_{U \cap V} \circ f^{-1}|_{U \cap V}$. De plus, étant donnés

\begin{enumerate}

\item un objet $E$ de $D(\mathcal{O}_X)$,
\item un morphisme $a : A \to E|_U$ de $D(\mathcal{O}_U)$,
\item un morphisme $b : B \to E|_V$ de $D(\mathcal{O}_V)$, 

\end{enumerate}

tels que
$$
a|_{U \cap V}  = b|_{U \cap V} \circ c.
$$
il existe un morphisme $F \to E$ dans $D(\mathcal{O}_X)$
dont la restriction à $U$ est $a \circ f$
et dont la restriction à $V$ est $b \circ g$.

\end{lemma}

\begin{proof}

Notons $j_U$, $j_V$, $j_{U \cap V}$ les immersions ouvertes correspondantes. Choisissons un triangle distingué
$$
F \to Rj_{U, *}A \oplus Rj_{V, *}B \to Rj_{U \cap V, *}(B|_{U \cap V})
\to F[1]
$$
où l'application $Rj_{V, *}B \to Rj_{U \cap V, *}(B|_{U \cap V})$ est l'application évidente et où
$Rj_{U, *}A \to Rj_{U \cap V, *}(B|_{U \cap V})$
est la composition de
$Rj_{U, *}A \to Rj_{U \cap V, *}(A|_{U \cap V})$
avec $Rj_{U \cap V, *}c$. En restreignant à $U$, nous obtenons
$$
F|_U \to A \oplus (Rj_{V, *}B)|_U \to (Rj_{U \cap V, *}(B|_{U \cap V}))|_U
\to F|_U[1]
$$
Notons $j : U \cap V \to U$. La compatibilité de la restriction aux ouverts avec la cohomologie montre que
$(Rj_{V, *}B)|_U$ et $(Rj_{U \cap V, *}(B|_{U \cap V}))|_U$
sont tous deux canoniquement isomorphes à $Rj_*(B|_{U \cap V})$. La deuxième flèche du dernier diagramme admet donc une section, et le morphisme $F|_U \to A$ est un isomorphisme. De même, le morphisme $F|_V \to B$ est un isomorphisme. L'existence du morphisme $F \to E$ résulte de la suite de Mayer--Vietoris pour $\Hom$ ; voir le lemme \ref{lemma-mayer-vietoris-hom}.

\end{proof}

\begin{lemma}

\label{lemma-vanishing-and-glueing}

Soit $f : (X, \mathcal{O}_X) \to (Y, \mathcal{O}_Y)$ un morphisme d'espaces annelés. Soit $\mathcal{B}$ une base de la topologie de $Y$.

\begin{enumerate}

\item Supposons que $K$ soit un objet de $D(\mathcal{O}_X)$ tel que, pour $V \in \mathcal{B}$, nous ayons $H^i(f^{-1}(V), K) = 0$ pour $i < 0$. Alors les faisceaux de cohomologie de $Rf_*K$ sont nuls en degrés négatifs,
$H^i(f^{-1}(V), K) = 0$ pour $i < 0$ pour tout ouvert $V \subset Y$, et la règle $V \mapsto H^0(f^{-1}V, K)$ définit un faisceau sur $Y$.
\item Supposons que $K, L$ soient des objets de $D(\mathcal{O}_X)$ tels que, pour $V \in \mathcal{B}$, nous ayons
$\Ext^i(K|_{f^{-1}V}, L|_{f^{-1}V}) = 0$ pour $i < 0$. Alors $\Ext^i(K|_{f^{-1}V}, L|_{f^{-1}V}) = 0$ pour $i < 0$
pour tout ouvert $V \subset Y$, et la règle $V \mapsto \Hom(K|_{f^{-1}V}, L|_{f^{-1}V})$ définit un faisceau sur $Y$.

\end{enumerate}

\end{lemma}

\begin{proof}
Le lemme \ref{lemma-unbounded-describe-higher-direct-images} affirme que $H^i(Rf_*K)$ est le faisceau associé au préfaisceau $V \mapsto H^i(f^{-1}(V), K) = H^i(V, Rf_*K)$. Les hypothèses de (1) impliquent que les faisceaux de cohomologie de $Rf_*K$ sont nuls en degrés $< 0$. Pour tout ouvert $V \subset Y$, le groupe de cohomologie $H^i(V, Rf_*K)$ est donc nul pour $i < 0$ et égal à $H^0(V, H^0(Rf_*K))$ pour $i = 0$. Cela démontre (1).

\medskip\noindent

Pour démontrer (2), appliquons (1) au complexe $R\SheafHom(K, L)$ et utilisons le lemme \ref{lemma-section-RHom-over-U} pour effectuer la traduction.

\end{proof}

\begin{situation}

\label{situation-locally-given}

Soit $(X, \mathcal{O}_X)$ un espace annelé. On se donne

\begin{enumerate}

\item une collection d'ouverts $\mathcal{B}$ de $X$,
\item pour $U \in \mathcal{B}$ un objet $K_U$ en $D(\mathcal{O}_U)$,
\item pour $V \subset U$ avec $V, U \in \mathcal{B}$ un isomorphisme
$\rho^U_V : K_U|_V \to K_V$ en $D(\mathcal{O}_V)$,

\end{enumerate}

de sorte que, pour $W \subset V \subset U$ avec $U, V, W$ dans
$\mathcal{B}$, on ait $\rho^U_W = \rho^V_W \circ \rho ^U_V|_W$.

\end{situation}

\noindent
Nous ne pourrons rien démontrer à ce sujet sans hypothèses supplémentaires. Un cas intéressant est celui où $\mathcal{B}$ est une base de la topologie de $X$. Un autre est celui où l'on dispose d'un morphisme $f : X \to Y$ d'espaces topologiques et où les éléments de $\mathcal{B}$ sont les images réciproques des éléments d'une base de la topologie de $Y$.

\medskip\noindent
Dans la situation \ref{situation-locally-given}, une {\it solution} est un couple $(K, \rho_U)$, où $K$ est un objet de $D(\mathcal{O}_X)$ et où les $\rho_U : K|_U \to K_U$, pour $U \in \mathcal{B}$, sont des isomorphismes tels que $\rho^U_V \circ \rho_U|_V = \rho_V$ pour tous $V \subset U$, $U, V \in \mathcal{B}$. Dans certains cas, les solutions sont uniques.

\begin{lemma}

\label{lemma-uniqueness}
Dans la situation \ref{situation-locally-given}, supposons que
\begin{enumerate}
\item $X = \bigcup_{U \in \mathcal{B}} U$ et pour $U, V \in \mathcal{B}$ nous avons $U \cap V = \bigcup_{W \in \mathcal{B}, W \subset U \cap V} W$,
\item pour tout $U \in \mathcal{B}$, nous ayons $\Ext^i(K_U, K_U) = 0$ pour $i < 0$.
\end{enumerate}
Si une solution $(K, \rho_U)$ existe, elle est unique à un unique isomorphisme près ; de plus, $\Ext^i(K, K) = 0$ pour $i < 0$.
\end{lemma}

\begin{proof}

Soient $(K, \rho_U)$ et $(K', \rho'_U)$ deux solutions. Soit $f : X \to Y$ l'application continue construite dans Topologie, lemme \ref{topology-lemma-create-map-from-subcollection}. Posons $\mathcal{O}_Y = f_*\mathcal{O}_X$. Alors $K, K'$ et $\mathcal{B}$ satisfont les hypothèses de la partie (2) du lemme \ref{lemma-vanishing-and-glueing}. Nous obtenons donc l'annulation des Ext négatifs de $K$, et la règle
$$
V \longmapsto \Hom(K|_{f^{-1}V}, K'|_{f^{-1}V})
$$
définit un faisceau sur $Y$. Comme $(K, \rho_U)$ et $(K', \rho'_U)$ sont deux solutions, les applications
$$
(\rho'_U)^{-1} \circ \rho_U : K|_U \longrightarrow K'|_U
$$
sur $U = f^{-1}(f(U))$ coïncident sur les intersections. Nous obtenons donc une unique section globale du faisceau ci-dessus ; elle définit l'isomorphisme recherché
$K \to K'$, compatible avec toutes les données.

\end{proof}

\begin{remark}

\label{remark-uniqueness}

Conservons les notations et les hypothèses du lemme \ref{lemma-uniqueness}. Soient $U, V \in \mathcal{B}$. Soit $\mathcal{B}'$ l'ensemble des éléments de $\mathcal{B}$ contenus dans $U \cap V$. Alors
$$
(\{K_{U'}\}_{U' \in \mathcal{B}'},
\{\rho_{V'}^{U'}\}_{V' \subset U'\text{ avec }U', V' \in \mathcal{B}'})
$$
est un système sur l'espace annelé $U \cap V$
satisfaisant aux hypothèses du lemme \ref{lemma-uniqueness}. De plus, $(K_U|_{U \cap V}, \rho^U_{U'})$ et
$(K_V|_{U \cap V}, \rho^V_{U'})$ sont deux solutions de ce système. Le lemme fournit donc un unique isomorphisme
$$
\rho_{U, V} : K_U|_{U \cap V} \longrightarrow K_V|_{U \cap V}
$$
tel que, pour tout $U' \subset U \cap V$, $U' \in \mathcal{B}$, le diagramme
$$
\xymatrix{
K_U|_{U'} \ar[rr]_{\rho_{U, V}|_{U'}} \ar[rd]_{\rho^U_{U'}} & &
K_V|_{U'} \ar[ld]^{\rho^V_{U'}} \\
& K_{U'}
}
$$
soit commutatif. Choisissons un troisième élément $W \in \mathcal{B}$. Nous obtenons des isomorphismes
$\rho_{U, W} : K_U|_{U \cap W} \to K_W|_{U \cap W}$ et
$\rho_{V, W} : K_U|_{V \cap W} \to K_W|_{V \cap W}$ satisfaisant des propriétés analogues à celles de $\rho_{U, V}$. Enfin, nous avons
$$
\rho_{U, W}|_{U \cap V \cap W} =
\rho_{V, W}|_{U \cap V \cap W} \circ \rho_{U, V}|_{U \cap V \cap W}
$$
Cela résulte de l'unicité dans le lemme, car les deux membres de l'égalité sont l'unique isomorphisme compatible avec les applications $\rho^U_{U''}$ et $\rho^W_{U''}$
pour $U'' \subset U \cap V \cap W$, $U'' \in \mathcal{B}$. Nous omettons quelques détails mineurs. La famille $(K_U, \rho_{U, V})$ est une donnée de descente dans la catégorie dérivée pour le recouvrement ouvert
$\mathcal{U} : X = \bigcup_{U \in \mathcal{B}} U$ de $X$. Dans ce langage, chercher une solution revient à chercher l'« effectivité de la donnée de descente ».

\end{remark}

\begin{lemma}

\label{lemma-solution-in-finite-case}

Dans la situation \ref{situation-locally-given}, supposons que

\begin{enumerate}

\item  $X = U_1 \cup \ldots \cup U_n$ avec $U_i \in \mathcal{B}$,
\item pour $U, V \in \mathcal{B}$, nous ayons
$U \cap V = \bigcup_{W \in \mathcal{B}, W \subset U \cap V} W$,
\item pour tout $U \in \mathcal{B}$, nous ayons $\Ext^i(K_U, K_U) = 0$
pour $i < 0$.

\end{enumerate}

Alors une solution existe et elle est unique à un unique isomorphisme près.

\end{lemma}

\begin{proof}
L'unicité a été établie dans le lemme \ref{lemma-uniqueness}. Démontrons le lemme par récurrence sur $n$. Le cas $n = 1$ est immédiat.

\medskip\noindent

Le cas $n = 2$. Considérons l'isomorphisme
$\rho_{U_1, U_2} : K_{U_1}|_{U_1 \cap U_2} \to K_{U_2}|_{U_1 \cap U_2}$
construit dans la remarque \ref{remark-uniqueness}. Par le lemme \ref{lemma-glue}, nous obtenons un objet $K$ de $D(\mathcal{O}_X)$
et des isomorphismes $\rho_{U_1} : K|_{U_1} \to K_{U_1}$ et
$\rho_{U_2} : K|_{U_2} \to K_{U_2}$ compatibles avec $\rho_{U_1, U_2}$. Soit $U \in \mathcal{B}$. Nous allons construire un isomorphisme
$\rho_U : K|_U \to K_U$ et laisserons au lecteur le soin de vérifier que $(K, \rho_U)$ est une solution. Considérons l'ensemble $\mathcal{B}'$
des éléments de $\mathcal{B}$ contenus dans $U \cap U_1$ ou dans
$U \cap U_2$. Alors $(K_U, \rho^U_{U'})$ est une solution du système
$(\{K_{U'}\}_{U' \in \mathcal{B}'},
\{\rho_{V'}^{U'}\}_{V' \subset U'\text{ avec }U', V' \in \mathcal{B}'})$
sur l'espace annelé $U$. Nous affirmons que $(K|_U, \tau_{U'})$ est une autre solution, où
$\tau_{U'}$, pour $U' \in \mathcal{B}'$, est défini comme suit : si $U' \subset U_1$, nous prenons la composée
$$
K|_{U'} \xrightarrow{\rho_{U_1}|_{U'}}
K_{U_1}|_{U'} \xrightarrow{\rho^{U_1}_{U'}}
K_{U'}
$$
et, si $U' \subset U_2$, la composée
$$
K|_{U'} \xrightarrow{\rho_{U_2}|_{U'}}
K_{U_2}|_{U'} \xrightarrow{\rho^{U_2}_{U'}}
K_{U'}.
$$
Pour vérifier qu'il s'agit d'une solution, utilisons la propriété de l'application $\rho_{U_1, U_2}$
décrite dans la remarque \ref{remark-uniqueness} et la compatibilité de
$\rho_{U_1}$ et $\rho_{U_2}$ avec $\rho_{U_1, U_2}$. Le lemme \ref{lemma-uniqueness} fournit alors un unique isomorphisme $K|_{U'} \to K_{U'}$ compatible avec les applications $\tau_{U'}$ et
$\rho^U_{U'}$ pour $U' \in \mathcal{B}'$.

\medskip\noindent

Le cas $n > 2$. Considérons le sous-espace ouvert
$X' = U_1 \cup \ldots \cup U_{n - 1}$ et soit $\mathcal{B}'$ l'ensemble des éléments de $\mathcal{B}$ contenus dans $X'$. Nous obtenons un système
$(\{K_U\}_{U \in \mathcal{B}'}, \{\rho_V^U\}_{U, V \in \mathcal{B}'})$
sur l'espace annelé $X'$, auquel nous pouvons appliquer l'hypothèse de récurrence. Nous obtenons une solution $(K_{X'}, \rho^{X'}_U)$. Considérons alors la famille
$\mathcal{B}^* = \mathcal{B} \cup \{X'\}$ des ouvertures de $X$ et nous voyons que nous obtenons un système
$(\{K_U\}_{U \in \mathcal{B}^*},
\{\rho_V^U\}_{V \subset U\text{ avec }U, V \in \mathcal{B}^*})$. Ce nouveau système satisfait encore la condition (3), par le lemme \ref{lemma-uniqueness} appliqué à la solution $K_{X'}$. Pour ce système, $X = X' \cup U_n$, ce qui nous ramène au cas $n = 2$ traité ci-dessus.

\end{proof}

\begin{lemma}

\label{lemma-glueing-increasing-union}

Soit $X$ un espace annelé. Soit $E$ un ensemble bien ordonné et soit
$$
X = \bigcup\nolimits_{\alpha \in E} W_\alpha
$$
être un recouvrement ouvert avec $W_\alpha \subset W_{\alpha + 1}$
et $W_\alpha = \bigcup_{\beta < \alpha} W_\beta$ si $\alpha$ n'est pas un successeur. Soit $K_\alpha$ un objet de $D(\mathcal{O}_{W_\alpha})$
tel que $\Ext^i(K_\alpha, K_\alpha) = 0$ pour $i < 0$. Supposons donnés des isomorphismes
$\rho_\beta^\alpha :  K_\alpha|_{W_\beta} \to K_\beta$ en
$D(\mathcal{O}_{W_\beta})$ pour tous $\beta < \alpha$ avec
$\rho_\gamma^\alpha = \rho_\gamma^\beta \circ \rho^\alpha_\beta|_{W_\gamma}$
pour $\gamma < \beta < \alpha$. Alors il existe un objet
$K$ en $D(\mathcal{O}_X)$ et les isomorphismes
$K|_{W_\alpha} \to K_\alpha$ pour $\alpha \in E$
compatible avec les isomorphismes $\rho_\beta^\alpha$.

\end{lemma}

\begin{proof}

Dans cette preuve, $\alpha, \beta, \gamma, \ldots$ désignent des éléments de $E$. Choisissons un complexe K-injectif
$I_\alpha^\bullet$ sur $W_\alpha$ représentant $K_\alpha$. Pour $\beta < \alpha$, notons $j_{\beta, \alpha} : W_\beta \to W_\alpha$
le morphisme d'inclusion. Par récurrence transfinie, nous allons construire, pour tout
$\beta < \alpha$, un morphisme de complexes
$$
\tau_{\beta, \alpha} :
(j_{\beta, \alpha})_!I_\beta^\bullet
\longrightarrow
I_\alpha^\bullet
$$
représentant l'adjoint de l'inverse de l'isomorphisme
$\rho^\alpha_\beta : K_\alpha|_{W_\beta} \to K_\beta$. Nous effectuerons de plus cette construction de sorte que, pour
$\gamma < \beta < \alpha$, nous ayons
$$
\tau_{\gamma, \alpha} = \tau_{\beta, \alpha} \circ
(j_{\beta, \alpha})_!\tau_{\gamma, \beta}
$$
comme morphismes de complexes. Supposons en effet les $\tau_{\gamma, \beta}$ déjà construits
et se composant correctement pour tous $\gamma < \beta < \alpha$. Si $\alpha = \alpha' + 1$ est un successeur, choisissons un morphisme de complexes
$$
(j_{\alpha', \alpha})_!I_{\alpha'}^\bullet \to I_\alpha^\bullet
$$
adjoint de l'inverse de l'isomorphisme
$\rho^\alpha_{\alpha'} : K_\alpha|_{W_{\alpha'}} \to K_{\alpha'}$
(ce qui est possible puisque $I_\alpha^\bullet$ est K-injectif) et, pour tout $\beta < \alpha'$, posons
$$
\tau_{\beta, \alpha} = \tau_{\alpha', \alpha} \circ
(j_{\alpha', \alpha})_!\tau_{\beta, \alpha'}
$$
Si $\alpha$ n'est pas un successeur, considérons sur $W_\alpha$ le complexe
$$
C^\bullet = \colim_{\beta < \alpha} (j_{\beta, \alpha})_!I_\beta^\bullet
$$
(colimite terme à terme), où les morphismes de transition sont les $\tau_{\beta', \beta}$ pour
$\beta' < \beta < \alpha$. Nous affirmons que $C^\bullet$
représente $K_\alpha$. En effet, pour $\beta < \alpha$, la restriction de la coprojection $(j_{\beta, \alpha})_!I_\beta^\bullet \to C^\bullet$
donne un morphisme
$$
\sigma_\beta : I_\beta^\bullet \longrightarrow C^\bullet|_{W_\beta}
$$
qui est un quasi-isomorphisme : si $x \in W_\beta$, le passage aux germes donne
$$
(C^\bullet)_x =
\colim_{\beta' < \alpha}
\left((j_{\beta', \alpha})_!I_{\beta'}^\bullet\right)_x =
\colim_{\beta \leq \beta' < \alpha} (I_{\beta'}^\bullet)_x
\longleftarrow
(I_\beta^\bullet)_x
$$
qui est un quasi-isomorphisme. Nous avons utilisé ici le fait que le passage aux germes commute aux colimites, que les colimites filtrantes sont exactes et que les morphismes $(I_\beta^\bullet)_x \to (I_{\beta'}^\bullet)_x$
sont des quasi-isomorphismes pour $\beta \leq \beta' < \alpha$. Ainsi, $(C^\bullet, \sigma_\beta^{-1})$ est une solution du système $(\{K_\beta\}_{\beta < \alpha},
\{\rho^\beta_{\beta'}\}_{\beta' < \beta < \alpha})$. Comme $(K_\alpha, \rho^\alpha_\beta)$ est une autre solution, nous obtenons un unique isomorphisme $\sigma : K_\alpha \to C^\bullet$
dans $D(\mathcal{O}_{W_\alpha})$, compatible avec tous nos morphismes ; voir le lemme \ref{lemma-solution-in-finite-case}
(c'est ici que nous utilisons l'annulation des groupes d'Ext négatifs). Choisissons un morphisme de complexes $\tau : C^\bullet \to I_\alpha^\bullet$
représentant $\sigma$. Posons alors
$$
\tau_{\beta, \alpha} = \tau|_{W_\beta} \circ \sigma_\beta
$$
pour obtenir les morphismes recherchés. Enfin, prenons pour $K$ l'objet de la catégorie dérivée représenté par le complexe
$$
K^\bullet = \colim_{\alpha \in E} (W_\alpha \to X)_!I_\alpha^\bullet
$$
où les morphismes de transition sont les morphismes $\tau_{\beta, \alpha}$ soigneusement construits pour $\beta < \alpha$. Le même raisonnement que ci-dessus montre que, pour tout $\alpha$,
la restriction de la coprojection détermine un isomorphisme
$$
K|_{W_\alpha} \longrightarrow K_\alpha
$$
compatible avec les morphismes donnés $\rho^\alpha_\beta$.

\end{proof}

\noindent
Une récurrence transfinie permet de démontrer le résultat dans le cas général.

\begin{theorem}
[Lemme de recollement de BBD]

\label{theorem-glueing-bbd-general}

\begin{reference}

Cas particulier de \cite[théorème 3.2.4]{BBD}, sans hypothèse de bornitude.

\end{reference}

Dans la situation \ref{situation-locally-given}, supposons que

\begin{enumerate}

\item  $X = \bigcup_{U \in \mathcal{B}} U$,
\item pour $U, V \in \mathcal{B}$, nous ayons
$U \cap V = \bigcup_{W \in \mathcal{B}, W \subset U \cap V} W$,
\item pour tout $U \in \mathcal{B}$, nous ayons $\Ext^i(K_U, K_U) = 0$
pour $i < 0$.

\end{enumerate}

Alors il existe un objet $K$ de $D(\mathcal{O}_X)$
et des isomorphismes $\rho_U : K|_U \to K_U$ dans $D(\mathcal{O}_U)$ pour
$U \in \mathcal{B}$, tels que $\rho^U_V \circ \rho_U|_V = \rho_V$
pour tous $V \subset U$ avec $U, V \in \mathcal{B}$. Le couple $(K, \rho_U)$ est unique à un unique isomorphisme près.

\end{theorem}

\begin{proof}
Un couple $(K, \rho_U)$ est appelé une solution dans le texte ci-dessus. L'unicité résulte du lemme \ref{lemma-uniqueness}. Si $X$ admet un recouvrement fini par des éléments de $\mathcal{B}$ (par exemple si $X$ est quasi-compact), le théorème résulte du lemme \ref{lemma-solution-in-finite-case}. Dans le cas général, nous raisonnons de la même manière, en utilisant une récurrence transfinie et le lemme \ref{lemma-glueing-increasing-union}.

\medskip\noindent

Utilisons d'abord la récurrence transfinie pour choisir un ouvert $W_\alpha \subset X$
pour tout ordinal $\alpha$. Posons $W_0 = \emptyset$. Si $\alpha = \beta + 1$ est un successeur, ou bien $W_\beta = X$
et nous posons $W_\alpha = X$, ou bien $W_\beta \not = X$ et nous posons
$W_\alpha = W_\beta \cup U_\alpha$ où
$U_\alpha \in \mathcal{B}$ n'est pas contenu dans $W_\beta$. Si $\alpha$ est un ordinal limite, posons
$W_\alpha = \bigcup_{\beta < \alpha} W_\beta$. Alors, pour $\alpha$ assez grand, on a $W_\alpha = X$. Pour tout $\alpha$, l'ouvert $W_\alpha$ est une réunion d'éléments de $\mathcal{B}$. Par conséquent, si
$\mathcal{B}_\alpha = \{U \in \mathcal{B}, U \subset W_\alpha\}$, alors
$$
S_\alpha = (\{K_U\}_{U \in \mathcal{B}_\alpha},
\{\rho_V^U\}_{V \subset U\text{ avec }U, V \in \mathcal{B}_\alpha})
$$
est un système comme dans le lemme \ref{lemma-uniqueness} sur l'espace annelé $W_\alpha$.

\medskip\noindent
Nous allons montrer par récurrence transfinie que, pour tout $\alpha$, le système $S_\alpha$ admet une solution. Cela démontrera le théorème, car pour $\alpha$ assez grand ce système est celui de l'énoncé.

\medskip\noindent
Supposons que $\alpha = \beta + 1$ soit un ordinal successeur. (Ce cas a déjà été traité dans la preuve du lemme précédent, mais nous répétons l'argument par souci de clarté.) Rappelons qu'alors $W_\alpha = W_\beta \cup U_\alpha$ pour un certain $U_\alpha \in \mathcal{B}$. Par hypothèse de récurrence, nous disposons d'une solution $(K_{W_\beta}, \{\rho^{W_\beta}_U\}_{U \in \mathcal{B}_\beta})$ du système $S_\beta$. Considérons la famille $\mathcal{B}_\alpha^* = \mathcal{B}_\alpha \cup \{W_\beta\}$ d'ouverts de $W_\alpha$ ; elle fournit un système $(\{K_U\}_{U \in \mathcal{B}_\alpha^*},
\{\rho_V^U\}_{V \subset U\text{ avec }U, V \in \mathcal{B}_\alpha^*})$. Ce nouveau système satisfait encore la condition (3), par le lemme \ref{lemma-uniqueness} appliqué à la solution $K_{W_\beta}$. Comme $W_\alpha = W_\beta \cup U_\alpha$, nous sommes ramenés au cas traité dans le lemme \ref{lemma-solution-in-finite-case}.

\medskip\noindent

Supposons que $\alpha$ soit un ordinal limite. Rappelons qu'alors
$W_\alpha = \bigcup_{\beta < \alpha} W_\beta$. Pour $\beta < \alpha$, soit
$(K_{W_\beta}, \{\rho^{W_\beta}_U\}_{U \in \mathcal{B}_\beta})$
la solution de $S_\beta$. Pour $\gamma < \beta < \alpha$, la restriction
$K_{W_\beta}|_{W_\gamma}$, munie des morphismes
$\rho^{W_\beta}_U$, $U \in \mathcal{B}_\gamma$
est une solution de $S_\gamma$. Par unicité, nous obtenons des isomorphismes uniques
$\rho_{W_\gamma}^{W_\beta} : K_{W_\beta}|_{W_\gamma} \to K_{W_\gamma}$
compatibles avec les morphismes $\rho^{W_\beta}_U$ et $\rho^{W_\gamma}_U$
pour $U \in \mathcal{B}_\gamma$. Ces morphismes se composent correctement, c'est-à-dire que $\rho_{W_\delta}^{W_\gamma} \circ \rho_{W_\gamma}^{W_\beta}|_{W_\delta}
= \rho^{W_\delta}_{W_\beta}$ pour $\delta < \gamma < \beta < \alpha$. Nous pouvons donc appliquer le lemme \ref{lemma-glueing-increasing-union}
(notons que les Ext négatifs de
$K_{W_\beta}$ sont nuls par le lemme \ref{lemma-uniqueness} appliqué à la solution $K_{W_\beta}$) afin d'obtenir $K_{W_\alpha}$ et des isomorphismes
$$
\rho_{W_\beta}^{W_\alpha} :
K_{W_\alpha}|_{W_\beta}
\longrightarrow
K_{W_\beta}
$$
compatibles avec les morphismes $\rho_{W_\gamma}^{W_\beta}$ pour
$\gamma < \beta < \alpha$.

\medskip\noindent

Pour montrer que $K_{W_\alpha}$ est une solution, il reste à construire des isomorphismes $\rho_U^{W_\alpha} : K_{W_\alpha}|_U \to K_U$ pour
$U \in \mathcal{B}_\alpha$, satisfaisant certaines compatibilités. Choisissons pour $\rho_U^{W_\alpha}$ l'unique morphisme tel que, pour $\beta < \alpha$ et tout $V \in \mathcal{B}_\beta$
avec $V \subset U$, le diagramme
$$
\xymatrix{
K_{W_\alpha}|_V \ar[r]_{\rho_U^{W_\alpha}|_V}
\ar[d]_{\rho_{W_\beta}^{W_\alpha}|_V}
& K_U|_V \ar[d]^{\rho_U^V} \\
K_{W_\beta} \ar[r]^{\rho_V^{W_\beta}}
& K_V
}
$$
soit commutatif. Ce choix est possible parce que
$$
(\{K_V\}_{V \subset U, V \in \mathcal{B}_\beta\text{ pour un certain }\beta < \alpha},
\{\rho_V^{V'}\}_{V \subset V'\text{ avec }V, V' \subset U
\text{ et }V, V' \in \mathcal{B}_\beta\text{ pour un certain }\beta < \alpha})
$$
est un système comme dans le lemme \ref{lemma-uniqueness} sur l'espace annelé $U$
et que $(K_U, \rho^U_V)$ et
$(K_{W_\alpha}|_U,  \rho_V^{W_\beta}\circ \rho_{W_\beta}^{W_\alpha}|_V)$
en sont deux solutions. Cela donne l'existence et l'unicité. Nous omettons de vérifier que ces morphismes satisfont les compatibilités voulues (il ne s'agit que d'une vérification formelle).

\end{proof}



\section{Complexes strictement parfaits}

\label{section-strictly-perfect}

\noindent
Les complexes strictement parfaits de modules serviront plus loin à définir les notions de complexe pseudo-cohérent et de complexe parfait. Les voici.

\begin{definition}

\label{definition-strictly-perfect}
Soit $(X, \mathcal{O}_X)$ un espace annelé et soit $\mathcal{E}^\bullet$ un complexe de $\mathcal{O}_X$-modules. On dit que $\mathcal{E}^\bullet$ est {\it strictement parfait} si $\mathcal{E}^i$ est nul sauf pour un nombre fini d'indices $i$ et si $\mathcal{E}^i$ est un facteur direct d'un $\mathcal{O}_X$-module libre de type fini pour tout $i$.
\end{definition}

\noindent
Attention : puisque nous ne supposons pas que $X$ est un espace localement annelé, un facteur direct d'un $\mathcal{O}_X$-module libre de type fini n'est pas nécessairement localement libre de type fini.

\begin{lemma}

\label{lemma-cone}

Le cône sur un morphisme de complexes strictement parfaits est strictement parfait.

\end{lemma}

\begin{proof}
Cela résulte immédiatement des définitions.
\end{proof}

\begin{lemma}

\label{lemma-tensor}

Le complexe total associé au produit tensoriel de deux complexes strictement parfaits est strictement parfait.

\end{lemma}

\begin{proof}
Démonstration omise.
\end{proof}

\begin{lemma}

\label{lemma-strictly-perfect-pullback}

Soit $f : (X, \mathcal{O}_X) \to (Y, \mathcal{O}_Y)$
un morphisme d'espaces annelés. Si $\mathcal{F}^\bullet$ est un complexe strictement parfait de $\mathcal{O}_Y$-modules, alors
$f^*\mathcal{F}^\bullet$ est un complexe strictement parfait de
$\mathcal{O}_X$-modules.

\end{lemma}

\begin{proof}

L'image réciproque d'un module libre de type fini est libre de type fini.
Le foncteur $f^*$ est additif, donc préserve les facteurs directs. Le lemme en résulte.

\end{proof}

\begin{lemma}

\label{lemma-local-lift-map}
Soit $(X, \mathcal{O}_X)$ un espace annelé. Étant donné un diagramme de $\mathcal{O}_X$-modules dont les flèches pleines sont fixées $$
\xymatrix{
\mathcal{E} \ar@{..>}[dr] \ar[r] & \mathcal{F} \\
& \mathcal{G} \ar[u]_p
}
$$ où $\mathcal{E}$ est un facteur direct d'un $\mathcal{O}_X$-module libre de type fini et où $p$ est surjectif, il existe localement sur $X$ une flèche pointillée rendant le diagramme commutatif.
\end{lemma}

\begin{proof}

Nous pouvons supposer $\mathcal{E} = \mathcal{O}_X^{\oplus n}$ pour un certain $n$. Trouver la flèche pointillée revient alors à relever dans $\Gamma(X, \mathcal{F})$ les images des éléments de base. C'est possible localement par la caractérisation des morphismes surjectifs de faisceaux (Faisceaux, section \ref{sheaves-section-exactness-points}).

\end{proof}

\begin{lemma}

\label{lemma-local-homotopy}

Soit $(X, \mathcal{O}_X)$ un espace annelé.

\begin{enumerate}

\item Soit $\alpha : \mathcal{E}^\bullet \to \mathcal{F}^\bullet$
un morphisme de complexes de $\mathcal{O}_X$-modules, avec $\mathcal{E}^\bullet$ strictement parfait et $\mathcal{F}^\bullet$
acyclique. Alors $\alpha$ est localement homotope à zéro sur $X$.
\item Soit $\alpha : \mathcal{E}^\bullet \to \mathcal{F}^\bullet$
être un morphisme de complexes de $\mathcal{O}_X$-modules avec $\mathcal{E}^\bullet$ strictement parfait, $\mathcal{E}^i = 0$
pour $i < a$ et $H^i(\mathcal{F}^\bullet) = 0$ pour $i \geq a$. Alors $\alpha$ est localement homotope à zéro sur $X$.

\end{enumerate}

\end{lemma}

\begin{proof}

La première assertion résulte de la seconde ; il suffit donc de démontrer (2). Procédons par récurrence sur la longueur du complexe
$\mathcal{E}^\bullet$. Si $\mathcal{E}^\bullet \cong \mathcal{E}[-n]$
pour un facteur direct $\mathcal{E}$ d'un
$\mathcal{O}_X$-module libre de type fini et un entier $n \geq a$, le résultat découle du lemme \ref{lemma-local-lift-map} et du fait que
$\mathcal{F}^{n - 1} \to \Ker(\mathcal{F}^n \to \mathcal{F}^{n + 1})$
est surjective par l'annulation supposée de $H^n(\mathcal{F}^\bullet)$. Si $\mathcal{E}^i$ est nul sauf pour $i \in [a, b]$, nous disposons de la suite exacte scindée de complexes
$$
0 \to \mathcal{E}^b[-b] \to \mathcal{E}^\bullet \to
\sigma_{\leq b - 1}\mathcal{E}^\bullet \to 0
$$
qui détermine un triangle distingué dans
$K(\mathcal{O}_X)$, et donc une suite exacte
$$
\Hom_{K(\mathcal{O}_X)}(
\sigma_{\leq b - 1}\mathcal{E}^\bullet, \mathcal{F}^\bullet)
\to
\Hom_{K(\mathcal{O}_X)}(\mathcal{E}^\bullet, \mathcal{F}^\bullet)
\to
\Hom_{K(\mathcal{O}_X)}(\mathcal{E}^b[-b], \mathcal{F}^\bullet)
$$
par les axiomes des catégories triangulées. La composée
$\mathcal{E}^b[-b] \to \mathcal{F}^\bullet$ est localement homotope à zéro ; nous pouvons donc supposer que notre morphisme provient d'un élément du membre de gauche de la suite exacte affichée ci-dessus. Cet élément est localement nul par l'hypothèse de récurrence.

\end{proof}

\begin{lemma}

\label{lemma-lift-through-quasi-isomorphism}
Soit $(X, \mathcal{O}_X)$ un espace annelé. Étant donné un diagramme de complexes de $\mathcal{O}_X$-modules dont les flèches pleines sont fixées $$
\xymatrix{
\mathcal{E}^\bullet \ar@{..>}[dr] \ar[r]_\alpha & \mathcal{F}^\bullet \\
& \mathcal{G}^\bullet \ar[u]_f
}
$$ où $\mathcal{E}^\bullet$ est strictement parfait, $\mathcal{E}^j = 0$ pour $j < a$, et où $H^j(f)$ est un isomorphisme pour $j > a$ et une surjection pour $j = a$, il existe localement sur $X$ une flèche pointillée rendant le diagramme commutatif à homotopie près.
\end{lemma}

\begin{proof}

Nos hypothèses sur $f$ impliquent que les faisceaux de cohomologie du cône $C(f)^\bullet$ sont nuls en degrés $\geq a$. Le lemme \ref{lemma-local-homotopy} garantit donc l'existence d'un recouvrement ouvert $X = \bigcup U_i$ tel que la composée
$\mathcal{E}^\bullet \to \mathcal{F}^\bullet \to C(f)^\bullet$
soit homotope à zéro sur $U_i$. Comme
$$
\mathcal{G}^\bullet \to \mathcal{F}^\bullet \to C(f)^\bullet \to
\mathcal{G}^\bullet[1]
$$
se restreint en un triangle distingué dans $K(\mathcal{O}_{U_i})$,
nous pouvons relever $\alpha|_{U_i}$, à homotopie près, en un morphisme
$\alpha_i : \mathcal{E}^\bullet|_{U_i} \to \mathcal{G}^\bullet|_{U_i}$
comme désiré.

\end{proof}

\begin{lemma}

\label{lemma-local-actual}

Soit $(X, \mathcal{O}_X)$ un espace annelé. Soient $\mathcal{E}^\bullet$, $\mathcal{F}^\bullet$ des complexes de $\mathcal{O}_X$-modules, avec $\mathcal{E}^\bullet$ strictement parfait.

\begin{enumerate}

\item Pour tout élément
$\alpha \in \Hom_{D(\mathcal{O}_X)}(\mathcal{E}^\bullet, \mathcal{F}^\bullet)$
il existe un recouvrement ouvert $X = \bigcup U_i$ tel que
$\alpha|_{U_i}$ soit représenté par un morphisme de complexes
$\alpha_i : \mathcal{E}^\bullet|_{U_i} \to \mathcal{F}^\bullet|_{U_i}$.
\item Étant donné un morphisme de complexes
$\alpha : \mathcal{E}^\bullet \to \mathcal{F}^\bullet$
dont l'image dans le groupe
$\Hom_{D(\mathcal{O}_X)}(\mathcal{E}^\bullet, \mathcal{F}^\bullet)$
est nulle, il existe un recouvrement ouvert $X = \bigcup U_i$ tel que
$\alpha|_{U_i}$ soit homotope à zéro.

\end{enumerate}

\end{lemma}

\begin{proof}
Démonstration de (1). Par construction de la catégorie dérivée, il existe un quasi-isomorphisme $f : \mathcal{F}^\bullet \to \mathcal{G}^\bullet$ et un morphisme de complexes $\beta : \mathcal{E}^\bullet \to \mathcal{G}^\bullet$ tels que $\alpha = f^{-1}\beta$. Le résultat découle donc du lemme \ref{lemma-lift-through-quasi-isomorphism}. Nous omettons la démonstration de (2).
\end{proof}

\begin{lemma}

\label{lemma-Rhom-strictly-perfect}
Soit $(X, \mathcal{O}_X)$ un espace annelé. Soient $\mathcal{E}^\bullet$, $\mathcal{F}^\bullet$ des complexes de $\mathcal{O}_X$-modules, avec $\mathcal{E}^\bullet$ strictement parfait. Alors le Hom interne $R\SheafHom(\mathcal{E}^\bullet, \mathcal{F}^\bullet)$ est représenté par le complexe $\mathcal{H}^\bullet$ de termes $$
\mathcal{H}^n =
\bigoplus\nolimits_{n = p + q}
\SheafHom_{\mathcal{O}_X}(\mathcal{E}^{-q}, \mathcal{F}^p)
$$ et de différentielle décrite dans la section \ref{section-hom-complexes}.
\end{lemma}

\begin{proof}

Choisissons un quasi-isomorphisme $\mathcal{F}^\bullet \to \mathcal{I}^\bullet$
vers un complexe K-injectif. Soit $(\mathcal{H}')^\bullet$ le complexe de termes
$$
(\mathcal{H}')^n =
\prod\nolimits_{n = p + q}
\SheafHom_{\mathcal{O}_X}(\mathcal{E}^{-q}, \mathcal{I}^p)
$$
qui représente $R\SheafHom(\mathcal{E}^\bullet, \mathcal{F}^\bullet)$
par la construction de la section \ref{section-internal-hom}. Il suffit de montrer que le morphisme
$$
\mathcal{H}^\bullet \longrightarrow (\mathcal{H}')^\bullet
$$
est un quasi-isomorphisme. Pour tout ouvert $U \subset X$, on constate que
$$
H^0(\mathcal{H}^\bullet(U)) =
\Hom_{K(\mathcal{O}_U)}(\mathcal{E}^\bullet|_U, \mathcal{I}^\bullet|_U)
\to
H^0((\mathcal{H}')^\bullet(U)) =
\Hom_{D(\mathcal{O}_U)}(\mathcal{E}^\bullet|_U, \mathcal{I}^\bullet|_U)
$$
Par le lemme \ref{lemma-local-actual}, le faisceau associé au préfaisceau
$U \mapsto H^0(\mathcal{H}^\bullet(U))$
est égal au faisceau associé au préfaisceau
$U \mapsto H^0((\mathcal{H}')^\bullet(U))$. Le même argument s'applique aux autres faisceaux de cohomologie. Ainsi, $\mathcal{H}^\bullet$
est quasi-isomorphe à $(\mathcal{H}')^\bullet$, ce qui démontre le lemme.

\end{proof}

\begin{lemma}

\label{lemma-Rhom-strictly-perfect-K-flat}
Dans la situation du lemme \ref{lemma-Rhom-strictly-perfect}, si $\mathcal{F}^\bullet$ est K-plat, alors $\mathcal{H}^\bullet$ est K-plat.
\end{lemma}

\begin{proof}

Notons que $\mathcal{H}^\bullet$ est simplement le complexe de Hom
$\SheafHom^\bullet(\mathcal{E}^\bullet, \mathcal{F}^\bullet)$
car la bornitude du complexe strictement parfait
$\mathcal{E}^\bullet$ assure que les produits intervenant dans la définition du complexe de Hom sont des sommes directes. Soit $\mathcal{K}^\bullet$ un complexe acyclique de
$\mathcal{O}_X$-modules. Considérons le morphisme
$$
\gamma :
\text{Tot}(\mathcal{K}^\bullet \otimes
\SheafHom^\bullet(\mathcal{E}^\bullet, \mathcal{F}^\bullet))
\longrightarrow
\SheafHom^\bullet(\mathcal{E}^\bullet,
\text{Tot}(\mathcal{K}^\bullet \otimes \mathcal{F}^\bullet))
$$
du lemme \ref{lemma-diagonal-better}. Comme $\mathcal{F}^\bullet$ est K-plat, le complexe $\text{Tot}(\mathcal{K}^\bullet \otimes \mathcal{F}^\bullet)$
est acyclique ; par le lemme \ref{lemma-local-actual}
(ou, si l'on préfère, le lemme \ref{lemma-Rhom-strictly-perfect}), le but de $\gamma$ est donc lui aussi acyclique. Pour démontrer le lemme, il suffit alors de montrer que $\gamma$ est un isomorphisme de complexes. Procédons par récurrence sur la longueur du complexe $\mathcal{E}^\bullet$. Si cette longueur est $\leq 1$, alors $\mathcal{E}^\bullet$
est un facteur direct de $\mathcal{O}_X^{\oplus n}[k]$ pour certains
$n \geq 0$ et $k \in \mathbf{Z}$, et le résultat est immédiat. Si la longueur est $> 1$, on la réduit en considérant la suite exacte courte de complexes, scindée terme à terme,
$$
0 \to \sigma_{\geq a + 1} \mathcal{E}^\bullet \to
\mathcal{E}^\bullet \to
\sigma_{\leq a} \mathcal{E}^\bullet \to 0
$$
pour un $a \in \mathbf{Z}$ convenable ; voir Homologie, section \ref{homology-section-truncations}. Alors $\gamma$ s'insère dans un morphisme de suites exactes courtes de complexes, scindées terme à terme. Par récurrence, $\gamma$ est un isomorphisme pour
$\sigma_{\geq a + 1} \mathcal{E}^\bullet$
et $\sigma_{\leq a} \mathcal{E}^\bullet$ ; le résultat en découle pour
$\mathcal{E}^\bullet$. Nous omettons quelques détails.

\end{proof}

\begin{lemma}

\label{lemma-Rhom-complex-of-direct-summands-finite-free}
Soit $(X, \mathcal{O}_X)$ un espace annelé. Soient $\mathcal{E}^\bullet$, $\mathcal{F}^\bullet$ des complexes de $\mathcal{O}_X$-modules tels que
\begin{enumerate}
\item $\mathcal{F}^n = 0$ pour $n \ll 0$,
\item $\mathcal{E}^n = 0$ pour $n \gg 0$, et
\item $\mathcal{E}^n$ est isomorphe à un facteur direct d'un $\mathcal{O}_X$-module libre de type fini.
\end{enumerate}
Alors le Hom interne $R\SheafHom(\mathcal{E}^\bullet, \mathcal{F}^\bullet)$ est représenté par le complexe $\mathcal{H}^\bullet$ de termes $$
\mathcal{H}^n =
\bigoplus\nolimits_{n = p + q}
\SheafHom_{\mathcal{O}_X}(\mathcal{E}^{-q}, \mathcal{F}^p)
$$ et de différentielle décrite dans la section \ref{section-internal-hom}.
\end{lemma}

\begin{proof}

Choisissons un quasi-isomorphisme $\mathcal{F}^\bullet \to \mathcal{I}^\bullet$
où $\mathcal{I}^\bullet$ est un complexe borné inférieurement d'objets injectifs. Notons que $\mathcal{I}^\bullet$ est K-injectif (Catégories dérivées, lemme
\ref{derived-lemma-bounded-below-injectives-K-injective}). La construction de la section \ref{section-internal-hom}
montre que
$R\SheafHom(\mathcal{E}^\bullet, \mathcal{F}^\bullet)$ est représenté par le complexe $(\mathcal{H}')^\bullet$ avec des termes
$$
(\mathcal{H}')^n =
\prod\nolimits_{n = p + q}
\SheafHom_{\mathcal{O}_X}(\mathcal{E}^{-q}, \mathcal{I}^p) =
\bigoplus\nolimits_{n = p + q}
\SheafHom_{\mathcal{O}_X}(\mathcal{E}^{-q}, \mathcal{I}^p)
$$
(l'égalité venant de ce qu'il n'y a qu'un nombre fini de termes non nuls). Notons que $\mathcal{H}^\bullet$ est le complexe total associé au bicomplexe de termes
$\SheafHom_{\mathcal{O}_X}(\mathcal{E}^{-q}, \mathcal{F}^p)$
et de même pour $(\mathcal{H}')^\bullet$. Le morphisme naturel $(\mathcal{H}')^\bullet \to \mathcal{H}^\bullet$ provient d'un morphisme de bicomplexes.
Pour montrer qu'il s'agit d'un quasi-isomorphisme, nous pouvons donc utiliser la suite spectrale d'un bicomplexe (Homologie, lemme \ref{homology-lemma-first-quadrant-ss})
$$
{}'E_1^{p, q} =
H^p(\SheafHom_{\mathcal{O}_X}(\mathcal{E}^{-q}, \mathcal{F}^\bullet))
$$
convergeant vers $H^{p + q}(\mathcal{H}^\bullet)$, et de même pour
$(\mathcal{H}')^\bullet$. Pour achever la démonstration, il suffit de montrer que $\mathcal{F}^\bullet \to \mathcal{I}^\bullet$
induit un isomorphisme
$$
H^p(\SheafHom_{\mathcal{O}_X}(\mathcal{E}, \mathcal{F}^\bullet))
\longrightarrow
H^p(\SheafHom_{\mathcal{O}_X}(\mathcal{E}, \mathcal{I}^\bullet))
$$
sur les faisceaux de cohomologie dès que $\mathcal{E}$ est un facteur direct d'un $\mathcal{O}_X$-module libre de type fini. Comme c'est clair lorsque $\mathcal{E}$
est libre de type fini, le résultat en découle.

\end{proof}





\section{Modules pseudo-cohérents}

\label{section-pseudo-coherent}

\noindent
Dans cette section, nous discutons des complexes pseudo-cohérents.

\begin{definition}

\label{definition-pseudo-coherent}

Soit $(X, \mathcal{O}_X)$ un espace annelé. Soit $\mathcal{E}^\bullet$
un complexe de $\mathcal{O}_X$-modules et soit $m \in \mathbf{Z}$.

\begin{enumerate}

\item On dit que $\mathcal{E}^\bullet$ est {\it $m$-pseudo-cohérent}
s'il existe un recouvrement ouvert $X = \bigcup U_i$ et, pour tout $i$,
un morphisme de complexes
$\alpha_i : \mathcal{E}_i^\bullet \to \mathcal{E}^\bullet|_{U_i}$
où $\mathcal{E}_i^\bullet$ est strictement parfait sur $U_i$ et
$H^j(\alpha_i)$ est un isomorphisme pour $j > m$ et $H^m(\alpha_i)$
est surjectif.
\item On dit que $\mathcal{E}^\bullet$ est {\it pseudo-cohérent}
s'il est $m$-pseudo-cohérent pour tout $m$.
\item On dit qu'un objet $E$ de $D(\mathcal{O}_X)$ est
{\it $m$-pseudo-cohérent} (resp.\ {\it pseudo-cohérent}) s'il peut être représenté par un complexe $m$-pseudo-cohérent (resp.\ pseudo-cohérent) de $\mathcal{O}_X$-modules.

\end{enumerate}

\end{definition}

\noindent
Si $X$ est quasi-compact, tout objet $m$-pseudo-cohérent de $D(\mathcal{O}_X)$ appartient à $D^-(\mathcal{O}_X)$. Ce n'est plus nécessairement vrai lorsque $X$ n'est pas quasi-compact.

\begin{lemma}

\label{lemma-pseudo-coherent-independent-representative}

Soit $(X, \mathcal{O}_X)$ un espace annelé et soit $E$ un objet de $D(\mathcal{O}_X)$.

\begin{enumerate}

\item S'il existe un recouvrement ouvert $X = \bigcup U_i$, des complexes strictement parfaits $\mathcal{E}_i^\bullet$ sur $U_i$ et des morphismes $\alpha_i : \mathcal{E}_i^\bullet \to E|_{U_i}$ dans
$D(\mathcal{O}_{U_i})$ avec $H^j(\alpha_i)$ un isomorphisme pour $j > m$
et $H^m(\alpha_i)$ surjectif, alors $E$ est $m$-pseudo-cohérent.
\item Si $E$ est $m$-pseudo-cohérent, alors tout complexe représentant
$E$ est $m$-pseudo-cohérent.

\end{enumerate}

\end{lemma}

\begin{proof}

Soit $\mathcal{F}^\bullet$ un complexe quelconque représentant $E$,
et soient $X = \bigcup U_i$ et
$\alpha_i : \mathcal{E}_i^\bullet \to E|_{U_i}$ comme dans (1). Montrons que $\mathcal{F}^\bullet$ est $m$-pseudo-cohérent comme complexe ; cela démontrera simultanément (1) et (2). Par le lemme \ref{lemma-local-actual},
après avoir raffiné le recouvrement ouvert $X = \bigcup U_i$, nous pouvons
représenter les morphismes $\alpha_i$ par des morphismes de complexes
$\alpha_i : \mathcal{E}_i^\bullet \to \mathcal{F}^\bullet|_{U_i}$. Par hypothèse
les $H^j(\alpha_i)$ sont des isomorphismes pour $j > m$, et $H^m(\alpha_i)$
est surjectif ; ainsi $\mathcal{F}^\bullet$ est
$m$-pseudo-cohérent.

\end{proof}

\begin{lemma}

\label{lemma-pseudo-coherent-pullback}
Soit $f : (X, \mathcal{O}_X) \to (Y, \mathcal{O}_Y)$ un morphisme d'espaces annelés. Soit $E$ un objet de $D(\mathcal{O}_Y)$. Si $E$ est $m$-pseudo-cohérent, alors $Lf^*E$ est $m$-pseudo-cohérent.
\end{lemma}

\begin{proof}

Représentons $E$ par un complexe $\mathcal{E}^\bullet$ de $\mathcal{O}_Y$-modules et choisissons un recouvrement ouvert $Y = \bigcup V_i$ ainsi que des morphismes
$\alpha_i : \mathcal{E}_i^\bullet \to \mathcal{E}^\bullet|_{V_i}$
comme dans la définition \ref{definition-pseudo-coherent}. Posons $U_i = f^{-1}(V_i)$. Par le lemme \ref{lemma-pseudo-coherent-independent-representative},
il suffit de montrer que $Lf^*\mathcal{E}^\bullet|_{U_i}$ est
$m$-pseudo-cohérent. Choisissons un triangle distingué
$$
\mathcal{E}_i^\bullet \to
\mathcal{E}^\bullet|_{V_i} \to
C \to
\mathcal{E}_i^\bullet[1]
$$
L'hypothèse sur $\alpha_i$ signifie exactement que les faisceaux de cohomologie
$H^j(C)$ sont nuls pour tout $j \geq m$. Notons $f_i : U_i \to V_i$
la restriction de $f$. Notons que $Lf^*\mathcal{E}^\bullet|_{U_i} = 
Lf_i^*(\mathcal{E}|_{V_i})$. En appliquant $Lf_i^*$, nous obtenons le triangle distingué
$$
Lf_i^*\mathcal{E}_i^\bullet \to
Lf_i^*\mathcal{E}|_{V_i} \to
Lf_i^*C \to
Lf_i^*\mathcal{E}_i^\bullet[1]
$$
La construction de $Lf_i^*$ comme foncteur dérivé à gauche montre que
$H^j(Lf_i^*C) = 0$ pour $j \geq m$ (par Catégories dérivées, lemme
\ref{derived-lemma-negative-vanishing}). Par conséquent, $H^j(Lf_i^*\alpha_i)$ est un isomorphisme pour $j > m$ et $H^m(Lf^*\alpha_i)$ est surjectif.
D'autre part, $Lf_i^*\mathcal{E}_i^\bullet = f_i^*\mathcal{E}_i^\bullet$, qui est strictement parfait par le lemme \ref{lemma-strictly-perfect-pullback}. Cela conclut la démonstration.

\end{proof}

\begin{lemma}

\label{lemma-cone-pseudo-coherent}

Soit $(X, \mathcal{O}_X)$ un espace annelé et soit $m \in \mathbf{Z}$. Soit $(K, L, M, f, g, h)$ un triangle distingué dans $D(\mathcal{O}_X)$.

\begin{enumerate}

\item Si $K$ est $(m + 1)$-pseudo-cohérent et $L$ est $m$-pseudo-cohérent, alors $M$ est $m$-pseudo-cohérent.
\item Si $K$ et $M$ sont $m$-pseudo-cohérents, alors $L$ est $m$-pseudo-cohérent.
\item Si $L$ est $(m + 1)$-pseudo-cohérent et $M$
est $m$-pseudo-cohérent, alors $K$ est $(m + 1)$-pseudo-cohérent.

\end{enumerate}

\end{lemma}

\begin{proof}

Démonstration de (1). Choisissons un recouvrement ouvert $X = \bigcup U_i$ et des morphismes $\alpha_i : \mathcal{K}_i^\bullet \to K|_{U_i}$ dans $D(\mathcal{O}_{U_i})$
avec $\mathcal{K}_i^\bullet$ strictement parfait et $H^j(\alpha_i)$
isomorphes pour $j > m + 1$ et surjectifs pour $j = m + 1$. Nous pouvons remplacer $\mathcal{K}_i^\bullet$ par
$\sigma_{\geq m + 1}\mathcal{K}_i^\bullet$
et donc nous pouvons supposer que $\mathcal{K}_i^j = 0$
pour $j < m + 1$. Après avoir raffiné le recouvrement ouvert, choisissons des morphismes $\beta_i : \mathcal{L}_i^\bullet \to L|_{U_i}$ dans $D(\mathcal{O}_{U_i})$
avec $\mathcal{L}_i^\bullet$ strictement parfait, tels que
$H^j(\beta)$ soit un isomorphisme pour $j > m$ et surjectif pour $j = m$. Par le lemme \ref{lemma-lift-through-quasi-isomorphism},
nous pouvons, après un nouveau raffinement, trouver des morphismes de complexes
$\gamma_i : \mathcal{K}^\bullet \to \mathcal{L}^\bullet$
tels que les diagrammes
$$
\xymatrix{
K|_{U_i} \ar[r] & L|_{U_i} \\
\mathcal{K}_i^\bullet \ar[u]^{\alpha_i} \ar[r]^{\gamma_i} &
\mathcal{L}_i^\bullet \ar[u]_{\beta_i}
}
$$
sont commutatifs dans $D(\mathcal{O}_{U_i})$ (il faut pour cela représenter les morphismes $\alpha_i$, $\beta_i$ et $K|_{U_i} \to L|_{U_i}$
par de véritables morphismes de complexes ; nous omettons quelques détails). Le cône $C(\gamma_i)^\bullet$ est strictement parfait (lemme \ref{lemma-cone}). La commutativité du diagramme implique l'existence d'un morphisme de triangles distingués
$$
(\mathcal{K}_i^\bullet, \mathcal{L}_i^\bullet, C(\gamma_i)^\bullet)
\longrightarrow
(K|_{U_i}, L|_{U_i}, M|_{U_i}).
$$
Il résulte du morphisme induit entre les suites exactes longues de cohomologie et de Homologie, lemmes \ref{homology-lemma-four-lemma} et
\ref{homology-lemma-five-lemma}
que $C(\gamma_i)^\bullet \to M|_{U_i}$ induit un isomorphisme en cohomologie dans les degrés $> m$ et une surjection en degré $m$. Ainsi, $M$ est $m$-pseudo-cohérent par le lemme \ref{lemma-pseudo-coherent-independent-representative}.

\medskip\noindent

Les assertions (2) et (3) suivent de (1) par rotation du triangle distingué.

\end{proof}

\begin{lemma}

\label{lemma-tensor-pseudo-coherent}

Soit $(X, \mathcal{O}_X)$ un espace annelé et soient $K, L$ des objets de $D(\mathcal{O}_X)$.

\begin{enumerate}

\item Si $K$ est $n$-pseudo-cohérent et $H^i(K) = 0$ pour $i > a$
et si $L$ est $m$-pseudo-cohérent et $H^j(L) = 0$ pour $j > b$, alors
$K \otimes_{\mathcal{O}_X}^\mathbf{L} L$ est $t$-pseudo-cohérent avec $t = \max(m + a, n + b)$.
\item Si $K$ et $L$ sont pseudo-cohérents, alors
$K \otimes_{\mathcal{O}_X}^\mathbf{L} L$ est pseudo-cohérent.

\end{enumerate}

\end{lemma}

\begin{proof}
Démonstration de (1). En remplaçant $X$ par les membres d'un recouvrement ouvert, nous pouvons supposer qu'il existe des complexes strictement parfaits $\mathcal{K}^\bullet$ et $\mathcal{L}^\bullet$ et des morphismes $\alpha : \mathcal{K}^\bullet \to K$ et $\beta : \mathcal{L}^\bullet \to L$, tels que $H^i(\alpha)$ soit un isomorphisme pour $i > n$ et surjectif pour $i = n$, et que $H^i(\beta)$ soit un isomorphisme pour $i > m$ et surjectif pour $i = m$. Alors le morphisme $$
\alpha \otimes^\mathbf{L} \beta :
\text{Tot}(\mathcal{K}^\bullet \otimes_{\mathcal{O}_X} \mathcal{L}^\bullet)
\to K \otimes_{\mathcal{O}_X}^\mathbf{L} L
$$ induit des isomorphismes sur les faisceaux de cohomologie en degré $i$ pour $i > t$ et une surjection pour $i = t$. Cela découle de la suite spectrale des Tor (détails omis).

\medskip\noindent

Démonstration de (2). En remplaçant d'abord $X$ par les membres d'un recouvrement ouvert, nous nous ramenons au cas où $K$ et $L$ sont bornés supérieurement. L'énoncé découle alors immédiatement de (1).

\end{proof}

\begin{lemma}

\label{lemma-summands-pseudo-coherent}
Soit $(X, \mathcal{O}_X)$ un espace annelé et soit $m \in \mathbf{Z}$. Si $K \oplus L$ est $m$-pseudo-cohérent (resp.\ pseudo-cohérent) dans $D(\mathcal{O}_X)$, alors $K$ et $L$ le sont aussi.
\end{lemma}

\begin{proof}

Supposons que $K \oplus L$ soit $m$-pseudo-cohérent. Après avoir remplacé $X$ par les membres d'un recouvrement ouvert, nous pouvons supposer que $K \oplus L \in D^-(\mathcal{O}_X)$, donc que
$L \in D^-(\mathcal{O}_X)$. Notons qu'il existe un triangle distingué
$$
(K \oplus L, K \oplus L, L \oplus L[1]) =
(K, K, 0) \oplus (L, L, L \oplus L[1])
$$
Voir Catégories dérivées, lemme \ref{derived-lemma-direct-sum-triangles}. Par le lemme \ref{lemma-cone-pseudo-coherent},
nous voyons que $L \oplus L[1]$ est $m$-pseudo-cohérent. Ainsi, $L[1] \oplus L[2]$ est lui aussi $m$-pseudo-cohérent. Par récurrence, $L[n] \oplus L[n + 1]$ est $m$-pseudo-cohérent. Comme $L$ est borné supérieurement, $L[n]$ est $m$-pseudo-cohérent pour $n$ assez grand. En remontant alors les triangles distingués
$$
(L[n], L[n] \oplus L[n - 1], L[n - 1])
$$
nous concluons que $L[n - 1], L[n - 2], \ldots, L$ sont $m$-pseudo-cohérents, comme souhaité.

\end{proof}

\begin{lemma}

\label{lemma-complex-pseudo-coherent-modules}
Soit $(X, \mathcal{O}_X)$ un espace annelé et soit $m \in \mathbf{Z}$. Soit $\mathcal{F}^\bullet$ un complexe de $\mathcal{O}_X$-modules (localement) borné supérieurement, tel que $\mathcal{F}^i$ soit $(m - i)$-pseudo-cohérent pour tout $i$. Alors $\mathcal{F}^\bullet$ est $m$-pseudo-cohérent.
\end{lemma}

\begin{proof}
Démonstration omise. Indication : utiliser le lemme \ref{lemma-cone-pseudo-coherent} et les troncatures comme dans la démonstration de Compléments sur l'algèbre, lemme \ref{more-algebra-lemma-complex-pseudo-coherent-modules}.
\end{proof}

\begin{lemma}

\label{lemma-cohomology-pseudo-coherent}
Soit $(X, \mathcal{O}_X)$ un espace annelé, soit $m \in \mathbf{Z}$ et soit $E$ un objet de $D(\mathcal{O}_X)$. Si $E$ est (localement) borné supérieurement et si $H^i(E)$ est $(m - i)$-pseudo-cohérent pour tout $i$, alors $E$ est $m$-pseudo-cohérent.
\end{lemma}

\begin{proof}
Démonstration omise. Indication : utiliser le lemme \ref{lemma-cone-pseudo-coherent} et les troncatures comme dans la démonstration de Compléments sur l'algèbre, lemme \ref{more-algebra-lemma-cohomology-pseudo-coherent}.
\end{proof}

\begin{lemma}

\label{lemma-finite-cohomology}

Soit $(X, \mathcal{O}_X)$ un espace annelé, soit $K$ un objet de $D(\mathcal{O}_X)$ et soit $m \in \mathbf{Z}$.

\begin{enumerate}

\item Si $K$ est $m$-pseudo-cohérent et $H^i(K) = 0$
pour $i > m$, alors $H^m(K)$ est un $\mathcal{O}_X$-module de type fini.
\item Si $K$ est $m$-pseudo-cohérent et $H^i(K) = 0$
pour $i > m + 1$, alors $H^{m + 1}(K)$ est un
$\mathcal{O}_X$-module de présentation finie.

\end{enumerate}

\end{lemma}

\begin{proof}

(1). Nous pouvons travailler localement sur $X$. Supposons donc qu'il existe un complexe strictement parfait $\mathcal{E}^\bullet$ et un morphisme
$\alpha : \mathcal{E}^\bullet \to K$ induisant un isomorphisme en cohomologie dans les degrés $> m$ et une surjection en degré $m$. Il suffit de démontrer le résultat pour $\mathcal{E}^\bullet$. Soit $n$ le plus grand entier tel que $\mathcal{E}^n \not = 0$. Si $n = m$, alors $H^m(\mathcal{E}^\bullet)$ est un quotient de
$\mathcal{E}^n$ et le résultat est clair. Si $n > m$, alors $\mathcal{E}^{n - 1} \to \mathcal{E}^n$ est surjectif puisque
$H^n(E^\bullet) = 0$. Par le lemme \ref{lemma-local-lift-map},
nous pouvons localement trouver une section de cette surjection et écrire
$\mathcal{E}^{n - 1} = \mathcal{E}' \oplus \mathcal{E}^n$. Il suffit donc de prouver le résultat pour le complexe
$(\mathcal{E}')^\bullet$ qui est le même que $\mathcal{E}^\bullet$
à ceci près qu'il contient $\mathcal{E}'$ en degré $n - 1$ et $0$ en degré $n$. On conclut par récurrence sur $n$.

\medskip\noindent
Démonstration de (2). Nous pouvons travailler localement sur $X$. Supposons donc qu'il existe un complexe strictement parfait $\mathcal{E}^\bullet$ et un morphisme $\alpha : \mathcal{E}^\bullet \to K$ induisant un isomorphisme en cohomologie dans les degrés $> m$ et une surjection en degré $m$. Comme dans la démonstration de (1), nous pouvons nous ramener au cas où $\mathcal{E}^i = 0$ pour $i > m + 1$. Alors $H^{m + 1}(K) \cong H^{m + 1}(\mathcal{E}^\bullet) =
\Coker(\mathcal{E}^m \to \mathcal{E}^{m + 1})$, qui est de présentation finie.
\end{proof}

\begin{lemma}

\label{lemma-n-pseudo-module}
Soit $(X, \mathcal{O}_X)$ un espace annelé et soit $\mathcal{F}$ un faisceau de $\mathcal{O}_X$-modules.
\begin{enumerate}
\item $\mathcal{F}$, considéré comme objet de $D(\mathcal{O}_X)$, est $0$-pseudo-cohérent si et seulement si $\mathcal{F}$ est un $\mathcal{O}_X$-module de type fini, et
\item $\mathcal{F}$, considéré comme objet de $D(\mathcal{O}_X)$, est $(-1)$-pseudo-cohérent si et seulement si $\mathcal{F}$ est un $\mathcal{O}_X$-module de présentation finie.
\end{enumerate}

\end{lemma}

\begin{proof}

Utiliser le lemme \ref{lemma-finite-cohomology}
pour démontrer les implications dans un sens, et le lemme \ref{lemma-cohomology-pseudo-coherent} pour la réciproque.

\end{proof}





\section{Dimension de Tor}

\label{section-tor}

\noindent
Dans cette section, nous examinons de plus près les résolutions par modules plats.

\begin{definition}

\label{definition-tor-amplitude}

Soit $(X, \mathcal{O}_X)$ un espace annelé. Soit $E$ un objet de $D(\mathcal{O}_X)$ et soient $a, b \in \mathbf{Z}$ avec $a \leq b$.

\begin{enumerate}

\item On dit que $E$ a une {\it amplitude de Tor dans $[a, b]$}
si $H^i(E \otimes_{\mathcal{O}_X}^\mathbf{L} \mathcal{F}) = 0$
pour tous $\mathcal{O}_X$-modules $\mathcal{F}$ et tous $i \not \in [a, b]$.
\item On dit que $E$ a une {\it dimension de Tor finie}
s'il a une amplitude de Tor dans $[a, b]$ pour certains $a, b$.
\item On dit que $E$ a {\it localement une dimension de Tor finie}
s'il existe un recouvrement ouvert $X = \bigcup U_i$ tel que
$E|_{U_i}$ ait une dimension de Tor finie pour tout $i$.

\end{enumerate}

Un $\mathcal{O}_X$-module $\mathcal{F}$ a une {\it dimension de Tor $\leq d$}
si $\mathcal{F}[0]$, considéré comme un objet de $D(\mathcal{O}_X)$, a une amplitude de Tor dans $[-d, 0]$.

\end{definition}

\noindent
Notons que si $E$, comme dans la définition, a une dimension de Tor finie, alors $E$ appartient à $D^b(\mathcal{O}_X)$, comme on le voit en prenant $\mathcal{F} = \mathcal{O}_X$ dans la définition ci-dessus.

\begin{lemma}

\label{lemma-last-one-flat}
Soit $(X, \mathcal{O}_X)$ un espace annelé. Soit $\mathcal{E}^\bullet$ un complexe borné supérieurement de $\mathcal{O}_X$-modules plats, d'amplitude de Tor contenue dans $[a, b]$. Alors $\Coker(d_{\mathcal{E}^\bullet}^{a - 1})$ est un $\mathcal{O}_X$-module plat.
\end{lemma}

\begin{proof}

Comme $\mathcal{E}^\bullet$ est un complexe borné supérieurement de modules plats, nous avons
$\mathcal{E}^\bullet \otimes_{\mathcal{O}_X} \mathcal{F} =
\mathcal{E}^\bullet \otimes_{\mathcal{O}_X}^{\mathbf{L}} \mathcal{F}$
pour tout $\mathcal{O}_X$-module $\mathcal{F}$. Par conséquent, pour tout $\mathcal{O}_X$-module $\mathcal{F}$, la suite
$$
\mathcal{E}^{a - 2} \otimes_{\mathcal{O}_X} \mathcal{F} \to
\mathcal{E}^{a - 1} \otimes_{\mathcal{O}_X} \mathcal{F} \to
\mathcal{E}^a \otimes_{\mathcal{O}_X} \mathcal{F}
$$
est exacte au milieu. Comme
$\mathcal{E}^{a - 2} \to \mathcal{E}^{a - 1} \to \mathcal{E}^a \to
\Coker(d^{a - 1}) \to 0$
est une résolution plate, il en résulte que
$\text{Tor}_1^{\mathcal{O}_X}(\Coker(d^{a - 1}), \mathcal{F}) = 0$
pour tout $\mathcal{O}_X$-module $\mathcal{F}$. Cela signifie que
$\Coker(d^{a - 1})$ est plat ; voir le lemme \ref{lemma-flat-tor-zero}.

\end{proof}

\begin{lemma}

\label{lemma-tor-amplitude}
Soit $(X, \mathcal{O}_X)$ un espace annelé, soit $E$ un objet de $D(\mathcal{O}_X)$ et soient $a, b \in \mathbf{Z}$ avec $a \leq b$. Les conditions suivantes sont équivalentes :
\begin{enumerate}
\item $E$ a une amplitude de Tor dans $[a, b]$.
\item $E$ est représenté par un complexe $\mathcal{E}^\bullet$ de $\mathcal{O}_X$-modules plats avec $\mathcal{E}^i = 0$ pour $i \not \in [a, b]$.
\end{enumerate}

\end{lemma}

\begin{proof}
Si (2) est vérifiée, nous pouvons calculer $E \otimes_{\mathcal{O}_X}^\mathbf{L} \mathcal{F} =
\mathcal{E}^\bullet \otimes_{\mathcal{O}_X} \mathcal{F}$, et (1) est immédiate.

\medskip\noindent

Supposons (1). Nous pouvons représenter $E$ par un complexe borné supérieurement de $\mathcal{O}_X$-modules plats $\mathcal{K}^\bullet$ ; voir la section \ref{section-flat}. Soit $n$ le plus grand entier tel que $\mathcal{K}^n \not = 0$. Si $n > b$, alors $\mathcal{K}^{n - 1} \to \mathcal{K}^n$ est surjectif puisque
$H^n(\mathcal{K}^\bullet) = 0$. Comme $\mathcal{K}^n$ est plat,
$\Ker(\mathcal{K}^{n - 1} \to \mathcal{K}^n)$ est plat (Modules, lemme \ref{modules-lemma-flat-ses}). Nous pouvons donc remplacer $\mathcal{K}^\bullet$ par
$\tau_{\leq n - 1}\mathcal{K}^\bullet$. Par récurrence sur $n$, nous nous ramenons ainsi au cas où $K^\bullet$ est un complexe de
$\mathcal{O}_X$-modules avec $\mathcal{K}^i = 0$ pour $i > b$.

\medskip\noindent

Posons $\mathcal{E}^\bullet = \tau_{\geq a}\mathcal{K}^\bullet$. Tout est clair, sauf la platitude de $\mathcal{E}^a$, qui résulte immédiatement du lemme \ref{lemma-last-one-flat}
et des définitions.

\end{proof}

\begin{lemma}

\label{lemma-tor-amplitude-pullback}
Soit $f : (X, \mathcal{O}_X) \to (Y, \mathcal{O}_Y)$ un morphisme d'espaces annelés et soit $E$ un objet de $D(\mathcal{O}_Y)$. Si $E$ a une amplitude de Tor dans $[a, b]$, alors $Lf^*E$ a une amplitude de Tor dans $[a, b]$.
\end{lemma}

\begin{proof}
Supposons que $E$ ait une amplitude de Tor dans $[a, b]$. Par le lemme \ref{lemma-tor-amplitude}, nous pouvons représenter $E$ par un complexe $\mathcal{E}^\bullet$ de $\mathcal{O}$-modules plats avec $\mathcal{E}^i = 0$ pour $i \not \in [a, b]$. Alors $Lf^*E$ est représenté par $f^*\mathcal{E}^\bullet$. Par Modules, lemme \ref{modules-lemma-base-change-flat}, les modules $f^*\mathcal{E}^i$ sont plats. Le lemme \ref{lemma-tor-amplitude} montre donc que $Lf^*E$ a une amplitude de Tor dans $[a, b]$.
\end{proof}

\begin{lemma}

\label{lemma-tor-amplitude-stalk}
Soit $(X, \mathcal{O}_X)$ un espace annelé, soit $E$ un objet de $D(\mathcal{O}_X)$ et soient $a, b \in \mathbf{Z}$ avec $a \leq b$. Les conditions suivantes sont équivalentes :
\begin{enumerate}
\item $E$ a une amplitude de Tor dans $[a, b]$.
\item pour tout $x \in X$, l'objet $E_x$ de $D(\mathcal{O}_{X, x})$ a une amplitude de Tor dans $[a, b]$.
\end{enumerate}

\end{lemma}

\begin{proof}
Prendre le germe en $x$ revient à prendre l'image réciproque par le morphisme d'espaces annelés $(x, \mathcal{O}_{X, x}) \to (X, \mathcal{O}_X)$. L'implication (1) $\Rightarrow$ (2) résulte donc du lemme \ref{lemma-tor-amplitude-pullback}. Réciproquement, le passage aux germes commute aux produits tensoriels (Modules, lemme \ref{modules-lemma-stalk-tensor-product}). Ainsi, $$
(E \otimes_{\mathcal{O}_X}^\mathbf{L} \mathcal{F})_x =
E_x \otimes_{\mathcal{O}_{X, x}}^\mathbf{L} \mathcal{F}_x
$$ D'autre part, le passage aux germes est exact, donc $$
H^i(E \otimes_{\mathcal{O}_X}^\mathbf{L} \mathcal{F})_x =
H^i((E \otimes_{\mathcal{O}_X}^\mathbf{L} \mathcal{F})_x) =
H^i(E_x \otimes_{\mathcal{O}_{X, x}}^\mathbf{L} \mathcal{F}_x)
$$ Or un faisceau $H^i(E \otimes_{\mathcal{O}_X}^\mathbf{L} \mathcal{F})$ est nul si et seulement si tous ses germes sont nuls (Modules, lemme \ref{modules-lemma-abelian}). Ainsi, (2) implique (1).
\end{proof}
\begin{lemma}

\label{lemma-cone-tor-amplitude}

Soit $(X, \mathcal{O}_X)$ un espace annelé. Soit $(K, L, M, f, g, h)$ un triangle distingué dans $D(\mathcal{O}_X)$ et soient $a, b \in \mathbf{Z}$.

\begin{enumerate}

\item Si $K$ a une amplitude de Tor dans $[a + 1, b + 1]$ et
$L$ une amplitude de Tor dans $[a, b]$, alors $M$ a une amplitude de Tor dans $[a, b]$.
\item Si $K$ et $M$ ont une amplitude de Tor dans $[a, b]$, alors
$L$ a une amplitude de Tor dans $[a, b]$.
\item Si $L$ a une amplitude de Tor dans $[a + 1, b + 1]$
et $M$ une amplitude de Tor dans $[a, b]$, alors
$K$ a une amplitude de Tor dans $[a + 1, b + 1]$.

\end{enumerate}

\end{lemma}

\begin{proof}

Démonstration omise. Indication : cela résulte de la suite exacte longue de cohomologie associée à un triangle distingué et du fait que
$- \otimes_{\mathcal{O}_X}^{\mathbf{L}} \mathcal{F}$
préserve les triangles distingués. L'assertion (2) est la plus facile à démontrer ; les autres s'en déduisent par translation.

\end{proof}

\begin{lemma}

\label{lemma-tensor-tor-amplitude}
Soit $(X, \mathcal{O}_X)$ un espace annelé et soient $K, L$ des objets de $D(\mathcal{O}_X)$. Si $K$ a une amplitude de Tor dans $[a, b]$ et $L$ une amplitude de Tor dans $[c, d]$, alors $K \otimes_{\mathcal{O}_X}^\mathbf{L} L$ a une amplitude de Tor dans $[a + c, b + d]$.
\end{lemma}

\begin{proof}

Démonstration omise. Indication : utiliser la suite spectrale des Tor.

\end{proof}

\begin{lemma}

\label{lemma-summands-tor-amplitude}
Soit $(X, \mathcal{O}_X)$ un espace annelé et soient $a, b \in \mathbf{Z}$. Pour des objets $K$, $L$ de $D(\mathcal{O}_X)$, si $K \oplus L$ a une amplitude de Tor dans $[a, b]$, alors $K$ et $L$ l'ont également.
\end{lemma}

\begin{proof}

Cela résulte immédiatement de l'additivité des foncteurs Tor.

\end{proof}





\section{Complexes parfaits}

\label{section-perfect}

\noindent
Dans cette section, nous discutons des propriétés des complexes parfaits sur les espaces annélés.

\begin{definition}

\label{definition-perfect}
Soit $(X, \mathcal{O}_X)$ un espace annelé et soit $\mathcal{E}^\bullet$ un complexe de $\mathcal{O}_X$-modules. On dit que $\mathcal{E}^\bullet$ est {\it parfait} s'il existe un recouvrement ouvert $X = \bigcup U_i$ tel que, pour tout $i$, il existe un quasi-isomorphisme de complexes $\mathcal{E}_i^\bullet \to \mathcal{E}^\bullet|_{U_i}$, où $\mathcal{E}_i^\bullet$ est un complexe strictement parfait de $\mathcal{O}_{U_i}$-modules. Un objet $E$ de $D(\mathcal{O}_X)$ est {\it parfait} s'il peut être représenté par un complexe parfait de $\mathcal{O}_X$-modules.
\end{definition}

\noindent
Si $X$ est quasi-compact, tout objet parfait de $D(\mathcal{O}_X)$ appartient à $D^b(\mathcal{O}_X)$. Ce n'est plus nécessairement vrai lorsque $X$ n'est pas quasi-compact.

\begin{lemma}

\label{lemma-perfect-independent-representative}
Soit $(X, \mathcal{O}_X)$ un espace annelé et soit $E$ un objet de $D(\mathcal{O}_X)$.
\begin{enumerate}
\item S'il existe un recouvrement ouvert $X = \bigcup U_i$ et des complexes strictement parfaits $\mathcal{E}_i^\bullet$ sur $U_i$ tels que $\mathcal{E}_i^\bullet$ représente $E|_{U_i}$ dans $D(\mathcal{O}_{U_i})$, alors $E$ est parfait.
\item Si $E$ est parfait, alors tout complexe représentant $E$ est parfait.
\end{enumerate}

\end{lemma}

\begin{proof}
La démonstration est identique à celle du lemme \ref{lemma-pseudo-coherent-independent-representative}.
\end{proof}

\begin{lemma}

\label{lemma-perfect-on-locally-ringed}

Soit $(X, \mathcal{O}_X)$ un espace annelé et soit $E$ un objet de
$D(\mathcal{O}_X)$. Supposons que tous les germes $\mathcal{O}_{X, x}$
soient des anneaux locaux. Alors les conditions suivantes sont équivalentes :

\begin{enumerate}

\item  $E$ est parfait,
\item il existe un recouvrement ouvert $X = \bigcup U_i$ tel que
$E|_{U_i}$ puisse être représenté par un complexe borné de $\mathcal{O}_{U_i}$-modules localement libres de type fini, et
\item il existe un recouvrement ouvert $X = \bigcup U_i$ tel que
$E|_{U_i}$ puisse être représenté par un complexe borné de $\mathcal{O}_{U_i}$-modules libres de type fini.

\end{enumerate}

\end{lemma}

\begin{proof}
Cela résulte du lemme \ref{lemma-perfect-independent-representative} et du fait que, sur $X$, tout facteur direct d'un module libre de type fini est localement libre de type fini. Voir Modules, lemme \ref{modules-lemma-direct-summand-of-locally-free-is-locally-free}.
\end{proof}

\begin{lemma}

\label{lemma-perfect-precise}
Soit $(X, \mathcal{O}_X)$ un espace annelé, soit $E$ un objet de $D(\mathcal{O}_X)$ et soient $a \leq b$ des entiers. Si $E$ a une amplitude de Tor dans $[a, b]$ et est $(a - 1)$-pseudo-cohérent, alors $E$ est parfait.
\end{lemma}

\begin{proof}

Après avoir remplacé $X$ par les membres d'un recouvrement ouvert, nous pouvons supposer qu'il existe un complexe strictement parfait $\mathcal{E}^\bullet$ et un morphisme
$\alpha : \mathcal{E}^\bullet \to E$ tel que $H^i(\alpha)$ soit un isomorphisme pour $i \geq a$. Remplaçons $\mathcal{E}^\bullet$ par
$\sigma_{\geq a - 1}\mathcal{E}^\bullet$ et choisissons un triangle distingué
$$
\mathcal{E}^\bullet \to E \to C \to \mathcal{E}^\bullet[1]
$$
L'annulation des faisceaux de cohomologie de $E$ et de $\mathcal{E}^\bullet$,
ainsi que l'hypothèse sur $\alpha$, donnent $C \cong \mathcal{K}[2 - a]$ avec
$\mathcal{K} = \Ker(\mathcal{E}^{a - 1} \to \mathcal{E}^a)$. Soit $\mathcal{F}$ un $\mathcal{O}_X$-module. En appliquant $- \otimes_{\mathcal{O}_X}^\mathbf{L} \mathcal{F}$,
l'hypothèse selon laquelle $E$ a une amplitude de Tor dans $[a, b]$
implique que $\mathcal{K} \otimes_{\mathcal{O}_X} \mathcal{F} \to
\mathcal{E}^{a - 1} \otimes_{\mathcal{O}_X} \mathcal{F}$ a l'image
$\Ker(\mathcal{E}^{a - 1} \otimes_{\mathcal{O}_X} \mathcal{F}
\to \mathcal{E}^a \otimes_{\mathcal{O}_X} \mathcal{F})$. Il s'ensuit que $\text{Tor}_1^{\mathcal{O}_X}(\mathcal{E}', \mathcal{F}) = 0$
où $\mathcal{E}' = \Coker(\mathcal{E}^{a - 1} \to \mathcal{E}^a)$. Ainsi, $\mathcal{E}'$ est plat (lemme \ref{lemma-flat-tor-zero}). Le module $\mathcal{E}'$ est donc localement facteur direct d'un module libre de type fini par Modules, lemme \ref{modules-lemma-flat-locally-finite-presentation}. Localement, le complexe
$$
\mathcal{E}' \to \mathcal{E}^{a + 1} \to \ldots \to \mathcal{E}^b
$$
est quasi-isomorphe à $E$ ; ainsi, $E$ est parfait.

\end{proof}

\begin{lemma}

\label{lemma-perfect}
Soit $(X, \mathcal{O}_X)$ un espace annelé et soit $E$ un objet de $D(\mathcal{O}_X)$. Les conditions suivantes sont équivalentes :
\begin{enumerate}
\item $E$ est parfait, et
\item $E$ est pseudo-cohérent et a localement une dimension finie de Tor.
\end{enumerate}

\end{lemma}

\begin{proof}
Supposons (1). Par définition, il existe un recouvrement ouvert $X = \bigcup U_i$ tel que $E|_{U_i}$ soit représenté par un complexe strictement parfait. Ainsi, $E$ est pseudo-cohérent (c'est-à-dire $m$-pseudo-cohérent pour tout $m$) par le lemme \ref{lemma-pseudo-coherent-independent-representative}. De plus, tout facteur direct d'un module libre de type fini est plat ; le lemme \ref{lemma-tor-amplitude} montre donc que $E|_{U_i}$ a une dimension de Tor finie. Ainsi, (2) est vérifiée.

\medskip\noindent
Supposons (2). Après avoir remplacé $X$ par les membres d'un recouvrement ouvert, nous pouvons supposer qu'il existe des entiers $a \leq b$ tels que $E$ ait une amplitude de Tor dans $[a, b]$. Puisque $E$ est $m$-pseudo-cohérent pour tout $m$, nous concluons par le lemme \ref{lemma-perfect-precise}.
\end{proof}

\begin{lemma}

\label{lemma-perfect-pullback}
Soit $f : (X, \mathcal{O}_X) \to (Y, \mathcal{O}_Y)$ un morphisme d'espaces annelés et soit $E$ un objet de $D(\mathcal{O}_Y)$. Si $E$ est parfait dans $D(\mathcal{O}_Y)$, alors $Lf^*E$ est parfait dans $D(\mathcal{O}_X)$.
\end{lemma}

\begin{proof}

Cela résulte des lemmes \ref{lemma-perfect},
\ref{lemma-tor-amplitude-pullback},
\ref{lemma-pseudo-coherent-pullback}. (Une autre démonstration consiste à reprendre celle du lemme \ref{lemma-pseudo-coherent-pullback}.)

\end{proof}

\begin{lemma}

\label{lemma-two-out-of-three-perfect}

Soit $(X, \mathcal{O}_X)$ un espace annelé et soit $(K, L, M, f, g, h)$
un triangle distingué dans $D(\mathcal{O}_X)$. Si deux des trois objets
$K, L, M$ sont parfaits, alors le troisième l'est également.

\end{lemma}

\begin{proof}

Première démonstration : combiner les lemmes \ref{lemma-perfect}, \ref{lemma-cone-pseudo-coherent} et
\ref{lemma-cone-tor-amplitude}. Deuxième démonstration (esquisse) : supposons $K$ et $L$ parfaits. Après avoir remplacé
$X$ par les membres d'un recouvrement ouvert, nous pouvons supposer que $K$ et $L$
sont représentés par des complexes strictement parfaits $\mathcal{K}^\bullet$
et $\mathcal{L}^\bullet$. Après un nouveau remplacement de $X$ par les membres d'un recouvrement ouvert, nous pouvons supposer que le morphisme $K \to L$ est représenté par un morphisme de complexes $\alpha : \mathcal{K}^\bullet \to \mathcal{L}^\bullet$ ; voir le lemme \ref{lemma-local-actual}. Alors $M$ est isomorphe au cône de $\alpha$, qui est strictement parfait par le lemme \ref{lemma-cone}.

\end{proof}

\begin{lemma}

\label{lemma-tensor-perfect}
Soit $(X, \mathcal{O}_X)$ un espace annelé. Si $K, L$ sont des objets parfaits de $D(\mathcal{O}_X)$, alors $K \otimes_{\mathcal{O}_X}^\mathbf{L} L$ l'est également.
\end{lemma}

\begin{proof}
Cela résulte des lemmes \ref{lemma-perfect}, \ref{lemma-tensor-pseudo-coherent} et \ref{lemma-tensor-tor-amplitude}.
\end{proof}

\begin{lemma}

\label{lemma-summands-perfect}
Soit $(X, \mathcal{O}_X)$ un espace annelé. Si $K \oplus L$ est un objet parfait de $D(\mathcal{O}_X)$, alors $K$ et $L$ le sont également.
\end{lemma}

\begin{proof}
Cela résulte des lemmes \ref{lemma-perfect}, \ref{lemma-summands-pseudo-coherent} et \ref{lemma-summands-tor-amplitude}.
\end{proof}

\begin{lemma}

\label{lemma-pushforward-perfect}
Soit $(X, \mathcal{O}_X)$ un espace annelé, soit $j : U \to X$ un sous-espace ouvert et soit $E$ un objet parfait de $D(\mathcal{O}_U)$ dont les faisceaux de cohomologie sont supportés sur une partie fermée $T \subset U$, avec $j(T)$ fermé dans $X$. Alors $Rj_*E$ est un objet parfait de $D(\mathcal{O}_X)$.
\end{lemma}

\begin{proof}
La perfection d'un complexe est une propriété locale sur $X$. Il suffit donc de vérifier que les restrictions de $Rj_*E$ à $U$ et à $V = X \setminus j(T)$ sont parfaites. Or $Rj_*E|_U = E$ est parfait, tandis que $Rj_*E|_V = 0$ puisque $E|_{U \setminus T} = 0$.
\end{proof}

\begin{lemma}

\label{lemma-perfect-max-open-coh-loc-free}
Soit $(X, \mathcal{O}_X)$ un espace annelé et soit $E$ un objet parfait de $D(\mathcal{O}_X)$. Supposons que tous les germes $\mathcal{O}_{X, x}$ soient des anneaux locaux. Alors l'ensemble $$
U =
\{x \in X \mid
H^i(E)_x\text{ est un }
\mathcal{O}_{X, x}\text{-module libre de type fini pour tout }i\in \mathbf{Z}\}
$$ est ouvert dans $X$ et constitue le plus grand ouvert $U \subset X$ tel que $H^i(E)|_U$ soit localement libre de type fini pour tout $i \in \mathbf{Z}$.
\end{lemma}

\begin{proof}

Notons que si $V \subset X$ est un ouvert tel que $H^i(E)|_V$
soit localement libre de type fini pour tout $i \in \mathbf{Z}$, alors $V \subset U$. Soit $x \in U$. Montrons qu'un voisinage ouvert de $x$
est contenu dans $U$ et que $H^i(E)$ y est localement libre de type fini pour tout $i$. Cela achèvera la démonstration. Au cours de celle-ci, nous pouvons remplacer un nombre fini de fois
$X$ par un voisinage ouvert de $x$. Nous pouvons donc supposer que $E$ est représenté par un complexe strictement parfait $\mathcal{E}^\bullet$. Supposons que
$\mathcal{E}^i = 0$ pour $i \not \in [a, b]$. Procédons par récurrence sur $b - a$. Le module
$H^b(E) = \Coker(d^{b - 1} : \mathcal{E}^{b - 1} \to \mathcal{E}^b)$
est de présentation finie. Comme $H^b(E)_x$ est libre de type fini, le module $H^b(E)$ est libre de type fini sur un voisinage ouvert de $x$, par Modules, lemme \ref{modules-lemma-finite-presentation-stalk-free}. Après avoir remplacé $X$ par un voisinage ouvert éventuellement plus petit, nous pouvons donc supposer que nous disposons d'une décomposition en somme directe
$\mathcal{E}^b = \Im(d^{b - 1}) \oplus H^b(E)$ et
$H^b(E)$ est libre de type fini ; voir le lemme \ref{lemma-local-lift-map}. En répétant le même argument, nous pouvons supposer que $\mathcal{E}^{b - 1} = \Ker(d^{b - 1}) \oplus \Im(d^{b - 1})$. Le complexe
$\mathcal{E}^a \to \ldots \to \mathcal{E}^{b - 2} \to \Ker(d^{b - 1})$
est strictement parfait et représente un objet parfait $E'$ tel que $H^i(E) = H^i(E')$ pour $i \not = b$. L'hypothèse de récurrence permet de conclure.

\end{proof}





\section{Duaux}

\label{section-duals}

\noindent
Dans cette section, nous caractérisons les objets dualisables de la catégorie des complexes et de la catégorie dérivée. En particulier, nous verrons qu'un objet de $D(\mathcal{O}_X)$ admet un dual si et seulement s'il est parfait (cela résulte de l'exemple \ref{example-dual-derived} et du lemme \ref{lemma-left-dual-derived}).

\begin{lemma}

\label{lemma-symmetric-monoidal-cat-complexes}

Soit $(X, \mathcal{O}_X)$ un espace annelé. La catégorie des complexes de $\mathcal{O}_X$-modules, munie du produit tensoriel défini par
$\mathcal{F}^\bullet \otimes \mathcal{G}^\bullet =
\text{Tot}(\mathcal{F}^\bullet \otimes_{\mathcal{O}_X} \mathcal{G}^\bullet)$
est une catégorie monoïdale symétrique (pour les règles de signes, voir Compléments sur l'algèbre, section \ref{more-algebra-section-sign-rules}).

\end{lemma}

\begin{proof}
Démonstration omise. Indications : comme unité $\mathbf{1}$, on prend le complexe égal à $\mathcal{O}_X$ en degré $0$ et nul dans les autres degrés, avec les isomorphismes évidents $\text{Tot}(\mathbf{1} \otimes_{\mathcal{O}_X} \mathcal{G}^\bullet) =
\mathcal{G}^\bullet$ et $\text{Tot}(\mathcal{F}^\bullet \otimes_{\mathcal{O}_X} \mathbf{1}) =
\mathcal{F}^\bullet$. Pour démontrer le lemme, il faut vérifier la commutativité de divers diagrammes ; voir Catégories, définitions \ref{categories-definition-monoidal-category} et \ref{categories-definition-symmetric-monoidal-category}. Dans chaque cas, les vérifications sont immédiates.
\end{proof}

\begin{example}

\label{example-dual}

Soit $(X, \mathcal{O}_X)$ un espace annelé. Soit $\mathcal{F}^\bullet$
un complexe localement borné de $\mathcal{O}_X$-modules tel que chaque
$\mathcal{F}^n$ soit localement un facteur direct d'un $\mathcal{O}_X$-module libre de type fini. Autrement dit, il existe un recouvrement ouvert
$X = \bigcup U_i$ tel que $\mathcal{F}^\bullet|_{U_i}$ soit un complexe strictement parfait. Considérons le complexe
$$
\mathcal{G}^\bullet = \SheafHom^\bullet(\mathcal{F}^\bullet, \mathcal{O}_X)
$$
comme dans la section \ref{section-hom-complexes}. Soient
$$
\eta :
\mathcal{O}_X
\to
\text{Tot}(\mathcal{F}^\bullet \otimes_{\mathcal{O}_X} \mathcal{G}^\bullet)
\quad\text{et}\quad
\epsilon :
\text{Tot}(\mathcal{G}^\bullet \otimes_{\mathcal{O}_X} \mathcal{F}^\bullet)
\to
\mathcal{O}_X
$$
les morphismes $\eta = \sum \eta_n$ et $\epsilon = \sum \epsilon_n$,
où $\eta_n : \mathcal{O}_X \to
\mathcal{F}^n \otimes_{\mathcal{O}_X} \mathcal{G}^{-n}$
et
$\epsilon_n : \mathcal{G}^{-n} \otimes_{\mathcal{O}_X} \mathcal{F}^n
\to \mathcal{O}_X$ sont comme dans Modules, exemple \ref{modules-example-dual}. Alors le triplet $\mathcal{G}^\bullet, \eta, \epsilon$
est un dual à gauche de $\mathcal{F}^\bullet$ au sens de Catégories, définition \ref{categories-definition-dual}. Nous omettons la vérification des égalités
$(1 \otimes \epsilon) \circ (\eta \otimes 1) = \text{id}_{\mathcal{F}^\bullet}$
et
$(\epsilon \otimes 1) \circ (1 \otimes \eta) =
\text{id}_{\mathcal{G}^\bullet}$. Comparer avec Compléments sur l'algèbre, lemme \ref{more-algebra-lemma-left-dual-complex}.

\end{example}

\begin{lemma}

\label{lemma-left-dual-complex}
Soit $(X, \mathcal{O}_X)$ un espace annelé et soit $\mathcal{F}^\bullet$ un complexe de $\mathcal{O}_X$-modules. Si $\mathcal{F}^\bullet$ admet un dual à gauche dans la catégorie monoïdale des complexes de $\mathcal{O}_X$-modules (Catégories, définition \ref{categories-definition-dual}), alors $\mathcal{F}^\bullet$ est un complexe localement borné dont les termes sont localement des facteurs directs de $\mathcal{O}_X$-modules libres de type fini, et son dual à gauche est celui construit dans l'exemple \ref{example-dual}.
\end{lemma}

\begin{proof}

Par l'unicité des duaux à gauche (Catégories, remarque \ref{categories-remark-left-dual-adjoint}), l'assertion finale résultera de ce que $\mathcal{F}^\bullet$
est bien de la forme annoncée. Soit le triplet $\mathcal{G}^\bullet, \eta, \epsilon$ un dual à gauche. Écrivons $\eta = \sum \eta_n$ et $\epsilon = \sum \epsilon_n$,
où $\eta_n : \mathcal{O}_X \to
\mathcal{F}^n \otimes_{\mathcal{O}_X} \mathcal{G}^{-n}$
et
$\epsilon_n : \mathcal{G}^{-n} \otimes_{\mathcal{O}_X} \mathcal{F}^n
\to \mathcal{O}_X$. Comme
$(1 \otimes \epsilon) \circ (\eta \otimes 1) = \text{id}_{\mathcal{F}^\bullet}$
et
$(\epsilon \otimes 1) \circ (1 \otimes \eta) = \text{id}_{\mathcal{G}^\bullet}$
par Catégories, définition \ref{categories-definition-dual}, nous obtenons immédiatement
$(1 \otimes \epsilon_n) \circ (\eta_n \otimes 1) = \text{id}_{\mathcal{F}^n}$
et
$(\epsilon_n \otimes 1) \circ (1 \otimes \eta_n) =
\text{id}_{\mathcal{G}^{-n}}$. Ainsi, $\mathcal{F}^n$ est localement un facteur direct d'un $\mathcal{O}_X$-module libre de type fini, par Modules, lemme \ref{modules-lemma-left-dual-module}. Comme la somme $\eta = \sum \eta_n$ est localement finie,
$\mathcal{F}^\bullet$ est localement borné.

\end{proof}

\begin{lemma}

\label{lemma-internal-hom-evaluate-isom}
Soit $(X, \mathcal{O}_X)$ un espace annelé et soient $K, L, M \in D(\mathcal{O}_X)$. Si $K$ est parfait, alors le morphisme $$
R\SheafHom(L, M) \otimes_{\mathcal{O}_X}^\mathbf{L} K
\longrightarrow
R\SheafHom(R\SheafHom(K, L), M)
$$ du lemme \ref{lemma-internal-hom-evaluate} est un isomorphisme.
\end{lemma}

\begin{proof}
Comme ce morphisme est défini globalement et que la formation de chacun de ses deux membres commute à la localisation (voir le lemme \ref{lemma-restriction-RHom-to-U}), nous pouvons travailler localement sur $X$. Nous pouvons donc supposer que $K$ est représenté par un complexe strictement parfait $\mathcal{E}^\bullet$.

\medskip\noindent
Si $K_1 \to K_2 \to K_3$ est un triangle distingué dans $D(\mathcal{O}_X)$, nous obtenons des triangles distingués $$
R\SheafHom(L, M) \otimes_{\mathcal{O}_X}^\mathbf{L} K_1 \to
R\SheafHom(L, M) \otimes_{\mathcal{O}_X}^\mathbf{L} K_2 \to
R\SheafHom(L, M) \otimes_{\mathcal{O}_X}^\mathbf{L} K_3
$$ et $$
R\SheafHom(R\SheafHom(K_1, L), M) \to
R\SheafHom(R\SheafHom(K_2, L), M) \to
R\SheafHom(R\SheafHom(K_3, L), M)
$$ Voir la section \ref{section-flat} et le lemme \ref{lemma-RHom-triangulated}. La flèche du lemme \ref{lemma-internal-hom-evaluate} est fonctorielle en $K$ ; nous obtenons donc un morphisme entre ces triangles distingués. Ainsi, si le résultat vaut pour $K_1$ et $K_3$, il vaut pour $K_2$ par Catégories dérivées, lemme \ref{derived-lemma-third-isomorphism-triangle}.

\medskip\noindent
En combinant les remarques ci-dessus avec les triangles distingués $$
\sigma_{\geq n}\mathcal{E}^\bullet \to \mathcal{E}^\bullet \to
\sigma_{\leq n - 1}\mathcal{E}^\bullet
$$ associés aux troncatures brutales, nous nous ramenons au cas où $K$ est un facteur direct d'un $\mathcal{O}_X$-module libre de type fini placé en un certain degré. La compatibilité évidente du problème avec les sommes directes (comme ci-dessus) et les décalages nous ramène au cas où $K = \mathcal{O}_X^{\oplus n}$ pour un certain entier $n$. Ce cas est clair.
\end{proof}

\begin{lemma}

\label{lemma-dual-perfect-complex}
Soit $(X, \mathcal{O}_X)$ un espace annelé et soit $K$ un objet parfait de $D(\mathcal{O}_X)$. Alors $K^\vee = R\SheafHom(K, \mathcal{O}_X)$ est également parfait et $(K^\vee)^\vee \cong K$. Il existe des isomorphismes fonctoriels $$
M \otimes^\mathbf{L}_{\mathcal{O}_X} K^\vee = R\SheafHom(K, M)
$$ et $$
H^0(X, M \otimes^\mathbf{L}_{\mathcal{O}_X} K^\vee) =
\Hom_{D(\mathcal{O}_X)}(K, M)
$$ pour tout $M$ dans $D(\mathcal{O}_X)$.
\end{lemma}

\begin{proof}
Par le lemme \ref{lemma-internal-hom-evaluate}, il existe une application canonique $$
K = R\SheafHom(\mathcal{O}_X, \mathcal{O}_X)
\otimes_{\mathcal{O}_X}^\mathbf{L} K \longrightarrow
R\SheafHom(R\SheafHom(K, \mathcal{O}_X), \mathcal{O}_X) =
(K^\vee)^\vee
$$ qui est un isomorphisme par le lemme \ref{lemma-internal-hom-evaluate-isom}. Pour vérifier les autres assertions, nous utiliserons sans autre mention le fait que la formation du Hom interne commute à la restriction aux ouverts (lemme \ref{lemma-restriction-RHom-to-U}). La perfection de $K^\vee$ peut se vérifier localement sur $X$. Par le lemme \ref{lemma-dual}, il suffit, pour l'assertion finale, de vérifier que le morphisme (\ref{equation-eval}) $$
M \otimes^\mathbf{L}_{\mathcal{O}_X} K^\vee
\longrightarrow
R\SheafHom(K, M)
$$ est un isomorphisme. Cette propriété est elle aussi locale sur $X$. Il suffit donc de démontrer ces deux assertions lorsque $K$ est représenté par un complexe strictement parfait.

\medskip\noindent

Supposons que $K$ soit représenté par le complexe strictement parfait
$\mathcal{E}^\bullet$. Il résulte alors du lemme \ref{lemma-Rhom-strictly-perfect}
que $K^\vee$ est représenté par le complexe dont les termes sont
$(\mathcal{E}^{-n})^\vee =
\SheafHom_{\mathcal{O}_X}(\mathcal{E}^{-n}, \mathcal{O}_X)$
en degré $n$. Comme $\mathcal{E}^{-n}$ est un facteur direct d'un $\mathcal{O}_X$-module libre de type fini, il en va de même de $(\mathcal{E}^{-n})^\vee$. Ainsi, $K^\vee$ est lui aussi représenté par un complexe strictement parfait, et nous voyons que $K^\vee$ est parfait. Pour voir que (\ref{equation-eval}) est un isomorphisme, représentons
$M$ par un complexe $\mathcal{F}^\bullet$. Par le lemme \ref{lemma-Rhom-strictly-perfect}, le complexe
$R\SheafHom(K, M)$ est représenté par le complexe avec des termes
$$
\bigoplus\nolimits_{n = p + q}
\SheafHom_{\mathcal{O}_X}(\mathcal{E}^{-q}, \mathcal{F}^p)
$$
D'autre part, l'objet $M \otimes^\mathbf{L}_{\mathcal{O}_X} K^\vee$
est représenté par le complexe avec des termes
$$
\bigoplus\nolimits_{n = p + q}
\mathcal{F}^p \otimes_{\mathcal{O}_X} (\mathcal{E}^{-q})^\vee
$$
L'assertion selon laquelle (\ref{equation-eval}) est un isomorphisme se ramène donc à celle selon laquelle le morphisme canonique
$$
\mathcal{F}
\otimes_{\mathcal{O}_X}
\SheafHom_{\mathcal{O}_X}(\mathcal{E}, \mathcal{O}_X)
\longrightarrow
\SheafHom_{\mathcal{O}_X}(\mathcal{E}, \mathcal{F})
$$
est un isomorphisme lorsque $\mathcal{E}$ est un facteur direct d'un $\mathcal{O}_X$-module libre de type fini et $\mathcal{F}$ un $\mathcal{O}_X$-module quelconque. Cela résulte immédiatement de l'énoncé correspondant lorsque
$\mathcal{E}$ est libre de type fini.

\end{proof}

\begin{lemma}

\label{lemma-symmetric-monoidal-derived}

Soit $(X, \mathcal{O}_X)$ un espace annelé. La catégorie dérivée
$D(\mathcal{O}_X)$ est une catégorie monoïdale symétrique dont le produit tensoriel est le produit tensoriel dérivé, muni des contraintes usuelles d'associativité et de commutativité (pour les règles de signes, voir Compléments sur l'algèbre, section \ref{more-algebra-section-sign-rules}).

\end{lemma}

\begin{proof}
Démonstration omise. Comparer avec le lemme \ref{lemma-symmetric-monoidal-cat-complexes}.
\end{proof}

\begin{example}

\label{example-dual-derived}

Soit $(X, \mathcal{O}_X)$ un espace annelé et soit $K$ un objet parfait de $D(\mathcal{O}_X)$. Posons $K^\vee = R\SheafHom(K, \mathcal{O}_X)$
comme dans le lemme \ref{lemma-dual-perfect-complex}. Alors le morphisme
$$
K \otimes_{\mathcal{O}_X}^\mathbf{L} K^\vee \longrightarrow R\SheafHom(K, K)
$$
est un isomorphisme (par ce lemme). Notons
$$
\eta :
\mathcal{O}_X
\longrightarrow
K \otimes_{\mathcal{O}_X}^\mathbf{L} K^\vee
$$
le morphisme envoyant $1$ sur la section correspondant à
$\text{id}_K$ sous l'isomorphisme ci-dessus.
$$
\epsilon : 
K^\vee
\otimes_{\mathcal{O}_X}^\mathbf{L} K
\longrightarrow
\mathcal{O}_X
$$
le morphisme d'évaluation (on peut par exemple le construire à l'aide du lemme \ref{lemma-internal-hom-composition}). Alors
le triplet $K^\vee, \eta, \epsilon$ est un dual à gauche de $K$ au sens de Catégories, définition \ref{categories-definition-dual}. Nous omettons de vérifier que
$(1 \otimes \epsilon) \circ (\eta \otimes 1) = \text{id}_K$
et
$(\epsilon \otimes 1) \circ (1 \otimes \eta) =
\text{id}_{K^\vee}$.

\end{example}

\begin{lemma}

\label{lemma-left-dual-derived}
Soit $(X, \mathcal{O}_X)$ un espace annelé et soit $M$ un objet de $D(\mathcal{O}_X)$. Si $M$ admet un dual à gauche dans la catégorie monoïdale $D(\mathcal{O}_X)$ (Catégories, définition \ref{categories-definition-dual}), alors $M$ est parfait et son dual à gauche est celui construit dans l'exemple \ref{example-dual-derived}.
\end{lemma}

\begin{proof}
Soit $x \in X$. Il suffit de trouver un voisinage ouvert $U$ de $x$ tel que la restriction de $M$ à $U$ soit un complexe parfait. Au cours de la démonstration, nous pouvons donc remplacer un nombre fini de fois $X$ par un voisinage ouvert de $x$. Soit le triplet $N, \eta, \epsilon$ un dual à gauche.

\medskip\noindent

Nous utiliserons plusieurs fois l'argument suivant. Choisissons un complexe $\mathcal{M}^\bullet$
de $\mathcal{O}_X$-modules représentant $M$, ainsi qu'un complexe K-plat
$\mathcal{N}^\bullet$ représentant $N$ et dont les termes sont des
$\mathcal{O}_X$-modules plats ; voir le lemme \ref{lemma-K-flat-resolution}. Considérons le morphisme
$$
\eta : \mathcal{O}_X \to
\text{Tot}(\mathcal{M}^\bullet \otimes_{\mathcal{O}_X} \mathcal{N}^\bullet)
$$
Après avoir rétréci $X$, on peut trouver un entier $N$ et, pour
$i = 1, \ldots, N$, des entiers $n_i \in \mathbf{Z}$ et des sections
$f_i$ et $g_i$ de $\mathcal{M}^{n_i}$ et $\mathcal{N}^{-n_i}$
tels que
$$
\eta(1) = \sum\nolimits_i f_i \otimes g_i
$$
Soit $\mathcal{K}^\bullet \subset \mathcal{M}^\bullet$ un sous-complexe quelconque de $\mathcal{O}_X$-modules contenant les sections $f_i$
pour $i = 1, \ldots, N$. Comme
$\text{Tot}(\mathcal{K}^\bullet \otimes_{\mathcal{O}_X} \mathcal{N}^\bullet)
\subset
\text{Tot}(\mathcal{M}^\bullet \otimes_{\mathcal{O}_X} \mathcal{N}^\bullet)$
par platitude des modules $\mathcal{N}^n$, le morphisme $\eta$ se factorise par
$$
\tilde \eta :
\mathcal{O}_X \to
\text{Tot}(\mathcal{K}^\bullet \otimes_{\mathcal{O}_X} \mathcal{N}^\bullet)
$$
En notant $K$ l'objet de $D(\mathcal{O}_X)$ représenté par
$\mathcal{K}^\bullet$, nous obtenons un diagramme commutatif
$$
\xymatrix{
M \ar[rr]_-{\eta \otimes 1} \ar[rrd]_{\tilde \eta \otimes 1} & &
M \otimes^\mathbf{L} N \otimes^\mathbf{L} M
\ar[r]_-{1 \otimes \epsilon} &
M \\
& &
K \otimes^\mathbf{L} N \otimes^\mathbf{L} M
\ar[u] \ar[r]^-{1 \otimes \epsilon} &
K \ar[u]
}
$$
Comme la composée de la ligne supérieure est l'identité de $M$,
nous concluons que $M$ est un facteur direct de $K$ dans $D(\mathcal{O}_X)$.

\medskip\noindent
Comme première application de cet argument, choisissons le sous-complexe $\mathcal{K}^\bullet = \sigma_{\geq a} \tau_{\leq b}\mathcal{M}^\bullet$ avec $a < n_i < b$ pour $i = 1, \ldots, N$. Ainsi, $M$ est un facteur direct, dans $D(\mathcal{O}_X)$, d'un complexe borné ; nous pouvons donc supposer que $M$ appartient à $D^b(\mathcal{O}_X)$. (Rappelons que la procédure ci-dessus impose de rétrécir $X$.)

\medskip\noindent
Puisque $M$ appartient à $D^b(\mathcal{O}_X)$, nous pouvons choisir pour $\mathcal{M}^\bullet$ un complexe borné supérieurement de modules plats (par Modules, lemme \ref{modules-lemma-module-quotient-flat}, et Catégories dérivées, lemme \ref{derived-lemma-subcategory-left-resolution}). Dans l'argument ci-dessus, nous pouvons alors prendre $\mathcal{K}^\bullet = \sigma_{\geq a}\mathcal{M}^\bullet$ avec $a < n_i$ pour $i = 1, \ldots, N$. Nous pouvons donc supposer que $M$ est un facteur direct, dans $D(\mathcal{O}_X)$, d'un complexe borné de modules plats. En particulier, $M$ a une amplitude de Tor finie.

\medskip\noindent

Supposons que $M$ ait une amplitude de Tor dans $[a, b]$. Si $M$ est $m$-pseudo-cohérent, montrons qu'après avoir rétréci $X$ nous pouvons supposer que $M$ est
$(m - 1)$-pseudo-cohérent. Le lemme \ref{lemma-perfect-precise} achèvera alors la démonstration, puisque
$M$ est de toute façon $(b + 1)$-pseudo-cohérent. Après avoir rétréci $X$, supposons qu'il existe un complexe strictement parfait $\mathcal{E}^\bullet$ et un morphisme $\alpha : \mathcal{E}^\bullet \to M$
dans $D(\mathcal{O}_X)$, tel que $H^i(\alpha)$ soit un isomorphisme pour
$i > m$ et surjectif pour $i = m$. Nous pouvons et allons supposer que $\mathcal{E}^i = 0$ pour $i < m$. Choisissons un triangle distingué
$$
\mathcal{E}^\bullet \to M \to L \to \mathcal{E}^\bullet[1]
$$
Notons que $H^i(L) = 0$ pour $i \geq m$. Nous pouvons donc représenter
$L$ par un complexe $\mathcal{L}^\bullet$ avec $\mathcal{L}^i = 0$
pour $i \geq m$. Le morphisme $L \to \mathcal{E}^\bullet[1]$
est représenté par un morphisme de complexes
$\mathcal{L}^\bullet \to \mathcal{E}^\bullet[1]$
qui est nul dans tous les degrés sauf en degré $m - 1$,
où nous obtenons un morphisme $\mathcal{L}^{m - 1} \to \mathcal{E}^m$ ; voir Catégories dérivées, lemme \ref{derived-lemma-negative-exts}. Alors $M$ est représenté par le complexe
$$
\mathcal{M}^\bullet :
\ldots \to
\mathcal{L}^{m - 2} \to
\mathcal{L}^{m - 1} \to
\mathcal{E}^m \to
\mathcal{E}^{m + 1} \to \ldots
$$
Appliquons à ce complexe la discussion du deuxième paragraphe afin d'obtenir des sections $f_i$ de $\mathcal{M}^{n_i}$ pour $i = 1, \ldots, N$. Pour $n < m$, soit $\mathcal{K}^n \subset \mathcal{L}^n$
le $\mathcal{O}_X$-sous-module engendré par les sections
$f_i$ telles que $n_i = n$ et par les $d(f_i)$ telles que $n_i = n - 1$. Pour $n \geq m$, posons $\mathcal{K}^n = \mathcal{E}^n$. Nous disposons manifestement d'un morphisme de triangles distingués
$$
\xymatrix{
\mathcal{E}^\bullet \ar[r] &
\mathcal{M}^\bullet \ar[r] &
\mathcal{L}^\bullet \ar[r] &
\mathcal{E}^\bullet[1] \\
\mathcal{E}^\bullet \ar[r] \ar[u] &
\mathcal{K}^\bullet \ar[r] \ar[u] &
\sigma_{\leq m - 1}\mathcal{K}^\bullet \ar[r] \ar[u] &
\mathcal{E}^\bullet[1] \ar[u]
}
$$
où tous les morphismes sont ceux indiqués ci-dessus. Notons $K$ l'objet de $D(\mathcal{O}_X)$ correspondant au complexe
$\mathcal{K}^\bullet$. Par les arguments du deuxième paragraphe de la démonstration, nous obtenons un morphisme $s : M \to K$ dans $D(\mathcal{O}_X)$ tel que la composée
$M \to K \to M$ est l'identité sur $M$. Nous ne savons pas que le diagramme
$$
\xymatrix{
\mathcal{E}^\bullet \ar[r] &
\mathcal{K}^\bullet \ar@{=}[r] &
K \\
\mathcal{E}^\bullet \ar[u]^{\text{id}} \ar[r]^i &
\mathcal{M}^\bullet \ar@{=}[r] &
M \ar[u]_s
}
$$
soit commutatif, mais nous savons qu'il le devient après composition avec le morphisme $K \to M$. Par le lemme \ref{lemma-local-actual}, après avoir rétréci $X$, nous pouvons supposer que $s \circ i$ est représenté par un morphisme de complexes
$\sigma : \mathcal{E}^\bullet \to \mathcal{K}^\bullet$. Par le même lemme, nous pouvons supposer que la composée de $\sigma$
avec l'inclusion $\mathcal{K}^\bullet \subset \mathcal{M}^\bullet$
est homotope à zéro par une homotopie
$\{h^i : \mathcal{E}^i \to \mathcal{M}^{i - 1}\}$. Ainsi, après avoir remplacé $\mathcal{K}^{m - 1}$ par
$\mathcal{K}^{m - 1} + \Im(h^m)$ (notons qu'après cette opération, $\mathcal{K}^{m - 1}$ reste engendré par un nombre fini de sections globales), nous voyons que
$\sigma$ lui-même est homotope à zéro ! Il existe donc un diagramme commutatif dont les flèches pleines sont fixées
$$
\xymatrix{
\mathcal{E}^\bullet \ar[r] &
M \ar[r] &
\mathcal{L}^\bullet \ar[r] &
\mathcal{E}^\bullet[1] \\
\mathcal{E}^\bullet \ar[r] \ar[u] &
K \ar[r] \ar[u] &
\sigma_{\leq m - 1}\mathcal{K}^\bullet \ar[r] \ar[u] &
\mathcal{E}^\bullet[1] \ar[u] \\
\mathcal{E}^\bullet \ar[r] \ar[u] &
M \ar[r] \ar[u]^s &
\mathcal{L}^\bullet \ar[r] \ar@{..>}[u] &
\mathcal{E}^\bullet[1] \ar[u]
}
$$
Les axiomes des catégories triangulées fournissent une flèche pointillée complétant le diagramme. En considérant les faisceaux de cohomologie en degré $m - 1$, nous obtenons
$$
\xymatrix{
H^{m - 1}(M) \ar[r] &
H^{m - 1}(\mathcal{L}^\bullet) \ar[r] &
H^m(\mathcal{E}^\bullet) \\
H^{m - 1}(K) \ar[r] \ar[u] &
H^{m - 1}(\sigma_{\leq m - 1}\mathcal{K}^\bullet) \ar[r] \ar[u] &
H^m(\mathcal{E}^\bullet) \ar[u] \\
H^{m - 1}(M) \ar[r] \ar[u] &
H^{m - 1}(\mathcal{L}^\bullet) \ar[r] \ar[u] &
H^m(\mathcal{E}^\bullet) \ar[u]
}
$$
Comme les composées verticales sont l'identité dans les colonnes de gauche et de droite, nous concluons que la composée verticale
$H^{m - 1}(\mathcal{L}^\bullet) \to
H^{m - 1}(\sigma_{\leq m - 1}\mathcal{K}^\bullet) \to
H^{m - 1}(\mathcal{L}^\bullet)$ au milieu est surjective ! En particulier, $H^{m - 1}(\sigma_{\leq m - 1}\mathcal{K}^\bullet) \to
H^{m - 1}(\mathcal{L}^\bullet)$ est surjectif. En utilisant le morphisme induit entre les suites exactes longues de faisceaux de cohomologie provenant du morphisme de triangles ci-dessus, une chasse au diagramme montre que $H^i(K) \to H^i(M)$ est un isomorphisme pour $i \geq m$ et surjectif pour $i = m - 1$. Par construction, nous pouvons choisir $r \geq 0$ et une surjection
$\mathcal{O}_X^{\oplus r} \to \mathcal{K}^{m - 1}$. Alors la composée
$$
(\mathcal{O}_X^{\oplus r} \to \mathcal{E}^m \to
\mathcal{E}^{m + 1} \to \ldots ) \longrightarrow
K \longrightarrow M
$$
induit un isomorphisme sur les faisceaux de cohomologie en degrés $\geq m$ et une surjection en degré $m - 1$, ce qui achève la démonstration.

\end{proof}





\section{Divers}

\label{section-misc}

\noindent
Quelques résultats qui ne trouvent pas leur place ailleurs.

\begin{lemma}

\label{lemma-colim-and-lim-of-duals}
Soit $(X, \mathcal{O}_X)$ un espace annelé. Soit $(K_n)_{n \in \mathbf{N}}$ un système d'objets parfaits de $D(\mathcal{O}_X)$. Soit $K = \text{hocolim} K_n$ la colimite dérivée (Catégories dérivées, définition \ref{derived-definition-derived-colimit}). Alors, pour tout objet $E$ de $D(\mathcal{O}_X)$, nous avons $$
R\SheafHom(K, E) = R\lim E \otimes^\mathbf{L}_{\mathcal{O}_X} K_n^\vee
$$ où $(K_n^\vee)$ est le système projectif des duaux des complexes parfaits.
\end{lemma}

\begin{proof}
Par le lemme \ref{lemma-dual-perfect-complex}, nous avons $R\lim E \otimes^\mathbf{L}_{\mathcal{O}_X} K_n^\vee =
R\lim R\SheafHom(K_n, E)$, qui s'insère dans le triangle distingué $$
R\lim R\SheafHom(K_n, E) \to
\prod R\SheafHom(K_n, E) \to
\prod R\SheafHom(K_n, E)
$$ Comme $K$ s'insère de même dans le triangle distingué $\bigoplus K_n \to \bigoplus K_n \to K$, il suffit de montrer que $\prod R\SheafHom(K_n, E) = R\SheafHom(\bigoplus K_n, E)$. Cela résulte formellement de (\ref{equation-internal-hom}) et du fait que le produit tensoriel dérivé commute aux sommes directes.
\end{proof}

\begin{lemma}

\label{lemma-ext-composition-is-cup}
Soit $(X, \mathcal{O}_X)$ un espace annelé. Soient $K$ et $E$ des objets de $D(\mathcal{O}_X)$, avec $E$ parfait. Le diagramme $$
\xymatrix{
H^0(X, K \otimes_{\mathcal{O}_X}^\mathbf{L} E^\vee) \times H^0(X, E)
\ar[r] \ar[d] &
H^0(X, K \otimes_{\mathcal{O}_X}^\mathbf{L} E^\vee
\otimes_{\mathcal{O}_X}^\mathbf{L} E) \ar[d] \\
\Hom_X(E, K) \times H^0(X, E) \ar[r] &
H^0(X, K)
}
$$ est commutatif ; la flèche horizontale supérieure est le cup-produit, la flèche verticale de droite utilise $\epsilon : E^\vee \otimes_{\mathcal{O}_X}^\mathbf{L} E \to \mathcal{O}_X$ (exemple \ref{example-dual-derived}), la flèche verticale de gauche utilise le lemme \ref{lemma-dual-perfect-complex}, et la flèche horizontale inférieure est la flèche évidente.
\end{lemma}

\begin{proof}
Nous abrégerons $\otimes = \otimes_{\mathcal{O}_X}^\mathbf{L}$ et $\mathcal{O} = \mathcal{O}_X$. Nous identifierons $E$ et $K$ à $R\SheafHom(\mathcal{O}, E)$ et $R\SheafHom(\mathcal{O}, K)$, et $E^\vee$ à $R\SheafHom(E, \mathcal{O})$.

\medskip\noindent
Soient $\xi \in H^0(X, K \otimes E^\vee)$ et $\eta \in H^0(X, E)$. Notons $\tilde \xi : \mathcal{O} \to K \otimes E^\vee$ et $\tilde \eta : \mathcal{O} \to E$ les morphismes correspondants dans $D(\mathcal{O})$. Par le lemme \ref{lemma-second-cup-equals-first}, le cup-produit $\xi \cup \eta$ correspond à $\tilde \xi \otimes \tilde \eta : \mathcal{O} \to
K \otimes E^\vee \otimes E$.

\medskip\noindent

Nous affirmons que le morphisme $\xi' : E \to K$ correspondant à $\xi$ par le lemme \ref{lemma-dual-perfect-complex} est la composée
$$
E = \mathcal{O} \otimes E
\xrightarrow{\tilde \xi \otimes 1_E}
K \otimes E^\vee \otimes E
\xrightarrow{1_K \otimes \epsilon}
K
$$
La construction du lemme \ref{lemma-dual-perfect-complex}
utilise le morphisme d'évaluation (\ref{equation-eval}), lui-même construit à l'aide de l'identification de $E$ à
$R\SheafHom(\mathcal{O}, E)$ et de la composition
$\underline{\circ}$ construite dans le lemme \ref{lemma-internal-hom-composition}. Ainsi, $\xi'$ est la composée

\begin{align*}
E = \mathcal{O} \otimes
R\SheafHom(\mathcal{O}, E)
& \xrightarrow{\tilde \xi \otimes 1}
R\SheafHom(\mathcal{O}, K) \otimes
R\SheafHom(E, \mathcal{O}) \otimes
R\SheafHom(\mathcal{O}, E) \\
& \xrightarrow{\underline{\circ} \otimes 1}
R\SheafHom(E, K) \otimes R\SheafHom(\mathcal{O}, E) \\
& \xrightarrow{\underline{\circ}}
R\SheafHom(\mathcal{O}, K) = K
\end{align*}

L'affirmation résulte immédiatement de ce qui précède, de l'associativité de la composition $\underline{\circ}$ construite dans le lemme \ref{lemma-internal-hom-composition} (insérer ici une référence future), et du fait que $\epsilon$
est définie comme la composition
$\underline{\circ} : E^\vee \otimes E \to \mathcal{O}$ dans l'exemple \ref{example-dual-derived}.

\medskip\noindent

Au vu des deux paragraphes précédents, l'énoncé du lemme revient à dire que
$(1_K \otimes \epsilon) \circ (\tilde \xi \otimes \tilde \eta)$
est égal à
$(1_K \otimes \epsilon) \circ (\tilde \xi \otimes 1_E)
\circ (1_\mathcal{O} \otimes \tilde \eta)$
ce qui est immédiat.

\end{proof}

\begin{lemma}

\label{lemma-pullback-internal-hom}

Soit $h : X \to Y$ un morphisme d'espaces annelés. Soient $K, M$ des objets de $D(\mathcal{O}_Y)$. Le morphisme canonique
$$
Lh^*R\SheafHom(K, M) \longrightarrow R\SheafHom(Lh^*K, Lh^*M)
$$
de la remarque \ref{remark-prepare-fancy-base-change}
est un isomorphisme dans les cas suivants :

\begin{enumerate}

\item  $K$ est parfait,
\item  $h$ est plat, $K$ est pseudo-cohérent, et $M$ est (localement) borné inférieurement,
\item  $\mathcal{O}_X$ a dimension finie de Tor sur $h^{-1}\mathcal{O}_Y$,
$K$ est pseudo-cohérent, et $M$ est (localement) borné inférieurement,

\end{enumerate}

\end{lemma}

\begin{proof}

(1). La question est locale sur $Y$ ; nous pouvons donc supposer que
$K$ est représenté par un complexe strictement parfait
$\mathcal{E}^\bullet$ ; voir la section \ref{section-perfect}. Choisissons un complexe K-plat $\mathcal{F}^\bullet$ représentant $M$. Le lemme \ref{lemma-Rhom-strictly-perfect} montre que
$R\SheafHom(K, L)$ est représenté par le complexe
$\mathcal{H}^\bullet =
\SheafHom^\bullet(\mathcal{E}^\bullet, \mathcal{F}^\bullet)$
avec des termes $\mathcal{H}^n = \bigoplus\nolimits_{n = p + q}
\SheafHom_{\mathcal{O}_X}(\mathcal{E}^{-q}, \mathcal{F}^p)$. Par la construction de $Lh^*$ dans la section \ref{section-derived-pullback},
nous voyons que $Lh^*K$ est représenté par le complexe strictement parfait
$h^*\mathcal{E}^\bullet$ (lemme \ref{lemma-strictly-perfect-pullback}). De même, l'objet $Lh^*M$ est représenté par le complexe $h^*\mathcal{F}^\bullet$. Enfin, l'objet $Lh^*R\SheafHom(K, M)$
est représenté par $h^*\mathcal{H}^\bullet$, puisque
$\mathcal{H}^\bullet$ est K-plat par le lemme \ref{lemma-Rhom-strictly-perfect-K-flat}. Pour achever la démonstration, il suffit donc de montrer que
$h^*\mathcal{H}^\bullet =
\SheafHom^\bullet(h^*\mathcal{E}^\bullet, h^*\mathcal{F}^\bullet)$. Pour cela, il suffit de noter que
$h^*\SheafHom(\mathcal{E}, \mathcal{F}) =
\SheafHom(h^*\mathcal{E}, \mathcal{F})$
dès que $\mathcal{E}$ est un facteur direct d'un
$\mathcal{O}_X$-module libre de type fini.

\medskip\noindent

(2). Puisque $h$ est plat, on peut calculer $Lh^*$
en appliquant simplement $h^*$ à tout complexe de $\mathcal{O}_Y$-modules. En particulier, $H^i(Lh^*K) = h^*H^i(K)$ pour tout $i \in \mathbf{Z}$. Supposons que $H^i(M) = 0$ pour $i < a$. Soit $K' \to K$ un morphisme de $D(\mathcal{O}_Y)$ induisant un isomorphisme
$H^i(K') \to H^i(K)$ pour tout $i \geq b$. Alors les morphismes correspondants
$$
R\SheafHom(K, M) \to R\SheafHom(K', M)
$$
et
$$
R\SheafHom(Lh^*K, Lh^*M) \to R\SheafHom(Lh^*K', Lh^*M)
$$
sont des isomorphismes sur les faisceaux de cohomologie en degrés $< a - b$ (détails omis). Ainsi, pour montrer que le morphisme de l'énoncé du lemme induit un isomorphisme sur les faisceaux de cohomologie en degrés $< a - b$, il suffit de démontrer le résultat pour $K'$ dans ces degrés. Comme dans la démonstration de (1), la question est locale sur $Y$. Nous pouvons donc supposer que $K$ est représenté par un complexe strictement parfait ; voir la section \ref{section-pseudo-coherent}. Nous sommes ainsi ramenés au cas (1).

\medskip\noindent

(3). La démonstration est la même que celle de (2), sauf que l'on utilise le fait que $Lh^*$ est de dimension cohomologique bornée afin d'obtenir l'annulation voulue. Nous omettons les détails.

\end{proof}

\begin{lemma}

\label{lemma-stalk-internal-hom}
Soit $X$ un espace annelé. Soient $K, M$ des objets de $D(\mathcal{O}_X)$. Soit $x \in X$. Le morphisme canonique $$
R\SheafHom(K, M)_x \longrightarrow
R\Hom_{\mathcal{O}_{X, x}}(K_x, M_x)
$$ est un isomorphisme dans les cas suivants
\begin{enumerate}
\item $K$ est parfait,
\item $K$ est pseudo-cohérent et $M$ est (localement) borné inférieurement.
\end{enumerate}

\end{lemma}

\begin{proof}
Soit $Y = \{x\}$ l'espace annelé réduit à un point, muni du faisceau structural donné par $\mathcal{O}_{X, x}$. On applique alors le lemme \ref{lemma-pullback-internal-hom} au morphisme d'inclusion plat $Y \to X$.
\end{proof}







\section{Objets inversibles dans la catégorie dérivée}
\label{section-invertible-D-or-R}

\noindent
Dans cette section, nous caractérisons les objets inversibles de la catégorie dérivée d'un espace annelé, aussi bien lorsque les germes du faisceau structural sont des anneaux locaux que lorsqu'ils ne le sont pas.

\begin{lemma}

\label{lemma-category-summands-finite-free}
Soit $(X, \mathcal{O}_X)$ un espace annelé. Posons $R = \Gamma(X, \mathcal{O}_X)$. La catégorie des $\mathcal{O}_X$-modules qui sont des facteurs directs de $\mathcal{O}_X$-modules libres de type fini est équivalente à la catégorie des $R$-modules projectifs de type fini.
\end{lemma}

\begin{proof}
Observons qu'un $R$-module projectif de type fini est la même chose qu'un facteur direct d'un $R$-module libre de type fini. L'équivalence est donnée par le foncteur $\mathcal{E} \mapsto
\Gamma(X, \mathcal{E})$. Le foncteur inverse est donné par la construction de Modules, lemme \ref{modules-lemma-construct-quasi-coherent-sheaves}.
\end{proof}

\begin{lemma}

\label{lemma-invertible-derived}

Soit $(X, \mathcal{O}_X)$ un espace annelé. Soit $M$ un objet de $D(\mathcal{O}_X)$. Les conditions suivantes sont équivalentes :

\begin{enumerate}

\item  $M$ est inversible dans $D(\mathcal{O}_X)$, voir Catégories, définition \ref{categories-definition-invertible},
\item il existe une décomposition localement finie en produit direct
$$
\mathcal{O}_X = \prod\nolimits_{n \in \mathbf{Z}} \mathcal{O}_n
$$
et, pour chaque $n$, il existe un $\mathcal{O}_n$-module inversible
$\mathcal{H}^n$ (Modules, définition \ref{modules-definition-invertible}) et $M = \bigoplus \mathcal{H}^n[-n]$ en $D(\mathcal{O}_X)$.

\end{enumerate}

Si (1) et (2) sont vérifiées, $M$ est un objet parfait de $D(\mathcal{O}_X)$. Si
$\mathcal{O}_{X, x}$ est un anneau local pour tout $x \in X$, ces conditions sont également équivalentes à :

\begin{enumerate}

\item[(3)] il existe un recouvrement ouvert $X = \bigcup U_i$
et, pour chaque $i$, un entier $n_i$ tel que $M|_{U_i}$
soit représenté par un $\mathcal{O}_{U_i}$-module inversible placé en degré $n_i$.

\end{enumerate}

\end{lemma}

\begin{proof}

Supposons (2). Considérons l'objet $R\SheafHom(M, \mathcal{O}_X)$
et le morphisme de composition
$$
R\SheafHom(M, \mathcal{O}_X) \otimes_{\mathcal{O}_X}^\mathbf{L} M \to
\mathcal{O}_X
$$
Pour montrer qu'il s'agit d'un isomorphisme, nous pouvons travailler localement. Nous pouvons donc supposer que $\mathcal{O}_X = \prod_{a \leq n \leq b} \mathcal{O}_n$
et $M = \bigoplus_{a \leq n \leq b} \mathcal{H}^n[-n]$. Alors il suffit de montrer que
$$
R\SheafHom(\mathcal{H}^m, \mathcal{O}_X)
\otimes_{\mathcal{O}_X}^\mathbf{L} \mathcal{H}^n
$$
est nul si $n \not = m$ et égal à $\mathcal{O}_n$ si $n = m$. Le cas $n \not = m$ résulte du fait que $\mathcal{O}_n$ et
$\mathcal{O}_m$ sont des $\mathcal{O}_X$-algèbres plates telles que
$\mathcal{O}_n \otimes_{\mathcal{O}_X} \mathcal{O}_m = 0$. En utilisant la structure locale des $\mathcal{O}_X$-modules inversibles (Modules, lemme \ref{modules-lemma-invertible}) et en travaillant localement, l'isomorphisme dans le cas $n = m$ s'obtient immédiatement ; nous omettons les détails. Puisque $D(\mathcal{O}_X)$ est monoïdale symétrique, nous concluons que $M$ est inversible.

\medskip\noindent
Supposons (1). La description donnée en (2) fournit un candidat pour $\mathcal{O}_n$, à savoir $\SheafHom_{\mathcal{O}_X}(H^n(M), H^n(M))$. Si ces faisceaux d'anneaux forment une famille localement finie et si $\mathcal{O}_X = \prod \mathcal{O}_n$, nous obtenons immédiatement la décomposition en somme directe $M = \bigoplus H^n(M)[-n]$ à l'aide des idempotents de $\mathcal{O}_X$ issus de la décomposition en produit. Cela montre que, pour démontrer (2), nous pouvons travailler localement sur $X$.

\medskip\noindent

Choisissons un objet $N$ de $D(\mathcal{O}_X)$
et un isomorphisme
$M \otimes_{\mathcal{O}_X}^\mathbf{L} N \cong \mathcal{O}_X$. Soit $x \in X$. Alors $N$ est un dual à gauche de $M$ dans la catégorie monoïdale
$D(\mathcal{O}_X)$, et nous concluons que $M$ est parfait par le lemme \ref{lemma-left-dual-derived}. Par symétrie,
$N$ est parfait. Après avoir remplacé $X$ par un voisinage ouvert de $x$, nous pouvons supposer que $M$ et $N$ sont représentés par des complexes strictement parfaits $\mathcal{E}^\bullet$ et $\mathcal{F}^\bullet$. Alors $M \otimes_{\mathcal{O}_X}^\mathbf{L} N$ est représenté par
$\text{Tot}(\mathcal{E}^\bullet \otimes_{\mathcal{O}_X} \mathcal{F}^\bullet)$. Après avoir de nouveau rétréci $X$, nous pouvons supposer que les isomorphismes réciproques
$\mathcal{O}_X \to M \otimes_{\mathcal{O}_X}^\mathbf{L} N$ et
$M \otimes_{\mathcal{O}_X}^\mathbf{L} N \to \mathcal{O}_X$
sont donnés par des morphismes de complexes
$$
\alpha : \mathcal{O}_X \to
\text{Tot}(\mathcal{E}^\bullet \otimes_{\mathcal{O}_X} \mathcal{F}^\bullet)
\quad\text{et}\quad
\beta :
\text{Tot}(\mathcal{E}^\bullet \otimes_{\mathcal{O}_X} \mathcal{F}^\bullet)
\to \mathcal{O}_X
$$
Voir le lemme \ref{lemma-local-actual}. Alors $\beta \circ \alpha = 1$
comme morphismes de complexes et $\alpha \circ \beta = 1$ comme morphisme dans $D(\mathcal{O}_X)$. Après avoir rétréci $X$,
nous pouvons supposer que la composée $\alpha \circ \beta$ est homotope à $1$
par une homotopie $\theta$ de composantes
$$
\theta^n :
\text{Tot}^n(\mathcal{E}^\bullet \otimes_{\mathcal{O}_X} \mathcal{F}^\bullet)
\to
\text{Tot}^{n - 1}(
\mathcal{E}^\bullet \otimes_{\mathcal{O}_X} \mathcal{F}^\bullet)
$$
par le même lemme que précédemment. Posons $R = \Gamma(X, \mathcal{O}_X)$. Par le lemme \ref{lemma-category-summands-finite-free},
nous obtenons :

\begin{enumerate}

\item  $M^\bullet = \Gamma(X, \mathcal{E}^\bullet)$ est un complexe borné de $R$-modules projectifs de type fini,
\item  $N^\bullet = \Gamma(X, \mathcal{F}^\bullet)$ est un complexe borné de $R$-modules projectifs de type fini,
\item  $\alpha$ et $\beta$ correspondent aux morphismes de complexes
$a : R \to \text{Tot}(M^\bullet \otimes_R N^\bullet)$ et
$b : \text{Tot}(M^\bullet \otimes_R N^\bullet) \to R$,
\item  $\theta^n$ correspond à un morphisme
$h^n : \text{Tot}^n(M^\bullet \otimes_R N^\bullet) \to
\text{Tot}^{n - 1}(M^\bullet \otimes_R N^\bullet)$,
\item  $b \circ a = 1$ et $b \circ a - 1 = dh + hd$,

\end{enumerate}

Il en résulte que $M^\bullet$ et $N^\bullet$ définissent des objets réciproquement inverses de $D(R)$. Par Compléments sur l'algèbre, lemme \ref{more-algebra-lemma-invertible-derived},
nous obtenons une décomposition en produit $R = \prod_{a \leq n \leq b} R_n$
et des $R_n$-modules inversibles $H^n$ tels que $M^\bullet \cong \bigoplus_{a \leq n \leq b} H^n[-n]$. Cet isomorphisme dans $D(R)$ se relève en un morphisme
$$
\bigoplus H^n[-n] \longrightarrow M^\bullet
$$
de complexes, car chaque $H^n$ est projectif comme $R$-module. De même, en appliquant une nouvelle fois le lemme \ref{lemma-category-summands-finite-free}, nous obtenons un morphisme
$$
\bigoplus H^n \otimes_R \mathcal{O}_X[-n] \to \mathcal{E}^\bullet
$$
qui est un isomorphisme dans $D(\mathcal{O}_X)$. En posant
$\mathcal{O}_n = R_n \otimes_R \mathcal{O}_X$, nous concluons que (2) est vérifiée.

\medskip\noindent

Si tous les germes de $\mathcal{O}_X$ sont locaux, l'équivalence de (2) et (3) se démontre immédiatement. Nous omettons les détails.

\end{proof}





\section{Objets compacts}

\label{section-compact}

\noindent
Dans cette section, nous étudions les objets compacts de la catégorie dérivée des modules sur un espace annelé. Rappelons que les objets compacts sont définis dans Catégories dérivées, définition \ref{derived-definition-compact-object}. Sur les espaces annelés convenables, les objets parfaits sont compacts.

\begin{lemma}

\label{lemma-when-jshriek-compact}
Soit $X$ un espace annelé. Soit $j : U \to X$ l'inclusion d'un ouvert. Le $\mathcal{O}_X$-module $j_!\mathcal{O}_U$ est un objet compact de $D(\mathcal{O}_X)$ s'il existe un entier $d$ tel que
\begin{enumerate}
\item $H^p(U, \mathcal{F}) = 0$ pour tous les $p > d$, et
\item les foncteurs $\mathcal{F} \mapsto H^p(U, \mathcal{F})$
commutent avec les sommes directes.
\end{enumerate}

\end{lemma}

\begin{proof}

Supposons (1) et (2). Comme
$\Hom(j_!\mathcal{O}_U, \mathcal{F}) = \mathcal{F}(U)$
par Faisceaux, lemme \ref{sheaves-lemma-j-shriek-modules},
nous avons $\Hom(j_!\mathcal{O}_U, K) = R\Gamma(U, K)$ pour
$K$ dans $D(\mathcal{O}_X)$. Il suffit donc de montrer que $R\Gamma(U, -)$ commute aux sommes directes.
La première hypothèse signifie que le foncteur
$F = H^0(U, -)$ est de dimension cohomologique finie. En outre, la seconde implique que toute somme directe de modules injectifs est $F$-acyclique. Soit la famille $K_i$ d'objets de $D(\mathcal{O}_X)$. Choisissons des représentants K-injectifs $I_i^\bullet$ à termes injectifs pour les $K_i$ ; voir Injectifs, théorème
\ref{injectives-theorem-K-injective-embedding-grothendieck}. Puisque nous pouvons calculer $RF$ en appliquant $F$ à tout complexe d'objets acycliques (Catégories dérivées, lemme \ref{derived-lemma-unbounded-right-derived}) et que $\bigoplus K_i$ est représenté par $\bigoplus I_i^\bullet$
(Injectifs, lemme \ref{injectives-lemma-derived-products}), nous concluons que $R\Gamma(U, \bigoplus K_i)$ est représenté par
$\bigoplus H^0(U, I_i^\bullet)$. Ainsi, $R\Gamma(U, -)$ commute aux sommes directes, comme souhaité.

\end{proof}

\begin{lemma}

\label{lemma-perfect-is-compact}

Soit $X$ un espace annelé. Supposons que l'espace topologique sous-jacent à $X$ possède les propriétés suivantes :

\begin{enumerate}

\item  $X$ est quasi-compact,
\item il existe une base d'ouverts quasi-compacts, et
\item l'intersection de deux ouverts quasi-compacts est quasi-compact.

\end{enumerate}

Soit $K$ un objet parfait de $D(\mathcal{O}_X)$. Alors :

\begin{enumerate}

\item[(a)] $K$ est un objet compact de $D^+(\mathcal{O}_X)$
au sens suivant : si $M = \bigoplus_{i \in I} M_i$ est borné inférieurement, alors $\Hom(K, M) = \bigoplus_{i \in I} \Hom(K, M_i)$.
\item[b)] Si $X$ est de dimension cohomologique finie, c'est-à-dire s'il existe un $d$ tel que $H^i(X, \mathcal{F}) = 0$ pour $i > d$, alors
$K$ est un objet compact de $D(\mathcal{O}_X)$.

\end{enumerate}

\end{lemma}

\begin{proof}
Soit $K^\vee$ le dual de $K$ ; voir le lemme \ref{lemma-dual-perfect-complex}. Nous avons alors $$
\Hom_{D(\mathcal{O}_X)}(K, M) =
H^0(X, K^\vee \otimes_{\mathcal{O}_X}^\mathbf{L} M)
$$ fonctoriellement en $M$ dans $D(\mathcal{O}_X)$. Puisque $K^\vee \otimes_{\mathcal{O}_X}^\mathbf{L} -$ commute aux sommes directes, il suffit de montrer que $R\Gamma(X, -)$ commute aux sommes directes considérées.

\medskip\noindent
Démonstration de (b). Comme $R\Gamma(X, K) = R\Hom(\mathcal{O}_X, K)$ et que $H^p(X, -)$ commute aux sommes directes par le lemme \ref{lemma-quasi-separated-cohomology-colimit}, il s'agit d'un cas particulier du lemme \ref{lemma-when-jshriek-compact}.

\medskip\noindent

Démonstration de (a). Soient $\mathcal{I}_i$, $i \in I$, une famille de
$\mathcal{O}_X$-modules injectifs. Par le lemme \ref{lemma-quasi-separated-cohomology-colimit},
nous avons
$$
H^p(X, \bigoplus\nolimits_{i \in I} \mathcal{I}_i) =
\bigoplus\nolimits_{i \in I} H^p(X, \mathcal{I}_i) = 0
$$
pour tout $p$. Maintenant, si $M = \bigoplus M_i$ est comme en (a), il existe un $a \in \mathbf{Z}$ tel que $H^n(M_i) = 0$
pour $n < a$. Nous pouvons donc choisir des complexes de $\mathcal{O}_X$-modules injectifs
$\mathcal{I}_i^\bullet$ représentant les $M_i$
avec $\mathcal{I}_i^n = 0$ pour $n < a$ ; voir Catégories dérivées, lemme \ref{derived-lemma-injective-resolutions-exist}. Par Injectifs, lemme \ref{injectives-lemma-derived-products},
le complexe somme directe $\bigoplus \mathcal{I}_i^\bullet$
représente $M$. Par acyclicité de Leray (Catégories dérivées, lemme \ref{derived-lemma-leray-acyclicity}), nous voyons que
$$
R\Gamma(X, M) = \Gamma(X, \bigoplus \mathcal{I}_i^\bullet) =
\bigoplus \Gamma(X, \bigoplus \mathcal{I}_i^\bullet) =
\bigoplus R\Gamma(X, M_i)
$$
comme souhaité.

\end{proof}

















\section{Formule de projection}

\label{section-projection-formula}

\noindent
Dans cette section, nous rassemblons plusieurs variantes de la formule de projection. La version la plus élémentaire est le lemme \ref{lemma-projection-formula}. Après l'avoir énoncée et démontrée, nous étudions une version plus générale faisant intervenir des complexes parfaits.

\begin{lemma}

\label{lemma-injective-tensor-finite-locally-free}
Soit $X$ un espace annelé. Soit $\mathcal{I}$ un $\mathcal{O}_X$-module injectif. Soit $\mathcal{E}$ un $\mathcal{O}_X$-module. Supposons que $\mathcal{E}$ soit localement libre de type fini sur $X$ ; voir Modules, définition \ref{modules-definition-locally-free}. Alors $\mathcal{E} \otimes_{\mathcal{O}_X} \mathcal{I}$ est un $\mathcal{O}_X$-module injectif.
\end{lemma}

\begin{proof}
Sous les hypothèses du lemme, nous avons en effet $$
\Hom_{\mathcal{O}_X}(\mathcal{F},
\mathcal{E} \otimes_{\mathcal{O}_X} \mathcal{I})
=
\Hom_{\mathcal{O}_X}(
\mathcal{F} \otimes_{\mathcal{O}_X} \mathcal{E}^\vee, \mathcal{I})
$$ où $\mathcal{E}^\vee = \SheafHom_{\mathcal{O}_X}(\mathcal{E}, \mathcal{O}_X)$ est le dual de $\mathcal{E}$, lui aussi localement libre de type fini. Puisque le produit tensoriel par un module localement libre de type fini est un foncteur exact, on conclut par Homologie, lemme \ref{homology-lemma-characterize-injectives}.
\end{proof}

\begin{lemma}

\label{lemma-projection-formula}
Soit $f : X \to Y$ un morphisme d'espaces annelés. Soit $\mathcal{F}$ un $\mathcal{O}_X$-module. Soit $\mathcal{E}$ un $\mathcal{O}_Y$-module. Supposons que $\mathcal{E}$ soit localement libre de type fini sur $Y$ ; voir Modules, définition \ref{modules-definition-locally-free}. Alors il existe des isomorphismes $$
\mathcal{E} \otimes_{\mathcal{O}_Y} R^qf_*\mathcal{F}
\longrightarrow
R^qf_*(f^*\mathcal{E} \otimes_{\mathcal{O}_X} \mathcal{F})
$$ pour tout $q \geq 0$. En fait, il existe un isomorphisme $$
\mathcal{E} \otimes_{\mathcal{O}_Y} Rf_*\mathcal{F}
\longrightarrow
Rf_*(f^*\mathcal{E} \otimes_{\mathcal{O}_X} \mathcal{F})
$$ dans $D^{+}(Y)$, fonctoriel en $\mathcal{F}$.
\end{lemma}

\begin{proof}
Choisissons une résolution injective $\mathcal{F} \to \mathcal{I}^\bullet$ sur $X$. Notons que $f^*\mathcal{E}$ est lui aussi localement libre de type fini ; nous obtenons donc une résolution $$
f^*\mathcal{E} \otimes_{\mathcal{O}_X} \mathcal{F}
\longrightarrow
f^*\mathcal{E} \otimes_{\mathcal{O}_X} \mathcal{I}^\bullet
$$ qui est injective par le lemme \ref{lemma-injective-tensor-finite-locally-free}. En appliquant $f_*$, nous obtenons $$
Rf_*(f^*\mathcal{E} \otimes_{\mathcal{O}_X} \mathcal{F})
=
f_*(f^*\mathcal{E} \otimes_{\mathcal{O}_X} \mathcal{I}^\bullet).
$$ Le lemme en résulte donc s'il est possible de montrer que $f_*(f^*\mathcal{E} \otimes_{\mathcal{O}_X} \mathcal{F}) =
\mathcal{E} \otimes_{\mathcal{O}_Y} f_*(\mathcal{F})$ fonctoriellement en le $\mathcal{O}_X$-module $\mathcal{F}$. C'est clair lorsque $\mathcal{E} = \mathcal{O}_Y^{\oplus n}$, et le cas général s'en déduit en travaillant localement sur $Y$. Nous omettons les détails.
\end{proof}
\noindent

Soit $f : X \to Y$ un morphisme d'espaces annelés. Soient $E \in D(\mathcal{O}_X)$ et $K \in D(\mathcal{O}_Y)$. Sans aucune hypothèse supplémentaire, il existe un morphisme

\begin{equation}
\label{equation-projection-formula-map}
Rf_*E \otimes^\mathbf{L}_{\mathcal{O}_Y} K
\longrightarrow
Rf_*(E \otimes^\mathbf{L}_{\mathcal{O}_X} Lf^*K)
\end{equation}

Il s'agit du morphisme adjoint au morphisme canonique
$$
Lf^*(Rf_*E \otimes^\mathbf{L}_{\mathcal{O}_Y} K) =
Lf^*Rf_*E \otimes^\mathbf{L}_{\mathcal{O}_X} Lf^*K
\longrightarrow
E \otimes^\mathbf{L}_{\mathcal{O}_X} Lf^*K
$$
provenant du morphisme $Lf^*Rf_*E \to E$ et des lemmes
\ref{lemma-pullback-tensor-product} et \ref{lemma-adjoint}. Une version raisonnablement générale de la formule de projection est la suivante.

\begin{lemma}

\label{lemma-projection-formula-perfect}
Soit $f : X \to Y$ un morphisme d'espaces annelés. Soient $E \in D(\mathcal{O}_X)$ et $K \in D(\mathcal{O}_Y)$. Si $K$ est parfait, alors $$
Rf_*E \otimes^\mathbf{L}_{\mathcal{O}_Y} K =
Rf_*(E \otimes^\mathbf{L}_{\mathcal{O}_X} Lf^*K)
$$ dans $D(\mathcal{O}_Y)$.
\end{lemma}

\begin{proof}

Pour vérifier que (\ref{equation-projection-formula-map}) est un isomorphisme, nous pouvons travailler localement sur $Y$, c'est-à-dire chercher un recouvrement ouvert
$\{V_j \to Y\}$
tel que le morphisme se restreigne en un isomorphisme sur chaque $V_j$. Par définition des objets parfaits, nous pouvons donc supposer que $K$ est représenté par un complexe strictement parfait de $\mathcal{O}_Y$-modules. Notons que, de manière tout à fait générale, l'énoncé vaut pour
$K = K_1 \oplus K_2$ si et seulement s'il vaut pour
$K_1$ et $K_2$. Nous pouvons donc supposer que $K$ est un complexe fini de $\mathcal{O}_Y$-modules libres de type fini. Dans ce cas, un argument simple faisant intervenir les troncatures brutales ramène l'énoncé au cas où $K$ est représenté par un $\mathcal{O}_Y$-module libre de type fini. Comme l'énoncé est invariant par passage aux facteurs directs finis en la variable $K$, il suffit de le démontrer pour $K = \mathcal{O}_Y[n]$,
cas dans lequel il est trivial.

\end{proof}

\noindent
Voici un cas où la formule de projection est vraie en généralité complète.

\begin{lemma}

\label{lemma-projection-formula-closed-immersion}
Soit $f : X \to Y$ un morphisme d'espaces annelés tel que $f$ soit un homéomorphisme sur une partie fermée. Alors (\ref{equation-projection-formula-map}) est toujours un isomorphisme.
\end{lemma}

\begin{proof}

Puisque $f$ est un homéomorphisme sur une partie fermée, le foncteur $f_*$
est exact (Modules, lemme \ref{modules-lemma-i-star-exact}). Ainsi,
$Rf_*$ se calcule en appliquant $f_*$ à n'importe quel complexe représentant l'objet considéré. Choisissons un complexe K-plat $\mathcal{K}^\bullet$ de $\mathcal{O}_Y$-modules représentant $K$, ainsi qu'un complexe quelconque $\mathcal{E}^\bullet$
de $\mathcal{O}_X$-modules représentant $E$. Alors
$Lf^*K$ est représenté par $f^*\mathcal{K}^\bullet$, qui est un complexe K-plat de $\mathcal{O}_X$-modules (lemme \ref{lemma-pullback-K-flat}). Ainsi, le membre de droite de (\ref{equation-projection-formula-map}) est représenté par
$$
f_*\text{Tot}(\mathcal{E}^\bullet
\otimes_{\mathcal{O}_X} f^*\mathcal{K}^\bullet)
$$
Par le même raisonnement, le membre de gauche est représenté par
$$
\text{Tot}(f_*\mathcal{E}^\bullet \otimes_{\mathcal{O}_Y} \mathcal{K}^\bullet)
$$
Puisque $f_*$ commute aux sommes directes (Modules, lemme \ref{modules-lemma-i-star-right-adjoint}), il suffit de montrer que
$$
f_*(\mathcal{E} \otimes_{\mathcal{O}_X} f^*\mathcal{K}) =
f_*\mathcal{E} \otimes_{\mathcal{O}_Y} \mathcal{K}
$$
pour tout $\mathcal{O}_X$-module $\mathcal{E}$ et $\mathcal{O}_Y$-module
$\mathcal{K}$. Vérifions cette égalité sur les germes. Soit $y \in Y$. Si $y \not \in f(X)$, les germes des deux membres sont nuls. Si $y = f(x)$, il faut montrer que
$$
\mathcal{E}_x \otimes_{\mathcal{O}_{X, x}}
(\mathcal{O}_{X, x} \otimes_{\mathcal{O}_{Y, y}} \mathcal{F}_y) =
\mathcal{E}_x \otimes_{\mathcal{O}_{Y, y}} \mathcal{F}_y
$$
(en utilisant Faisceaux, lemme \ref{sheaves-lemma-stalks-closed-pushforward}
et lemme \ref{sheaves-lemma-stalk-pullback-modules}). Cette égalité est vraie, ce qui démontre le lemme.

\end{proof}

\begin{remark}

\label{remark-compatible-with-diagram}

Le morphisme (\ref{equation-projection-formula-map}) est compatible avec le morphisme de changement de base de la remarque \ref{remark-base-change} au sens suivant. Supposons que
$$
\xymatrix{
X' \ar[r]_{g'} \ar[d]_{f'} &
X \ar[d]^f \\
Y' \ar[r]^g &
Y
}
$$
soit un diagramme commutatif d'espaces annelés. Soient $E \in D(\mathcal{O}_X)$ et $K \in D(\mathcal{O}_Y)$. Alors le diagramme
$$
\xymatrix{
Lg^*(Rf_*E \otimes^\mathbf{L}_{\mathcal{O}_Y} K) \ar[r]_p \ar[d]_t &
Lg^*Rf_*(E \otimes^\mathbf{L}_{\mathcal{O}_X} Lf^*K) \ar[d]_b \\
Lg^*Rf_*E \otimes^\mathbf{L}_{\mathcal{O}_{Y'}} Lg^*K \ar[d]_b &
Rf'_*L(g')^*(E \otimes^\mathbf{L}_{\mathcal{O}_X} Lf^*K) \ar[d]_t \\
Rf'_*L(g')^*E \otimes^\mathbf{L}_{\mathcal{O}_{Y'}} Lg^*K \ar[rd]_p &
Rf'_*(L(g')^*E \otimes^\mathbf{L}_{\mathcal{O}_{Y'}} L(g')^*Lf^*K) \ar[d]_c \\
& Rf'_*(L(g')^*E \otimes^\mathbf{L}_{\mathcal{O}_{Y'}} L(f')^*Lg^*K)
}
$$
est commutatif. Ici, les flèches étiquetées $t$ sont obtenues en appliquant le lemme \ref{lemma-pullback-tensor-product}, celles étiquetées $b$ en appliquant la remarque \ref{remark-base-change}, celles étiquetées $p$
en appliquant (\ref{equation-projection-formula-map}), et
$c$ vient de $L(g')^* \circ Lf^* = L(f')^* \circ Lg^*$. Nous omettons la vérification.

\end{remark}















\section{Un opérateur introduit par Berthelot et Ogus}

\label{section-eta}

\noindent
Cette section continue la discussion commencée dans Compléments sur l'algèbre, section \ref{more-algebra-section-eta}. Nous encourageons vivement le lecteur à lire cette section en premier.

\begin{lemma}

\label{lemma-invertible-ideal-sheaf}

Soit $(X, \mathcal{O}_X)$ un espace annelé.
Soit $\mathcal{I} \subset \mathcal{O}_X$ un faisceau d'idéaux. Considérons les deux conditions suivantes :

\begin{enumerate}

\item pour tout $x \in X$, il existe un voisinage ouvert
$U \subset X$ de $x$ et un élément $f \in \mathcal{I}(U)$ tels que
$\mathcal{I}|_U = \mathcal{O}_U \cdot f$ et
$f : \mathcal{O}_U \to \mathcal{O}_U$ est injectif, et
\item  $\mathcal{I}$ est inversible comme $\mathcal{O}_X$-module.

\end{enumerate}

Alors (1) implique (2), et la réciproque est vraie si tous les germes
$\mathcal{O}_{X, x}$ du faisceau structural sont des anneaux locaux.

\end{lemma}

\begin{proof}
Démonstration omise. Indication : utiliser Modules, lemme \ref{modules-lemma-invertible-is-locally-free-rank-1}.
\end{proof}

\begin{situation}

\label{situation-eta}
Soit $(X, \mathcal{O}_X)$ un espace annelé. Soit $\mathcal{I} \subset \mathcal{O}_X$ un faisceau d'idéaux satisfaisant à la condition (1) du lemme \ref{lemma-invertible-ideal-sheaf}\footnote{La discussion de cette section se généralise au cas où l'on suppose seulement que $\mathcal{I}$ est un $\mathcal{O}_X$-module inversible au sens de Modules, section \ref{modules-section-invertible}.}.
\end{situation}

\begin{lemma}

\label{lemma-I-torsion-free}

Dans la situation \ref{situation-eta},
soit $\mathcal{F}$ un $\mathcal{O}_X$-module. Les conditions suivantes sont équivalentes :

\begin{enumerate}

\item le sous-faisceau $\mathcal{F}[\mathcal{I}] \subset \mathcal{F}$
des sections annulées par $\mathcal{I}$ est nul,
\item le sous-faisceau $\mathcal{F}[\mathcal{I}^n]$ est nul pour tout $n \geq 1$,
\item l'application de multiplication
$\mathcal{I} \otimes_{\mathcal{O}_X} \mathcal{F} \to \mathcal{F}$
est injectif,
\item pour tout ouvert $U \subset X$ tel que
$\mathcal{I}|_U = \mathcal{O}_U \cdot f$
pour un certain $f \in \mathcal{I}(U)$,
le morphisme $f : \mathcal{F}|_U \to \mathcal{F}|_U$ est injectif,
\item pour tout $x \in X$ et tout générateur $f$ de l'idéal
$\mathcal{I}_x \subset \mathcal{O}_{X, x}$, l'élément $f$
est un non-diviseur de zéro sur le germe $\mathcal{F}_x$.

\end{enumerate}

\end{lemma}

\begin{proof}
Démonstration omise.
\end{proof}

\noindent
Dans la situation \ref{situation-eta}, soit $\mathcal{F}$ un $\mathcal{O}_X$-module. Si les conditions équivalentes du lemme \ref{lemma-I-torsion-free} sont vérifiées, nous dirons que $\mathcal{F}$ est {\it sans $\mathcal{I}$-torsion}. Dans ce cas, pour tout $i \in \mathbf{Z}$, nous noterons $$
\mathcal{I}^i\mathcal{F} =
\mathcal{I}^{\otimes i} \otimes_{\mathcal{O}_X} \mathcal{F}
$$ de sorte que nous ayons des inclusions $$
\ldots \subset
\mathcal{I}^{i + 1}\mathcal{F} \subset
\mathcal{I}^i\mathcal{F} \subset
\mathcal{I}^{i - 1}\mathcal{F} \subset \ldots
$$ Les modules $\mathcal{I}^i\mathcal{F}$ sont localement, mais pas nécessairement globalement, isomorphes à $\mathcal{F}$ comme $\mathcal{O}_X$-modules.

\medskip\noindent

Soit $\mathcal{F}^\bullet$ un complexe sans $\mathcal{I}$-torsion de $\mathcal{O}_X$-modules, de différentielles $d^i : \mathcal{F}^i \to \mathcal{F}^{i + 1}$. Nous définissons alors $\eta_\mathcal{I}\mathcal{F}^\bullet$
comme le complexe dont les termes sont

\begin{align*}
(\eta_\mathcal{I}\mathcal{F})^i
& =
\Ker\left(
d^i, -1 :
\mathcal{I}^i\mathcal{F}^i \oplus \mathcal{I}^{i + 1}\mathcal{F}^{i + 1}
\to
\mathcal{I}^i\mathcal{F}^{i + 1}
\right) \\
& =
\Ker\left(d^i :
\mathcal{I}^i\mathcal{F}^i
\to
\mathcal{I}^i\mathcal{F}^{i + 1}/
\mathcal{I}^{i + 1}\mathcal{F}^{i + 1}
\right)
\end{align*}

et dont la différentielle est induite par $d^i$. Autrement dit, une section locale
$s$ de $(\eta_\mathcal{I}\mathcal{F})^i$ est la même chose qu'une section locale
$s$ de $\mathcal{I}^i\mathcal{F}^i$ telle que son image $d^i(s)$
dans $\mathcal{I}^i\mathcal{F}^{i + 1}$ appartienne au sous-faisceau
$\mathcal{I}^{i + 1}\mathcal{F}^{i + 1}$. Observons que $\eta_\mathcal{I}\mathcal{F}^\bullet$
est encore un complexe de modules sans $\mathcal{I}$-torsion.

\medskip\noindent
Soit $a^\bullet : \mathcal{F}^\bullet \to \mathcal{G}^\bullet$ un morphisme de complexes sans $\mathcal{I}$-torsion de $\mathcal{O}_X$-modules. Nous obtenons alors un morphisme de complexes $$
\eta_\mathcal{I} a^\bullet :
\eta_\mathcal{I}\mathcal{F}^\bullet
\longrightarrow
\eta_\mathcal{I}\mathcal{G}^\bullet
$$ induit par les morphismes $\mathcal{I}^i\mathcal{F}^i \to \mathcal{I}^i\mathcal{G}^i$. Le lecteur vérifiera que l'on obtient ainsi un endofoncteur de la catégorie des complexes sans $\mathcal{I}$-torsion de $\mathcal{O}_X$-modules.

\medskip\noindent

Si $a^\bullet, b^\bullet : \mathcal{F}^\bullet \to \mathcal{G}^\bullet$
sont deux morphismes de complexes sans $\mathcal{I}$-torsion de $\mathcal{O}_X$-modules et si $h = \{h^i : \mathcal{F}^i \to \mathcal{G}^{i - 1}\}$ est une homotopie entre $a^\bullet$ et $b^\bullet$, nous définissons
$\eta_\mathcal{I}h$ comme la famille de morphismes
$(\eta_\mathcal{I}h)^i : (\eta_\mathcal{I}\mathcal{F})^i \to
(\eta_\mathcal{I}\mathcal{G})^{i - 1}$
qui envoie $x$ sur $h^i(x)$ ; cette définition a un sens, car si
$x$ est une section locale de $\mathcal{I}^i\mathcal{F}^i$,
alors $h^i(x)$ est une section locale de $\mathcal{I}^i\mathcal{G}^{i - 1}$,
qui est bien contenue dans $(\eta_\mathcal{I}\mathcal{G})^{i - 1}$. Le lecteur vérifiera que $\eta_\mathcal{I}h$ est une homotopie entre $\eta_\mathcal{I}a^\bullet$ et $\eta_\mathcal{I}b^\bullet$. En définitive, nous obtenons un foncteur
$$
\eta_f :
K(\mathcal{I}\text{-sans torsion }\mathcal{O}_X\text{-modules})
\longrightarrow
K(\mathcal{I}\text{-sans torsion }\mathcal{O}_X\text{-modules})
$$
sur la catégorie d'homotopie (Catégories dérivées, section \ref{derived-section-homotopy}) de la catégorie additive des modules sans
$\mathcal{I}$-torsion sur $\mathcal{O}_X$. Le foncteur $\eta_\mathcal{I}$ n'est en aucun sens un foncteur exact de catégories triangulées ; comparer avec Compléments sur l'algèbre, exemple \ref{more-algebra-example-eta-not-distinguished}.

\begin{lemma}

\label{lemma-eta-stalks}
Dans la situation \ref{situation-eta}, soit $\mathcal{F}^\bullet$ un complexe sans $\mathcal{I}$-torsion de $\mathcal{O}_X$-modules. Pour $x \in X$, choisissons un générateur $f \in \mathcal{I}_x$. Alors le germe $(\eta_\mathcal{I}\mathcal{F}^\bullet)_x$ est canoniquement isomorphe au complexe $\eta_f\mathcal{F}^\bullet_x$ construit dans Compléments sur l'algèbre, section \ref{more-algebra-section-eta}.
\end{lemma}

\begin{proof}
Démonstration omise.
\end{proof}

\begin{lemma}

\label{lemma-eta-first-property}

Dans la situation \ref{situation-eta},
soit $\mathcal{F}^\bullet$ un complexe sans
$\mathcal{I}$-torsion de $\mathcal{O}_X$-modules. Il existe un isomorphisme canonique
$$
\mathcal{I}^{\otimes i} \otimes_{\mathcal{O}_X}
\left(
H^i(\mathcal{F}^\bullet)/H^i(\mathcal{F}^\bullet)[\mathcal{I}]
\right)
\longrightarrow H^i(\eta_\mathcal{I}\mathcal{F}^\bullet)
$$
des faisceaux de cohomologie.

\end{lemma}

\begin{proof}

Nous définissons un morphisme
$$
\mathcal{I}^{\otimes i} \otimes_{\mathcal{O}_X} H^i(\mathcal{F}^\bullet)
\longrightarrow
H^i(\eta_\mathcal{I}\mathcal{F}^\bullet)
$$
comme suit. Soit $g$ une section locale de $\mathcal{I}^{\otimes i}$
et soit $\overline{s}$ une section locale de $H^i(\mathcal{F}^\bullet)$. Alors $\overline{s}$ est localement la classe d'une section locale $s$ de
$\Ker(d^i : \mathcal{F}^i \to \mathcal{F}^{i + 1})$. Nous envoyons $g \otimes \overline{s}$ sur la section locale
$gs$ de $(\eta_\mathcal{I}\mathcal{F})^i \subset \mathcal{I}^i\mathcal{F}$. Bien entendu, $gs$ appartient au noyau de $d^i$ sur
$\eta_\mathcal{I}\mathcal{F}^\bullet$ et définit donc une section locale de $H^i(\eta_\mathcal{I}\mathcal{F}^\bullet)$. Il est immédiat que cette définition ne dépend pas des choix. Nous affirmons que ce morphisme se factorise par l'isomorphisme donné dans le lemme. Nous pouvons le vérifier sur les germes ; le lemme \ref{lemma-eta-stalks} ramène alors l'assertion au résultat de Compléments sur l'algèbre, lemme \ref{more-algebra-lemma-eta-first-property}.

\end{proof}

\begin{lemma}

\label{lemma-eta-qis}
Dans la situation \ref{situation-eta}, soit $\mathcal{F}^\bullet \to \mathcal{G}^\bullet$ un morphisme de complexes sans $\mathcal{I}$-torsion de $\mathcal{O}_X$-modules. Alors le morphisme induit $\eta_\mathcal{I}\mathcal{F}^\bullet \to \eta_\mathcal{I}\mathcal{G}^\bullet$ est lui aussi un quasi-isomorphisme.
\end{lemma}

\begin{proof}

Cela résulte de la compatibilité des isomorphismes du lemme \ref{lemma-eta-first-property}
avec les morphismes de complexes.

\end{proof}

\begin{lemma}

\label{lemma-Leta}
Dans la situation \ref{situation-eta}, il existe un foncteur additif\footnote{Attention : ce foncteur n'est pas exact, c'est-à-dire qu'il ne transforme pas les triangles distingués en triangles distingués.} $L\eta_\mathcal{I} : D(\mathcal{O}_X) \to D(\mathcal{O}_X)$ tel que, si $M$ dans $D(\mathcal{O}_X)$ est représenté par un complexe $\mathcal{F}^\bullet$ sans $\mathcal{I}$-torsion de $\mathcal{O}_X$-modules, alors $L\eta_\mathcal{I}M = \eta_\mathcal{I}\mathcal{F}^\bullet$. Il en va de même pour les morphismes.
\end{lemma}

\begin{proof}

Notons $\mathcal{T} \subset \textit{Mod}(\mathcal{O}_X)$ la sous-catégorie pleine des modules sans $\mathcal{I}$-torsion sur $\mathcal{O}_X$. Nous avons l'inclusion correspondante
$$
K(\mathcal{T})
\quad\subset\quad
K(\textit{Mod}(\mathcal{O}_X)) = K(\mathcal{O}_X)
$$
de $K(\mathcal{T})$ comme sous-catégorie triangulée pleine de $K(\mathcal{O}_X)$. Soit $S \subset \text{Arrows}(K(\mathcal{T}))$ l'ensemble des quasi-isomorphismes. Nous allons appliquer Catégories dérivées, lemme \ref{derived-lemma-localization-subcategory},
pour montrer que l'application
$$
S^{-1}K(\mathcal{T}) \longrightarrow D(\mathcal{O}_X)
$$
est une équivalence de catégories triangulées. D'après ce lemme, il suffit de démontrer que, pour tout complexe $\mathcal{G}^\bullet$ de
$\mathcal{O}_X$-modules, il existe un quasi-isomorphisme $\mathcal{F}^\bullet \to \mathcal{G}^\bullet$
avec $\mathcal{F}^\bullet$
un complexe sans $\mathcal{I}$-torsion de $\mathcal{O}_X$-modules. Par le lemme \ref{lemma-K-flat-resolution}, nous pouvons trouver un quasi-isomorphisme
$\mathcal{F}^\bullet \to \mathcal{G}^\bullet$ de sorte que le complexe
$\mathcal{F}^\bullet$ soit K-plat (ce dont nous ne nous servirons pas) et formé de $\mathcal{O}_X$-modules plats $\mathcal{F}^i$. Par la troisième caractérisation du lemme \ref{lemma-I-torsion-free},
un $\mathcal{O}_X$-module plat est un module sans
$\mathcal{I}$-torsion parmi les $\mathcal{O}_X$-modules, ce qui achève cette étape.

\medskip\noindent

Ces préliminaires étant établis, nous pouvons définir $L\eta_f$. En effet, la discussion qui suit le lemme \ref{lemma-I-torsion-free}
nous a déjà fourni un foncteur bien défini
$$
K(\mathcal{T}) \xrightarrow{\eta_f} K(\mathcal{T}) \to
K(\mathcal{O}_X) \to D(\mathcal{O}_X)
$$
qui, d'après le lemme \ref{lemma-eta-qis}, envoie les quasi-isomorphismes sur des quasi-isomorphismes. Ce foncteur se factorise donc par
$S^{-1}K(\mathcal{T}) = D(\mathcal{O}_X)$, en vertu de Catégories, lemme \ref{categories-lemma-properties-left-localization}.

\end{proof}

\noindent

Dans la situation \ref{situation-eta}, construisons les opérateurs de Bockstein. Observons d'abord qu'il existe un diagramme commutatif
$$
\xymatrix{
0 \ar[r] &
\mathcal{I}^{i + 1} \ar[r] \ar[d] &
\mathcal{I}^i \ar[r] \ar[d] &
\mathcal{I}^i/\mathcal{I}^{i + 1}
\ar[r] \ar@{=}[d] &
0 \\
0 \ar[r] &
\mathcal{I}^{i + 1}/\mathcal{I}^{i + 2} \ar[r] &
\mathcal{I}^i/\mathcal{I}^{i + 2} \ar[r] &
\mathcal{I}^i/\mathcal{I}^{i + 1} \ar[r] &
0
}
$$
dont les lignes sont des suites exactes courtes de $\mathcal{O}_X$-modules. Soit $M$ un objet de $D(\mathcal{O}_X)$. En tensorisant le diagramme ci-dessus par $M$, on obtient un morphisme
$$
\xymatrix{
M \otimes^\mathbf{L} \mathcal{I}^{i + 1} \ar[r] \ar[d] &
M \otimes^\mathbf{L} \mathcal{I}^i \ar[r] \ar[d] &
M \otimes^\mathbf{L} \mathcal{I}^i/\mathcal{I}^{i + 1} \ar[d]^{\text{id}} \\
M \otimes^\mathbf{L} \mathcal{I}^{i + 1}/\mathcal{I}^{i + 2} \ar[r] &
M \otimes^\mathbf{L} \mathcal{I}^i/\mathcal{I}^{i + 2} \ar[r] &
M \otimes^\mathbf{L} \mathcal{I}^i/\mathcal{I}^{i + 1}
}
$$
de triangles distingués. La suite exacte longue des faisceaux de cohomologie associée au triangle inférieur détermine en particulier un
{\it opérateur de Bockstein}

$$
\beta = \beta^i :
H^i(M \otimes^\mathbf{L} \mathcal{I}^i/\mathcal{I}^{i + 1})
\longrightarrow
H^{i + 1}(M \otimes^\mathbf{L} \mathcal{I}^{i + 1}/\mathcal{I}^{i + 2})
$$
pour tout $i \in \mathbf{Z}$. Pour un usage ultérieur, notons que le diagramme commutatif ci-dessus fournit une factorisation

\begin{equation}
\label{equation-factorization-bockstein}
\vcenter{
\xymatrix{
H^i(M \otimes^\mathbf{L} \mathcal{I}^i/\mathcal{I}^{i + 1})
\ar[r]_\delta \ar[rd]_\beta &
H^{i + 1}(M \otimes^\mathbf{L} \mathcal{I}^{i + 1}) \ar[d] \\
&
H^{i + 1}(M \otimes^\mathbf{L} \mathcal{I}^{i + 1}/\mathcal{I}^{i + 2})
}
}
\end{equation}

de l'opérateur de Bockstein, où $\delta$ est l'opérateur de bord provenant du triangle distingué supérieur du diagramme commutatif ci-dessus. Nous obtenons un complexe

\begin{equation}
\label{equation-complex-bocksteins}
H^\bullet(M/\mathcal{I}) =
\left[
\begin{matrix}
\ldots \\
\downarrow \\
H^{i - 1}(M \otimes^\mathbf{L} \mathcal{I}^{i - 1}/\mathcal{I}^i) \\
\downarrow \beta \\
H^i(M \otimes^\mathbf{L} \mathcal{I}^i/\mathcal{I}^{i + 1}) \\
\downarrow \beta \\
H^{i + 1}(M \otimes^\mathbf{L} \mathcal{I}^{i + 1}/\mathcal{I}^{i + 2}) \\
\downarrow \\
\ldots
\end{matrix}
\right]
\end{equation}

c'est-à-dire que $\beta \circ \beta = 0$. En effet, nous pouvons le vérifier sur les germes et le déduire alors du résultat algébrique correspondant démontré dans Compléments sur l'algèbre, section \ref{more-algebra-section-eta}. Autre démonstration : les suites exactes courtes
$0 \to \mathcal{I}^{i + 1}/\mathcal{I}^{i + 2}
\to \mathcal{I}^i/\mathcal{I}^{i + 2}
\to \mathcal{I}^i/\mathcal{I}^{i + 1} \to 0$
définissent des morphismes
$b^i : \mathcal{I}^i/\mathcal{I}^{i + 1} \to
(\mathcal{I}^{i + 1}/\mathcal{I}^{i + 2})[1]$
en $D(\mathcal{O}_X)$
qui induisent les morphismes $\beta$ ci-dessus en tensorisant par $M$
et en prenant les faisceaux de cohomologie. On montre ensuite que la composée
$b^{i + 1}[1] \circ b^i : \mathcal{I}^i/\mathcal{I}^{i + 1} \to
(\mathcal{I}^{i + 1}/\mathcal{I}^{i + 2})[1] \to
(\mathcal{I}^{i + 2}/\mathcal{I}^{i + 3})[2]$
est nulle dans $D(\mathcal{O}_X)$ en appliquant le critère de Catégories dérivées, lemme \ref{derived-lemma-cup-ext-1-zero},
et en utilisant le fait que le module $\mathcal{I}^i/\mathcal{I}^{i + 3}$
est une extension de $\mathcal{I}^{i + 1}/\mathcal{I}^{i + 3}$
par $\mathcal{I}^i/\mathcal{I}^{i + 1}$.

\begin{lemma}

\label{lemma-eta-second-property}

Dans la situation \ref{situation-eta}, soit $M$ un objet de
$D(\mathcal{O}_X)$. Il existe un isomorphisme canonique
$$
L\eta_\mathcal{I}M \otimes^\mathbf{L} \mathcal{O}_X/\mathcal{I}
\longrightarrow
H^\bullet(M/\mathcal{I})
$$
dans $D(\mathcal{O}_X)$, où le membre de droite est le complexe (\ref{equation-complex-bocksteins}).

\end{lemma}

\begin{proof}
Par la construction de $L\eta_\mathcal{I}$ dans le lemme \ref{lemma-eta-qis}, nous pouvons supposer que $M$ est représenté par un complexe sans $\mathcal{I}$-torsion de $\mathcal{O}_X$-modules $\mathcal{F}^\bullet$. Alors $L\eta_\mathcal{I}M$ est représenté par le complexe $\eta_\mathcal{I}\mathcal{F}^\bullet$, lui aussi sans $\mathcal{I}$-torsion comme complexe de $\mathcal{O}_X$-modules. Ainsi, $L\eta_\mathcal{I}M \otimes^\mathbf{L} \mathcal{O}_X/\mathcal{I}$ est représenté par le complexe $\eta_\mathcal{I}\mathcal{F}^\bullet \otimes \mathcal{O}_X/\mathcal{I}$. De même, le complexe $H^\bullet(M/\mathcal{I})$ a pour termes $H^i(\mathcal{F}^\bullet \otimes \mathcal{I}^i/\mathcal{I}^{i + 1})$.

\medskip\noindent

Soit $f$ un générateur local de $\mathcal{I}$. Soit $s$ une section locale de $(\eta_\mathcal{I}\mathcal{F})^i$. Nous pouvons écrire $s = f^is'$ pour une section locale $s'$ de
$\mathcal{F}^i$ et, de même, $d^i(s) = f^{i + 1}t$ pour une section locale $t$ de $\mathcal{F}^{i + 1}$. Ainsi, $d^i$ envoie $f^is'$
sur zéro dans $\mathcal{F}^{i + 1} \otimes \mathcal{I}^i/\mathcal{I}^{i + 1}$. Nous pouvons donc envoyer $s$ sur la classe de $f^is'$ dans
$H^i(\mathcal{F}^\bullet \otimes \mathcal{I}^i/\mathcal{I}^{i + 1})$. Cette règle définit un morphisme
$$
(\eta_\mathcal{I}\mathcal{F})^i \otimes \mathcal{O}_X/\mathcal{I}
\longrightarrow
H^i(\mathcal{F}^\bullet \otimes \mathcal{I}^i/\mathcal{I}^{i + 1})
$$
de $\mathcal{O}_X$-modules. Un calcul montre que ces morphismes sont compatibles aux différentielles (essentiellement parce que $\beta$
envoie la classe de $f^is'$ sur celle de $f^{i + 1}t$), d'où un morphisme de complexes représentant la flèche de l'énoncé du lemme.

\medskip\noindent

Pour achever la démonstration, observons que la construction du paragraphe précédent coïncide sur les germes avec les morphismes construits dans Compléments sur l'algèbre, lemme
\ref{more-algebra-lemma-eta-second-property}
ce qui permet de conclure.

\end{proof}

\begin{lemma}

\label{lemma-eta-tensor-invertible}
Dans la situation \ref{situation-eta}, soit $\mathcal{F}^\bullet$ un complexe sans $\mathcal{I}$-torsion de $\mathcal{O}_X$-modules. Soit $\mathcal{L}$ un $\mathcal{O}_X$-module inversible. Alors $\eta_\mathcal{I}(\mathcal{F}^\bullet \otimes \mathcal{L}) =
(\eta_\mathcal{I}\mathcal{F}^\bullet) \otimes \mathcal{L}$.
\end{lemma}

\begin{proof}
C'est immédiat par construction.
\end{proof}

\begin{lemma}

\label{lemma-eta-cohomology-locally-free}
Dans la situation \ref{situation-eta}, soit $M$ un objet de $D(\mathcal{O}_X)$. Soit $x \in X$ tel que $\mathcal{O}_{X, x}$ soit non nul. Si $H^i(M)_x$ est libre de type fini sur $\mathcal{O}_{X, x}$, alors $H^i(L\eta_\mathcal{I}M)_x$ est libre de type fini sur $\mathcal{O}_{X, x}$ et de même rang.
\end{lemma}

\begin{proof}
Soit en effet $f \in \mathcal{O}_{X, x}$ un générateur du germe $\mathcal{I}_x$. Alors $f$ est un non-diviseur de zéro dans $\mathcal{O}_{X, x}$, d'où $H^i(M)_x[f] = 0$. Par le lemme \ref{lemma-eta-first-property}, $H^i(L\eta_\mathcal{I}M)_x$ est donc isomorphe à $\mathcal{I}^i_x \otimes_{\mathcal{O}_{X, x}} H^i(M)_x$, qui est libre de même rang, comme souhaité.
\end{proof}









\begin{multicols}{2}[\section{Autres chapitres}]
\noindent
Préliminaires
\begin{enumerate}
\item \hyperref[introduction-section-phantom]{Introduction}
\item \hyperref[conventions-section-phantom]{Conventions}
\item \hyperref[sets-section-phantom]{Théorie des ensembles}
\item \hyperref[categories-section-phantom]{Catégories}
\item \hyperref[topology-section-phantom]{Topologie}
\item \hyperref[sheaves-section-phantom]{Faisceaux sur les espaces}
\item \hyperref[sites-section-phantom]{Sites et faisceaux}
\item \hyperref[stacks-section-phantom]{Champs}
\item \hyperref[fields-section-phantom]{Corps}
\item \hyperref[algebra-section-phantom]{Algèbre commutative}
\item \hyperref[brauer-section-phantom]{Groupes de Brauer}
\item \hyperref[homology-section-phantom]{Algèbre homologique}
\item \hyperref[derived-section-phantom]{Catégories dérivées}
\item \hyperref[simplicial-section-phantom]{Méthodes simpliciales}
\item \hyperref[more-algebra-section-phantom]{Compléments d'algèbre}
\item \hyperref[smoothing-section-phantom]{Lissage des morphismes d'anneaux}
\item \hyperref[modules-section-phantom]{Faisceaux de modules}
\item \hyperref[sites-modules-section-phantom]{Modules sur les sites}
\item \hyperref[injectives-section-phantom]{Objets injectifs}
\item \hyperref[cohomology-section-phantom]{Cohomologie des faisceaux}
\item \hyperref[sites-cohomology-section-phantom]{Cohomologie sur les sites}
\item \hyperref[dga-section-phantom]{Algèbre différentielle graduée}
\item \hyperref[dpa-section-phantom]{Algèbre à puissances divisées}
\item \hyperref[sdga-section-phantom]{Faisceaux différentiels gradués}
\item \hyperref[hypercovering-section-phantom]{Hyperrecouvrements}
\end{enumerate}
Schémas
\begin{enumerate}
\setcounter{enumi}{25}
\item \hyperref[schemes-section-phantom]{Schémas}
\item \hyperref[constructions-section-phantom]{Constructions de schémas}
\item \hyperref[properties-section-phantom]{Propriétés des schémas}
\item \hyperref[morphisms-section-phantom]{Morphismes de schémas}
\item \hyperref[coherent-section-phantom]{Cohomologie des schémas}
\item \hyperref[divisors-section-phantom]{Diviseurs}
\item \hyperref[limits-section-phantom]{Limites de schémas}
\item \hyperref[varieties-section-phantom]{Variétés}
\item \hyperref[topologies-section-phantom]{Topologies sur les schémas}
\item \hyperref[descent-section-phantom]{Descente}
\item \hyperref[perfect-section-phantom]{Catégories dérivées de schémas}
\item \hyperref[more-morphisms-section-phantom]{Compléments sur les morphismes}
\item \hyperref[flat-section-phantom]{Compléments sur la platitude}
\item \hyperref[groupoids-section-phantom]{Schémas en groupoïdes}
\item \hyperref[more-groupoids-section-phantom]{Compléments sur les schémas en groupoïdes}
\item \hyperref[etale-section-phantom]{Morphismes étales de schémas}
\end{enumerate}
Thèmes de théorie des schémas
\begin{enumerate}
\setcounter{enumi}{41}
\item \hyperref[chow-section-phantom]{Homologie de Chow}
\item \hyperref[intersection-section-phantom]{Théorie de l'intersection}
\item \hyperref[pic-section-phantom]{Schémas de Picard des courbes}
\item \hyperref[weil-section-phantom]{Théories cohomologiques de Weil}
\item \hyperref[adequate-section-phantom]{Modules adéquats}
\item \hyperref[dualizing-section-phantom]{Complexes dualisants}
\item \hyperref[duality-section-phantom]{Dualité pour les schémas}
\item \hyperref[discriminant-section-phantom]{Discriminants et différentes}
\item \hyperref[derham-section-phantom]{Cohomologie de de Rham}
\item \hyperref[local-cohomology-section-phantom]{Cohomologie locale}
\item \hyperref[algebraization-section-phantom]{Géométrie algébrique et formelle}
\item \hyperref[curves-section-phantom]{Courbes algébriques}
\item \hyperref[resolve-section-phantom]{Résolution des surfaces}
\item \hyperref[models-section-phantom]{Réduction semi-stable}
\item \hyperref[functors-section-phantom]{Foncteurs et morphismes}
\item \hyperref[equiv-section-phantom]{Catégories dérivées de variétés}
\item \hyperref[pione-section-phantom]{Groupes fondamentaux de schémas}
\item \hyperref[etale-cohomology-section-phantom]{Cohomologie étale}
\item \hyperref[crystalline-section-phantom]{Cohomologie cristalline}
\item \hyperref[proetale-section-phantom]{Cohomologie pro-étale}
\item \hyperref[relative-cycles-section-phantom]{Cycles relatifs}
\item \hyperref[more-etale-section-phantom]{Compléments de cohomologie étale}
\item \hyperref[trace-section-phantom]{La formule des traces}
\end{enumerate}
Espaces algébriques
\begin{enumerate}
\setcounter{enumi}{64}
\item \hyperref[spaces-section-phantom]{Espaces algébriques}
\item \hyperref[spaces-properties-section-phantom]{Propriétés des espaces algébriques}
\item \hyperref[spaces-morphisms-section-phantom]{Morphismes d'espaces algébriques}
\item \hyperref[decent-spaces-section-phantom]{Espaces algébriques décents}
\item \hyperref[spaces-cohomology-section-phantom]{Cohomologie des espaces algébriques}
\item \hyperref[spaces-limits-section-phantom]{Limites d'espaces algébriques}
\item \hyperref[spaces-divisors-section-phantom]{Diviseurs sur les espaces algébriques}
\item \hyperref[spaces-over-fields-section-phantom]{Espaces algébriques sur des corps}
\item \hyperref[spaces-topologies-section-phantom]{Topologies sur les espaces algébriques}
\item \hyperref[spaces-descent-section-phantom]{Descente et espaces algébriques}
\item \hyperref[spaces-perfect-section-phantom]{Catégories dérivées d'espaces}
\item \hyperref[spaces-more-morphisms-section-phantom]{Compléments sur les morphismes d'espaces}
\item \hyperref[spaces-flat-section-phantom]{Platitude sur les espaces algébriques}
\item \hyperref[spaces-groupoids-section-phantom]{Groupoïdes dans les espaces algébriques}
\item \hyperref[spaces-more-groupoids-section-phantom]{Compléments sur les groupoïdes dans les espaces}
\item \hyperref[bootstrap-section-phantom]{Amorçage}
\item \hyperref[spaces-pushouts-section-phantom]{Sommes amalgamées d'espaces algébriques}
\end{enumerate}
Thèmes de géométrie
\begin{enumerate}
\setcounter{enumi}{81}
\item \hyperref[spaces-chow-section-phantom]{Groupes de Chow des espaces}
\item \hyperref[groupoids-quotients-section-phantom]{Quotients de groupoïdes}
\item \hyperref[spaces-more-cohomology-section-phantom]{Compléments sur la cohomologie des espaces}
\item \hyperref[spaces-simplicial-section-phantom]{Espaces simpliciaux}
\item \hyperref[spaces-duality-section-phantom]{Dualité pour les espaces}
\item \hyperref[formal-spaces-section-phantom]{Espaces algébriques formels}
\item \hyperref[restricted-section-phantom]{Algébrisation des espaces formels}
\item \hyperref[spaces-resolve-section-phantom]{Résolution des surfaces revisitée}
\end{enumerate}
Théorie des déformations
\begin{enumerate}
\setcounter{enumi}{89}
\item \hyperref[formal-defos-section-phantom]{Théorie formelle des déformations}
\item \hyperref[defos-section-phantom]{Théorie des déformations}
\item \hyperref[cotangent-section-phantom]{Le complexe cotangent}
\item \hyperref[examples-defos-section-phantom]{Problèmes de déformation}
\end{enumerate}
Champs algébriques
\begin{enumerate}
\setcounter{enumi}{93}
\item \hyperref[algebraic-section-phantom]{Champs algébriques}
\item \hyperref[examples-stacks-section-phantom]{Exemples de champs}
\item \hyperref[stacks-sheaves-section-phantom]{Faisceaux sur les champs algébriques}
\item \hyperref[criteria-section-phantom]{Critères de représentabilité}
\item \hyperref[artin-section-phantom]{Axiomes d'Artin}
\item \hyperref[quot-section-phantom]{Espaces de Quot et de Hilbert}
\item \hyperref[stacks-properties-section-phantom]{Propriétés des champs algébriques}
\item \hyperref[stacks-morphisms-section-phantom]{Morphismes de champs algébriques}
\item \hyperref[stacks-limits-section-phantom]{Limites de champs algébriques}
\item \hyperref[stacks-cohomology-section-phantom]{Cohomologie des champs algébriques}
\item \hyperref[stacks-perfect-section-phantom]{Catégories dérivées de champs}
\item \hyperref[stacks-introduction-section-phantom]{Introduction aux champs algébriques}
\item \hyperref[stacks-more-morphisms-section-phantom]{Compléments sur les morphismes de champs}
\item \hyperref[stacks-geometry-section-phantom]{La géométrie des champs}
\end{enumerate}
Thèmes de théorie des espaces de modules
\begin{enumerate}
\setcounter{enumi}{107}
\item \hyperref[moduli-section-phantom]{Champs de modules}
\item \hyperref[moduli-curves-section-phantom]{Espaces de modules de courbes}
\end{enumerate}
Divers
\begin{enumerate}
\setcounter{enumi}{109}
\item \hyperref[examples-section-phantom]{Exemples}
\item \hyperref[exercises-section-phantom]{Exercices}
\item \hyperref[guide-section-phantom]{Guide bibliographique}
\item \hyperref[desirables-section-phantom]{Desiderata}
\item \hyperref[coding-section-phantom]{Style de programmation}
\item \hyperref[obsolete-section-phantom]{Obsolète}
\item \hyperref[fdl-section-phantom]{Licence de documentation libre GNU}
\item \hyperref[index-section-phantom]{Index généré automatiquement}
\end{enumerate}
\end{multicols}


\bibliography{my} \bibliographystyle{amsalpha}

\end{document}
